\documentclass[10pt, letterpaper]{article}

% Inhaltsverzeichnis für Pakettypen (nur für Übersicht im Header, wird nicht im Dokument angezeigt)
% 1. Seitenlayout und Ränder
% 2. Sprache und Zeichensatz
% 3. Mathematik und Theorem-Umgebungen
% 4. Eigene Makros
% 5. Diagramme und Grafiken
% 6. Tabellen und Aufzählungen
% 7. Inhaltsverzeichnis
% 8. Abschnittsüberschriften
% 9. Abstrakt-Umgebung
% 10. Todos/Notizen
% 11. Rahmen/Box-Umgebungen
% 12. Python-Integration
% 13. Literaturverwaltung
% 14. Hyperlinks
% 15. Absatzeinstellungen
% 16. Umgebungen
% 17  Graphik
% 18  Extra
% 00. Titel und Autor

\usepackage{slashed}

% --- 1. Seitenlayout und Ränder ---
\usepackage[margin=3cm]{geometry}

% --- 2. Sprache und Zeichensatz ---
\usepackage[english]{babel}
\usepackage[T1]{fontenc}
\usepackage[utf8]{inputenc}

% --- 3. Mathematik und Theorem-Umgebungen ---
\usepackage{amsmath, amssymb, amsthm}
\usepackage{mathrsfs}
\DeclareMathOperator{\WF}{WF}

% --- 4. Eigene Makros ---
\usepackage{xcolor}
\newcommand{\SKP}{\langle\cdot,\cdot\rangle}
\newcommand{\R}{\mathbb{R}}
\newcommand{\N}{\mathbb{N}}
\newcommand{\Q}{\mathbb{Q}}
\newcommand{\Z}{\mathbb{Z}}
\newcommand{\C}{\mathbb{C}}
\newcommand{\entwurf}[1]{\textcolor{red}{#1}}

% --- 5. Diagramme und Grafiken ---
\usepackage{graphicx}
\graphicspath{{../../Library/Images/}}
\usepackage[export]{adjustbox}
\usepackage{tikz}
\usetikzlibrary{decorations.pathreplacing, arrows.meta, positioning}
\usepackage{tikz-cd}

% --- 6. Tabellen und Aufzählungen ---
\usepackage{enumitem}
\setlist[itemize]{left=0.5cm}

\newenvironment{romanenum}[1][]
  {%
    \ifx&#1&
    \else
      \textbf{#1}\quad
    \fi
    \begin{enumerate}[label=\roman*)]
  }
  {%
    \end{enumerate}%
  }

% --- 7. Inhaltsverzeichnis ---
\usepackage{tocloft}
\renewcommand{\cftsecfont}{\footnotesize}
\renewcommand{\cftsubsecfont}{\footnotesize}
\renewcommand{\cftsubsubsecfont}{\footnotesize}
\renewcommand{\cftsecpagefont}{\footnotesize}
\renewcommand{\cftsubsecpagefont}{\footnotesize}
\renewcommand{\cftsubsubsecpagefont}{\footnotesize}
\usepackage{etoc}

% --- 8. Abschnittsüberschriften ---
\usepackage{titlesec}
\titleformat{\section}{\normalfont\large\bfseries}{\thesection}{1em}{}
\titleformat{\subsection}{\normalfont\normalsize\bfseries}{\thesubsection}{0.5em}{}
\titleformat{\subsubsection}{\normalfont\normalsize\bfseries}{\thesubsubsection}{0.5em}{}
\setcounter{secnumdepth}{4}

% --- 9. Abstrakt-Umgebung ---
\usepackage{changepage}
\renewenvironment{abstract}
  {
    \begin{adjustwidth}{1.5cm}{1.5cm}
    \small
    \textsc{Abstract. –}%
  }
  {
    \end{adjustwidth}
  }

% --- 10. Todos/Notizen ---
\usepackage{todonotes}
% Hellblau definieren
\definecolor{hellblau}{rgb}{0.3,0.6,1}
\newenvironment{note}{\par\begingroup\color{hellblau}\itshape}{\par\endgroup}

% --- 11. Rahmen/Box-Umgebungen ---
\usepackage{mdframed}
\usepackage{tcolorbox}
\colorlet{shadecolor}{gray!25}

\newenvironment{customTheorem}
  {\vspace{10pt}%
   \begin{mdframed}[
     backgroundcolor=gray!20,
     linewidth=0pt,
     innertopmargin=10pt,
     innerbottommargin=10pt,
     skipabove=\dimexpr\topsep+\ht\strutbox\relax,
     skipbelow=\topsep,
   ]}
  {\end{mdframed}
   \vspace{10pt}%
  }

% --- 12. Python-Integration ---
% (Deaktiviert in dieser Version, aktiviere bei Bedarf)
% \usepackage{pythontex}
% \usepackage[makestderr]{pythontex}

% --- 13. Literaturverwaltung ---
\usepackage{csquotes}
\usepackage[backend=biber, style=alphabetic, citestyle=alphabetic]{biblatex}
\addbibresource{bibliography.bib}

% --- 14. Hyperlinks ---
\usepackage{hyperref}
\hypersetup{
  colorlinks   = true,
  urlcolor     = blue,
  linkcolor    = blue,
  citecolor    = blue,
  frenchlinks  = true
}

% --- 15. Absatzeinstellungen ---
\usepackage[parfill]{parskip}
\sloppy

% --- 16. Umgebungen ---
\usepackage{thmtools}

\newcommand{\CustomHeading}[3]{%
  \par\medskip\noindent%
  \textbf{#1 #2} \textnormal{(#3)}.\enskip%
}

\newenvironment{DEF}[2]{\begin{unitbox}\CustomHeading{Definition}{#1}{#2}}{\end{unitbox}}
\newenvironment{PROP}[2]{\begin{unitbox}\CustomHeading{Proposition}{#1}{#2}}{\end{unitbox}}
\newenvironment{THEO}[2]{\begin{unitbox}\CustomHeading{Theorem}{#1}{#2}}{\end{unitbox}}
\newenvironment{LEM}[2]{\begin{unitbox}\CustomHeading{Lemma}{#1}{#2}}{\end{unitbox}}
\newenvironment{KORO}[2]{\begin{unitbox}\CustomHeading{Corollar}{#1}{#2}}{\end{unitbox}}
\newenvironment{REM}[2]{\begin{unitbox}\CustomHeading{Remark}{#1}{#2}}{\end{unitbox}}
\newenvironment{EXA}[2]{\begin{unitbox}\CustomHeading{Example}{#1}{#2}}{\end{unitbox}}
\newenvironment{STUD}[2]{\begin{unitbox}\CustomHeading{Study}{#1}{#2}}{\end{unitbox}}
\newenvironment{CONC}[2]{\begin{unitbox}\CustomHeading{Concept}{#1}{#2}}{\end{unitbox}}
\newenvironment{OTH}[2]{\begin{unitbox}\CustomHeading{Other}{#1}{#2}}{\end{unitbox}}
\newenvironment{EXE}[2]{\begin{unitbox}\CustomHeading{Exercise}{#1}{#2}}{\end{unitbox}}
\newenvironment{MOT}[2]{\begin{unitbox}\CustomHeading{Motivation}{#1}{#2}}{\end{unitbox}}
\newenvironment{PROOF}[2]{\begin{unitbox}\CustomHeading{Proof}{#1}{#2}}{\end{unitbox}}

% --- Unit Umgebung für Source-Inhalte ---
\usepackage{mdframed}
\newmdenv[
  linewidth=1pt,
  topline=false,
  bottomline=false,
  rightline=false,
  leftmargin=0cm,
  rightmargin=0cm,
  skipabove=10pt,
  skipbelow=10pt,
  innertopmargin=0.5\baselineskip,
  innerbottommargin=0.5\baselineskip,
  backgroundcolor=gray!10,
  linecolor=gray
]{unitbox}

\newenvironment{unit}[1]
  {\begin{unitbox}\textbf{Unit #1}\par\smallskip}
  {\end{unitbox}}

% --- 17. ... ---

% --- 18. Extras ---
\usepackage{stmaryrd}
\usepackage{bbold}  % falls du \mathbb{1} nutzen willst

% --- 00. Titel und Autor ---
\title{Mein Titel}
\author{Tim Jaschik}
\date{\today}

\begin{document}

\maketitle
\rule{\textwidth}{0.5pt}
\begin{abstract}
Kurze Beschreibung …
\end{abstract}
\rule{\textwidth}{0.5pt}
\vspace{0.5cm}

\tableofcontents

\pagebreak


\section {Einleitung}
Riemanns fundamentaler Idee, Funktionentheorie nicht nur auf der komplexen Ebene oder Gebieten in der komplexen Ebene zu studieren, sondern auch kompliziertere Objekte - heutzutage Riemannsche Flächen genannte - zuzulassen, sind viele Bücher gewidmet. Ausgehend vom dynamischen System eines polygonalen Billardtisches treten Riemannsche Flächen mit einer Zusatzstruktur auf, es entstehen sogenannte flache Flächen.

Quellen: Das Material über Riemannsche Flächen ist eine Zusammenstellung aus den Büchern von Forster [FG81], Lamotke [Lam09] und Reyssat [Rey89]. Über flache Flächen ist derzeit kein Lehrbuch verfügbar. Der Abschnitt über affine Gruppen stammt aus [Vor96] und [GJ00]. Grundbegriffe aus der Funktionentheorie in verschiedenen Abschnitten stammen aus dem Buch von Freitag und Busam [FB06].

Voraussetzungen sind Grundbegriffe über komplexe Zahlen wie komplexe Differenzierbarkeit bis hin zum Cauchy-Integralsatz. Grundbegriffe der Topologie und Gruppentheorie sind im Anhang zusammengefasst.

Dieses Skript entstand aus einer 4 -std. Vorlesung im WiSe 10/11, in der die Kapitel 2, Abschnitt3, Abschnitt4, Abschnitt7, die Einleitung von Abschnitt6, die Hälfte von Abschnitt 8, Abschnitt 9 sowie Abschnitt 10 behandelt wurden. In einer 4-std. Vorlesung im WiSe 13/14 wurden Abschnitt 2, Abschnitt 3, Abschnitt 4, Abschnitt 5, Abschnitt 8, Abschnitt 9 und in Kurzform Abschnitt 7 diskutiert.

\subsection {Motivation 1: Die Riemannsche Fläche des Logarithmus}
Als erstes Beispiel zur Motivation von Riemannschen Flächen geben wir ein Beispiel, das Riemann ursprünglich (ca. 1851) zur Definition bewogen hat.

Die Potenzreihe $\sum_{n=0}^{\infty} \frac{z^{n}}{n!}$ ist für alle $z \in \mathbb{C}$ absolut konvergent und die Funktion $\exp (z)=\sum_{n=0}^{\infty} \frac{z^{n}}{n!}$ heißt komplexe Exponentialfunktion. Es gilt bekanntlich für $z, w \in \mathbb{C}$ und $t \in \mathbb{R}$

$$
\exp (z+w)=\exp (z) \cdot \exp (w), \quad \exp (i z)=\cos z+i \sin (z)
$$

Was ist das Problem, wenn wir definieren wollen: „Sei log $z$ die Umkehrfunktion der Exponentialfunktion"? Es soll natürlich

$$
\exp (\log (z))=\log (\exp (z))=z
$$

gelten. Wir müssen die Existenz eines solchen Wertes $\log z$ zeigen und auch, dass dieser Wert eindeutig bestimmt ist. Nach (1.1) ist für $z \in \mathbb{C}$

$$
\exp (z)=1 \Leftrightarrow \cos z=1 \text { und } \operatorname{Re} z=0
$$

also genau dann, wenn $z \in 2 \pi i \mathbb{Z}$. Folglich ist Eindeutigkeit nicht gegeben. Da allgemeiner $\exp (z)=\exp (2 \pi i+z)$ gilt, beschränken wir uns auf den Streifen

$$
S=\{z \in \mathbb{C}:-\pi<\operatorname{Im} z \leq \pi\}
$$

Wiederum mit Hilfe von (1.1) rechnet man nach, dass die Einschränkung von $\exp (\cdot)$ auf $S$ injektiv ist. Bezeichnet $S^{\circ}=S \backslash\{x+i \cdot \pi: x \in \mathbb{R}\}$ den Streifen ohne den oberen Rand, so ist

$$
\left.\exp \right|_{S^{\circ}}:\left\{\begin{array}{rll}
S^{\circ} & \longrightarrow & \mathbb{C}^{-}=\mathbb{C} \backslash \mathbb{R}_{\leq 0} \\
z & \longmapsto & \exp (z)
\end{array}\right.
$$

eine Bijektion. Folglich ist

$$
\log (z): \mathbb{C}^{-} \longrightarrow S^{\circ} \subseteq \mathbb{C}
$$

als Umkehrfunktion der Exponentialfunktion wohldefiniert. Man nennt diese Abbildung den Hauptzweig des Logarithmus. Aber $\log (z)$ kann nicht zu einer stetigen Abbildung $\mathbb{C}^{*} \longrightarrow \mathbb{C}$ fortgesetzt werden, denn

$$
\lim _{\substack{z \rightarrow-1 \\ \operatorname{Im}(z)>0}} \log (z)=+\pi i, \quad \text { aber } \quad \lim _{\substack{z \rightarrow-1 \\ \operatorname{Im}(z)<0}} \log (z)=-\pi i
$$

Für alle $t \in \mathbb{R}$ bezeichnen wir mit $S_{t}$ den verschobenen Streifen und mit $S_{t}^{\circ}$ sein Inneres,\\
$S_{t}=\{z \in \mathbb{C}:-\pi+t<\operatorname{Im}(z) \leq \pi+t\} \quad$ und $\quad S_{t}^{\circ}=S_{t} \backslash\{x+i(\pi+t), x \in \mathbb{R}\}$.\\
Das Bild von $\left.\exp \right|_{S_{t}^{0}}$ ist nun $e^{i t} \mathbb{C}^{-}=\mathbb{C} \backslash e^{i t} \mathbb{R}_{\leq 0}$ und $\left.\exp (\cdot)\right|_{S_{t}}$ ist wiederum eine Bijektion. Wir bezeichnen mit ${ }_{t} \log (\cdot)$ die Umkehrfunktion von $\left.\exp (\cdot)\right|_{S_{t}^{\circ}}$.

\begin{figure}[h]
\begin{center}
  \includegraphics[scale=0.2]{2025_10_06_374d4cfa8b1948eb4c41g-005}
 
\caption{Abbildung 1.1: Ein anderer Zweig \$\{ }\_\{t\} \textbackslash log (\textbackslash cdot)\$ des Logarithmus\}\end{center}
\end{figure}

Warum sollte man den Hauptzweig des Logarithmus bevorzugen? Riemanns Idee war, dass alle Zweige ${ }_{t} \log (\cdot)$ des Logarithmus für $t \in \mathbb{R}$ gleichberechtigt auftreten sollen. Wir setzen also alle geschlitzten Ebenen $e^{i t} \mathbb{C}^{-}$zusammen und definieren den Logarithmus dort.

Dazu fixieren wir $t_{2}<t_{1}$ mit $t_{1}-t_{2}<2 \pi$. Dann ist

$$
S_{t_{1}, t_{2}}=S_{t_{1}}^{\circ} \cap S_{t_{2}}^{\circ}=\left\{z \in \mathbb{C}:-\pi+t_{1}<\operatorname{Im}(z)<\pi+t_{2}\right\}
$$

nicht leer. Das Bild der Exponentialabbildung eingeschränkt auf $S_{t_{1}, t_{2}}$ ist der Sektor

$$
\mathbb{C}_{t_{1}, t_{2}}=\left\{w=r \cdot(\cos \varphi+i \sin \varphi):-\pi+t_{1}<\varphi<\pi+t_{2}, r>0\right\} .
$$

\begin{figure}[h]
\begin{center}
  \includegraphics[scale=0.2]{2025_10_06_374d4cfa8b1948eb4c41g-006}

\caption{Abbildung 1.2: Die Verklebungssektoren \$\textbackslash mathbb\{C}\textit{\{t}\{1\}, t\_\{2\}\}\$\}\end{center}
\end{figure}

Wir beobachten, dass die Einschränkungen der Zweige $t_{1} \log$ und $_{t_{2}} \log$ des Logarithmus

$$
\left.t_{1} \log \right|_{S_{t_{1}, t_{2}}} \text { und }\left.t_{2} \log \right|_{S_{t_{1}, t_{2}}}: \mathbb{C}_{t_{1}, t_{2}} \longrightarrow \mathbb{C}
$$

auf $S_{t_{1}, t_{2}}$ übereinstimmen. Folglich erhalten wir durch Verkleben der geschlitzten Ebenen $e^{i t} \cdot \mathbb{C}^{-}$entlang der Sektoren $\mathbb{C}_{t_{1}, t_{2}}$ ein neues Objekt, das lokal aus Teilgebieten von $\mathbb{C}$ zusammengesetzt ist. Deswegen können wir auf diesem neuen Objekt alles aus der Funktionentheorie gewohnt auch machen, also z.B. komplex differentieren. Zudem ist auf diesem Objekt „komplizierte" Funktion „komplexer Logarithmus" wohldefininiert. Diese Objekt ist unser erstes Beispiel einer Riemannschen Fläche, welche nicht ein Gebiet in $\mathbb{C}$ ist.

\subsection {Motivation 2: Polygonale Billardtische ([ZK])}
Will man die Trajektorien einer Billard-Kugel auf einem polygonalen - nicht notwendig rechteckigen - Billiardtisch studieren, stößt man auf das Problem, dass die Trajektorien nur stückweise Geraden sind und sich „wirr" selbst kreuzen. Reflektiert man statt der Kugel an der Bande die Bande selbst, bewegt sich die Kugel auf gerader Linie - allerdings auf einem viel komplizierteren Objekt. Denn sobald die Kugel erneut auf die Bande trifft, wird man wieder die Bande reflektieren müssen.

Jeden einzelnen Billardtisch fassen wir als Teil der komplexen Ebene auf. Das neue Objekt entsteht durch „Verkleben" des ursprünglichen polygonalen Billardtisches. Was soll „Verkleben" dabei genau bedeuten? Und wieso kann

\begin{figure}[h]
\begin{center}
  \includegraphics[scale=0.2]{2025_10_06_374d4cfa8b1948eb4c41g-007}
 
\caption{Abbildung 1.3: Entfalten eines Billiardtisches}
\end{center}
\end{figure}

man auf dem verklebten Objekt mehr über Billardtrajektorien aussagen als auf dem ursrprünglichen Billardtisch?

Die Fläche, die durch Entfalten des Billardtischs entsteht ist auch eine Riemannsche Fläche, sogar eine sehr spezielle Sorte davon, welche flache Fläche genannt wird.

Die Entfaltungskonstruktion taucht erstmal ca. 1975 (z.B. in [ZK75]) auf. Die resultierenden flachen Flächen sind ein aktuelles Forschungsgebiet der Riemannschen Flächen. Wir gehen im Abschnitt 8 darauf ein. Eine neuerer Übersichtsartikel hierzu ist [MT02].

\subsection {Ziele}
Von der Verklebefläche, auf der der Logarithmus definiert ist hat man in naürlicher Weise eine Abbildung auf die komplexe Ebene ohne den Nullpunkt. Diese Abbildung ist das erste Beispiel einer Überlagerung. Überlagerungen sind ein wichtige Methode, aus einer Riemannschen Fläche eine neue zu gewinnen. Welche Überlagerungen eine Riemannsche Fläche zulässt misst die Fundamentalgruppe. Diesen beiden Begriffen ist der Abschnitt 3 gewidmet, der etwas Topologie und Gruppentheorie voraussetzt.\\
Es gibt immer eine 'oberste' Überlagerung, die sogenannte universelle Überlagerung. Bei der Vielzahl von Riemannschen Flächen ist es erstaunlich, dass es nur drei Möglichkeiten (bis auf Isomorphie) für die universelle Überlagerung gibt. Dies ist ein Ziel von Abschnitt 7, in welchem der Laplace-Operator und damit Analysis eine zentrale Rolle spielt.\\
Die andere zentrale Frage über Riemannsche Flächen ist die nach der Existenz von Abbildungen. Die ist besonders für kompakte Riemannsche Flächen interessant. Die Bestimmung des Raums dieser Abbildungen ist unter dem Namen Satz von Riemann-Roch bekannt. Es gibt zwei Zugänge hierzu, zum einen mittels harmonischer Funktionen in Abschnitt 7, zum anderen mittels Garbenkohomologie basierend auf Abschnitt 5. Dieser algebraische Zugang hat auch viele andere Vorteile, liefert aber nicht den Uniformisierungssatz. Beide Zugänge setzen zunächst das Studium von Differenzialformen in Ab-\\
schnitt 4 voraus. Die beiden Zugänge zum Satz von Riemann-Roch werden in Abschnitt 6 vorgestellt.

Der Abschnitt 8 greift das zweite motivierende Beispiel auf und führt systematisch Riemannsche Flächen mit einer Zusatzstruktur ein, wie sie obige Entfaltungsflächen automatisch haben.


\pagebreak


\section {Riemannsche Flächen}
Im folgenden Abschnitt werden wir genau erklären, was der in der Einleitung benutzte Begriff des Verklebens formal bedeutet.

\subsection {Karten, Atlanten, komplexe Strukturen}
Riemannsche Flächen (wie auch Mannigfaltigkeiten) bestehen lokal aus der bekannten komplexen Ebene (oder $\mathbb{R}^{n}$ oder $\mathbb{C}^{n}$ ). Die Sturktur der Mannigfaltigkeit wird dadurch bestimmt, welche 'Güte' die Übergangsfunktionen zwischen den einzelnen lokalen Bausteinen haben.

\begin{figure}[h]
\begin{center}
  \includegraphics[scale=0.2]{2025_10_06_374d4cfa8b1948eb4c41g-008}
 
\caption{Abbildung 2.1: Kartenwechsel}
\end{center}
\end{figure}

Definition 2.1 (Karten, Atlanten, komplexe Strukturen) 


Sei $X$ ein topologischer Raum. Eine Karte $(U, h)$ von $X$ ist ein Homöomorphismus $h: U \rightarrow h(U) \subseteq \mathbb{C}$ von einer offenen Teilmenge $U \subseteq X$ auf eine offene Teilmenge von $\mathbb{C}$. Ein Atlas $A=\left\{\left(U_{i}, h_{i}\right)\right\}$ von $X$ ist eine Menge von Karten, sodass die Vereinigung der $U_{i}$ den topologischen Raum X überdeckt. Zu zwei Karten ( $U_{i}, h_{i}$ ) und ( $U_{j}, h_{j}$ ) heißt
$$
h_{j} \circ h_{i}^{-1}: h_{i}\left(U_{i} \cap U_{j}\right) \longrightarrow h_{j}\left(U_{i} \cap U_{j}\right)
$$
die Kartenwechselabbildung. Der Atlas A heißt komplexer Atlas, falls für alle Paare ( $i, j$ ) die Kartenwechselabbildung biholomorph ist. Zwei Atlanten $A= \left\{\left(U_{i}, h_{i}\right)\right\}$ und $B=\left\{\left(V_{i}, g_{i}\right)\right\}$ heißen äquivalent, falls für alle $\left(U_{i}, h_{i}\right) \in A$ und für alle $\left(V_{j}, g_{j}\right) \in B$ die Kartenwechselabbildung $g_{j} \circ h_{i}^{-1}$ biholomorph ist.

Eine Äquivalenzklasse von Atlanten heißt komplexe Struktur auf $X$. 



\begin{DEF}{RS-M25-01-01}{Komplexe Struktur auf einem topologischen Raum}
Sei $X$ ein topologischer Raum.
\begin{enumerate}
\item Eine \textbf{Karte} von $X$ ist ein Homöomorphismus 
\[
h:U\longrightarrow h(U)\subseteq\mathbb{C},
\]
wobei $U\subseteq X$ offen ist. Das Paar $(U,h)$ heißt dann \textbf{lokale Karte} von $X$.

\item Eine \textbf{Überdeckung} $A=\{(U_i,h_i)\}_{i\in I}$ von $X$ durch Karten heißt \textbf{Atlas}, falls
\[
X=\bigcup_{i\in I}U_i.
\]
Für zwei Karten $(U_i,h_i)$ und $(U_j,h_j)$ heißt die Abbildung
\[
h_j\circ h_i^{-1}: h_i(U_i\cap U_j)\longrightarrow h_j(U_i\cap U_j)
\]
\textbf{Kartenwechselabbildung}.

\item Der Atlas $A$ heißt \textbf{komplexer Atlas}, falls für alle $i,j\in I$ die Kartenwechselabbildungen
\[
h_j\circ h_i^{-1}
\]
\emph{biholomorph} sind (d.\,h. holomorph mit holomorpher Inversen).

\item Zwei Atlanten 
\[
A=\{(U_i,h_i)\}_{i\in I}
\quad\text{und}\quad 
B=\{(V_j,g_j)\}_{j\in J}
\]
heißen \textbf{äquivalent}, falls jede Kartenwechselabbildung zwischen einer Karte aus $A$ und einer Karte aus $B$
\[
g_j\circ h_i^{-1}: h_i(U_i\cap V_j)\longrightarrow g_j(U_i\cap V_j)
\]
biholomorph ist.

\item Eine \textbf{komplexe Struktur} auf $X$ ist eine Äquivalenzklasse komplexer Atlanten unter dieser Äquivalenzrelation. 

Ein Paar $(X,\mathcal{A})$, bestehend aus einem topologischen Raum $X$ und einer komplexen Struktur $\mathcal{A}$, heißt \textbf{komplexe Mannigfaltigkeit (der Dimension 1)}.
\end{enumerate}
\end{DEF}



Zu einem gegebenen Atlas $A=\left\{\left(U_{i}, h_{i}\right)\right\}$ von $X$ und $x \in X$ sei $(U, h)$ eine Karte einer Umgebung von $x$. Dann ist auch $(U, \widetilde{h})$ mit $\widetilde{h}(z)=h(z)-h(x)$ eine Karte bei $x$. Der Atlas $A \cup(U, \widetilde{h})$ ist offenbar zu $A$ äquivalent, da die Verschiebung um $h(x)$ eine holomorphe Abbildung ist. Führt man dies für alle $x \in X$ durch, so sieht man, dass in einer komplexe Struktur stets eine Karte enthalten ist, die einen gegebenen Punkt auf Null abbildet. Dies ist für lokale Rechnung mit Potenzreihen sehr nützlich.



\begin{LEM}{RS-M25-01-02}{Lokale Nullkarten in einer komplexen Struktur}
Sei $A=\{(U_i,h_i)\}$ ein (komplexer) Atlas auf einem topologischen Raum $X$, und sei $x\in X$. 
\begin{enumerate}
\item Ist $(U,h)$ eine Karte mit $x\in U$, so definiert
\[
\widetilde{h}:U\longrightarrow\mathbb{C},\qquad \widetilde{h}(z):=h(z)-h(x)
\]
ebenfalls eine Karte bei $x$. Dabei gilt $\widetilde{h}(x)=0$.
\item Der erweiterte Atlas $A\cup\{(U,\widetilde{h})\}$ ist zum ursprünglichen Atlas $A$ äquivalent, da die Abbildung
\[
\phi:\mathbb{C}\longrightarrow\mathbb{C},\qquad \phi(w)=w-h(x)
\]
biholomorph ist und somit auch die Kartenwechselabbildung $h\circ \widetilde{h}^{-1}=\phi^{-1}$ biholomorph bleibt.
\end{enumerate}
\textbf{Folgerung.}  
Zu jedem Punkt $x\in X$ existiert in jeder komplexen Struktur eine Karte $(U,\widetilde{h})$, die $x$ auf $0$ abbildet.  
Solche „Nullkarten“ sind insbesondere für lokale Rechnungen mittels Potenzreihen oder Taylorentwicklungen äußerst nützlich.
\end{LEM}



Definition 2.2 (Riemannsche Fläche) 

Eine Riemannsche Fläche ist ein hausdorffscher topologischer Raum mit einer komplexen Struktur. Eine Riemannsche Fläche heißt kompakt, wenn der zugrundeliegende topologische Raum kompakt ist.


\begin{DEF}{RS-M25-01-03}{Riemannsche Fläche}
Eine \textbf{Riemannsche Fläche} ist ein \textit{hausdorffscher topologischer Raum} $X$, der mit einer \textbf{komplexen Struktur} versehen ist, d.\,h. mit einer Äquivalenzklasse komplexer Atlanten $\mathcal{A}$, deren Karten lokale Homöomorphismen auf offene Teilmengen von $\mathbb{C}$ sind und deren Kartenwechselabbildungen biholomorph sind.  
Das Paar $(X,\mathcal{A})$ heißt dann komplexe Mannigfaltigkeit der komplexen Dimension~$1$.  

Eine Riemannsche Fläche heißt \textbf{kompakt}, wenn der zugrunde liegende topologische Raum $X$ kompakt ist.
\end{DEF}



>> Riemannische Fläche durch komplex konjugierter Atlant


\begin{DEF}{RS-M25-01-04}{Komplex konjugierte Riemannsche Fläche}
Sei $(X,\mathcal{A})$ eine Riemannsche Fläche mit komplexem Atlas 
$\mathcal{A}=\{(U_i,h_i)\}_{i\in I}$, wobei jede Kartenabbildung 
$h_i:U_i\to h_i(U_i)\subset\mathbb{C}$ holomorph ist.  

\textbf{(a)} Die \emph{komplex konjugierte Riemannsche Fläche} 
$\overline{X}$ ist definiert als derselbe topologische Raum $X$, versehen 
mit dem \emph{komplex konjugierten Atlas} 
\[
\overline{\mathcal{A}}:=\{(U_i,\overline{h_i})\}_{i\in I},
\]
wobei $\overline{h_i}(p):=\overline{h_i(p)}$ für $p\in U_i$ gilt.  

\textbf{(b)} Die Kartenwechselabbildungen von $\overline{\mathcal{A}}$ 
sind genau die komplex konjugierten derjenigen von $\mathcal{A}$ und 
damit antiholomorph. Somit trägt $\overline{X}$ eine wohldefinierte 
komplexe Struktur, die als „konjugierte Struktur“ zu der von $X$ gilt.

\textbf{(c)} Zwischen $X$ und $\overline{X}$ existiert die 
\emph{kanonische Identität} 
\[
\operatorname{id}_X : X \longrightarrow \overline{X},
\]
die eine \emph{antiholomorphe Abbildung} ist.
\end{DEF}



Beispiele 2.3 


Die komplexe Ebene $\mathbb{C}$, die obere Halbebene $\mathbb{H}=\{z \in \mathbb{C}$ : $\operatorname{Im}(z)>0\}$ und die Einheitskreisscheibe $\triangle=\{z:|z|<1\}$ sind Riemannsche Flächen, gegeben durch eine Karte, 

ebenso $\mathbb{C}^{*}=\mathbb{C} \backslash\{0\}$ und $\triangle^{*}=\triangle \backslash\{0\}$ die punktierte Ebene (bzw. punktierte Einheitskreisscheibe).


\begin{EXA}{RS-M25-01-05}{Grundlegende Beispiele Riemannscher Flächen}
\begin{enumerate}
\item \textbf{Komplexe Ebene.}  
Die komplexe Ebene $\mathbb{C}$ trägt die natürliche komplexe Struktur, die durch die \emph{globale Karte}
\[
(\mathbb{C}, \operatorname{id}_{\mathbb{C}})
\]
gegeben ist. Damit ist $\mathbb{C}$ eine Riemannsche Fläche (sogar eine komplexe Mannigfaltigkeit der Dimension $1$).

\item \textbf{Obere Halbebene.}  
Die Menge
\[
\mathbb{H} := \{ z \in \mathbb{C} \mid \operatorname{Im}(z) > 0 \}
\]
ist eine offene Teilmenge von $\mathbb{C}$ und wird mit der \emph{restriktierten Karte}
\[
(\mathbb{H}, \operatorname{id}_{\mathbb{C}}|_{\mathbb{H}})
\]
zu einer Riemannschen Fläche. Die komplexe Struktur ist somit diejenige, die von der Standardstruktur auf $\mathbb{C}$ induziert wird.

\item \textbf{Einheitskreisscheibe.}  
Ebenso ist
\[
\triangle := \{ z \in \mathbb{C} \mid |z| < 1 \}
\]
eine offene Teilmenge von $\mathbb{C}$ und trägt über die Karte
\[
(\triangle, \operatorname{id}_{\mathbb{C}}|_{\triangle})
\]
die induzierte komplexe Struktur. Auch $\triangle$ ist daher eine Riemannsche Fläche.
\end{enumerate}
\textbf{Bemerkung.}  
Diese drei Beispiele sind lokal biholomorph zueinander — insbesondere sind $\mathbb{H}$ und $\triangle$ durch eine Möbiustransformation biholomorph, also komplex-äquivalente Modelle derselben lokalen Geometrie.
\end{EXA}



>> Unverträglichkeit von $\mapsto x+iy$ und $\mapsto x-iy$ bei $\triangle$


\begin{CONC}{RS-M25-01-06}{Unverträglichkeit konjugierter Kartenabbildungen}
Betrachten wir die Einheitskreisscheibe 
\[
\triangle = \{ z \in \mathbb{C} \mid |z| < 1 \}
\]
mit der Standardkarte 
\[
h : \triangle \to \mathbb{C}, \quad h(x + i y) = x + i y.
\]
Ersetzt man $h$ durch die „spiegelbildliche“ Karte
\[
\tilde{h}(x + i y) = x - i y = \overline{z},
\]
so beschreibt $\tilde{h}$ denselben topologischen Raum, aber mit \emph{komplex konjugierter Orientierung}.  

Die beiden Karten $(\triangle, h)$ und $(\triangle, \tilde{h})$ sind \emph{nicht kompatibel} im Sinne der Definition einer komplexen Struktur, denn die Kartenwechselabbildung
\[
\tilde{h} \circ h^{-1} : h(\triangle) \to \tilde{h}(\triangle), 
\quad z \longmapsto \overline{z}
\]
ist \emph{nicht holomorph}, sondern \emph{antiholomorph}.  

\textbf{Folge.}  
Die Strukturen, die durch $h$ und $\tilde{h}$ definiert werden, sind zueinander \emph{komplex konjugiert}, aber nicht äquivalent als komplexe Atlanten.  
Sie definieren also zwei verschiedene komplexe Strukturen auf demselben topologischen Raum~$\triangle$ — die eine mit der „gewöhnlichen“ Orientierung, die andere mit der entgegengesetzten.
\end{CONC}



Beispiel 2.4 (Der Torus) 

Als zugrundeliegender topologischer Raum sei $X= S^{1} \times S^{1}$ gegeben, wobei $S^{1}=\{z \in \mathbb{C}:|z|=1\}$. Zu $w_{1}, w_{2} \in \mathbb{C}^{*}$ mit $\frac{w_{2}}{w_{1}} \notin \mathbb{R}$ definieren wir eine Riemannsche Fläche wie folgt. Sei (für $t_{1}, t_{2} \in \mathbb{R}$ )
$$
q:\left\{\begin{array}{ccc}
\mathbb{C} & \longrightarrow & S^{1} \times S^{1} \\
z=t_{1} \cdot w_{1}+t_{2} \cdot w_{2} & \longmapsto & \left(\exp \left(2 \pi i t_{1}\right), \exp \left(2 \pi i_{2}\right)\right)
\end{array}\right.
$$
Die Abbildung $q$ ist surjektiv und jeder Punkt $z \in \mathbb{C}$ besitzt eine Umgebung $U_{z}$, sodass $U_{z} \cap\left(w+U_{z}\right)=\emptyset$ für alle $w \in \mathbb{Z} w_{1}+\mathbb{Z} w_{2}$ mit $w \neq 0$ gilt. Folglich ist $\left.q\right|_{U_{z}}$ ein Homöomorphismus und die Abbildungen $\left(\left.q\right|_{U_{z}}\right)^{-1}$ bilden einen Atlas von $S^{1} \times S^{1}$.

Ist $q\left(U_{1}\right) \cap q\left(U_{2}\right) \neq \emptyset$ für zwei solche Umgebungen $U_{1}, U_{2}$, so ist die Kartenwechselabbildung $\left(\left.q\right|_{U_{1}}\right)^{-1} \circ\left(\left.q\right|_{U_{2}}\right)$ gegeben durch $z \mapsto z+w$ für ein $w \in \mathbb{Z} w_{1}+\mathbb{Z} w_{2}$, also ein Biholomorphismus. Die Riemannsche Fläche $X$ mit diesen Karten heißt der Torus mit Perioden $w_{1}$ und $w_{2}$.
$$
\includegraphics[scale=0.2]{2025_10_06_374d4cfa8b1948eb4c41g-010}
$$



\begin{EXA}{RS-M25-01-08}{$2$-Torus als Riemannsche Fläche mit Perioden}
Als zugrundeliegender topologischer Raum betrachten wir das Produkt
\[
X := S^{1} \times S^{1}, \qquad \text{wobei } S^{1} = \{ z \in \mathbb{C} \mid |z| = 1 \}.
\]
Seien nun $w_{1}, w_{2} \in \mathbb{C}^{*}$ so gewählt, dass $\frac{w_{2}}{w_{1}} \notin \mathbb{R}$ gilt — das heißt, $w_{1}$ und $w_{2}$ spannen ein \emph{Gitter}
\[
\Lambda := \mathbb{Z} w_{1} + \mathbb{Z} w_{2}
\]
in der komplexen Ebene $\mathbb{C}$ auf.

\textbf{(a) Definition der Quotientenabbildung.}  
Wir definieren eine surjektive Abbildung
\[
q : \mathbb{C} \longrightarrow S^{1} \times S^{1}, \qquad 
q(t_{1} w_{1} + t_{2} w_{2}) := \left( e^{2\pi i t_{1}}, \, e^{2\pi i t_{2}} \right)
\quad \text{für } t_{1}, t_{2} \in \mathbb{R}.
\]
Zwei Punkte $z, z' \in \mathbb{C}$ werden genau dann auf denselben Punkt in $S^{1} \times S^{1}$ abgebildet, wenn $z - z' \in \Lambda$ gilt.

\textbf{(b) Lokale Eigenschaften.}  
Für jedes $z \in \mathbb{C}$ existiert eine offene Umgebung $U_{z} \subset \mathbb{C}$, sodass
\[
U_{z} \cap (w + U_{z}) = \varnothing
\quad \text{für alle } w \in \Lambda \setminus \{0\}.
\]
Damit ist $\left.q\right|_{U_{z}} : U_{z} \to q(U_{z})$ ein Homöomorphismus.  
Die inversen Abbildungen $\bigl( \left.q\right|_{U_{z}} \bigr)^{-1}$ bilden einen Atlas von Karten auf $X = S^{1} \times S^{1}$.

\textbf{(c) Kartenwechselabbildungen.}  
Falls $q(U_{1}) \cap q(U_{2}) \neq \varnothing$, dann gibt es $w \in \Lambda$ mit
\[
\left( \left.q\right|_{U_{1}} \right)^{-1} \circ \left( \left.q\right|_{U_{2}} \right)(z)
= z + w.
\]
Dies ist eine holomorphe Abbildung, da Translationen in $\mathbb{C}$ biholomorph sind.  
Somit sind alle Kartenwechselabbildungen biholomorph, und der Atlas definiert eine komplexe Struktur auf $S^{1} \times S^{1}$.

\textbf{(d) Ergebnis.}  
Die so erhaltene Riemannsche Fläche
\[
X \cong \mathbb{C} / \Lambda
\]
heißt \emph{Torus mit Perioden $w_{1}$ und $w_{2}$}.  
Jede Wahl linear unabhängiger Perioden $w_{1}, w_{2}$ mit $\frac{w_{2}}{w_{1}} \notin \mathbb{R}$ liefert eine komplexe Struktur auf dem topologischen Torus~$S^{1} \times S^{1}$.

\begin{center}
  \includegraphics[scale=0.2]{2025_10_06_374d4cfa8b1948eb4c41g-010}
\end{center}

\textbf{Bemerkung.}  
Die entstehende Fläche $\mathbb{C}/\Lambda$ kann äquivalent als Quotient mit der \emph{Quotiententopologie} beschrieben werden, wobei offene Mengen $U \subseteq \mathbb{C}/\Lambda$ genau diejenigen sind, deren Urbilder $\pi^{-1}(U)$ offen in $\mathbb{C}$ sind. Diese Struktur macht den Torus zu einer kompakten Riemannschen Fläche vom Geschlecht~$1$.
\end{EXA}


>> Quotiententopologie auf Torus


\begin{DEF}{RS-M25-01-07}{Quotiententopologie auf dem Torus}
Sei $\pi : \mathbb{R}^2 \to \mathbb{R}^2 / \mathbb{Z}^2$ die kanonische Projektion, welche Punkte identifiziert, die sich um ganzzahlige Gitterverschiebungen unterscheiden:
\[
(x, y) \sim (x', y') \quad \Longleftrightarrow \quad (x - x', y - y') \in \mathbb{Z}^2 .
\]
\textbf{(a) Definition der Topologie.}  
Auf dem Quotientenraum
\[
\mathbb{T}^2 := \mathbb{R}^2 / \mathbb{Z}^2
\]
wird die \emph{Quotiententopologie} definiert durch:
\[
U \subseteq \mathbb{T}^2 \text{ ist offen } \; \Longleftrightarrow \; 
\pi^{-1}(U) \subseteq \mathbb{R}^2 \text{ ist offen.}
\]
Damit wird $\pi$ eine stetige und offene Abbildung, und $\mathbb{T}^2$ ein kompakter, hausdorffscher topologischer Raum.

\textbf{(b) Interpretation.}  
Anschaulich entsteht $\mathbb{T}^2$ durch Identifikation der gegenüberliegenden Ränder des Einheitsquadrats:
\[
[0,1] \times [0,1] \; \text{mit den Randpunkten identifiziert gemäß} \;
(x,0) \sim (x,1), \quad (0,y) \sim (1,y).
\]
\textbf{(c) Komplexe Struktur.}  
Die Projektion $\pi$ ist lokal ein Homöomorphismus, sodass $\mathbb{T}^2$ mit der von $\mathbb{R}^2 \simeq \mathbb{C}$ induzierten Kartenstruktur eine \emph{Riemannsche Fläche} wird.  
In komplexer Schreibweise kann man auch $\mathbb{T}^2 \cong \mathbb{C}/\Lambda$ betrachten, wobei $\Lambda \subset \mathbb{C}$ ein Gitter ist (etwa $\Lambda = \mathbb{Z} + i \mathbb{Z}$).

\textbf{Bemerkung.}  
Die Quotiententopologie ist hierbei entscheidend, da sie sicherstellt, dass lokale Eigenschaften (z.\,B. Offenheit, Stetigkeit, Holomorphie) über die Projektion $\pi$ mit der Struktur auf $\mathbb{C}$ kompatibel bleiben.
\end{DEF}



>> Alternative Konstruktion von Karten durch Bild gleich Torus minus parametrisierten Meridian $e^{2\pi it_1}$ und Longitude $e^{2\pi it_2}$


>> Translationsatlas


\begin{DEF}{RS-M25-01-09}{Translationsatlas}
Sei $\Lambda = \mathbb{Z} w_{1} + \mathbb{Z} w_{2} \subset \mathbb{C}$ ein Gitter, erzeugt von zwei linear unabhängigen Perioden $w_{1}, w_{2} \in \mathbb{C}$ mit $\frac{w_{2}}{w_{1}} \notin \mathbb{R}$.  
Dann definiert die kanonische Projektion
\[
\pi : \mathbb{C} \longrightarrow \mathbb{C}/\Lambda, \qquad \pi(z) = z + \Lambda
\]
eine Riemannsche Fläche, deren Karten durch lokale Inverse von $\pi$ gegeben sind.

\textbf{(a) Lokale Karten.}  
Für jedes $z \in \mathbb{C}$ existiert eine Umgebung $U_{z}$, so dass
\[
U_{z} \cap (U_{z} + w) = \varnothing \quad \text{für alle } w \in \Lambda \setminus \{0\}.
\]
Die Einschränkung $\pi|_{U_{z}}$ ist ein Homöomorphismus auf ihr Bild, und die Inversen
\[
\phi_{z} := \left(\pi|_{U_{z}}\right)^{-1} : \pi(U_{z}) \longrightarrow U_{z}
\]
bilden zusammen einen Atlas
\[
\mathcal{A} = \{ (\pi(U_{z}), \phi_{z}) \mid z \in \mathbb{C} \},
\]
der die komplexe Struktur auf $\mathbb{C}/\Lambda$ bestimmt.

\textbf{(b) Kartenwechsel.}  
Für zwei überlappende Karten $(\pi(U_{z_{1}}), \phi_{z_{1}})$ und $(\pi(U_{z_{2}}), \phi_{z_{2}})$ gilt:
\[
\phi_{z_{1}} \circ \phi_{z_{2}}^{-1}(w) = w + (z_{1} - z_{2}),
\]
also eine \emph{Translation} in $\mathbb{C}$.

\textbf{(c) Bezeichnung.}  
Ein Atlas, dessen Kartenwechselabbildungen durch komplexe Translationen gegeben sind,
\[
z \longmapsto z + c, \quad c \in \mathbb{C},
\]
heißt \emph{Translationsatlas}.  
Er ist automatisch ein komplexer Atlas, da Translationen biholomorph sind.

\textbf{(d) Beispiel.}  
Der Torus $\mathbb{C}/\Lambda$ trägt den kanonischen Translationsatlas, der direkt von der komplexen Ebene $\mathbb{C}$ induziert wird.  
Alle kompakten Riemannschen Flächen vom Geschlecht~$1$ können auf diese Weise beschrieben werden.
\end{DEF}



>> Universelle Abbildungs Eigenschaft


\begin{DEF}{RS-M25-01-10}{Universelle Abbildungseigenschaft der Quotientenabbildung}
Sei $X$ ein topologischer Raum, $\sim$ eine Äquivalenzrelation auf $X$ und 
\[
\pi : X \longrightarrow X/{\sim}, \qquad x \longmapsto [x]
\]
die kanonische Quotientenabbildung.  
Dann erfüllt $\pi$ die folgende \textbf{universelle Abbildungseigenschaft (UAE)}:

\textbf{(a) Aussage.}  
Sei $Y$ ein topologischer Raum und $f : X \to Y$ eine \emph{stetige} Abbildung, die \emph{äquivalenzverträglich} ist, d.\,h.
\[
x \sim x' \quad \Longrightarrow \quad f(x) = f(x').
\]
Dann existiert genau eine stetige Abbildung
\[
\tilde{f} : X/{\sim} \longrightarrow Y
\]
so dass das folgende Diagramm kommutiert:
\[
\begin{tikzcd}
X \arrow[rr, "f"] \arrow[d, "\pi"'] & & Y \\
X/{\sim} \arrow[urr, "\tilde{f}"', dashed]
\end{tikzcd}
\]
d.h.
\[
f = \tilde{f} \circ \pi.
\]

\textbf{(b) Charakterisierung der Quotiententopologie.}  
Die Topologie auf $X/{\sim}$ ist durch die Forderung eindeutig bestimmt, dass $\pi$ \emph{stetig} und \emph{universell initial} bezüglich dieser Eigenschaft ist.  
Konkret heißt das:
\[
U \subseteq X/{\sim} \text{ ist offen } 
\;\Longleftrightarrow\; 
\pi^{-1}(U) \subseteq X \text{ ist offen}.
\]

\textbf{(c) Anwendung.}  
Diese Eigenschaft erlaubt es, stetige Abbildungen $f : X \to Y$, die auf Äquivalenzklassen konstant sind, eindeutig zu „faktorisieren“ über den Quotientenraum $X/{\sim}$.  
In geometrischen Situationen — etwa bei der Konstruktion des Tori $\mathbb{C}/\Lambda$ oder anderer Riemannscher Flächen durch Identifikation — garantiert die UAE, dass der Quotient die „richtige“ (d.h. stetigste mögliche) Topologie trägt.

\textbf{(d) Beispiel.}  
Für das Gitter $\Lambda = \mathbb{Z}w_{1} + \mathbb{Z}w_{2} \subset \mathbb{C}$ sei $\pi : \mathbb{C} \to \mathbb{C}/\Lambda$ die Projektion.  
Ist $f : \mathbb{C} \to Y$ $\Lambda$-periodisch (d.\,h. $f(z+w) = f(z)$ für alle $w \in \Lambda$), so existiert nach der UAE genau eine stetige Abbildung
\[
\tilde{f} : \mathbb{C}/\Lambda \longrightarrow Y
\quad \text{mit} \quad f = \tilde{f} \circ \pi.
\]
\end{DEF}


Definition: Vekleben von topologischen Räumen

Seien für $i \in I$ topologische Räume $X_{i}$ gegeben und für $j \in J \subset \mathbb{N}$ seien topologische Räume $U_{j}$ zusammen mit Abbildungen $\varphi_{1}: U_{j} \rightarrow X_{i_{1}}$ und $\varphi_{2}: U_{j} \rightarrow X_{i_{2}}$ gegeben, wobei $i_{1}$ und $i_{2}$ von $j$ abhängen. Sind die Bilder von $\varphi_{1}$ und $\varphi_{2}$ offen und die Abbildungen $\varphi_{1}, \varphi_{2}$ für alle $j \in J$ Homöomorphismen auf ihr Bild, so nennt man den topologischen Raum $X=\left(\coprod_{i \in I} X_{i}\right) / \sim$, mit der Äquivalenzrelation
$$
\begin{aligned}
x_{1} \sim x_{2} \Leftrightarrow & \exists j \in J \quad \text { und } \quad y \in U_{j}: \\
x_{1}=\varphi_{1}(y), x_{2}=\varphi_{2}(y) & \quad \text { oder } \quad x_{2}=\varphi_{1}(y), x_{1}=\varphi_{2}(y)
\end{aligned}
$$
das Verkleben von $X_{i}$ längs der $U_{j}$. Wir versehen $X$ vermöge der Abbildung $\coprod_{i \in I} X_{i} \longrightarrow X$ mit der Quotiententopologie. 



\begin{DEF}{RS-M25-01-11}{Verkleben von topologischen Räumen}
Sei $I$ eine Indexmenge und seien für jedes $i \in I$ topologische Räume $X_{i}$ gegeben.  
Ferner sei $J \subseteq \mathbb{N}$ eine (endliche oder abzählbare) Indexmenge und für jedes $j \in J$ sei ein topologischer Raum $U_{j}$ zusammen mit zwei stetigen Abbildungen
\[
\varphi_{1}^{(j)} : U_{j} \longrightarrow X_{i_{1}(j)}, 
\qquad
\varphi_{2}^{(j)} : U_{j} \longrightarrow X_{i_{2}(j)},
\]
gegeben, wobei $i_{1}(j), i_{2}(j) \in I$ von $j$ abhängen dürfen.

\textbf{(a) Voraussetzungen.}  
Wir fordern:
\begin{itemize}
  \item Die Bilder $\varphi_{1}^{(j)}(U_{j}) \subset X_{i_{1}(j)}$ und $\varphi_{2}^{(j)}(U_{j}) \subset X_{i_{2}(j)}$ seien offen.
  \item Die Abbildungen $\varphi_{1}^{(j)}$ und $\varphi_{2}^{(j)}$ seien Homöomorphismen auf ihr jeweiliges Bild.
\end{itemize}

\textbf{(b) Konstruktion der Äquivalenzrelation.}  
Betrachte die disjunkte Vereinigung
\[
\bigsqcup_{i \in I} X_{i}
\]
und definiere eine Äquivalenzrelation $\sim$ durch
\[
x_{1} \sim x_{2}
\quad \Longleftrightarrow \quad
\exists j \in J \; \exists y \in U_{j} :
\bigl( x_{1} = \varphi_{1}^{(j)}(y),\, x_{2} = \varphi_{2}^{(j)}(y) \bigr)
\;\text{ oder }\;
\bigl( x_{2} = \varphi_{1}^{(j)}(y),\, x_{1} = \varphi_{2}^{(j)}(y) \bigr).
\]

\textbf{(c) Definition des Verklebungsraums.}  
Der \emph{verklebte Raum} wird definiert als der Quotient
\[
X := \left( \bigsqcup_{i \in I} X_{i} \right) \Big/ \sim
\]
versehen mit der \emph{Quotiententopologie}, d.\,h. eine Menge $U \subseteq X$ ist offen genau dann, wenn ihr Urbild unter der kanonischen Projektion
\[
\pi : \bigsqcup_{i \in I} X_{i} \longrightarrow X
\]
offen ist.

\textbf{(d) Bedeutung.}  
Man sagt, $X$ entsteht durch \emph{Verkleben der Räume $X_{i}$ entlang der Übergangsabbildungen $\varphi_{1}^{(j)}$ und $\varphi_{2}^{(j)}$}.  
Intuitiv werden dabei offene Teilmengen $\varphi_{1}^{(j)}(U_{j}) \subset X_{i_{1}(j)}$ und $\varphi_{2}^{(j)}(U_{j}) \subset X_{i_{2}(j)}$ mittels der Homöomorphismen $\varphi_{2}^{(j)} \circ (\varphi_{1}^{(j)})^{-1}$ identifiziert.

\textbf{(e) Beispiel.}  
Die Riemannsche Fläche des Torus kann als Verkleben der Kanten eines Einheitsquadrats $[0,1] \times [0,1] \subset \mathbb{C}$ konstruiert werden, indem man die gegenüberliegenden Seiten durch die Identifikationen
\[
(x,0) \sim (x,1), \qquad (0,y) \sim (1,y)
\]
verklebt.  
Die Quotiententopologie macht den entstehenden Raum homöomorph zu $\mathbb{C}/(\mathbb{Z} + i\mathbb{Z})$.
\end{DEF}

Die folgende Aussage ist unmittelbare Konsequenz der Definitionen.

Proposition 2.5 

Sind die $X_{i}$ und $U_{j}$ offene Teilmengen von $\mathbb{C}$, die Inklusionsabbildungen $\varphi_{1}$ und $\varphi_{2}$ für alle $j \in J$ biholomorph auf ihr Bild und der durch Verkleben entstehende Raum $X$ hausdorffsch, so ist $X$ eine Riemannsche Fläche.



\begin{EXA}{RS-M25-01-12}{Verkleben offener Teilmengen der Ebene liefert eine Riemannsche Fläche}
Seien $\{X_i\}_{i\in I}$ und $\{U_j\}_{j\in J}$ offene Teilmengen von $\C$. Für jedes $j\in J$ gebe es Indizes
$i_1(j),i_2(j)\in I$ und Inklusionen
\[
\varphi^{(j)}_1:U_j\hookrightarrow X_{i_1(j)},\qquad
\varphi^{(j)}_2:U_j\hookrightarrow X_{i_2(j)},
\]
so dass $\varphi^{(j)}_1$ und $\varphi^{(j)}_2$ Homöomorphismen auf ihre (offenen) Bilder sind und die \emph{Übergangskarten}
\[
\varphi^{(j)}_2\circ\bigl(\varphi^{(j)}_1\bigr)^{-1}:
\varphi^{(j)}_1(U_j)\longrightarrow \varphi^{(j)}_2(U_j)
\]
biholomorph sind. Sei $X$ der durch Verkleben entstehende Quotientraum
\(
X=\bigl(\bigsqcup_{i\in I}X_i\bigr)\big/\!\sim
\)
mit der Quotiententopologie, und sei $X$ hausdorffsch. Dann trägt $X$ eine komplexe Struktur; insbesondere ist $X$ eine Riemannsche Fläche.


\begin{proof}
Bezeichne die kanonische Quotientenabbildung mit
\(
\pi:\bigsqcup_{i\in I}X_i\to X.
\)
Für $x\in X_i$ wählen wir eine offene Umgebung $V\subset X_i$, die entweder disjunkt zu allen verklebten Teilmengen ist
oder über genau eine Verklebung $\varphi^{(j)}_1,\varphi^{(j)}_2$ mit einem offenen Teil von $X_{i_2(j)}$ identifiziert wird.
Nach der Definition des Verklebens ist dann die Einschränkung
\[
\pi|_V:V\longrightarrow \pi(V)
\]
ein Homöomorphismus. Definiere eine Karte bei $\pi(x)$ durch
\[
h_V := (\pi|_V)^{-1}:\ \pi(V)\longrightarrow V\subset\C.
\]
Die durch alle solchen $(\pi(V),h_V)$ gebildeten Karten überdecken $X$, denn jedes $x\in\bigsqcup X_i$ besitzt eine solche Umgebung
$V$.

Es bleibt zu zeigen, dass die Kartenwechselabbildungen biholomorph sind. Seien dazu $(\pi(V_1),h_{V_1})$ und
$(\pi(V_2),h_{V_2})$ zwei Karten mit nichtleerem Überlapp. Ist $V_1,V_2\subset X_i$ für denselben Index $i$, so ist
\[
h_{V_2}\circ h_{V_1}^{-1}
=\bigl(\pi|_{V_2}\bigr)^{-1}\circ \pi|_{V_1}
= \operatorname{id}_{V_1\cap V_2},
\]
also holomorph. Liegen $V_1\subset X_{i_1}$ und $V_2\subset X_{i_2}$ mit $i_1\neq i_2$, so kann der Überlapp nur aus Identifikationen
über die vorgegebenen Verklebungen stammen; auf einem geeigneten offenen Teil $W\subset V_1$ ist
\[
h_{V_2}\circ h_{V_1}^{-1}
=\bigl(\pi|_{V_2}\bigr)^{-1}\circ \pi|_{V_1}
= \varphi^{(j)}_2\circ\bigl(\varphi^{(j)}_1\bigr)^{-1}\big|_{W},
\]
also biholomorph nach Voraussetzung. Damit ist die Gesamtheit der Karten ein komplexer Atlas auf $X$.

Da $X$ per Annahme hausdorffsch ist, erhält $X$ damit die Struktur einer Riemannschen Fläche (im hier verwendeten Sinn:
hausdorffscher topologischer Raum mit komplexer Struktur).
\end{proof}
\end{EXA}



Beispiel 2.6  Riemannsche Zahlenkugel / Projektive Grade

Ausgangsflächen seien $X_{1}=\mathbb{C}$ und $X_{2}=\mathbb{C}$. Die Verklebefläche ist $U_{1}=\mathbb{C}^{*}$ mit
$$
\varphi_{1}:\left\{\begin{array}{rlc}
\mathbb{C}^{*} & \longrightarrow X_{1} \\
z & \longmapsto z
\end{array} \quad \text { und } \quad \varphi_{2}:\left\{\begin{array}{rll}
\mathbb{C}^{*} & \longmapsto X_{2} \\
z & \longmapsto \frac{1}{z}
\end{array}\right.\right.
$$
Die Voraussetzungen der Proposition 2.5 sind in diesem Beispiel erfüllt. Die Fläche $X$ heißt Riemannsche Zahlenkugel oder Projektive Gerade $\mathbb{P}_{\mathbb{C}}^{1}$. Die Menge $\mathbb{P}_{\mathbb{C}}^{1} \backslash X_{1}$ ist einpunktig. Diesen Punkt werden wir mit $\infty \in \mathbb{P}_{\mathbb{C}}^{1}$ bezeichnen.



\begin{EXA}{RS-M25-01-14}{Riemannsche Zahlenkugel / Projektive Gerade}
Wir nehmen zwei Kopien der Ebene $X_1=\C$ und $X_2=\C$ und kleben sie über
$U:=\C^\times$ mittels
\[
\varphi_1:U\to X_1,\ z\mapsto z,
\qquad
\varphi_2:U\to X_2,\ z\mapsto \frac{1}{z}.
\]
Die Übergangsfunktion auf dem Überlapp ist damit
\[
\varphi_2\circ\varphi_1^{-1}: \C^\times\longrightarrow \C^\times,\quad
z\longmapsto \frac{1}{z},
\]
also biholomorph. Nach Prop.~2.5 entsteht so ein hausdorffscher
topologischer Raum $X$ mit komplexer Struktur; wir schreiben
\[
\mathbb{P}^1_\C \;=\; X \;=\; (\C \sqcup \C)/\!\sim
\quad\text{(Riemannsche Zahlenkugel, komplexe projektive Gerade).}
\]

\textbf{Atlas und der Punkt $\infty$.}
Wir erhalten einen komplexen Atlas aus zwei Karten:
\[
(U_0,\; z_0:=\mathrm{id}_\C),\qquad (U_\infty,\; z_\infty:=\mathrm{id}_\C),
\]
wobei $U_0$ die erste Kopie $X_1\simeq\C$ und $U_\infty$ die zweite Kopie $X_2\simeq\C$
bezeichnet. Die Kartenwechselabbildung ist $z_\infty=\frac{1}{z_0}$ auf
$U_0\cap U_\infty\simeq\C^\times$.
Die Menge $\mathbb{P}^1_\C\setminus U_0$ besteht aus genau einem Punkt; wir nennen ihn
$\infty\in\mathbb{P}^1_\C$. Analog ist $\mathbb{P}^1_\C\setminus U_\infty=\{0\}$ in der $z_0$-Karte.

\textbf{Nachbarschaften von $\infty$.}
Eine Basis von Umgebungen von $\infty$ sind die Mengen
\[
\{\,p\in\mathbb{P}^1_\C \mid |z_\infty(p)|<R\,\}\quad (R>0),
\]
d.\,h. Urbilder beschränkter Mengen in der $z_\infty$-Koordinate; äquivalent:
$\{\,p\in\mathbb{P}^1_\C \mid |z_0(p)|>1/R\,\}\cup\{\infty\}$.

\textbf{Identifikation mit der Sphäre.}
Über die stereographische Projektion erhält man einen Homöomorphismus
\[
S^2 \;\xrightarrow{\ \cong\ }\; \mathbb{P}^1_\C,
\]
der $S^2$ mit der einen Zusatzstelle $\infty$ zur \emph{Einpunktkompaktifizierung}
der Ebene identifiziert. Damit ist $\mathbb{P}^1_\C$ kompakt, zusammenhängend und orientierbar.

\textbf{Holomorphe/meromorphe Funktionen.}
Jede holomorphe Funktion $\mathbb{P}^1_\C\to\C$ ist konstant (Kompaktheit $\Rightarrow$
großer Satz von Liouville). Meromorphe Funktionen auf $\mathbb{P}^1_\C$ sind genau die
rationalen Funktionen $R(z)=\frac{P(z)}{Q(z)}$ auf der $z_0$-Karte (inkl.~einer
Polstelle/Nullstelle an $\infty$ entsprechend dem Grad).

\textbf{Zusammenfassung.}
Die projektive Gerade $\mathbb{P}^1_\C$ ist die aus zwei Kopien von $\C$ über
$z\mapsto 1/z$ verklebte Riemannsche Fläche; sie besteht aus $\C$ plus einem
Punkt $\infty$ und ist (als topologischer Raum) mit $S^2$ homöomorph.
\end{EXA}




Beispiel 2.7 


Wir wählen $X_{1}, X_{2}$ und $U_{1,2}$ wie in Beispiel 2.4, aber $\varphi_{2}(z)=z$. Der resultierende topologische Raum ist nicht hausdorffsch, also keine Riemannsche Fläche.



\begin{EXA}{RS-M25-01-15}{Nicht‐Hausdorff durch triviales Verkleben}
Wir nehmen zwei Kopien der Ebene $X_1=\C$ und $X_2=\C$ und kleben sie wie in Bsp.~2.6
über $U=\C^\times$, aber diesmal mit \emph{identischer} Übergangsfunktion
\[
\varphi_1:U\to X_1,\ z\mapsto z,\qquad
\varphi_2:U\to X_2,\ z\mapsto z.
\]
Bezeichne den verklebten Raum mit $X:=(\C\sqcup\C)/\!\sim$. Es gibt zwei verschiedene Punkte
\[
p:=[0\in X_1],\qquad q:=[0\in X_2]\in X,
\]
die nicht identifiziert werden.

\textbf{Behauptung.} $X$ ist nicht hausdorffsch (also keine Riemannsche Fläche).

\textit{Beweis.}
Seien $U_p,U_q\subset X$ offene Umgebungen von $p$ bzw.\ $q$.
Dann existieren offene Kreisscheiben $D_r\subset X_1$ und $D_s'\subset X_2$ um $0$
mit $D_r\subset \pi^{-1}(U_p)$ und $D_s'\subset \pi^{-1}(U_q)$, wobei
$\pi:\C\sqcup\C\to X$ die Quotientenabbildung ist.
Für $0<\rho<\min\{r,s\}$ enthalten $D_r\setminus\{0\}$ und $D_s'\setminus\{0\}$ je
eine punktierte Kreisscheibe um $0$, und auf $U=\C^\times$ werden die Punkte via
$\varphi_1=\varphi_2=\mathrm{id}$ identifiziert. Also gilt
\[
\pi\big(D_r\setminus\{0\}\big)=\pi\big(D_s'\setminus\{0\}\big)\subset U_p\cap U_q.
\]
Insbesondere $U_p\cap U_q\neq\emptyset$. Damit lassen sich $p$ und $q$ nicht durch disjunkte
offene Umgebungen trennen; $X$ ist nicht hausdorffsch.
\qed
\end{EXA}




Beispiel 2.8  (Punktierter) Torus durch Verklebung

Wir betrachten nochmal den Torus mit Perioden 1 und $i$ aus Beispiel 2.4. Wir wollen ihn nun durch Verkleben konstruieren. Sei
$$
\begin{aligned}
& X_{1}=\{z=x+i y \in \mathbb{C}: x \in(0,1), y \in(0,1)\} \\
& X_{2}=\{z \in \mathbb{C}: x \in(0,1), y \in(-\varepsilon,+\varepsilon)\} \\
& X_{3}=\{z \in \mathbb{C}: x \in(-\varepsilon,+\varepsilon), y \in(0,1)\}
\end{aligned}
$$
wie in Abbildung 2.3 und $T^{*}$ die Verklebung von $\coprod_{i=1}^{3} X_{i}$ entlang der offenen Mengen $U_{2}^{+}, U_{2}^{-}, U_{3}^{+}, U_{3}^{-}$wie folgt.
$$
\includegraphics[scale=0.2]{2025_10_06_374d4cfa8b1948eb4c41g-011}
$$
Seien
$$
\begin{array}{lll}
U_{2}^{+}=X_{2} \cap \mathbb{H} & \text { mit } & \varphi_{1}: U_{2}^{+} \longrightarrow X_{1}, z \longmapsto z \\
U_{2}^{-}=X_{2} \cap\{\operatorname{Im} z<0\} & \text { mit } & \varphi_{1}: U_{2}^{-} \longrightarrow X_{1}, z \longmapsto i+z \\
U_{3}^{+}=X_{3} \cap\{\operatorname{Re} z>0\} & & \varphi_{1}: U_{3}^{+} \longrightarrow X_{1}, z \longmapsto z \\
U_{3}^{-}=X_{3} \cap\{\operatorname{Re} z<0\} & & \varphi_{1}: U_{3}^{-} \longrightarrow X_{1}, z \longmapsto 1+z
\end{array}
$$
und sei $\varphi_{2}$ jeweils die Inklusion $U_{i}^{+} \longrightarrow X_{i}$ bzw. $U_{i}^{-} \longrightarrow X_{i}$. Die Riemannsche Fläche $T^{*}$ wird punktierter Torus genannt.

Im Unterschied zu $S^{1} \times S^{1}$ ist $T^{*}$ nicht kompakt. Wir haben in der Ecke noch ein Loch gelassen, das wir nun stopfen. Sei schließlich $T$ die Verklebung von $T^{*}$ mit der Kreisscheibe $\Delta$ entlang $U_{1}=\Delta^{*}$, wobei wir $\varphi_{1}$ auf den offenen Vierteln
$$
U_{1, k}=\left\{z \in \Delta: \arg (z) \in\left((k-1) \frac{\pi}{2}, k \frac{\pi}{2}\right)\right\}, \quad k=1, \ldots, 4
$$
definieren durch
$$
\begin{array}{ll}
\varphi_{1}: U_{1,1} \longrightarrow X_{1}, z \longmapsto z & \varphi_{1}: U_{1,3} \longrightarrow X_{1}, z \longmapsto z+1+i \\
\varphi_{1}: U_{1,2} \longrightarrow X_{1}, z \longmapsto z+1 & \varphi_{1}: U_{1,4} \longrightarrow X_{1}, z \longmapsto z+i
\end{array}
$$
und auf den Achsen durch
$$
\begin{array}{ll}
\varphi_{1}:(0,1) \longrightarrow X_{2}, z \longmapsto z & \varphi_{1}:(0,1) i \longrightarrow X_{3}, z \longmapsto z \\
\varphi_{1}:(-1,0) \longrightarrow X_{2}, z \longmapsto z+1 & \varphi_{1}:(-1,0) i \longrightarrow X_{3}, z \longmapsto z+i
\end{array}
$$
und $\varphi_{2}: U_{1} \longrightarrow \Delta$ die Inklusionsabbildung ist.




\begin{EXA}{RS-M25-01-16}{(Punktierter) Torus via Verkleben}
Wir konstruieren den Torus mit Perioden $1$ und $i$ durch Verkleben von planaren
Gebieten mittels \emph{Translationskarten}.

\textbf{Schritt 1: Der punktierte Torus $T^*$.}
Sei
\[
\begin{aligned}
X_1&=\{x+iy\in\C: x\in(0,1),\,y\in(0,1)\},\\
X_2&=\{x+iy\in\C: x\in(0,1),\,y\in(-\varepsilon,\varepsilon)\},\\
X_3&=\{x+iy\in\C: x\in(-\varepsilon,\varepsilon),\,y\in(0,1)\},
\end{aligned}
\]
und definiere die offenen Teilmengen
\[
\begin{array}{lll}
U_2^+=X_2\cap\mathbb{H}, & \varphi_1:U_2^+\to X_1,\ z\mapsto z, & \varphi_2:U_2^+\hookrightarrow X_2,\\
U_2^-=X_2\cap\{\Im z<0\}, & \varphi_1:U_2^-\to X_1,\ z\mapsto i+z, & \varphi_2:U_2^-\hookrightarrow X_2,\\
U_3^+=X_3\cap\{\Re z>0\}, & \varphi_1:U_3^+\to X_1,\ z\mapsto z, & \varphi_2:U_3^+\hookrightarrow X_3,\\
U_3^-=X_3\cap\{\Re z<0\}, & \varphi_1:U_3^-\to X_1,\ z\mapsto 1+z, & \varphi_2:U_3^-\hookrightarrow X_3.
\end{array}
\]
Die jeweils angegebenen $\varphi_1$ sind Translationen (also biholomorph), die $\varphi_2$
sind Inklusionen. Wir verkleben die disjunkte Vereinigung $X_1\sqcup X_2\sqcup X_3$
entlang der Paare $(\varphi_1,\varphi_2)$ und versehen das Quotient
$T^*:=\bigl(\coprod_{i=1}^3 X_i\bigr)/\!\sim$ mit der Quotiententopologie.

\textit{Behauptung.} $T^*$ ist eine Riemannsche Fläche (hausdorffsch, mit
Translationsatlas) und nicht kompakt.

\textit{Begründung.} (i) \emph{Kompatibilität:} Alle Übergänge sind durch
Translationen $z\mapsto z+\omega$ mit $\omega\in\{0,1,i,1+i\}$ gegeben, also
biholomorph. Damit liefern die inversen Karten lokale Koordinaten und die
Kartenwechsel sind holomorph. (ii) \emph{Hausdorff:} Die Verklebungen erfolgen
entlang offener Teilmengen, und es werden nur Randpunkte mittels Translationen
identifiziert; die entstehende Äquivalenzrelation ist abgeschlossen in der disjunkten
Vereinigung, also ist der Quotient hausdorffsch (Standardargument aus der
Gluing‐Proposition). (iii) \emph{Nichtkompaktheit:} Es bleibt in der „Ecke”
(entspricht dem Gitterpunkt $0$ im Fundamentalparallelogramm) ein Loch; daher ist
$T^*$ topologisch ein Torus ohne einen Punkt und somit nicht kompakt.

\medskip
\textbf{Schritt 2: Auffüllen der Spitze — der (kompakte) Torus $T$.}
Wir füllen das Loch, indem wir die Einheitskreisscheibe $\Delta$ entlang des
Randes $\Delta^*=\Delta\setminus\{0\}$ ankleben.
Zerlege $\Delta^*$ in die vier offenen Quadranten
\[
U_{1,k}=\{z\in\Delta:\ \arg z\in((k-1)\tfrac{\pi}{2},k\tfrac{\pi}{2})\},\quad k=1,\dots,4,
\]
und definiere die Verklebung auf den Quadranten durch Translationen
\[
\begin{array}{ll}
\varphi_1:U_{1,1}\to X_1,\ z\mapsto z,&
\varphi_1:U_{1,2}\to X_1,\ z\mapsto z+1,\\[2pt]
\varphi_1:U_{1,3}\to X_1,\ z\mapsto z+1+i,&
\varphi_1:U_{1,4}\to X_1,\ z\mapsto z+i,
\end{array}
\]
sowie auf den Achsen (zur Kompatibilität mit $X_2,X_3$)
\[
\begin{array}{ll}
\varphi_1:(0,1)\to X_2,\ z\mapsto z, & \varphi_1:(0,1)i\to X_3,\ z\mapsto z,\\
\varphi_1:(-1,0)\to X_2,\ z\mapsto z+1, & \varphi_1:(-1,0)i\to X_3,\ z\mapsto z+i,
\end{array}
\]
und setze $\varphi_2:U_1=\Delta^*\hookrightarrow\Delta$ als Inklusion.
Die so entstandene Verklebung
\[
T:=\bigl(T^*\sqcup\Delta\bigr)\big/\!\sim
\]
ist wieder durch Translationen definiert, also liefert ein komplexes Atlas.

\textit{Behauptung.} $T$ ist eine kompakte Riemannsche Fläche und (biholomorph)
gleich dem Quotienten $\C/(\Z\oplus i\Z)$.

\textit{Begründung.} (i) \emph{Atlas:} Auch hier sind alle Übergangsfunktionen
Translationen und damit biholomorph, also ist $T$ eine Riemannsche Fläche.
(ii) \emph{Kompakt:} Topologisch erhält man ein Fundamentalviereck mit gegenüberliegenden
Rändern identifiziert—also den Torus. Das Auffüllen des fehlenden Punktes macht
den Raum kompakt. (iii) \emph{Identifikation mit einem Gitterquotienten:}
Die Karten stimmen mit dem Translationsatlas für $\C/(\Z+i\Z)$ überein; die natürliche
Quotientenabbildung $\C\to \C/(\Z+i\Z)$ ist lokal ein Holomorphismus und induziert
einen Bihomöomorphismus auf $T$.

\medskip
\textbf{Fazit.}
\begin{itemize}
\item $T^*$ ist der \emph{punktierte Torus}, eine nicht‐kompakte Riemannsche Fläche,
erhalten durch Verkleben von $X_1,X_2,X_3$ mittels Translationen.
\item $T$ entsteht aus $T^*$ durch Ankleben der Scheibe $\Delta$ am Rand $\Delta^*$;
$T$ ist kompakt und biholomorph zu $\C/(\Z\oplus i\Z)$.
\end{itemize}
\end{EXA}


Das ist auf den ersten Blick ziemlich umständlich. Wir werden diese Art von Verkleben im Abschnitt 8 über flache Flächen nochmal formalisieren und das Einkleben der 'Eckpunkte' automatisieren.

Wir haben nun zwei Arten gesehen Riemannsche Flächen zu konstruieren, durch eine Quotientabbildung, bei der das Objekt 'oben' bereits eine Riemannsche Fläche ist, und synthetische durch Verkleben. Die dritte wichtige Art beschreibt Riemannsche Flächen als Unterobjekte.


Definition (Punktierte) elliptische Kurve 

Wir betrachten dazu $Z=\mathbb{C}^{2} \cong \mathbb{R}^{4}$ mit der üblichen reellen Topologie. Dann sei zu vorgegebenen, paarweise verschiedenen $x_{1}, x_{2}, x_{3} \in \mathbb{C}$ der topologische Raum $X$ definiert durch
$$
X=\left\{(x, y) \in Z: y^{2}=f(x)\right\} \quad \text { mit } \quad f(x)=\left(x-x_{1}\right)\left(x-x_{2}\right)\left(x-x_{3}\right) .
$$
Für $x \notin B=\left\{x_{1}, x_{2}, x_{3}\right\}$ sei wählen wir eine Umgebung $V(x) \subset \mathbb{C}$, welche $B$ nicht trifft. Für $(x, y) \in X$ ist dann die Zusammenhangskomponente von $U(x)=\left\{\left(x^{\prime}, y^{\prime}\right) \in X: x^{\prime} \in V(x)\right\}$, welche ( $x, y$ ) enthält, eine Umgebung von ( $x, y$ ) in $X$ und die Projektion auf die erste Koordinate eine Karte. Für $x=x_{i}$ wollen wir die Projektion auf die zweite Koordinate zu einer Karte in einer Umgebung von $\left(x_{i}, 0\right) \in X$ machen. Dazu müssen wir eine Umgebung von ( $x_{i}, 0$ ) angeben, auf der die Abbildung ( $x, y$ ) $\mapsto y$ in Homöomorphismus ist. Da die $x_{i}$ paarweise verschieden sind, ist die Ableitung von $f(x)$ im Punkt $x=x_{1}$ von Null verschieden und die gewünschte Umgebung existiert nach dem Satz über implizite Funktionen. Man prüft nun noch nach, dass die Kartenwechselabbildungen in der Tat holomorph sind.

Eine so definierte Riemannsche Fläche wird (punktierte) elliptische Kurve genannt.



\begin{DEF}{RS-M25-01-17}{(Punktierte) elliptische Kurve}
Wir betrachten den Raum $Z=\C^2\cong\R^4$ mit der üblichen reellen Topologie.
Seien drei paarweise verschiedene Punkte $x_1,x_2,x_3\in\C$ gegeben und definiere
\[
X=\{(x,y)\in Z:\ y^2=f(x)\},\qquad f(x)=(x-x_1)(x-x_2)(x-x_3).
\]
Dann ist $X$ eine reelle zweidimensionale (also komplex eindimensionale)
Untermannigfaltigkeit von $\C^2$.

\textbf{Lokale Karten.}
\begin{enumerate}
\item Für $x\notin B:=\{x_1,x_2,x_3\}$ wählen wir eine Umgebung $V(x)\subset\C$, die $B$ nicht schneidet.
Definiere
\[
U(x)=\{(x',y')\in X:\ x'\in V(x)\}.
\]
Die Zusammenhangskomponente von $U(x)$, die $(x,y)$ enthält, sei $U_{(x,y)}$.
Dann ist die Projektion
\[
\pi_1:U_{(x,y)}\to V(x),\quad (x',y')\mapsto x'
\]
ein Homöomorphismus auf ihr Bild. Sie definiert eine Karte
\[
h_1:(U_{(x,y)},\C),\quad h_1(x',y')=x'.
\]
\item Für $x=x_i$ wollen wir die Projektion auf die zweite Koordinate verwenden.
Da $f'(x_i)\neq 0$, folgt aus dem Satz über implizite Funktionen, dass es eine
Umgebung $U_i\subset X$ von $(x_i,0)$ gibt, auf der
\[
\pi_2:U_i\to\C,\quad (x,y)\mapsto y
\]
ein Homöomorphismus auf ihr Bild ist.
Damit erhalten wir Karten
\[
h_2:(U_i,\C),\quad h_2(x,y)=y,\qquad i=1,2,3.
\]
\end{enumerate}

\textbf{Kartenwechsel.}
In den Überlappungsbereichen der obigen Karten gelten Übergangsfunktionen der Form
\[
h_1\circ h_2^{-1}: y\mapsto x_i+\frac{y^2}{f'(x_i)}+\dots,
\]
die holomorph sind. Somit ist der Atlas aus den obigen Karten ein komplexer Atlas.

\textbf{Fazit.}
Der Raum $X$ mit dieser komplexen Struktur ist eine Riemannsche Fläche. Entfernt man den Punkt
„im Unendlichen“, erhält man die \emph{punktierte elliptische Kurve}, während die kompakte
Ergänzung durch den Punkt $\infty$ die (vollständige) \emph{elliptische Kurve}
\[
E=\{(x,y)\in\mathbb{P}^1\times\mathbb{P}^1: y^2z=f(x,z)\}\subset\mathbb{P}^2
\]
liefert.
\end{DEF}




\pagebreak




\subsection {Holomorphe Abbildungen}


Nachdem wir nun die Objekte definiert haben, mit denen wir umgehen, folgt nun die Definition der Morphismen zwischen Riemannschen Flächen und damit auch, wann wir Riemannsche Flächen als gleich ansehen.

Definition 2.9 Holomorphe Abbildung

Eine Abbildung $f: X \rightarrow Y$ zwischen zwei Riemannschen Flächen heißt holomorph, falls $f: X \rightarrow Y$ stetig ist und für jede Karte $(U, g)$ von $X$ und jede Karte $(V, h)$ von $Y$ die Verkettung $h \circ f \circ g^{-1}$ : $g\left(U \cap f^{-1}(V)\right) \rightarrow \mathbb{C}$ holomorph ist. $f$ heißt biholomorph oder Isomorphismus; falls $f$ zudem bijektiv ist.


\begin{DEF}{RS-M25-02-01}{Holomorphe Abbildung zwischen Riemannschen Flächen}
Seien $X$ und $Y$ Riemannsche Flächen.  
Eine Abbildung
\[
f:X\longrightarrow Y
\]
heißt \emph{holomorph}, falls
\begin{enumerate}
\item $f$ stetig ist, und
\item für jede Karte $(U,g)$ von $X$ und jede Karte $(V,h)$ von $Y$ die verkettete Abbildung
\[
h\circ f\circ g^{-1}:\ g\bigl(U\cap f^{-1}(V)\bigr)\longrightarrow \C
\]
eine holomorphe Abbildung zwischen offenen Teilmengen der komplexen Ebene ist.
\end{enumerate}

Ist $f$ zusätzlich bijektiv und $f^{-1}$ ebenfalls holomorph, so heißt $f$ \emph{biholomorph}  
(oder \emph{Isomorphismus} der Riemannschen Flächen). In diesem Fall sind $X$ und $Y$ \emph{biholomorph äquivalent}.
\end{DEF}


Beispiel: Isomorphie zwischen Torus ($1,i$) und Torus ($1,i'$)

Es ist eine einfache Übung zu prüfen, dass der Torus mit Perioden $1, i$ aus Beispiel 2.4 zum Torus mit Perioden 1, $i$ aus Beispiel 2.8 isomorph ist. Dieser ist wiederum zum Torus mit Perioden $1, i+1$ isomorph.


\begin{EXA}{RS-M25-02-02}{Isomorphie zwischen Tori mit äquivalenten Perioden}
Wir betrachten die Riemannschen Flächen, die durch Gitterquotienten der komplexen Ebene
\[
T_{\tau}=\C/(\Z+\tau\Z)
\]
gegeben sind, wobei $\tau\in\C$ mit $\Im(\tau)>0$ eine Periode bezeichnet.

\textbf{Behauptung.}
Die Tori $T_i=\C/(\Z+i\Z)$ und $T_{i+1}=\C/(\Z+(i+1)\Z)$ sind biholomorph.

\textbf{Beweisidee.}
Die Abbildung
\[
f:\C\longrightarrow\C,\qquad f(z)=e^{\pi i/4}\,z
\]
ist linear und holomorph. Sie erfüllt
\[
f(\Z+i\Z)=e^{\pi i/4}(\Z+i\Z)=\Z+(i+1)\Z,
\]
da Multiplikation mit $e^{\pi i/4}$ das Gitter $\Z+i\Z$ in das Gitter $\Z+(i+1)\Z$ überführt. Somit induziert $f$ eine wohldefinierte Abbildung
\[
\tilde{f}:T_i\longrightarrow T_{i+1},\qquad \tilde{f}([z])=[f(z)].
\]
Da $f$ biholomorph ist und die Gitterabbildung respektiert, ist auch $\tilde{f}$ biholomorph.

\textbf{Folgerung.}
Insbesondere sind die in Beispiel~2.4 und Beispiel~2.8 konstruierten Tori mit Perioden $(1,i)$ und $(1,i+1)$ biholomorph isomorph.  
Beide beschreiben somit dieselbe komplexe Struktur auf dem Torus.
\end{EXA}



Definition: Automorphismen-Gruppe von Riemannischen Flächen

Biholomorphe Selbstabbildungen einer Riemannschen Fläche $X$ (auch Biholomorphismen oder Automorphismen von $X$ genannt) bilden eine Gruppe, die wir mit $\operatorname{Aut}(X)$ bezeichnen. 




\begin{DEF}{RS-M25-02-03}{Automorphismengruppe einer Riemannschen Fläche}
Sei $X$ eine Riemannsche Fläche.  
Eine \emph{biholomorphe Selbstabbildung} von $X$, d.\,h. eine bijektive Abbildung
\[
f:X\longrightarrow X
\]
mit $f$ und $f^{-1}$ holomorph, heißt \emph{Automorphismus} (oder \emph{Biholomorphismus}) von $X$.

Die Menge aller solchen Abbildungen
\[
\operatorname{Aut}(X):=\{\,f:X\to X\mid f\text{ biholomorph}\,\}
\]
bildet bezüglich der Komposition eine Gruppe, die \emph{Automorphismengruppe} der Riemannschen Fläche $X$ genannt wird.
\end{DEF}



Beispiel: Elliptische Involution

Im Beispiel der elliptischen Kurve am Ende von Abschnitt 2.1 ist $(x, y) \mapsto(x,-y)$ ein Automorphismus von $X$, der elliptische Involution genannt wird.



\begin{EXA}{RS-M25-02-04}{Elliptische Involution}
Betrachte die in Abschnitt~2.1 eingeführte elliptische Kurve
\[
X=\{(x,y)\in\C^2 \mid y^2=(x-x_1)(x-x_2)(x-x_3)\},
\]
wobei $x_1,x_2,x_3\in\C$ paarweise verschieden sind.  

Die Abbildung
\[
\iota:X\longrightarrow X,\qquad (x,y)\longmapsto(x,-y),
\]
ist offensichtlich bijektiv, stetig und erfüllt
\[
(-y)^2 = y^2 = (x-x_1)(x-x_2)(x-x_3),
\]
also $\iota(X)=X$.  

Da $\iota$ in lokalen Koordinaten durch $y\mapsto -y$ gegeben ist, ist sie holomorph und ihre Inverse ist $\iota$ selbst. Somit gilt $\iota\in\operatorname{Aut}(X)$.

\textbf{Definition.}
Diese Abbildung $\iota$ heißt \emph{elliptische Involution} der Kurve $X$.  
Sie ist ein nichttrivialer Automorphismus der elliptischen Kurve und erfüllt $\iota^2=\mathrm{id}_X$.
\end{EXA}


Es ist eine Konsequenz aus der Definition oben, dass die Umkehrabbildung einer biholomorphen Abbildung ebenfalls wieder holomorph ist. Das wird am Ende des Abschnitts klar sein, wenn wir die lokale Struktur von holomorphen Abbildungen beschrieben haben. Wesentliches Hilfsmittel dazu ist, dass holomorphe Funktionen stets in Potenzreihen entwickelbar sind (siehe FTDGL, Satz 7.7]).


Definition 2.10 Verzweigungsindex 

Sei $f: X \rightarrow Y$ eine holomorphe Abbildung zwischen Riemannschen Flächen, $a \in X$ und ( $U, g$ ) und ( $V, h$ ) Karten von Umgebungen von a bzw. von $f(a)$ mit $g(a)=0$ und $h(f(a))=0$. Der Verzweigungsindex von $f$ bei a ist definiert als
$$
e_{a}(f)=\min \left\{n: b_{n} \neq 0\right\}, \text { wobei } h \circ f \circ g^{-1}(z)=\sum_{n \geq 1} b_{n} z^{n}
$$
Wir zeigen, dass dies wohldefiniert ist. Sei $(\widetilde{U}, \widetilde{h})$ eine andere Karte im Bildbereich mit $\widetilde{h}(f(a))=0$, so ist $\widetilde{h} \circ h^{-1}$ eine biholomorphe Abbildung von ihrem Definitionsbereich auf ihren Bildbereich. Also ist $\widetilde{h} \circ h^{-1}=\sum_{n \geq 1} c_{n} z^{n}$ eine Potenzreihemit $c_{1} \neq 0$. Schreiben wir
$$
\left(\widetilde{h} \circ f \circ g^{-1}\right)(z)=\left(\widetilde{h} \circ h^{-1}\right)\left(h \circ f \circ g^{-1}\right)(z)=\sum_{n \geq 1} \widetilde{b}_{n} z^{n}
$$
und ist $b_{n}=0$ für $n=1, \ldots, k$, so ist auch $\widetilde{b}_{n}=0$ für $n=1 \ldots k$, und umgekehrt. Gleiches gilt für einen Kartenwechsel auf $X$.

Nach dem Identitätssatz ist $e_{a}(f)<\infty$, falls $f$ nicht in einer Umgebung von $a$ konstant ist. Aus diesem Satz folgt auch, dass falls $f$ in einer Umgebung von $a$ konstant ist, so ist $f$ auf der ganzen Zusammenhangskomponente von $X$, welche $a$ enthält konstant. 



\begin{DEF}{RS-M25-02-05}{Verzweigungsindex einer holomorphen Abbildung in Punkt}
Seien $f:X\to Y$ eine holomorphe Abbildung zwischen Riemannschen Flächen und $a\in X$.  
Wähle Karten $(U,g)$ um $a$ und $(V,h)$ um $f(a)$ mit $g(a)=0$ und $h(f(a))=0$.  
In diesen lokalen Koordinaten hat $f$ die Darstellung
\[
(h\circ f\circ g^{-1})(z)=\sum_{n\ge1} b_n z^n.
\]
Der \emph{Verzweigungsindex} von $f$ in $a$ ist definiert durch
\[
e_a(f):=\min\{\,n\in\mathbb{N}\mid b_n\neq0\,\}.
\]

\textbf{Bemerkungen.}
\begin{enumerate}
\item Diese Definition ist wohldefiniert, d.\,h. unabhängig von der Wahl der Karten:
Ist $(\widetilde{U},\widetilde{h})$ eine andere Karte um $f(a)$ mit $\widetilde{h}(f(a))=0$, so gilt
\[
\widetilde{h}\circ h^{-1}(z)=\sum_{n\ge1} c_n z^n,\quad c_1\neq0,
\]
und damit
\[
(\widetilde{h}\circ f\circ g^{-1})(z)
=(\widetilde{h}\circ h^{-1})\bigl(h\circ f\circ g^{-1}(z)\bigr)
=\sum_{n\ge1}\widetilde{b}_n z^n.
\]
Sind $b_1=\dots=b_k=0$, so gilt auch $\widetilde{b}_1=\dots=\widetilde{b}_k=0$, und umgekehrt.
Dasselbe gilt bei einem Kartenwechsel auf $X$.

\item Nach dem Identitätssatz der Funktionentheorie ist $e_a(f)<\infty$, falls $f$ nicht in einer Umgebung von $a$ konstant ist.  

\item Ist $f$ in einer Umgebung von $a$ konstant, so ist $f$ auf der gesamten Zusammenhangskomponente von $X$, die $a$ enthält, konstant.
\end{enumerate}
\end{DEF}



\begin{THEO}{RS-M25-02-06}{Identitätssatz der Funktionentheorie}
Sei $D\subseteq\C$ ein Gebiet (d.\,h.\ eine offene und zusammenhängende Menge) und  
$f,g:D\to\C$ zwei holomorphe Funktionen.  

Falls die Menge
\[
M:=\{\,z\in D\mid f(z)=g(z)\,\}
\]
einen Häufungspunkt in $D$ besitzt, so gilt
\[
f\equiv g \quad\text{auf } D.
\]

\textbf{Spezialfall:}
Ist $f:D\to\C$ holomorph und verschwindet auf einer Menge mit Häufungspunkt in $D$, so folgt $f\equiv0$ auf $D$.
\end{THEO}






Definition: Nicht verschwindende holomorphe Abbildungen

\begin{DEF}{RS-M25-02-07}{Nicht verschwindende holomorphe Abbildungen}
Eine holomorphe Abbildung $f: X \longrightarrow Y$ heißt nirgends konstant, falls $e_{a}(f)<\infty$ für alle $a \in X$ ist.
\end{DEF}



Wir benutzen Begriff der Ordnung, um eine Normalform für holomorphe Abbildungen herzuleiten.



Proposition 2.11 

\begin{PROP}{RS-M25-02-08}{Normalform von holomorphen Abbildungen}
Sei $f: X \rightarrow Y$ holomorph und $a \in X$ mit $n=e_{a}(f)<\infty$. Dann gibt es Karten $(U, g)$ von $X$ und $(V, h)$ von $Y$, so dass $h \circ f \circ g^{-1}=z^{n}$ ist.
\end{PROP}

>> $f$ nicht konstant


Beweis: 


\begin{PROOF}{RS-M25-02-09}{P: Normalform von holomorphen Abbildungen}
Wähle zunächst $\operatorname{Karten}(\widetilde{U}, \widetilde{g})$ und $(V, h)$ mit $\widetilde{g}(a)=0, h(f(a))=0$ und $f(\widetilde{U}) \subset V$ ). Wir wollen aus der Potenzreihe
$$
\left(h \circ f \circ \widetilde{g}^{-1}\right)(z)=\sum_{j \geq n} b_{j} z^{j}
$$
eine $n$-te Wurzel $\varphi(z)=z \sum_{i \geq 0} c_{i} z^{i}$ ziehen. Dann folgt aus dem Ansatz
$$
\sum_{j \geq n} b_{j} z^{j}=\left(z \sum_{i \geq 0} c_{i} z^{i}\right)^{n}
$$
dass $c_{0}$ eine $n$-te Wurzel von $b_{n}$ sein muss und dass sich die folgenden Terme rekursiv eindeutig auflösen lassen, beginnend mit $c_{1}=\frac{b_{n+1}}{n c_{0}^{n-1}}$. Die Funktion $\varphi$ ist also in allen Fällen biholomorph auf einer Umgebung von Null. Wir setzen $g=\varphi \circ \widetilde{g}$ und schränken $\widetilde{U}$ zu einer Umgebung $U$ von $a$ ein, sodass $g^{-1}$ auf $g(U)$ definiert ist. Aus der Definition der Umkehrabbildung folgt $\varphi^{-1}(z) \sum_{j \geq 0} c_{j}\left(\varphi^{-1}\right)^{j}(z)=z$ und aus
$$
h \circ f \circ g^{-1}(z)=\left(h \circ f \circ \widetilde{g}^{-1} \circ \varphi^{-1}\right)(z)=\left(\varphi^{-1}(z) \sum_{j \geq 0} c_{j}\left(\varphi^{-1}(z)\right)^{j}\right)^{n}=z^{n}
$$
folgt die Behauptung.

Es verbleibt noch zu zeigen, dass diese $n$-te Wurzel nicht nur eine formale Potenzreihe sondern in der Tat eine holomorphe Funktion auf einer Umgebung von Null ist (Übung).
\end{PROOF}

>> Koeffizienten Vergleich für $c_j$ für $u(z)=1+\sum_{j\geq 1}b_jz^j=(\sum_{j\geq 0}c_jz^j)^n$

>> $v(z)=c_j$-Potenzreihe hat positiven Konvergenzradius um $0$: $v(z)=\exp(1/n\log(u(z)))$




Korollar 2.12 (Offenheitssatz) 

\begin{KORO}{RS-M25-02-10}{Offenheitssatz}
Eine nirgends konstante, holomorphe Abbildung ist offen.
\end{KORO}

Beweis: 

\begin{PROOF}{RS-M25-02-11}{P: Offenheitssatz}
Für alle $n \in \mathbb{N}$ ist das Bild der offenen Kreisscheibe $\triangle_{r}$ unter $z \longmapsto z^{n}$ eine offene Kreisscheibe mit Radius $r^{n}$.
\end{PROOF}

Eine abgeschlossene Menge $A \subseteq X$ heißt analytisch, wenn jeder Punkt in $A$ eine Umgebung $U$ besitzt, sodass $U \cap A$ Nullstellenmenge einer holomorphen Abbildung $U \longrightarrow \mathbb{C}$ ist.\\
Eine Menge $M \subset X$ heißt lokal endlich, wenn für alle kompakten Mengen $K \subset X$ der Schnitt $K \cap M$ endlich ist. Wir verwenden oft, dass in einer lokal endlichen Menge $M$ jeder Punkt $m \in M$ eine offene Umgebung $U$ besitzt, sodass $U \cap M=\{m\}$ ist.



\begin{DEF}{RS-M25-02-12}{Analytische und lokal endliche Mengen}
Sei $X$ eine Riemannsche Fläche.
\begin{enumerate}[label=(\alph*)]
\item Eine abgeschlossene Teilmenge $A\subseteq X$ heißt \emph{analytisch},  
wenn zu jedem Punkt $p\in A$ eine Umgebung $U\subseteq X$ existiert,  
so dass
\[
U\cap A = \{\, z\in U \mid f(z)=0 \,\}
\]
für eine holomorphe Funktion $f:U\to\C$ gilt.  
Mit anderen Worten: lokal ist $A$ die Nullstellenmenge einer holomorphen Funktion.

\item Eine Teilmenge $M\subseteq X$ heißt \emph{lokal endlich},  
wenn für jede kompakte Teilmenge $K\subseteq X$ der Schnitt $K\cap M$ endlich ist.

\item Für lokal endliche Mengen gilt insbesondere:  
Zu jedem Punkt $m\in M$ existiert eine offene Umgebung $U\subseteq X$,  
so dass
\[
U\cap M=\{m\}.
\]
\end{enumerate}
\end{DEF}



Satz 2.13 

Sei $X$ zusammenhängend und $A$ analytisch. Dann ist $A=X$ oder $A$ ist lokal endlich in $X$.


\begin{PROP}{RS-M25-02-15}{Struktur analytischer Mengen auf Riemannschen Flächen}
Sei $X$ eine zusammenhängende Riemannsche Fläche und $A\subseteq X$ eine analytische Menge.  
Dann gilt:
\[
A = X \quad \text{oder} \quad A \text{ ist lokal endlich in } X.
\]
\end{PROP}

Beweis des Satzes 2.13. 

\begin{PROOF}{RS-M25-02-16}{P: Struktur analytischer Mengen auf Riemannschen Flächen}
Sei $M$ die Menge der Häufungspunkte von $A$. Diese ist abgeschlossen und $M \subset A$. Anwenden des Lemmas auf einen Punkt $a \in M$ ergibt, dass es eine offene Umgebung von $a$ gibt, die ganz in $M$ liegt. Also ist $M$ offen. Folglich ist $M=X$ oder $M=\emptyset$. Damit hat jeder Punkt in $a \in A$ eine offene Umgebung $U_{A}$ mit $A \cap U_{A}=\{a\}$ und damit ist $A$ lokal endlich.
\end{PROOF}


Lemma 2.14 

\begin{LEM}{RS-M25-02-13}{Lokales Verschwinden um Häufungspunkte der Nullstellenmenge}
Sei $U \subseteq \mathbb{C}$ ein Gebiet und $f: U \longrightarrow \mathbb{C}$ holomorph und $a \in U$ ein Häufungspunkt der Nullstellenmenge von $f$. Dann gibt es eine Umgebung von a, auf der $f$ identisch verschwindet.
\end{LEM}

Beweis: 

\begin{PROOF}{RS-M25-02-14}{P: Lokales Verschwinden um Häufungspunkte der Nullstellenmenge}
Wir leiten die Aussage aus dem bekannten Satz, dass holomorphe Funktionen lokal in Potenzreihen entwickelbar sind, her. Sei also
$$
f(z)=\sum_{n \geq 0} b_{n}(z-a)^{n}
$$
Wegen der Stetigkeit ist $f(a)=b_{0}=0$. Der Punkt $a \in U$ ist auch Häufungspunkt der Nullstellenmenge von
$$
\frac{f(z)}{z-a}=\sum_{n \geq 1} b_{n}(z-a)^{n-1}
$$
Folglich ist $b_{1}=0$ und induktiv folgt, dass alle $b_{n}$ verschwinden.
\end{PROOF}


Korollar 2.15 

Seien $f_{1}, f_{2}: X \longrightarrow Y$ zwei holomorphe Abbildungen und $X$ zusammenhängend. Wenn $A=\left\{a \in X: f_{1}(a)=f_{2}(a)\right\}$ einen Häufungspunkt besitzt, so ist $f_{1}=f_{2}$. Insbesondere ist jede nicht-konstante holomorphe Abbildung offen.

Beweis: Die Menge $A$ ist abgeschlossen und analytisch, die Behauptungen folgen unmittelbar aus dem Satz.



\begin{THEO}{RS-M25-02-17}{Identitätssatz und Offenheit holomorpher Abbildungen (Riemannsche Flächen)}
Seien $X,Y$ Riemannsche Flächen, $X$ zusammenhängend, und $f_1,f_2:X\to Y$ holomorph. Setze
\[
A:=\{\,a\in X \mid f_1(a)=f_2(a)\,\}.
\]
\begin{enumerate}
\item[(i)] Besitzt $A$ einen Häufungspunkt in $X$, so gilt $f_1\equiv f_2$ auf $X$.
\item[(ii)] Insbesondere ist jede nichtkonstante holomorphe Abbildung $f:X\to Y$ eine offene Abbildung.
\end{enumerate}
\end{THEO}

\begin{PROOF}{RS-M25-02-18}{P: Identitätssatz und Offenheit holomorpher Abbildungen (Riemannsche Flächen)}
\emph{Zu (i).} Für $p\in X$ wähle Karten $(U,g)$ bei $p$ und $(V,h)$ bei $f_1(p)=f_2(p)$ mit $g(p)=0$ und $h(f_i(p))=0$. Dann ist
\[
F:=h\circ f_1\circ g^{-1}-h\circ f_2\circ g^{-1}
\]
eine holomorphe Funktion auf $g(U)\subset\C$, und es gilt
\[
A\cap U \;=\;\{x\in U \mid F(g(x))=0\}.
\]
Also ist $A$ lokal die Nullstellenmenge einer holomorphen Funktion, somit analytisch und (da Gleichheitsmenge stetiger Abbildungen) abgeschlossen.  
Besitzt $A$ einen Häufungspunkt, so ist $A$ nach Proposition~2.16 gleich $X$ (denn eine analytische, abgeschlossene Teilmenge mit Häufungspunkt kann auf einer zusammenhängenden Riemannschen Fläche nicht lokal endlich sein, ohne ganz $X$ zu sein). Folglich $f_1=f_2$.

\smallskip
\emph{Zu (ii) (Offenheit).} Sei $f:X\to Y$ holomorph und nicht konstant. Für $p\in X$ wähle Karten $(U,g)$ bei $p$ und $(V,h)$ bei $f(p)$, $g(p)=0$, $h(f(p))=0$. Dann ist $F:=h\circ f\circ g^{-1}$ eine nichtkonstante holomorphe Funktion in einer Umgebung von $0\in\C$. Nach dem klassischen Offenheitssatz (bzw. dem lokalen Weierstraßschen Vorbereitungssatz/der lokalen Normalform $F(z)=z^m\cdot u(z)$ mit $u(0)\neq 0$) ist $F$ in $0$ offen. Also bildet $f$ $U$ offen auf $V$ ab. Da $p$ beliebig war, ist $f$ offen.
\end{PROOF}





\pagebreak


\section {Überlagerungen und Fundamentalgruppen}


Das erste Beispiel einer Überlagerung steckt implizit in der Einleitung. Ist $X$ die Riemannsche Fläche des Logarithmus, so definiert die Abbildung, die jeder Karte (bestehend aus einer geschlitzten Ebene) in die Ebene abbildet, eine wohldefinierte surjektive Abbildung $X \rightarrow \mathbb{C}^{*}$. Abbildungen, die so aussehen, werden wir als unverzweigte Überlagerungen bezeichnen. Der allgemeine Begriff der Überlagerung lässt lokal noch etwas mehr Flexibilität zu.

Fundamentalgruppen beschreiben, welche unverzweigten Überlagerungen eine Riemannsche Fläche zulässt. Sie sind die erste Invariante, die uns erlaubt, Riemannsche Flächen zu unterscheiden. Allerdings ist diese Invariante rein topologischer Natur und daher recht grob.

\subsection {Überlagerungen}
Seien $X$ und $Y$ Riemannsche Flächen. Wir setzen in diesem Abschnitt generall $X$ und $Y$ als nichtleer und zusammenhängend voraus. In Proposition 2.11 haben wir jede nicht-konstante holomorphe Abbildung $f: X \rightarrow Y$ lokal in $X$ in geeigneten Karten als $z \mapsto z^{n}$ beschrieben. 

Der im folgenden definierte Begriff der Überlagerung beschreibt holomorphe Abbildungen lokal aus der Sichtweise von $Y$. In manchen Quellen wird dieser verzweigte Überlagerung genannt. Der Begriff der Überlagerung aus der Topologie korrespondiert zu dem unten definierten Begriff einer unverzweigten Überlagerung.

\begin{figure}[h]
\begin{center}
  \includegraphics[scale=0.2]{2025_10_06_374d4cfa8b1948eb4c41g-016}
 
\caption{Abbildung 3.1: Eine (verzweigte) Überlagerung}
\end{center}
\end{figure}

Definition 3.1 (Eigentliche und endliche Abbildungen, Überlagerungen)

i) Eine holomorphe Abbildung $f: X \rightarrow Y$ heißt eigentlich, falls das Urbild jedes Kompaktums in $Y$ kompakt in $X$ ist.\\
ii) Eine holomorphe Abbildung $f$ heißt endlich, falls $f$ eigentlich ist und für jeden Punkt $y \in Y$ die Faser $f^{-1}(y)$ endlich ist.\\
iii) Eine holomorphe Abbildung $f$ heißt Überlagerung, wenn für alle $y \in Y$ ein $r=r(y)$ sowie Kartenumgebungen $\left(V, h: V \rightarrow \Delta_{r}\right)$ von $y$ und Kartenumgebungen ( $U_{i}, g_{i}$ ) von allen Punkten $x_{i} \in f^{-1}(y)$ existieren, sodass $h(y)=0$ und $g_{i}\left(x_{i}\right)=0$ ist, sodass $f\left(U_{i}\right)=V$ gilt und $f^{-1}(V)=\bigcup_{i \in I} U_{i}$ die disjunkte Vereinigung der $U_{i}$ ist, und $f$ in den Karten die Darstellung

$$
h \circ f \circ g_{i}^{-1}: \Delta_{r^{1 / n_{i}}} \rightarrow \Delta_{r}, \quad z \mapsto z^{n_{i}}
$$

mit einem $n_{i} \in \mathbb{N}$ hat.

Die Mengen $U_{i}$ und $V$ aus der Definition einer Überlagerung werden Standardkarten genannt.

In der Definition einer Überlagerung darf $n_{i}$ von $i$ und $y$ abhängen und $I$ ist dabei nicht notwendig endlich.


\begin{DEF}{COMPL-03-01}{Eigentliche, endliche und Überlagerungsabbildungen}
Seien $X$ und $Y$ Riemannsche Flächen und $f:X\to Y$ eine holomorphe Abbildung.
\begin{enumerate}\itemsep0.3em

\item \textbf{Eigentliche Abbildungen.}  
$f$ heißt \emph{eigentlich}, falls für jedes kompakte $K\subset Y$ das Urbild $f^{-1}(K)$ kompakt in $X$ ist.

\item \textbf{Endliche Abbildungen.}  
$f$ heißt \emph{endlich}, falls $f$ eigentlich ist und zusätzlich jede Faser
\[
f^{-1}(y)=\{x\in X:\ f(x)=y\}
\]
für alle $y\in Y$ endlich ist.

\item \textbf{Überlagerungen.}  
$f$ heißt \emph{Überlagerung}, wenn für jedes $y\in Y$ ein Radius $r=r(y)>0$ und lokale Kartenumgebungen
\[
(V,h:V\to\Delta_r)\quad\text{um }y,\qquad (U_i,g_i)\text{ um }x_i\in f^{-1}(y)
\]
existieren, so dass gilt:
\begin{itemize}
  \item $h(y)=0$ und $g_i(x_i)=0$ für alle $i$;
  \item $f(U_i)=V$ und $f^{-1}(V)=\bigsqcup_{i\in I}U_i$ ist die disjunkte Vereinigung der $U_i$;
  \item In diesen Karten gilt die Normalform
  \[
  h\circ f\circ g_i^{-1}:\ \Delta_{r^{1/n_i}}\to \Delta_r,\qquad z\mapsto z^{n_i},
  \]
  wobei $n_i\in\mathbb{N}$ von $i$ und $y$ abhängen darf.
\end{itemize}

\end{enumerate}
Die Mengen $U_i$ und $V$ heißen \emph{Standardkarten} der Überlagerung.
\end{DEF}


>> Beispiel für nicht eigentlich Abbildung: $L$-Fläche aus ÜB1 auf punktierten Torus


Die Nützlichkeit des Begriffs "eigentlich" zeigt sich an folgender Beobachtung. Ist $f: X \rightarrow Y$ eine holomorphe Abbildung und $X$ kompakt, so ist $f$ insbesondere eigentlich. Aber auch die Einschränkung einer solchen Abbildung $f: X \backslash f^{-1}(A) \rightarrow Y \backslash A$ für eine abgeschlossene Teilmenge $A \subset Y$ ist noch immer eigentlich, obwohl der Definitionsbereich nun nicht mehr kompakt ist.

\begin{REM}{AG-02-04}{Eigentliche Abbildungen und kompakte Urbilder}
Die Nützlichkeit des Begriffs der \emph{Eigentlichkeit} zeigt sich an folgender Beobachtung:

Sei \(f:X\to Y\) eine holomorphe Abbildung zwischen Riemannschen Flächen.  
Ist \(X\) kompakt, so ist \(f\) insbesondere \emph{eigentlich}, da für jede kompakte Teilmenge \(K\subseteq Y\) auch das Urbild \(f^{-1}(K)\subseteq X\) kompakt ist.

Bemerkenswerterweise bleibt diese Eigenschaft erhalten, wenn man \(f\) auf den komplementären offenen Teilbereich zu einer abgeschlossenen Teilmenge einschränkt:
\[
f|_{X\setminus f^{-1}(A)}: X\setminus f^{-1}(A)\longrightarrow Y\setminus A,
\qquad A\subseteq Y\ \text{abgeschlossen.}
\]
Obwohl der neue Definitionsbereich \(X\setminus f^{-1}(A)\) im Allgemeinen \emph{nicht mehr kompakt} ist,
bleibt die Abbildung weiterhin \emph{eigentlich}.
Dies illustriert, dass Eigentlichkeit eine stabilere Eigenschaft ist als Kompaktheit des Definitionsbereichs.
\end{REM}

Lemma 3.2 Jede eigentliche holomorphe Abbildung $f: X \rightarrow Y$ ist abgeschlossen.


Ist \(f:X\to Y\) \emph{eigentlich} und holomorph, so ist \(f\) eine
\emph{abgeschlossene Abbildung}, d.\,h.
für jede abgeschlossene Teilmenge \(A\subseteq X\) ist auch \(f(A)\subseteq Y\) abgeschlossen.


Beweis: 

Sei $A \subset X$ abgeschlossen und $b \in Y \backslash f(A)$. Sei $V$ eine Umgebung von $b$ mit kompakter Hülle $\bar{V}$. Dann ist $f^{-1}(\bar{V}) \cap A$ kompakt. Es gilt stets $f\left(f^{-1}(\bar{V}) \cap A\right)=\bar{V} \cap f(A)$ und diese Menge ist als stetiges Bild einer kompakten Menge wieder kompakt. Folglich ist $V \backslash(\bar{V} \cap f(A))$ eine Umgebung von $b$, die $f(A)$ nicht trifft.

Satz 3.3 Sei $f: X \rightarrow Y$ eine holomorphe Abbildung.\\
a) $f$ ist genau dann endlich, wenn $f$ eigentlich und nicht-konstant ist.\\
b) Ist $f$ endlich, so ist $f$ eine Überlagerung.\\
c) Ist $f$ eine Überlagerung und zudem für einen Punkt $y \in Y$ die Indexmenge $I=I_{y}$ aus der Definition endlich, so ist $f$ endlich und somit insbesondere eigentlich.


\begin{THEO}{AG-03-03}{Charakterisierung endlicher holomorpher Abbildungen zwischen Riemannschen Flächen}
Sei \(f:X\to Y\) eine holomorphe Abbildung zwischen zusammenhängenden,
nichtleeren Riemannschen Flächen.

\begin{enumerate}[label=\textbf{(\alph*)}]
\item
\(f\) ist genau dann \emph{endlich}, wenn \(f\) \emph{eigentlich} und
\emph{nicht-konstant} ist.

\item
Ist \(f\) endlich, so ist \(f\) eine \emph{Überlagerung}, d.\,h.\
jedes \(y\in Y\) besitzt eine Umgebung \(V\subseteq Y\), so dass
\(f^{-1}(V)\) eine endliche disjunkte Vereinigung offener Mengen
\(U_1,\dots,U_n\subseteq X\) mit
\[
f|_{U_j}:U_j\longrightarrow V
\]
biholomorph für alle \(j=1,\dots,n\) ist.

\item
Ist \(f\) eine Überlagerung und existiert ein Punkt \(y\in Y\),
für den die zugehörige Indexmenge \(I_y\) (der Blätter über \(y\))
endlich ist, so ist \(f\) \emph{endlich} und damit insbesondere \emph{eigentlich}.
\end{enumerate}
\end{THEO}

Beweis: 

a) Wenn $f$ endlich ist, dann ist für alle $y \in Y$ die Faser $f^{-1}(y)$ endlich, also insbesondere verschieden von $X$, da $X \neq \emptyset$. Damit ist $f$ eigentlich und nicht-konstant. Ist $f$ eigentlich, so ist $f^{-1}(y)$ kompakt. Die Faser ist aber auch eine analytische Menge und von $X$ verschieden, also lokal endlich und somit endlich.

b) Sei $y \in Y$. Da $f$ eigentlich und nicht konstant ist, und da $Y$ zusammenhängend und $X \neq \emptyset$ ist, muss $f$ surjektiv sein (Übung). Also ist $f^{-1}(y)= \left\{x_{1}, \ldots, x_{d}\right\}$. Wir wählen Karten ( $V, h$ ) um $y$ mit $h(y)=0$ und ( $U_{i}, g_{i}$ ) um $x_{i}$ mit $g\left(x_{i}\right)=0$ und können hierbei nach Verkleinern der Umgebungen $U_{i}$ ohne Einschränkung annehmen, dass $f\left(U_{i}\right) \subset V$ gilt und dass die $U_{i}$ disjunkt sind. Nach Proposition 2.11und deren Beweis finden wir für jedes $i$ eine Karte ( $\widetilde{U}_{i}, \widetilde{g}_{i}$ ) mit $\widetilde{U}_{i} \subset U_{i}$ so dass
$$
h \circ f \circ \widetilde{g}_{i}^{-1}: \widetilde{g}_{i}\left(\widetilde{U}_{i}\right) \rightarrow h(V), \quad z \mapsto z^{n_{i}}
$$
mit einem $n_{i}=e_{x_{i}}(f) \in \mathbb{N}$ gilt. Indem wir zunächst $V$ durch den Schnitt 
$$V \cap \bigcap_{i=1}^{d} f\left(\widetilde{U}_{i}\right)$$ 
und dann $\widetilde{U}_{i}$ durch $U_{i}=\widetilde{U}_{i} \cap f^{-1}(V)$ ersetzen, können wir $f\left(U_{i}\right)=V$ sicherstellen. Da $f$ eigentlich ist, gibt es zu jedem $i$ ein $r_{i}$, sodass $h^{-1}$ auf $\Delta_{r_{i}}$ und $g_{i}^{-1}$ auf $\Delta_{r_{i}^{1 / n} i}$ definiert ist. Wir nehmen $r$ als das Minimum der $r_{i}$ und schließlich $V=h^{-1}\left(\Delta_{r}\right)$ sowie $U_{i}=g_{i}^{-1}\left(\Delta_{r^{1 / n}}\right)$.

c) Wir zeigen zunächst, dass $f^{-1}(y)$ für jedes $y \in Y$ endlich ist. Dazu definieren wir
$$
\operatorname{deg}(f, y): Y \rightarrow \mathbb{N} \cup\{\infty\}, y \mapsto \sum_{x \in f^{-1}(y)} e_{x}(f) .
$$
Da $f$ eine Überlagerung ist, ist diese Abbildung lokal konstant, also konstant, da $Y$ zusammenhängend ist. Nach Voraussetzung ist die Konstante endlich, da $I=I_{y}$ in einem Punkt endlich ist. Es fehlt noch, die Eigentlichkeit zu zeigen. Sei dazu $K \subset Y$ kompakt. Für jedes $y \in Y$ gibt es Umgebungen $V_{y}$ und $U_{y, i}$ aus der Definition von Überlagerung. Da $Y$ lokal kompakt ist, finden wir auch offene Umgebungen $V_{y}^{\prime}$ von $y$, so dass $\overline{V_{y}^{\prime}}$ kompakt und in $V_{y}$ enthalten sind. Da $K$ kompakt ist, reichen endlich viele $V_{1}^{\prime}, \ldots, V_{r}^{\prime}$ um $K$ zu überdecken. Es reicht nun zu zeigen, dass $f^{-1}\left(\overline{V_{i}^{\prime}}\right)$ für alle $i=1, \ldots, r$ kompakt ist. Denn die endliche Vereinigung von Kompakta ist kompakt und 
$$f^{-1}(K) \subset \bigcup_{i=1}^{r} f^{-1}\left(\overline{V_{i}^{\prime}}\right)$$ 
ist als abgeschlossene Menge in einem Kompaktum ebenfalls kompakt. Dafür reicht es zu zeigen, dass 
$$\left.f\right|_{U_{i j}}: U_{i j} \rightarrow V_{i}$$ 
eigentlich ist. Dies kann man, da die Karten $h$ und $g_{i}$ Homöomorphismen sind, genauso gut für $h \circ f \circ g_{i}^{-1}$ nachprüfen. Die Abbildung $z \mapsto z^{n}$ ist aber eigentlich, denn das Urbild einer kompakten Menge ist beschränkt - da das Urbild von $B_{r}(0)$ in $B_{\sqrt[n]{r}}(0)$ liegt und abgeschlossen, also kompakt.



Die in obigem Beweis definierte Funktion $\operatorname{deg}(f)$ heißt Grad der endlichen Abbildung. Die Punkte $x \in X$ mit $e_{x}(f)>1$ heißen Verzweigungspunkte der Überlagerungsabbildung. Überlagerungsabbildungen ohne Verzweigungspunkte heißen unverzweigt.

\begin{DEF}{AG-03-04}{Grad und Verzweigungspunkte endlicher holomorpher Abbildungen}
Sei \(f:X\to Y\) eine endliche holomorphe Abbildung zwischen Riemannschen Flächen.
Dann ist für jedes \(y\in Y\) das Urbild \(f^{-1}(y)\) endlich.
Wir definieren die \emph{lokale Verzweigungszahl} \(e_x(f)\in\mathbb{N}\)
für \(x\in X\) gemäß der lokalen Normalform \(f(z)=z^{\,e_x(f)}\) (vgl.\ Proposition~2.11).

\medskip
\noindent
\textbf{Definition des Grades.}
Der \emph{Grad} der Abbildung \(f\) ist die Funktion
\[
\deg(f):Y\longrightarrow \mathbb{N}\cup\{\infty\},\qquad
\deg(f,y)=\sum_{x\in f^{-1}(y)} e_x(f).
\]
Ist \(f\) endlich, so ist die Summe für jedes \(y\in Y\) endlich,
und der Grad \(\deg(f)\) ist konstant auf ganz \(Y\).

\medskip
\noindent
\textbf{Verzweigungspunkte.}
Ein Punkt \(x\in X\) heißt \emph{Verzweigungspunkt} von \(f\),
falls seine lokale Verzweigungszahl größer als eins ist, also
\[
e_x(f)>1.
\]
Die Menge aller solcher Punkte wird mit
\[
\operatorname{Br}(f)=\{x\in X\mid e_x(f)>1\}
\]
bezeichnet.

\medskip
\noindent
\textbf{Unverzweigte Überlagerungen.}
Eine endliche Überlagerungsabbildung \(f:X\to Y\) heißt
\emph{unverzweigt}, wenn keine Verzweigungspunkte auftreten, d.\,h.\
wenn für alle \(x\in X\) gilt:
\[
e_x(f)=1.
\]
\end{DEF}


Beispiel 3.4 Die Abbildung $\exp (\cdot): \mathbb{C} \longrightarrow \mathbb{C}^{*}$ ist eine unverzweigte Überlagerung, aber keine endliche Abbildung.

\begin{EXA}{AG-03-04}{Exponentialabbildung als unverzweigte, nicht-endliche Überlagerung}
Betrachte die holomorphe Abbildung
\[
\exp:\C\longrightarrow\C^*,\qquad z\longmapsto e^{z}.
\]
Diese Abbildung ist eine \emph{unverzweigte Überlagerung}, denn für jedes
\(w\in\C^*\) gilt
\[
\exp^{-1}(w)=\{\log w+2\pi i k \mid k\in\Z\},
\]
und die Ableitung \(\exp'(z)=e^z\) verschwindet nirgends,
also existieren keine Verzweigungspunkte (\(e_z(\exp)=1\) für alle \(z\in\C\)).

Allerdings ist \(\exp\) \emph{nicht endlich}, da die Urbilder
\(\exp^{-1}(w)\) unendlich viele Punkte enthalten.
Insbesondere ist die Exponentialabbildung nicht eigentlich:
das Urbild einer kompakten Teilmenge von \(\C^*\)
ist im Allgemeinen nicht kompakt in \(\C\).
\end{EXA}

Die Abbildung $z \mapsto z^{2}$ is eine Überlagerung, da aber nicht unverzweigt ist. Der Nullpunkt ist ein Verzweigungspunkt.

\begin{EXA}{AG-03-05}{Quadratabbildung als verzweigte Überlagerung}
Die Abbildung
\[
f:\C\longrightarrow\C,\qquad z\longmapsto z^{2},
\]
ist eine \emph{Überlagerungsabbildung}, jedoch \emph{nicht unverzweigt}.

\medskip
\noindent
Für jedes \(w\in\C\setminus\{0\}\) gilt
\[
f^{-1}(w)=\{\sqrt{w},\,-\sqrt{w}\},
\]
und in einer Umgebung jedes solchen \(w\) ist \(f\) eine zweiblättrige
Überlagerung. Dagegen gilt im Punkt \(z=0\)
\[
f'(z)=2z \quad\Rightarrow\quad f'(0)=0,
\]
also verschwindet die Ableitung dort, und die lokale Form
\(f(z)=z^{2}\) zeigt, dass
\[
e_{0}(f)=2.
\]
Daher ist \(z=0\) ein \emph{Verzweigungspunkt} der Überlagerungsabbildung.

\medskip
\noindent
Die Abbildung \(z\mapsto z^{2}\) ist somit eine \emph{verzweigte Überlagerung}
vom Grad \(2\) mit genau einem Verzweigungspunkt bei \(z=0\).
\end{EXA}

\pagebreak

\subsection{Wege und Homotopien}
Der Begriff eines Weges lässt sich in offensichtlicher Weise von Gebieten in $\mathbb{C}$ auf topologische Räume, also insbesondere auf Riemannsche Flächen verallgemeinern.

Definition 3.5 (Wege) Eine stetige Abbildung $\gamma:[0,1] \longrightarrow X$ in einen topologischen Raum $X$ heißt Weg von $\gamma(\alpha)$ nach $\gamma(\beta)$. Der Raum $X$ heißt wegzusammenhängend, falls es zwischen je zwei Punkten $a, b \in X$ einen Weg von a nach $b$ gibt.

\begin{DEF}{AG-03-05}{Wege und Wegzusammenhang}
Sei \(X\) ein topologischer Raum.

\medskip
\noindent
Eine \emph{stetige Abbildung}
\[
\gamma:[0,1]\longrightarrow X
\]
heißt \emph{Weg in \(X\)} (oder \emph{Pfad}), wobei
\(\gamma(0)\) als \emph{Anfangspunkt} und \(\gamma(1)\) als \emph{Endpunkt}
des Weges bezeichnet wird.
Man sagt, \(\gamma\) ist ein \emph{Weg von} \(\gamma(0)\) \emph{nach} \(\gamma(1)\).

\medskip
\noindent
Ein topologischer Raum \(X\) heißt \emph{wegzusammenhängend},
falls es zu je zwei Punkten \(a,b\in X\)
einen Weg \(\gamma:[0,1]\to X\) mit
\[
\gamma(0)=a,\qquad \gamma(1)=b
\]
gibt.
\end{DEF}

Wege gibt es auf Riemannschen Flächen viele, zu viele um daraus eine interessante Invariante zu erhalten. Wir wollen nur Klassen von Wegen modulo der Relation „daran wackeln" betrachten.

Definition 3.6 (Homotopie) Zwei Wege $\gamma_{0}, \gamma_{1}: I=[0,1] \rightarrow X$ von a nach $b$ heißen homotop, falls es eine Schar von Wegen gibt, d.h. eine stetige Abbildung
$$
h: I \times I \rightarrow X, \gamma_{t}(s):=h(s, t)
$$
sodass $h(0, t)=a$ und $h(1, t)=b$ für alle $t \in I$ gilt und sodass $h(s, 0)=\gamma_{0}(s)$ und $h(s, 1)=\gamma_{1}(s)$ ist. $h$ wird Homotopie genannt.


\begin{DEF}{AG-03-06}{Homotopie von Wegen}
Sei \(X\) ein topologischer Raum und \(I=[0,1]\) das Einheitsintervall.

\medskip
\noindent
Zwei Wege \(\gamma_0,\gamma_1:I\to X\), die beide denselben Anfangs- und Endpunkt besitzen,
\[
\gamma_0(0)=\gamma_1(0)=a, \qquad \gamma_0(1)=\gamma_1(1)=b,
\]
heißen \emph{homotop} (relativ zu den Endpunkten),
wenn es eine stetige Abbildung
\[
h:I\times I\longrightarrow X,\qquad (s,t)\longmapsto h(s,t)
\]
gibt, sodass gilt:
\[
\begin{aligned}
h(0,t)&=a, & h(1,t)&=b \quad &\text{für alle } t\in I,\\
h(s,0)&=\gamma_0(s), & h(s,1)&=\gamma_1(s) \quad &\text{für alle } s\in I.
\end{aligned}
\]
Die Abbildung \(h\) heißt \emph{Homotopie zwischen \(\gamma_0\) und \(\gamma_1\)}.

\medskip
\noindent
Man schreibt in diesem Fall
\[
\gamma_0 \simeq \gamma_1 \quad (\text{rel } \{0,1\}).
\]
\end{DEF}


Lemma 3.7 „Homotop sein" ist eine Äquivalenzrelation.

\begin{LEM}{AG-03-07}{Homotopie als Äquivalenzrelation}
Die Relation „homotop sein“ auf der Menge der Wege in einem topologischen Raum \(X\)
mit festen Endpunkten ist eine \emph{Äquivalenzrelation}.
\medskip
\noindent
Das heißt, für alle Wege \(\gamma_0,\gamma_1,\gamma_2:[0,1]\to X\) mit
\(\gamma_i(0)=a\) und \(\gamma_i(1)=b\) gilt:
\begin{enumerate}
  \item \emph{(Reflexivität)} \(\gamma_0 \simeq \gamma_0\),
  \item \emph{(Symmetrie)} Aus \(\gamma_0 \simeq \gamma_1\) folgt \(\gamma_1 \simeq \gamma_0\),
  \item \emph{(Transitivität)} Aus \(\gamma_0 \simeq \gamma_1\) und \(\gamma_1 \simeq \gamma_2\)
        folgt \(\gamma_0 \simeq \gamma_2\).
\end{enumerate}
\end{LEM}

Beweis: Die Homotopien $h(s, t):=\gamma(s)$ und $\widetilde{h}(s, t)=h(s, 1-t)$ beweisen Reflexivität und Symmetrie. Ist $h$ eine Homotopie zwischen $\gamma_{0}$ und $\gamma_{1}$ und $h_{2}$ eine Homotopie zwischen $\gamma_{1}$ und $\gamma_{2}$, so ist
$$
h(s, t)= \begin{cases}h_{1}(s, 2 t) & \text { für } \quad 0 \leq t \leq \frac{1}{2} \\ h_{2}(s, 2 t-1) & \text { für } \quad \frac{1}{2}<t \leq 1\end{cases}
$$
eine Homotopie zwischen $\gamma_{0}$ und $\gamma_{2}$. Man beachte die Stetigkeit der Abbildung $h$ bei $t=1 / 2$. 



Ist $\gamma_{1}: I \rightarrow X$ ein Weg von $a$ nach $b$ und $\gamma_{2}: I \rightarrow X$ ein Weg von $b$ nach $c$, so definieren wir die Verkettung von Wegen als
$$
\gamma_{2} \cdot \gamma_{1}:\left\{\begin{array}{l}
I \longrightarrow X \\
s \longmapsto \begin{cases}\gamma_{1}(2 s) & \text { für } 0 \leq s \leq \frac{1}{2} \\
\gamma_{2}(2 s-1) & \text { für } \frac{1}{2}<s \leq 1 .\end{cases}
\end{array}\right.
$$



Lemma 3.8 Verkettung von Wegen ist assoziativ auf Homotopieklassen:

\begin{LEM}{AG-03-08}{Assoziativität der Wegverkettung bis auf Homotopie}
Seien \(X\) ein topologischer Raum und 
\[
\gamma_1:[0,1]\to X,\quad \gamma_2:[0,1]\to X,\quad \gamma_3:[0,1]\to X
\]
drei Wege mit gemeinsamen Endpunkten, sodass
\[
\gamma_1(1)=\gamma_2(0), \qquad \gamma_2(1)=\gamma_3(0).
\]
Dann gilt für ihre Homotopieklassen (relativ zu den Endpunkten)
\[
[\gamma_1 * \gamma_2] * \gamma_3 \;\simeq\; \gamma_1 * [\gamma_2 * \gamma_3],
\]
d.\,h. die Verkettung von Wegen ist auf Homotopieklassen \emph{assoziativ}.
\end{LEM}


Beweis: 

Sei $\omega_{0}=\gamma_{3} \cdot\left(\gamma_{2} \cdot \gamma_{1}\right)$ und $\omega_{1}=\left(\gamma_{3} \cdot \gamma_{2}\right) \cdot \gamma_{1}$.

Die beiden Wege beschreiben dieselbe Bildmenge, die Homotopie muss nur die Umparametrisierung durchführen. Die Homotopie
$$
h(s, t)= \begin{cases}\gamma_{1}(2(2-t) \cdot s) & 0 \leq s \leq \frac{1}{2(2-t)} \\ \gamma_{2}\left(\left(s-\frac{1}{2(2-t)}\right) \cdot \frac{2(2-t)(1+t)}{1+2 t-2 t^{2}}\right) & \frac{1}{2(2-t)}<s \leq 1-\frac{1}{2(1+t)} \\ \gamma_{3}\left(2 \cdot(1+t) \cdot\left(s-1+\frac{1}{2(1+t)}\right)\right) & 1-\frac{1}{2(1+t)}<s \leq 1\end{cases}
$$
leistet das Verlangte.


Definition 3.9 

Ein Weg $\gamma$ von a nach a heißt geschlossen oder Schleife. Der Weg $\gamma^{-1}(t)=\gamma(1-t)$ heißt umgekehrter Weg. Mit $\pi_{1}(X, a, b)$ bezeichnen wir die Menge der Homotopieklassen von Wegen von a nach b.

\begin{DEF}{AG-03-09}{Geschlossene Wege und Fundamentalgruppe}
Sei \(X\) ein topologischer Raum und \(a,b\in X\).

\begin{enumerate}
\item Ein Weg \(\gamma:[0,1]\to X\) mit \(\gamma(0)=a\) und \(\gamma(1)=a\) 
      heißt \emph{geschlossen} oder \emph{Schleife um \(a\)}.
\item Der \emph{umgekehrte Weg} zu einem Weg \(\gamma:[0,1]\to X\) ist definiert durch
      \[
      \gamma^{-1}(t)=\gamma(1-t), \qquad t\in[0,1].
      \]
      Er läuft denselben Pfad in entgegengesetzter Richtung.
\item Mit
      \[
      \pi_1(X,a,b)
      \]
      bezeichnen wir die Menge der Homotopieklassen (relativ zu den Endpunkten)
      von Wegen \(\gamma:[0,1]\to X\) mit \(\gamma(0)=a\), \(\gamma(1)=b\).
\end{enumerate}
Im Spezialfall \(a=b\) spricht man von der \emph{Fundamentalgruppe} von \(X\) im Punkt \(a\):
\[
\pi_1(X,a):=\pi_1(X,a,a).
\]
\end{DEF}


Definition und Proposition 

3.10 

Die Homotopieklassen von Wegen von a nach a bilden eine Gruppe
$$
\pi_{1}(X, a):=\pi_{1}(X, a, a),
$$
die Fundamentalgruppe von $X$ mit Basispunkt a. Ist b ein weiterer Punkt in $X$ und ist $X$ wegzusammenhängend, so ist $\pi_{1}(X, b)$ isomorph zu $\pi_{1}(X, a)$.

\begin{DEF}{AG-03-10}{Fundamentalgruppe und Basispunktunabhängigkeit}
\begin{enumerate}
\item
Die Menge der Homotopieklassen von geschlossenen Wegen (Schleifen) 
in einem topologischen Raum \(X\) mit Basispunkt \(a\in X\),
\[
\pi_1(X,a):=\pi_1(X,a,a),
\]
bildet eine \emph{Gruppe}, die \emph{Fundamentalgruppe von \(X\) im Punkt \(a\)} genannt wird.
\end{enumerate}

\begin{proof}[Gruppenstruktur]
Die Verknüpfung zweier Homotopieklassen \([\gamma_1],[\gamma_2]\in\pi_1(X,a)\)
wird durch die Verkettung der Wege definiert:
\[
[\gamma_1]\cdot[\gamma_2]:=[\gamma_1 * \gamma_2].
\]
Diese Operation ist wohldefiniert und assoziativ (nach Lemma~\ref{LEM:AG-03-08}).  
Das neutrale Element ist die konstante Schleife \(e_a(t)\equiv a\),
und das Inverse zu \([\gamma]\) ist \([\gamma^{-1}]\). 
\end{proof}

\begin{enumerate}\setcounter{enumi}{1}
\item
Sind \(a,b\in X\) zwei Punkte in einem \emph{wegzusammenhängenden} Raum \(X\),
so sind die Fundamentalgruppen \(\pi_1(X,a)\) und \(\pi_1(X,b)\) \emph{isomorph}.
\end{enumerate}
\end{DEF}

Beweis: 

Neutrales Element ist $\gamma(t):=a$, inverses Element ist der umgekehrte Weg. Ist $\eta$ ein Weg von $a$ nach $b$, so definiert
$$
\gamma \mapsto \eta \cdot \gamma \cdot \eta^{-1}: \pi_{1}(X, a) \rightarrow \pi_{1}(X, b)
$$
den gewünschten Isomorphismus. 



Definition 3.11 

(Einfach zusammenhängend) 

Ein zusammenhängender topologischer Raum $X$ heißt einfach zusammenhängend, falls $\pi_{1}(X, a)=\{i d\}$ für ein $a \in X$ ist.

\begin{DEF}{AG-03-11}{Einfach zusammenhängende Räume}
Ein zusammenhängender topologischer Raum \(X\) heißt \emph{einfach zusammenhängend},  
falls seine Fundamentalgruppe trivial ist, d.\,h. es existiert (und damit für alle) ein Punkt \(a\in X\) mit
\[
\pi_1(X,a)=\{\mathrm{id}\}.
\]
Äquivalent bedeutet dies: Jeder geschlossene Weg (Schleife) in \(X\) ist zu einer konstanten Schleife im selben Punkt homotop.
\end{DEF}

Proposition 3.12 Ist $X \subset \mathbb{C}$ sternförmig, so ist $X$ einfach zusammenhängend.


\begin{PROP}{AG-03-12}{Sternförmige Mengen sind einfach zusammenhängend}
Sei \(X \subset \mathbb{C}\) eine sternförmige Menge, d.\,h.\ es existiert ein Punkt \(x_0 \in X\), sodass  
für alle \(x \in X\) das Verbindungssegment
\[
[x_0,x]:=\{\, (1-t)x_0 + t x \mid t \in [0,1] \,\}
\subset X
\]
liegt.

Dann ist \(X\) \emph{einfach zusammenhängend}, d.\,h.
\[
\pi_1(X,x_0)=\{\mathrm{id}\}.
\]
\end{PROP}


Sei \(\gamma:[0,1]\to X\) eine geschlossene Schleife mit \(\gamma(0)=\gamma(1)=x_0\).
Wir definieren eine Homotopie, welche \(\gamma\) stetig auf die konstante Schleife am Basispunkt \(x_0\) zusammenzieht:
\[
H:[0,1]\times[0,1]\to X,\qquad H(s,t) := (1-t)x_0 + t\,\gamma(s).
\]
Da \(X\) sternförmig bezüglich \(x_0\) ist, liegt für jedes \(s,t\in[0,1]\) der Punkt \(H(s,t)\in X\).
Es gilt:
\[
H(s,0)=x_0,\quad H(s,1)=\gamma(s),\quad
H(0,t)=H(1,t)=x_0.
\]
Damit ist \(H\) eine Homotopie zwischen \(\gamma\) und der konstanten Schleife bei \(x_0\),
also ist jede Schleife in \(X\) nullhomotop.
Somit ist \(\pi_1(X,x_0)=\{\mathrm{id}\}\).



Beweis: 

Sei $a$ der Referenzpunkt des Sterns und $\gamma$ ein Weg. Dann ist
$$
h(s, t)=t \gamma(s)+(1-t) a
$$
eine Homotopie von $\gamma$ zum konstanten Weg am Punkt $a$. 




Oftmals ist es nützlich, einen Weg lokal zu modifizieren. Dazu verwendet man folgenden Zerlegungssatz.


Proposition 3.13 

Sei $\gamma: I \rightarrow X$ ein Weg und $U=\left\{U_{i}\right\}$ eine offene Überdeckung von $X$. Dann gibt es eine Zerlegung $0=s_{0}<s_{1}<\ldots<s_{n}=1$, sodass jeder Teilweg $\gamma\left(\left[s_{j}, s_{j+1}\right]\right)$ in einem $U_{i}$ enthalten ist.

\begin{LEM}{AG-03-13}{Zerlegung eines Weges unter offener Überdeckung}
Sei \(X\) ein topologischer Raum, \(\gamma:I=[0,1]\to X\) ein stetiger Weg und
\(\mathcal U=\{U_i\}_{i\in I_0}\) eine offene Überdeckung von \(X\).
Dann existiert eine Zerlegung
\[
0=s_0<s_1<\dots<s_n=1
\]
und Indizes \(i_0,\dots,i_{n-1}\in I_0\) mit
\[
\gamma\big([s_j,s_{j+1}]\big)\subset U_{i_j}\quad\text{für alle }j=0,\dots,n-1.
\]
\end{LEM}

\begin{proof}
Da \(\gamma\) stetig ist und \(\mathcal U\) eine offene Überdeckung von \(X\) ist,
bildet die Familie
\[
\mathcal V:=\{\gamma^{-1}(U_i)\subset[0,1]\mid i\in I_0\}
\]
eine offene Überdeckung des kompakten Intervalls \([0,1]\).
Nach Endlichkeitssatz der Kompaktheit (Heine--Borel) existiert eine
endliche Teilüberdeckung \(\gamma^{-1}(U_{i_1}),\dots,\gamma^{-1}(U_{i_m})\).

Da \([0,1]\) linear geordnet ist, können wir diese offenen Mengen durch
endlich viele offene Intervalle \(J_k=(a_k,b_k)\) mit
\(\overline{J_k}\subset\gamma^{-1}(U_{i_{\ell(k)}})\) überdecken und so wählen,
dass die Vereinigung \(\bigcup_k J_k\) bereits \([0,1]\) überdeckt
(benutze die Regularität von \(\mathbb R\) bzw.\ die Existenz lokaler Intervalle
um jeden Punkt und Kompaktheit zur Wahl einer endlichen, sogar
lokal-endlichen Familie mit \(\overline{J_k}\subset \gamma^{-1}(U_{i_{\ell(k)}})\)).

Ordne die endlich vielen Endpunkte \(\{a_k,b_k\}_k\) streng aufsteigend
und ergänze \(0\) und \(1\);
dies liefert eine Zerlegung \(0=s_0<s_1<\dots<s_n=1\).
Für jedes Teilintervall \([s_j,s_{j+1}]\) existiert (wegen der Überdeckung
durch die \(J_k\)) ein Index \(k(j)\) mit
\([s_j,s_{j+1}]\subset \overline{J_{k(j)}}\subset \gamma^{-1}(U_{i_{\ell(k(j))}})\).
Damit folgt
\(\gamma([s_j,s_{j+1}])\subset U_{i_{\ell(k(j))}}\), wie behauptet.
\end{proof}

\begin{REM}
Eine alternative, oft zitierte Argumentation benutzt das \emph{Lebesgue-Zahl-Lemma}:
Da \(\gamma([0,1])\) kompakt ist, besitzt die Überdeckung
\(\{U_i\cap \gamma([0,1])\}\) eine Lebesgue-Zahl \(\delta>0\).
Wählt man eine äquidistante Zerlegung mit Maschenweite \(<\delta\),
liegt jedes Bildsegment \(\gamma([s_j,s_{j+1}])\) vollständig in einem \(U_i\).
\end{REM}


Die Beweisidee veranschaulicht das Bild, die Details sind eine einfache Übung.
\begin{figure}[h]
\begin{center}
  \includegraphics[scale=0.2]{2025_10_06_374d4cfa8b1948eb4c41g-020}
 
\caption{Abbildung 3.3: Zerlegung in Teilwege}
\end{center}
\end{figure}




Korollar 3.14 Sei $A \subset X$ lokal endlich und $\gamma: I \rightarrow X$ ein Weg. Dann gibt es einen Weg $\widetilde{\gamma}: I \rightarrow X$ homotop zu $\gamma$ mit $\widetilde{\gamma}(I) \cap A=\phi$.

\begin{KORO}{AG-03-14}{Ausweichhomotopie um lokal endliche Teilmengen}
Sei $X$ eine Riemannsche Fläche, $A\subset X$ lokal endlich (d.\,h.\ jede kompakte Teilmenge von $X$ schneidet $A$ nur endlich oft), und
$\gamma:I=[0,1]\to X$ ein Weg. Dann existiert ein Weg $\widetilde\gamma:I\to X$, der zu $\gamma$ homotop ist, mit
\[
\widetilde\gamma(I)\cap A=\varnothing.
\]
\end{KORO}

Beweis: 

Überdecke $X$ durch Wahl von je einer Kreisscheibe $U_{a}$ für jeden Punkt $a \in A$ mit der Eigenschaft $U_{a} \cap A=\{a\}$ und durch das Komplement $X \backslash A$. Wir wenden nun den Satz an und nehmen durch Zusammenfassen von Teilwegen an, dass kein Teilweg $\gamma\left(\left[s_{k}, s_{k+1}\right]\right)$ ausschließlich in einem $U_{a}$ liegt. Folglich liegt kein Teilungspunkt in $A$. Jeder Teilweg in einem $U_{a}$ trifft maximal einen Punkt in $A$. Wir ersetzen ihn durch einen anderen Weg in $U_{a}$, der $A$ nicht trifft. Da $U_{a}$ sternförmig, also einfach zusammenhängend ist, sind diese Wege homotop.



Korollar 3.15 Die Zahlenkugel $\mathbb{P}_{\mathbb{C}}^{1}$ ist einfach zusammenhängend.

\begin{KORO}{AG-03-15}{Einfachzusammenhängendheit der Zahlenkugel $\mathbb{P}^1_\C$}
Die komplexe projektive Gerade $\mathbb{P}^1_\C$ (Riemannsche Zahlenkugel) ist einfach zusammenhängend.
\end{KORO}


Beweis: 

Wir können nach Korollar 3.14 annehmen, dass der Weg $\infty=\mathbb{P}_{\mathbb{C}}^{1} \backslash X_{1}$ vermeidet, also in $X_{1} \cong \mathbb{C}$ verläuft. Da $\mathbb{C}$ sternförmig ist, folgt die Behauptung aus Proposition 3.12.



Selbst die Bestimmung der einfachsten Fundamentalgruppe, die der punktierten Ebene oder punktierten Kreisscheibe, verwendet den Begriff der Liftung auf eine Überlagerung, und dies werden wir im folgenden Abschnitt allgemein untersuchen.



\pagebreak

\subsection{Unverzweigte Überlagerungen}
Wir wollen die Fundamentalgruppen von $X$ und $Y$ miteinander vergleichen, falls es eine unverzweigte Überlagerung von $X$ nach $Y$ gibt. In diesem Abschnitt seien $X$ und $Y$ stets zusammenhängend. Sei zunächst $f: X \rightarrow Y$ nur eine stetige Abbildung und $\gamma: I \rightarrow Y$ ein Weg. Ein Weg $\eta: I \rightarrow X$ heißt Liftung oder Hochhebung von $\gamma$, falls $\gamma=f \circ \eta$ gilt.

Die Abbildung $f$ heißt unbegrenzt, falls jeder Weg $\gamma:[0,1) \rightarrow X$ sich stetig nach 1 fortsetzen lässt, sofern dies schon für den Weg $f \circ \gamma:[0,1) \rightarrow Y$ gilt.


\begin{DEF}{CV-04-01}{Weghebung und Unbegrenztheit}
Seien $X,Y$ zusammenhängende topologische Räume und $f:X\to Y$ stetig.
\begin{enumerate}
\item Für einen Weg $\gamma:I\to Y$ heißt ein Weg $\eta:I\to X$ \emph{Lift} (oder \emph{Hochhebung}) von $\gamma$, falls $f\circ\eta=\gamma$ gilt.
\item Die Abbildung $f$ heißt \emph{unbegrenzt}, wenn für jeden Weg $\gamma:[0,1)\to X$ gilt:
Lässt sich $f\circ\gamma$ stetig in $t=1$ fortsetzen, so lässt sich auch $\gamma$ stetig in $t=1$ fortsetzen.
\end{enumerate}
\end{DEF}

Proposition 3.16 

Ist $f: X \rightarrow Y$ eine unverzweigte Überlagerung, so sind Liftungen von Wegen eindeutig durch den Anfangspunkt bestimmt, d.h. zwei Liftungen $\gamma_{1}, \gamma_{2}$ des Weges $\gamma: I \rightarrow Y$ stimmen überein, sobald sie an einem Punkt $c \in I$ übereinstimmen. Außerdem ist $f$ unbegrenzt.

\begin{PROP}{AG-03-16}{Eindeutigkeit von Weghebungen und Unbegrenztheit}
Sei $f:X\to Y$ eine unverzweigte Überlagerung zusammenhängender Riemannscher Flächen.

\begin{enumerate}
\item (\emph{Eindeutigkeit der Liftung})  
Sind $\gamma_1,\gamma_2:I\to X$ zwei Liftungen desselben Weges $\gamma:I\to Y$ mit  
\[
f\circ\gamma_1=f\circ\gamma_2=\gamma,
\]
und existiert ein Punkt $c\in I$ mit $\gamma_1(c)=\gamma_2(c)$, so gilt bereits
\[
\gamma_1=\gamma_2 \text{ auf } I.
\]

\item (\emph{Unbegrenztheit})  
Die Abbildung $f$ ist unbegrenzt, d.\,h.\ jeder Weg $\eta:[0,1)\to X$ lässt sich stetig in $t=1$ fortsetzen, sobald dies für $f\circ\eta$ möglich ist.
\end{enumerate}
\end{PROP}


Beweis: Zum Beweis der Eindeutigkeitsaussage betrachten wir die Menge $W=\left\{s \in I: \gamma_{1}(s)=\gamma_{2}(s)\right\} \subseteq I$, an der die Liftungen übereinstimmen. Diese ist nicht leer, abgeschlossen, da $X$ hausdorffsch ist und offen, da $f$ lokal ein Biholomorphismus ist und beide Wege Lifts desselben $\gamma$ sind. Da $I$ zusammenhängend ist, muss $W=I$ sein.

Für die zweite Aussage wähle eine Umgebung $V$ von $f \circ \gamma(1)$, sodass $f$ auf jeder Komponenten von $f^{-1}(V)$ die Identität ist. Dann ist die Fortsetzbarkeit klar.



Unbegrenztheit ist die wesentliche Voraussetzung, um allgemein Wege zu liften.


Proposition 3.17 Ist $f: X \rightarrow Y$ eine unverzweigte Überlagerung und $\eta: I \rightarrow Y$ ein Weg, so gibt es zu jedem $x \in X$ mit $f(x)=\eta(0)$ einen Weg $\gamma: I \rightarrow X$ mit $\gamma(0)=x$, welcher $\eta$ liftet, d.h. sodass $f \circ \gamma=\eta$.

\begin{PROP}{AG-03-17}{Existenz von Weghebungen für unverzweigte Überlagerungen}
Sei $f:X\to Y$ eine unverzweigte Überlagerung zusammenhängender Riemannscher Flächen
und $\eta:I=[0,1]\to Y$ ein stetiger Weg. Für jeden Punkt $x\in X$ mit
$f(x)=\eta(0)$ existiert ein Weg $\gamma:I\to X$ mit
\[
\gamma(0)=x \quad\text{und}\quad f\circ\gamma=\eta,
\]
d.\,h.\ $\gamma$ ist eine Liftung (Hochhebung) von $\eta$ durch $f$ mit vorgegebenem Anfangspunkt.
\end{PROP}


Beweis: Wir überdecken das Bild des Wegs $\eta$ mit $V_{i}$ wie in der Definition einer Überlagerung und zerlegen den Weg $\eta$ in Teilwege $\eta_{i}:\left[s_{i}, s_{i+1}\right] \rightarrow Y$, sodass $\eta_{i}\left(\left[s_{i}, s_{i+1}\right]\right) \subset V_{i+1}$. Wir zeigen die Existenz eines Lifts auf [ $0, s_{k}$ ] durch Induktion nach $k$. Für $k=1$ leistet dies die lokale Umkehrabbildung aus der Definition einer unverzweigten Überlagerung. Wenn wir $\gamma$ auf $\left[0, s_{k-1}\right]$ bereits konstruiert haben, so gibt es genau eine Komponente $U_{k}$ von $f^{-1}\left(V_{k}\right)$, die $\gamma\left(s_{k-1}\right)$ enthält. Sei also $\tau_{k}: V_{k} \rightarrow U_{k}$ die Umkehrabbildung mit dieser Komponente als Bildbereich. Dann setzen wir $\gamma$ auf $\left[s_{k-1}, s_{k}\right]$ durch $\tau_{k} \circ \eta$ fort.

Offensichtlich ist das Bild einer Homotopie unter einer stetigen Abbildung wieder eine Homotopie der Bildwege. Wesentlich zum Vergleich der Fundamentalgruppen ist nun folgender Liftungssatz für Homotopien.

Satz 3.18 (Liftungssatz) Sei $f: X \rightarrow Y$ eine unverzweigte Überlagerung. Zwei Wege $\gamma_{0}, \gamma_{1}: I \rightarrow X$ mit gleichem Anfangspunkt a sind homotop, sobald die Bilder $\eta_{i}=f \circ \gamma_{i}$ in $Y$ homotop sind.

Beweis: Sei $k: I^{2} \rightarrow Y$ eine Homotopie von $\eta_{0}$ nach $\eta_{1}$. Die Wege $\eta_{t}: I \rightarrow Y, \eta_{t}(s)=k(s, t)$ haben alle den Anfangspunkt $f(a)$, den Endpunkt $f\left(\gamma_{0}(1)\right)$ und lassen sich nach Proposition $3.17 \mathrm{zu} \gamma_{t}: I \rightarrow X$ liften. Es liegt also nahe,

$$
h: I^{2} \rightarrow X, \quad h(s, t):=\gamma_{t}(s)
$$

zu definieren Zu zeigen ist die Stetigkeit von $h$. Wir überdecken $\eta_{0}(I)$ mit $V_{i}$ aus der Definition der Überlagerung. Auf diesen $V_{i}$ hat $f$ eine stetigen Schnitt $\tau_{i}$ (d.h. $f \circ \tau_{i}=\mathrm{id}_{V_{i}}$ ) auf die Komponente, die $\gamma_{0}$ enthält. Zusammengesetzt erhalten wir somit einen stetigen Schnitt $\tau: \cup V_{i} \rightarrow X$, d.h. insbesondere gilt $f \circ \tau \circ \eta_{0}=\eta_{0}$. Da $I$ kompakt ist, genügen hierfür endlich viele offene Mengen $V_{i}$.

Zunächst gibt es ein $\varepsilon>0$, sodass $k(I \times[0, \varepsilon)) \subseteq \bigcup V_{i}$ wegen der Stetigkeit von $k$. Dann ist $\tau \circ k: I \times[0, \varepsilon) \rightarrow X$ ebenfalls stetig und stimmt dort mit h überein. Damit haben wir die Stetigkeit von $h$ auf $I \times[0, \varepsilon)$ gezeigt.

Angenommen $h$ ist nicht auf ganz $I \times I$ stetig. Dann gibt es ein minimales $t_{0}$, sodass $h$ auf $I \times\left[0, t_{0}\right]$ nicht stetig ist. Nach dem vorangehenden Argument ist $t_{0} \geq \varepsilon>0$. Sei also $\left(s_{0}, t_{0}\right)$ eine der Unstetigkeitsstellen von $h$. Wir nehmen nun eine Umgebung $V$ von $k\left(s_{0}, t_{0}\right)$ wie in der Definition einer Überlagerung her. Dann gibt es wieder einen stetigen Schnitt $\tau: V \rightarrow X$ auf die Komponente des Urbilds von $V$, die $h\left(t_{0}, s_{0}\right)$ enthält. Da $k$ stetig ist, ist $k\left(V_{s_{0}} \times V_{t_{0}}\right) \subset V$\\
für genügend kleine offene Umgebungen $V_{s_{0}}, V_{t_{0}} \subset I$ von $s_{0}$ bzw. von $t_{0}$. Wir betrachten $h$ nun als Liftung der Wege $u_{s}(t)=k(s, t)$ für $t \in\left[t_{0}-\delta, t_{0}\right]$ und $s \in V_{s_{0}}$ (d.h. wir vertauschen die Rollen von $s$ und $t$ ). Diese Wege haben nach der vorangehenden Proposition eine eindeutige Fortsetzung nach $t=t_{0}$. Da $\tau \circ k$ auf $V_{s_{0}} \times V_{t_{0}}$ einen Abbildung ist, die alle diese Wege liftet, muss $h$ mit $\tau \circ k$ übereinstimmen. Da $\tau \circ k$ stetig ist, folgt auch die Stetigkeit von $h, \mathrm{im}$ Widerspruch zur Annahme.

Korollar 3.19 Sei $f: X \rightarrow Y$ eine unverzweigte Überlagerung und $a \in X$. Dann ist der induzierte Homomorphismus $f_{*}: \pi_{1}(X, a) \rightarrow \pi_{1}(Y, f(a))$ injektiv.

Eine Anwendung ist ein topologisches Kriterium für Isomorphie von Überlagerungen. Wir nennen zwei unverzweigte Überlagerungen $f_{1}: X_{1} \rightarrow Y$ und $f_{2}: X_{2} \rightarrow Y$ isomorph, falls es eine biholomorphe Abbildung $\varphi: X_{1} \rightarrow X_{2}$ mit $f_{2} \circ \varphi=f_{1}$ gibt.

Satz 3.20 Zwei unverzweigte Überlagerungen $f_{1}: X_{1} \rightarrow Y$ und $f_{2}: X_{2} \rightarrow Y$ sind isomorph, falls es $a_{i} \in X_{i}$ und $b \in Y$ gibt mit $f_{i}\left(a_{i}\right)=b \in Y$ und

$$
\left(f_{1}\right)_{*}\left(\pi_{1}\left(X_{1}, a_{1}\right)\right)=\left(f_{2}\right)_{*}\left(\pi_{1}\left(X_{2}, a_{2}\right)\right)
$$

als Untergruppen von $\pi_{1}(Y, b)$.\\
Der Satz folgt aus zweimaliger Anwendung der folgenden Hilfsüberlegung.\\
Lemma 3.21 Sei $f: X \rightarrow Y$ eine unverzweigte Überlagerung und $g: Z \rightarrow Y$ eine stetige Abbildung mit $f(a)=g(c)=b$ und $g_{*}\left(\pi_{1}(Z, c)\right) \subseteq f_{*}\left(\pi_{1}(X, a)\right)$. Dann gibt es genau eine stetige Liftung $\widetilde{g}: Z \rightarrow X$ mit $f \circ \widetilde{g}=g$ und $\widetilde{g}(c)=a$.

Beweis: $\mathrm{Zu} z \in Z$ wähle einen Weg $\gamma$ von $c$ nach $z$ und lifte $g \circ \gamma$ zu einem Weg $\widetilde{\gamma}: I \rightarrow X$ mit $\widetilde{\gamma}(0)=a$. Sei $\widetilde{g}(z)$ der Endpunkt von $\widetilde{\gamma}$. Wir prüfen die Wohldefiniertheit. Ist $\gamma_{2}$ ein weiterer Weg von c nach z , so ist $\gamma_{2}^{-1} \cdot \gamma$ geschlossen, also $g_{*}\left(\gamma_{2}^{-1} \cdot \gamma\right)=f_{*}(\delta)$ in $\pi_{1}(Y, c)$, d.h.

$$
g \circ \gamma \sim\left(g \circ \gamma_{2}\right) \cdot(f \circ \delta)
$$

Die Lifte dieser beiden Wege nach $X$ haben den selben Startpunkt $a \in X$ und nach dem Homotopieliftungssatz denselben Endpunkt. Es ist noch zu zeigen, dass $\widetilde{g}$ stetig ist. Um dies beim Punkt $z_{0}$ zu zeigen, wählt man eine Umgebung $V$ von $g\left(z_{0}\right)$ wie in der Definition von Überlagerung, $U$ die Komponente von $f^{-1}(V)$, die den Endpunkt des Lifts eines Weges $\gamma_{0}$ von $c$ nach $z_{0}$ enthält und $W=g^{-1}(V)$. Für jeden Punkt $z$ in $W$ kann man einen Weg $\gamma$ von $c$ nach $z$ wählen, indem man $\gamma_{0}$ mit einem Weg in $W$ verkettet. Damit ist $\left.\widetilde{g}\right|_{W}=\left.f\right|_{U} ^{-1} \circ g$ und somit stetig.

Zum Beweis der Eindeutigkeit überdeckt man $g(\gamma(I))$ mit offenen Mengen wie in der Definition einer Überlagerung, und zerlegt $I$ in das Urbild dieser Mengen. Der Eindeutigkeitsbeweis erfolgt nun induktiv über die so erzeugten Teilintervalle wie auch schon in Proposition 3.17.

Beweis des Satzes: Das Lemma angewandt auf $f_{1}=f, f_{2}=g$ gibt eine stetige Abbildung $\widetilde{g}$ mit $f_{2}=f_{1} \circ \widetilde{g}$. Das gleiche Lemma angewandt auf $f_{2}=f$ und $f_{1}=g$ ergibt eine stetige Abbildung in die umgekehrte Richtung. Die Eindeutigkeitsaussage impliziert, dass diese Abbildungen zueinander invers sind. Betrachtet man die Konstruktion im vorigen Lemma genau, so sieht man, dass $\widetilde{g}$ die Verkettung von $g$ und einer lokalen Umkehrabbildung von $f$ ist. Da also $f$ und $g$ holomorph sind, ist auch $\widetilde{g}$ holomorph und somit biholomorph. $\square$

Definition 3.22 Eine unverzweigte Überlagerung $f: \widetilde{Y} \rightarrow Y$ heißt universell, falls es zu jeder unverzweigten Überlagerung $g: X \rightarrow Y$ mit zusammenhängendem Definitionsbereich $X$ und zu jeder Wahl von Basispunkten $c \in \widetilde{Y}$ und $a \in g^{-1}(f(c))$ genau eine holomorphe Abbildung $h: \widetilde{Y} \rightarrow X$ mit $h(c)=a$ und $g \circ h=f$ gibt.

Nach dem Satz 3.20 ist eine universelle Überlagerung (wenn sie existiert) bis auf Isomorphie eindeutig. Die Existenz werden wir später in Satz 3.43 zeigen.

Korollar 3.23 Jede einfach zusammenhängende, unverzweigte Überlagerung ist universell.

Die obere Halbebene $\mathbb{H}$ ist die universelle Überlagerung der punktierten Kreisscheibe $\triangle^{*}$, denn sie ist einfach zusammenhängend und $z \mapsto \exp (2 \pi i z)$ ist eine unverzweigte Überlagerungsabbildung, wie man durch Einschränken auf die Streifen $k<\operatorname{Re}(z)<k+1$ für $k \in \mathbb{Z}$ und $k \in \mathbb{Z}+\frac{1}{2}$ erkennt.


\pagebreak

\subsection {Fundamentalgruppenbestimmung I}

Von der punktierten Ebene bestimmen wir die Fundamentalgruppe mit Hilfe der Liftungseigenschaften aus dem vorigen Abschnitt direkt.

Proposition 3.24 Die Fundamentalgruppe von $\mathbb{C}^{*}$, von $\Delta^{*}$ und von $S^{1}$ ist jeweils isomorph zu $\mathbb{Z}$.

Beweis: Es genügt die Aussage für $S^{1}$ zu beweisen, denn ist $\gamma(t)$ ein Weg in $\mathbb{C}^{*}$, so ist $t \mapsto \gamma(t) /|\gamma(t)|$ ein dazu homotoper Weg in $S^{1}$, eine Homotopie ist $h(s, t)=\gamma(t) /|s \gamma(t)+(1-s)|$. In gleicher Weise vergleicht man die Fundamentalgruppe von $S^{1}$ und von $\frac{1}{2} S^{1}$ sowie $\frac{1}{2} S^{1}$ und $\Delta^{*}$.

Wir betrachen die unverzweigte Überlagerung $\mathbb{H} \rightarrow \Delta^{*}, z \mapsto \exp (2 \pi i z)$. Das Bild von $a=-\frac{1}{2 \pi} \log \left(\frac{1}{2}\right) i \in \mathbb{H}$ ist offenbar $\frac{1}{2} \in \Delta^{*}$. Sei $\eta \in \pi_{1}\left(\Delta^{*}, \frac{1}{2}\right)$. Dann gibt es genau einen Lift $\gamma$ von $\eta$ nach $\mathbb{H}$ mit Startpunkt $a$. Der Endpunkt des Weges $\eta$ hat die Gestalt $a+k$ für ein $k \in \mathbb{Z}$. Die Abbildung $\eta \mapsto k$ definiert offenbar einen surjektiven Homomorphismus $\pi_{1}\left(\Delta^{*}, \frac{1}{2}\right) \rightarrow \mathbb{Z}$. Ist $\eta$ im Kern, so ist der Lift von $\eta$ homotop zum trivialen Weg, da $\mathbb{H}$ einfach zusammenhängend ist. Das Bild dieser Homotopie überführt $\eta$ in den trivialen Weg. Damit haben wir die Bijektivität gezeigt.

Die Fundamentalgruppe von allen komplizierteren Riemannschen Flächen baut man aus dieser Proposition und dem folgenden Satz zusammen. Der Beweis hier folgt [Hat02].

Satz 3.25 (Seifert/van Kampen) Sei $X=U \cup V$ mit $U, V$ offen, zusammenhängend und wegzusammenhängend und sei auch $U \cap V$ zusammenhängend und wegzusammenhängend. Ist $a \in U \cap V$, so ist

$$
\pi_{1}(U \cup V, a)=\pi(U, a) *_{\pi_{1}(U \cap V, a)} \pi_{1}(V, a)
$$

ein amalgamiertes Produkt.\\
Beweis: Wir lassen den Basispunkt $a$ aller Fundamentalgruppen ab sofort in der Notation weg. Es seien $s_{1}: \pi_{1}(U \cap V) \rightarrow \pi_{1}(U)$ und $s_{2}: \pi_{1}(U \cap V) \rightarrow \pi_{1}(V)$ sowie $\alpha_{1}^{\prime}: \pi_{1}(U) \rightarrow \pi_{1}(X)$ und $\alpha_{2}^{\prime}: \pi_{1}(V) \rightarrow \pi_{1}(X)$ die natürlichen Abbildungen, die von der Inklusion induziert sind. Nach der universellen Eigenschaft des amalgamierten Produkts gibt es also einen eindeutigen Homomorphismus

$$
\varphi: \pi(U) *_{\pi_{1}(U \cap V)} \pi_{1}(V) \rightarrow \pi_{1}(X)
$$

Wir müssen zeigen, dass dieser surjektiv und injektiv ist.\\
Für die Surjektivität genügt es zu zeigen, dass wir jeden geschlossenen Weg $\gamma$ in $X$ mit Startpunkt $a$ als Verkettung von geschlossenen Wegen schreiben können, wovon jeder einzelne ganz in $U$ oder ganz in $V$ liegt. Dazu überdecken wir $I=[0,1]$ mit den offenen Mengen, bestehend aus den Zusammenhangskomponenten von $\gamma^{-1}(U)$ und $\gamma^{-1}(V)$. Da $I$ kompakt ist, genügen endliche viele solcher Mengen zur Überdeckung. Es gibt also, wie in Proposition 3.13 Zwischenpunkte $0=s_{0}<s_{1}<\ldots s_{n}=1$, sodass das Bild jedes Intervalls $\gamma\left(\left[s_{i}, s_{i+1}\right]\right)$ ganz in $U$ oder ganz in $V$ enthalten ist und $\gamma\left(\left[s_{i}, s_{i+1}\right]\right)$ genau dann ganz in $U$ enthalten ist, wenn $\gamma\left(\left[s_{i+1}, s_{i+2}\right]\right)$ ganz in $V$ enthalten ist. Für jedes $i=0, \ldots, n$ liegt $\gamma\left(s_{i}\right)$ dann in $U \cap V$ und wir wählen wir einen Hilfsweg $\eta_{i}$ von $a$ nach $\gamma\left(s_{i}\right)$, welcher ganz in $U \cap V$ enthalten ist. Dies ist aufgrund der Zusammenhangsvoraussetzung von $U \cap V$ möglich. Sei nun $\gamma_{i}=\left.\eta_{i} \cdot \gamma\right|_{\left[s_{i-1}, s_{i}\right]} \cdot \eta_{i-1}^{-1}$ für $i=1, \ldots, n$. Dann ist $\gamma=\prod_{i=1}^{n} \gamma_{n-i+1}$ und jedes der $\gamma_{i}$ liegt nach Konstruktion in einer der offenen Mengen $U$ oder $V$.

Wir fassen diesen Schritt nochmal zusammen. Eine Faktorisierung eines Weges $\gamma$ in $X$ ist eine formale Zerlegung $\gamma=\left[\gamma_{n}\right] \cdots\left[\gamma_{1}\right]$ in stetige Abbildungen $\gamma_{i}:\left[s_{i-1}, s_{i}\right] \rightarrow X$ mit $\gamma_{i}\left(s_{i}\right)=\gamma_{i}\left(s_{i-1}\right)$, sodass für jedes $i$ eine der Inklusionen $\operatorname{Bild}\left(\gamma_{i}\right) \subset U$ oder $\operatorname{Bild}\left(\gamma_{i}\right) \subset V$ gilt und sodass $\gamma$ homotop zur Verkettung $\gamma_{n} \circ \cdots \circ \gamma_{1}$ ist. Im ersten Schritt haben wir gezeigt, dass jeder Weg in $X$ eine Faktorisierung besitzt. Wir definieren auf der Menge aller Faktorisierungen von $\gamma$ eine Äquivalenzrelation, die von folgenden zwei Elementaroperationen erzeugt ist. Zum Einen definieren wir, dass falls in $\gamma=\left[\gamma_{n}\right] \cdots\left[\gamma_{1}\right]$ ist und $\gamma_{i}$ und $\gamma_{i+1}$ beide in $U$ oder beide in $V$ liegen, diese Teilwege zur Faktorisierung $\gamma=\left[\gamma_{n}\right] \cdots\left[\gamma_{i+1} \circ \gamma_{i}\right] \ldots\left[\gamma_{1}\right]$ zusammengefasst werden können. Des weiteren\\
erlauben wir, die Teilmengenzugehörigkeit eines Teilstücks $\gamma_{i}$ zu $U$ oder zu $V$ zu ändern, falls Bild $\left(\gamma_{i}\right) \subset U \cap V$ ist.

Zum Beweis der Injektivität ist nach Definition der Wortäquivalenz in einem amalgamierten Produkt zu zeigen, dass zwei Faktorisierungen $\gamma=\left[\gamma_{n}\right] \cdots\left[\gamma_{1}\right]$ und $\gamma=\left[\gamma_{n}\right] \cdots\left[\gamma_{1}\right]$ von in $X$ homotopen Wegen $\gamma$ und $\gamma^{\prime}$ äquivalent sind. Hierbei haben wir durch Übergang zu einer gemeinsamen Verfeinerung bereits angenommen, dass die Trennstellen beider Wege die gleichen $s_{i}$ sind.

Sei $h: I \times I \rightarrow X$ eine Homotopie von $\gamma$ zu $\gamma^{\prime}$. Wir behaupten, dass wir die kompakte Menge $I \times I$ mit abgeschlossenen achsenparallelen Rechtecken $R_{i j}=\left[s_{i}, s_{i+1}\right] \times\left[r_{j}, r_{j+1}\right]$ überdecken können, sodass jedes Rechteck ganz einer Zusammenhangskomponenten von $h^{-1}(U)$ und $h^{-1}(V)$ liegt und sodass die Trennpunkte der Faktorisierungen gerade die $s_{i}$ sind. Um dies einzusehen betrachten wir alle in $I \times I$ offenen achsenparallen Rechtecke die ganz in einer Komponente von $h^{-1}(U)$ oder $h^{-1}(V)$ liegen. Da $U$ und $V$ offen sind, überdeckt diese Kollektion ganz $I \times I$ und aufgrund der Kompaktheit genügt eine endliche Teilmenge. Jedes der offenen Rechtecke kann nun zu einem in $I \times I$ abgeschlossenen Rechteck verkleinert werden, sodass deren Vereinigung immer noch $I \times I$ überdeckt. Nach Übergang zu einer Verfeinerung sind all diese Rechtecke von der Form $\left[s_{i}, s_{i+1}\right] \times\left[r_{j}, r_{j+1}\right]$ und nach einer weiteren Verfeinerung die Trennpunkte unter den $s_{i}$, wie behauptet. Dabei eine $I=\{0, \ldots, n\}$ und $J=\{0, \ldots, m\}$ die Indexmengen für $i$ bzw. für $j$. Wir schreiben auch $R_{i+j n}$ für das Rechteck $R_{i j}$.

Seien $x_{i j}=\left(s_{i}, r_{j}\right)$ die Eckpunkte dieser Rechtecke und $\eta_{i j}$ Wege von $a$ nach $h\left(x_{i j}\right)$ ganz in $U$, falls $x_{i j} \in h^{-1}(U)$, und (zudem) ganz in $V$, falls $x_{i j} \in h^{-1}(V)$. Jeder Wege in $I \times I$ von einem Punkte in $\{0\} \times I$ zu einem Punkt in $\{1\} \times I$ definiert einen Weg in $X$, der bei $a$ beginnt und endet. Darunter sind inbesondere die Wege $\gamma^{(\ell)}$, die die Rechtecke $R_{1}, \ldots, R_{\ell}$ von $R_{\ell+1}, \ldots, R_{m n}$ trennt. Inbesondere ist $\gamma=\gamma^{(0)}$ und $\gamma^{\prime}=\gamma^{(m n)}$. Für jedes $\gamma^{(\ell)}$ betrachten wir die Faktorisierung entsprechend der Randkanten der Rechtecke, die $\gamma^{(\ell)}$ durchläuft, wobei wir die Wege mit den $\eta_{i j}$ prä- und postkomponieren, sodass jeder Teilweg in $a$ startet und endet. Es genügt also zu zeigen, dass die Faktorisierungen von $\gamma^{(\ell)}$ und $\gamma^{(\ell+1)}$ für alle $\ell$ äquivalent sind. Dies ist nun aber offensichtlich, indem wir den eine Homotopie zwischen den zwei Randwegen eines Rechtecks von ( $s_{i}, r_{j}$ ) nach $s_{i+1}, r_{j-1}$ ) verwenden, und gegebenenfalls zuvor die zweite Eigenschaft der Faktorisierungsäquivalenz, um die Wegzugehörigkeit von $U$ nach $V$ (oder umgekehrt) zu ändern.

Ist insbesondere $U \cap V$ einfach zusammenhängend, so ist $\pi_{1}(X)$ ein freies Produkt von $\pi_{1}(U)$ und $\pi_{1}(V)$.

Eine einfache Schleife um $a_{j} \in X$ mit Basispunkt $a$ ist ein Weg $\gamma: I \rightarrow X$ der folgenden Gestalt: Es gibt eine Karte $h: U \rightarrow \mathbb{C}$ mit $h\left(a_{j}\right)=0$, einen Punkt $\widetilde{a} \in U \backslash a_{j}$ in der Nähe von $a_{j}$, und einen Weg $\widetilde{\gamma}$ von $a$ nach $\widetilde{a}$, sodass $\gamma$ zur
Verkettung
$$
\widetilde{\gamma}^{-1} \cdot\left(h^{-1} \circ\left(t \mapsto e^{2 \pi i t} h(\widetilde{a})\right) \circ h\right) \cdot \widetilde{\gamma}
$$
homotop ist.



Satz 3.26 

Seien $a_{0} \ldots, a_{r} \in \mathbb{P}_{\mathbb{C}}^{1}$ paarweise verschieden mit $r \geq 1$. Es gibt einfache Schleifen $\gamma_{j}$ um die Punkte $a_{j}$ für $j=0, \ldots, r$, sodass
$$
\pi_{1}\left(\mathbb{P}_{\mathbb{C}}^{1} \backslash\left\{a_{0}, \ldots, a_{r}\right\}\right) \cong F_{\left\{\gamma_{1}, \ldots, \gamma_{r}\right\}} \quad \text { und } \quad \prod_{j=0}^{r} \gamma_{i}=i d
$$

Beweis: Durch Automorphismen von $\mathbb{P}_{\mathbb{C}}^{1}$ erreicht man ohne Einschränkung, dass $a_{0}=\infty$ und danach zudem dass die Realteile der $a_{i}$ verschieden sind (unter allen Abbildungen $z \mapsto e^{i \theta} z$ für $\theta \in \mathbb{R}$ ist diese Bedingung nur für endlich viele $\theta$ nicht erfüllt). Sei $R$ ein achsenparalleles Rechteck, das alle $a_{i}$, $i \geq 1$ enthält.

\begin{figure}[h]
\begin{center}
  \includegraphics[scale=0.2]{2025_10_06_374d4cfa8b1948eb4c41g-027}
 
\caption{Abbildung 3.4: Zerlegung der Punktierungen in zwei Teilmengen}
\end{center}
\end{figure}

Sei $x_{0} \in \partial R$ ein Basispunkt. Der Rand des Rechtecks ist ein einfacher Weg $\gamma_{0}$ um $\infty$ und es gilt

$$
\pi_{1}\left(\mathbb{P}_{\mathbb{C}}^{1} \backslash\left\{a_{0}, \ldots, a_{r}\right\}\right) \cong \pi_{1}\left(R \backslash\left\{a_{1}, \ldots, a_{r}\right\}\right.
$$

da alle Wege durch zentrische Streckung zu einem Weg in $R$ homotop sind. Wir beweisen die Aussage durch Induktion nach $r$. Für $r=1$ ist $\gamma_{1}=\gamma_{0}^{-1}$ eine gesuchte einfache Schleife um $a_{1}$. Für den Induktionsschritt wenden wir den Satz 3.25 von Seifert/van Kampen auf die Überdeckung von $R$ durch $U$ und $V$ (spezifiziert durch den im Bild markierten Realteil) an.


\pagebreak

\subsection {Triangulierung und kanonische Zerschneidung}
Ziel dieses Abschnittes ist es, jede kompakte Riemannsche Fläche topologisch aus einem relativ einfachen Polygon zu kontruieren. Hilfsmittel hierfür sind Triangulierungen, die wir nun definieren.

Sei $S$ das Dreieck mit Ecken $0,1, e^{2 \pi i / 3}$ in $\mathbb{C} \cong \mathbb{R}^{2}$. Ein 2-Simplex der Riemannschen Fläche $X$ ist eine stetige Abbildung $S \rightarrow X$, die eine Homöomorphismus auf das Bild ist. Die Bilder von $0,1, e^{2 \pi i / 3}$ werden Ecken des Simplexes genannt, die Bilder der Geradensegmente [ 0,1 ], $\left[1, e^{2 \pi i / 3}\right]$ und $\left[0, e^{2 \pi i / 3}\right]$ werden Kanten des Simplexes genannt, die Bilder der entsprechenden offenen Geradensegmente auch offene Kanten.

Eine Triangulierung von $X$ ist eine Kollektion von 2-Simplices $\alpha_{i}: S \rightarrow X$ mit $\bigcup_{i} \alpha_{i}(S)=X$ und folgenden Eigenschaften. Für jeden Punkt $a \in X$ gilt:\\
i) Liegt $a \in X$ im Bild von $\alpha_{i}$, aber nicht auf dessen Kante, so gibt es genau ein $i$ mit $a \in$ Bild $\left(\alpha_{i}\right)$. In diesem Fall ist $\alpha_{i}(S)$ eine Umgebung von $a$ in $X$.\\
ii) Liegt $a$ auf der offenen Kante $K$ eines Simplexes, so gibt es genau zwei Indices $i, j$ mit $K=\alpha_{i}(S) \cap \alpha_{j}(S)$ und $\alpha_{i}(S) \cup \alpha_{j}(S)$ ist Umgebung von $a$.\\
iii) Andernfalls gibt es eine endliche Zahl von Indices $i_{k}, k=1 \ldots n$, sodass $a$ eine Ecke von jedem $\alpha_{i_{k}}(S)$ ist, $\cup_{k=1}^{n} \alpha_{i_{k}}(S)$ ist eine Umgebung von $a$ und in geeigneter Nummerierung haben $\alpha_{i_{k}}(S)$ und $\alpha_{i_{k+1}}(S)$ sowie $\alpha_{i_{n}}(S)$ und $\alpha_{i_{1}}(S)$ genau eine Kante gemeinsam und alle anderen $\alpha_{i_{k}}(S)$, $\alpha_{i_{l}}(S)$ mit $k \neq l$ haben genau den Punkt $a$ gemeinsam.

Kurz gesagt ist es also bei einer Triangulierung nicht gestattet, dass der Eckpunkt eines Simplexes in der Mitte der Kante eines anderen Simplexes liegt.

\begin{DEF}{AG-03-18}{Triangulierung einer Riemannschen Fläche}
Sei 
\[
S \subset \C \cong \R^2
\]
das (abgeschlossene) Dreieck mit Ecken \(0,1,e^{2\pi i/3}\).
Eine stetige Abbildung
\[
\alpha:S\to X
\]
in eine Riemannsche Fläche \(X\) heißt \emph{\(2\)-Simplex} von \(X\), falls
\(\alpha\) ein Homöomorphismus auf ihr Bild ist. Die Bilder der Punkte
\(0,1,e^{2\pi i/3}\) heißen \emph{Ecken} des Simplexes, die Bilder der
Geradensegmente \([0,1], [1,e^{2\pi i/3}], [0,e^{2\pi i/3}]\) heißen
\emph{Kanten}, und die Bilder der entsprechenden offenen Segmente heißen
\emph{offene Kanten}.

Eine \emph{Triangulierung} von \(X\) ist eine Familie von \(2\)-Simplexes
\(\{\alpha_i:S\to X\}_{i\in I}\) mit
\[
X = \bigcup_{i\in I} \alpha_i(S),
\]
die folgende Eigenschaften erfüllt. Für jeden Punkt \(a\in X\) gilt:

\begin{enumerate}[label=\roman*)]
\item Liegt \(a\) im Inneren eines Simplexes, d.\,h.\ 
\(a\in \alpha_i(S)\) und \(a\) liegt nicht auf einer Kante von \(\alpha_i(S)\),
so existiert genau ein Index \(i\) mit \(a\in \alpha_i(S)\), und
\(\alpha_i(S)\) ist eine Umgebung von \(a\) in \(X\).

\item Liegt \(a\) auf einer offenen Kante \(K\) eines Simplexes,
so gibt es genau zwei Indizes \(i,j\) mit
\[
K = \alpha_i(S)\cap \alpha_j(S),
\]
und die Vereinigung \(\alpha_i(S)\cup \alpha_j(S)\) ist eine Umgebung von \(a\).

\item Andernfalls ist \(a\) eine Ecke von endlich vielen Simplexes:
es existieren Indizes \(i_1,\dots,i_n\) mit
\[
a \text{ Ecke von } \alpha_{i_k}(S) \quad (k=1,\dots,n),
\]
\(\bigcup_{k=1}^n \alpha_{i_k}(S)\) ist eine Umgebung von \(a\),
und in geeigneter Nummerierung haben
\(\alpha_{i_k}(S)\) und \(\alpha_{i_{k+1}}(S)\) (sowie
\(\alpha_{i_n}(S)\) und \(\alpha_{i_1}(S)\)) genau \emph{eine} Kante
gemeinsam, während für alle übrigen Paare \(k\neq \ell\) gilt, dass
\(\alpha_{i_k}(S)\) und \(\alpha_{i_\ell}(S)\) nur den Punkt \(a\) gemeinsam haben.
\end{enumerate}
\end{DEF}

\begin{REM}{AG-03-19}{Geometrische Intuition}
Die obigen Bedingungen schließen insbesondere Konfigurationen aus, in denen
eine Ecke eines Simplexes im \emph{Inneren} der Kante eines anderen Simplexes
liegt. Kanten treffen sich also nur vollständig (als ganze Kanten) oder
punktweise in gemeinsamen Ecken.
\end{REM}



Satz 3.27 

Jede kompakte Riemannsche Fläche besitzt eine Triangulierung.

\begin{PROP}{AG-03-27}{Triangulierbarkeit kompakter Riemannscher Flächen}
Jede kompakte Riemannsche Fläche besitzt eine Triangulierung.
\end{PROP}

\begin{proof}[Idee des Beweises]
Da die Fläche kompakt ist, besitzt sie eine endliche atlasartige Überdeckung durch holomorphe Karten. In jeder Karte können kleine Dreiecke in der komplexen Ebene gewählt werden, deren Bilder unter den Kartenabbildungen
Simplexes in der Fläche liefern. Durch geeignete Verfeinerung und Anpassung der Kanten über den Überlappungsbereichen erhält man eine endliche Familie solcher Simplexes, deren Vereinigung die gesamte Fläche überdeckt und die die Bedingungen einer Triangulierung erfüllen.
\end{proof}

>> Ausführlicherer Beweis: Geometric Topology in Dimensions 2 and 3



Beispiel 3.28 

Das folgende Bild erhält ein Beispiel für eine Triangulierung des Torus mit 18 Dreiecken.
\begin{figure}[h]
\begin{center}
  \includegraphics[scale=0.2]{2025_10_06_374d4cfa8b1948eb4c41g-028}
 
\caption{Abbildung 3.5: Triangulierung des Torus}
\end{center}
\end{figure}

Beweis: 

Wegen der vorausgesetzten Kompaktheit können wir $X$ mit endlich vielen Kreisscheiben überdecken und annehmen, dass keine dieser Scheiben ganz in einer anderen enthalten ist, d.h. $X=\bigcup_{i=1}^{n} D_{i}$ mit $D_{i}=\Delta_{\varepsilon_{i}}\left(a_{i}\right)$ für geeignete Punkte $a_{i} \in X$ und Radien $\varepsilon_{i} \in \mathbb{R}$ und $D_{i} \nsubseteq D_{j}$ für $i \neq j$. Ein topologisches Lemma (z.B. das Lemma auf S. 20 in [Rey89], Induktion nach $n$ und geeignetes leichtes Vergrößern der Kreisscheiben) erlaubt anzunehmen, dass $\bigcup_{i=1}^{n} \partial \triangle_{\varepsilon_{i}}\left(a_{i}\right)$ eine endliche Vereinigung von Bögen $A$ zwischen der Menge von Schnittpunkten

$$
S=\bigcup_{i<j} \partial \Delta_{\varepsilon_{i}}\left(a_{i}\right) \cap \partial \Delta_{\varepsilon_{j}}\left(a_{j}\right)
$$

ist. Falls es kein Paar $(i, j)$ gibt mit $\partial D_{i} \subseteq D_{j}$, so enthält der Rand eines jeden $D_{i}$ einen Punkt in $S$. Also ist $X$ durch Polygone mit Kanten in $A$ und Eckpunkten in $S$ zerlegt, welche durch Einfügen weiterer Kanten zu einer Triangulierung verfeinert werden kann. Andernfalls ist $X \cong S^{2}$, hat also auch eine Triangulierung (z.B. mit Kanten auf dem Äquator und in den Längenkreisen).

Triangulierungen kann man auch für topologische Flächen (mit Karten in $\mathbb{R}^{2}$ und Homomorphismen als Kartenwechselabbildungen) definieren. Triangulierbarkeit für Riemannsche Flächen ist auch ohne die Kompaktheitsvoraussetzung richtig, aber hierzu benötigt man die Methoden aus Abschnitt 7 . Im allgemeineren topologischen Kontext benötigt man eine topologische Gutartigkeitsvoraussetzung (Abzählbarkeit).

\begin{REM}{AG-03-28}{Triangulierbarkeit in weiteren Kategorien}
Triangulierungen lassen sich nicht nur in der Kategorie der
Riemannschen Flächen formulieren: Auch \emph{topologische Flächen},
also zweidimensionale Mannigfaltigkeiten mit Karten in $\mathbb{R}^2$
und Homöomorphismen als Kartenwechselabbildungen, können
trianguliert werden.

\medskip

Für Riemannsche Flächen gilt die Triangulierbarkeit sogar \emph{ohne}
Kompaktheitsannahme; der Beweis erfordert jedoch die analytischen
Methoden von Abschnitt~7 (insbesondere Normalformen holomorpher
Abbildungen und Kontrollargumente für lokale Gitterstrukturen).

\medskip

Im rein topologischen Rahmen ist die Situation subtiler:
Eine allgemeine topologische Fläche ist genau dann triangulierbar,
wenn sie eine \emph{topologische Abzählbarkeitsbedingung} erfüllt
(etwa „zweite Abzählbarkeit“). Diese Voraussetzung garantiert, daß
sich die lokale Struktur der Fläche durch eine zählbare Menge von
Karten kontrollieren lässt, was die Konstruktion kompatibler Simplexe
erst ermöglicht.

\medskip

Kurz gesagt:
\begin{itemize}
  \item Triangulierbarkeit kompakter Riemannscher Flächen ist einfach,
        weil Kompaktheit eine endliche Kontrolle ermöglicht.
  \item Nicht-kompakte Riemannsche Flächen sind ebenfalls triangulierbar,
        erfordern aber analytische Werkzeuge.
  \item Topologische Flächen benötigen für Triangulierbarkeit
        eine zusätzliche Gutartigkeitsvoraussetzung
        (etwa zweite Abzählbarkeit).
\end{itemize}
\end{REM}

Die Umkehrung des Satzes gilt auch. Dazu erinneren wir, dass eine topologische Fläche mit Karten in $\mathbb{R}^{2}$ orientierbar genannt wird, wenn die Kartenwechselabbildungen die Orientierung von $\mathbb{R}^{2}$ respektieren. Holomorphe Kartenwechselabbildungen haben diese Eigenschaft, weswegen Riemannsche Flächen stets orientierbar sind.


\begin{PROP}{AG-03-29}{Orientierbarkeit als Charakterisierung der Triangulierbarkeit}
Die Umkehrung des vorangegangenen Satzes gilt ebenfalls:  
Eine topologische Fläche ist triangulierbar genau dann, wenn sie
orientierbar ist.

\medskip

\noindent
Erinnerung:  
Eine topologische Fläche mit Karten $(U_i,\phi_i)$ in $\mathbb{R}^2$
heißt \emph{orientierbar}, wenn alle Kartenwechselabbildungen
\[
\phi_j\circ \phi_i^{-1} : \phi_i(U_i\cap U_j)\longrightarrow \phi_j(U_i\cap U_j)
\]
die Standardorientierung von $\mathbb{R}^2$ respektieren.  
Holomorphe Kartenwechsel sind stets orientierungserhaltend, sodass
jede Riemannsche Fläche automatisch eine orientierbare
topologische Fläche ist.

\medskip

\noindent
Damit gilt:
\begin{itemize}
  \item Jede Riemannsche Fläche ist orientierbar und somit triangulierbar.
  \item Umgekehrt besitzt jede orientierbare topologische Fläche eine
        kompatible Triangulierung.
\end{itemize}

\medskip

\noindent
Die Orientierbarkeit ist also die präzise topologische Bedingung, die
Triangulierungen in zwei Dimensionen ermöglicht.
\end{PROP}

Proposition 3.29 

Jede orientierbare triangulierbare topologische Fläche besitzt einen komplexen Atlas.


Beweis: Siehe z.B. Théorème 2.3 in Rey89.


\begin{PROP}{RF-03-29}{Komplexe Struktur auf orientierbaren triangulierbaren Flächen}
Jede orientierbare triangulierbare topologische Fläche besitzt einen komplexen Atlas, d.\,h.\ sie trägt eine Struktur einer Riemannschen Fläche.
\end{PROP}

\begin{proof}
Der Beweis findet sich z.\,B.\ in \textsc{Reynaud}, Théorème~2.3 in~\cite{Rey89}.
\end{proof}

Wir können nun das Hauptresultat des Abschnitts formulieren.


>> Kanonische Zerschneidung: Jede kompakte Riemannsche Fläche (nicht homöomorph zu $S^2$) besitzt Triangulierung mit kanonischer Zerschneidung (mit invariante $g$) :: Das ist 3.30



Satz 3.30 

Jede kompakte Riemannsche Fläche $X$ ist homöomorph zu $S^{2}$ oder (als topologischer Raum) zur Verklebung wie in Abbildung 3.6 eines Polygons mit $2 g$ Kanten, wobei Kanten mit gleichem Symbol (bis auf Potenz) unter Beachtung der Pfeilrichtung identifiziert werden.

\begin{THEO}{RF-03-30}{Topologische Klassifikation kompakter Riemannscher Flächen}
Jede kompakte Riemannsche Fläche $X$ ist (als topologischer Raum) homöomorph zu genau einem der folgenden Räume:
\begin{enumerate}
  \item der $2$-Sphäre $S^{2}$ (entspricht dem Fall $g=0$);
  \item der durch Verklebung eines $4g$-Polygons gemäß dem kanonischen Zerschneidung (in Abbildung~3.6 dargestellten Identifikationsschema) entstehenden Fläche, wobei jeweils die Kantenpaare mit gleichem Symbol (bis auf Potenz) unter Beachtung der Orientierungen miteinander identifiziert werden.
\end{enumerate}
Damit ist jede kompakte Riemannsche Fläche durch genau ein Geschlecht $g\ge 0$
klassifiziert, und topologisch homöomorph zu einer Fläche vom Geschlecht $g$.

Man beachte, dass bei der Zerschneidung wie in Abbildung 3.6 die Anfangsund Endpunkte aller Kanten zu einem einzigen Punkt in der Riemannschen Fläche $X$ identifiziert werden.
\end{THEO}


Beweis: 

Man pflastere einen zusammenhängenden und einfach zusammenhängenden Teil der Ebene mit den Dreiecken der Triangulierung. Man nummeriere Kanten am Rand des Polygons, die auf $X$ identifiziert waren mit den Kanten-Symbolen $a_{i}$ und $a_{i}^{-1}, i \in I_{K}$. Man identifiziere weitere Ecken, die dem gleichen Punkt in $X$ entsprechen, mit demselben Ecken-Symbol $P_{i}$ mit $i \in I_{E}$. Dieses Polygon wird mit den folgenden Modifikationen in Normalform gebracht.
\begin{enumerate}
  \item Falls es nur zwei Seiten $a$ und $a^{-1}$ gibt, so ist $X \sim S^{2}$.
  \item Falls es mehr als zwei Seiten gibt, sorgt man zunächst dafür, dass nur ein Ecksymbol, sagen wir $p$ auftritt. Denn gibt es noch ein weiteres Symbol $q$, so sorgt das Abtrennen und Neuverkleben in Figur (untenstehend) dafür, dass die Anzahl der Ecken mit Symbol $q$ um eins erniedrigt und die mit Symbol $p$ um eins erhöht wird.
\end{enumerate}
\begin{figure}[h]
\begin{center}
  \includegraphics[scale=0.2]{2025_10_06_374d4cfa8b1948eb4c41g-030}
 
\caption{Abbildung 3.6: Kanonische Zerschneidung}
\end{center}
\end{figure}
\begin{figure}[h]
\begin{center}
  \includegraphics[scale=0.2]{2025_10_06_374d4cfa8b1948eb4c41g-030(1)}
\caption{Abbildung 3.7: Schritt 2 von Satz 3.30}
\end{center}
\end{figure}
\begin{enumerate}
  \setcounter{enumi}{2}
  \item Benachbarte $a$ und $a^{-1}$ werden verklebt. (Die Aussage des Satzes ist topologisch, die Einbettung der Figur in die Ebene ist irrelevant.)
  \item Zu jeder Kante $a$ gibt es eine Kante $b$, sodass $a$ und $b$ verschachtelt sind, d.h. dass die zugehörigen Symbole in der zyklischen Reihenfolge $\ldots a \ldots b \ldots a^{-1} \ldots b^{-1} \ldots$ beim Umlauf der Kantensymbole am Rand auftreten. Denn andernfalls könnte man das Polygon wie in Abbildung 3.8 längs einer Kante $c$ von $a$ nach $a^{-1}$ in $P_{1}$ und $P_{2}$ zerschneiden und entsprechend der Seitenidentifikationen verkleben. Man erhält so eine berandete Riemannsche Fläche $X_{0}$ mit Rand $c$, sodass $X$ durch Verkleben entlang $c$ entsteht. Sei $V$ eine beliebige Umgebung von $P$ in $X$. In der Riemannschen Fläche $X$ zerfällt $V \backslash\{P\}$ in zwei Zusammenhangskomponenten. Dies ist ein Widerspruch und folglich kann die Riemannsche Fläche nicht die angenommene Gestalt haben.
  \item Zu jedem Paar verschachtelter Kanten kann man durch die Wiederverklebung wie skizziert in Abbildung 3.9 ein paar aufeinanderfolgender verschachtelter Kanten herstellen.
\end{enumerate}
\begin{figure}[h]
\begin{center}
  \includegraphics[scale=0.2]{2025_10_06_374d4cfa8b1948eb4c41g-031}
\caption{Abbildung 3.8: Schritt 4 von Satz 3.30}
\end{center}
\end{figure}
\begin{figure}[h]
\begin{center}
  \includegraphics[scale=0.2]{2025_10_06_374d4cfa8b1948eb4c41g-031(1)}
\caption{Abbildung 3.9: Schritt 5 von Satz 3.30}
\end{center}
\end{figure}
\begin{enumerate}
  \setcounter{enumi}{5}
  \item Die gestrichelten Abschnitte bleiben in Schritt 5 unberührt. Man kann also die Schritte 3 und 5 wiederholen, bis alle verschachtelten Paare aufeinanderfolgen.
\end{enumerate}



Bemerkung 3.31 Beim Wiederverkleben bilden Paare aufeinanderfolgender verschachtelter Kanten eine „Henkelzerschneidung" der Riemannschen Fläche, wie in Abbildung 3.10 für $g=2$ dargestellt.


\pagebreak

\subsection {Fundamentalgruppenbestimmung II}
Die Anzahl der Kanten der kanonischen Zerschneidung einer kompakten Riemannschen Fläche ist die erste fundamentale Invariante dieser Objekte.

Satz 3.32 Die kompakte Riemannsche Fläche $X$ habe eine Zerschneidung in ein $4 g$ gon wie in Satz 3.30 mit $g \geq 1$. Dann ist die Fundamentalgruppe von $X$ präsentiert durch

$$
\pi_{1}(X) \cong\left(\left\{a_{1}, \ldots, a_{g}, b_{1}, \ldots, b_{g}\right\} \mid \prod_{i=1}^{g}\left[a_{i}, b_{i}\right] .\right)
$$

\begin{figure}[h]
\begin{center}
  \includegraphics[scale=0.2]{2025_10_06_374d4cfa8b1948eb4c41g-032}
 
\caption{Abbildung 3.10: Henkelzerschneidung}
\end{center}
\end{figure}

Korollar 3.33 In diesem Fall ist $H_{1}(X, \mathbb{Z})=\pi_{1}(X)^{\mathrm{ab}} \cong \mathbb{Z}^{2 g}$. Insbesondere ist $g$ eine Invariante der kompakten Riemannschen Fläche $X$.

Definition 3.34 (Geschlecht) Diese Invariante $g=\frac{1}{2} \operatorname{Rang}\left(H_{1}(X, \mathbb{Z})\right)$ wird Geschlecht der Riemannschen Fläche genannt. Wir definieren $g\left(S^{2}\right)=0$.

Beweis des Satzes: Sei $q$ ein Punkt im Inneren des Polygons, $U=X \backslash\{q\}$ und $V$ sei $X$ ohne den Rand des $4 g$-gons. Dann ist $X=U \cup V$ und $U \cap V$ ist homöomorph zu $\Delta^{*}$, hat also nach Proposition 3.24 die Fundamentalgruppe $\pi_{1}(U \cap V) \cong \mathbb{Z}$ erzeugt von $\prod_{i=1}^{g}\left[a_{i}, b_{i}\right]$. Weiter ist $V$ homöomorph zu $\Delta$, also $\pi_{1}(V)=\{i d\}$. Zentrische Streckung definiert eine Homotopie $\pi_{1}(U)=\pi_{1}(\partial(4 g$-gon $))$, wobei der Rand des $4 g$-gons ein Bouquet $B_{2 g}$ von $2 g$ Kreisen ist.

\begin{figure}[h]
\begin{center}
  \includegraphics[scale=0.2]{2025_10_06_374d4cfa8b1948eb4c41g-032(1)}
 
\caption{Abbildung 3.11: Bouquet \$B\_\{2 g}\$\}\end{center}
\end{figure}

Nach dem Satz von Seifert/van Kampen ist

$$
\pi_{1}(X) \cong \pi_{1}\left(B_{2 g}\right) *_{\mathbb{Z} \cdot\left(\pi\left[a_{i}, b_{i}\right]\right)} \quad\{\mathrm{id}\} \cong \pi_{1}\left(B_{2 g}\right) /\left\langle\prod_{i=1}^{g}\left[a_{i}, b_{i}\right]\right\rangle .
$$

Per Induktion und wiederum mit Seifert/van Kampen folgt, dass $\pi_{1}\left(B_{2 g}\right)= F_{2 g}$ eine freie Gruppe erzeugt von $a_{1}, \ldots, a_{g}, b_{1}, \ldots, b_{g}$ ist.

Satz 3.35 Ist $X$ eine kompakte Fläche mit einer Triangulierung (oder allgemeiner: Zerlegung in zusammenhängende und einfach zusammenhängende Zellen), mit $F$ Flächen (oder Zellen), K Kanten und E Ecken, so gilt

$$
2-2 g=E-K+F
$$

Beweis: Man zeigt leicht der Reihe nach:\\
i) $E-K+F$ bleibt beim Verfeinern einer Triangulierung invariant.\\
ii) Diese Größe bleibt auch invariant, wenn man ein zusammenhängendes und einfach zusammenhängendes $n$-gon zu einer Triangulierung verfeinert.\\
iii) Die Formel ist richtig für die Triangulierung einer kanonischen Zerschneidung einer Riemannschen Fläche mit $g \geq 1$, die man erhält, in dem man die Eckpunkte und Seitenmitten mit einem Punkt im Innern des $4 g$-gons verbindet. Sie ist auch für eine Triangulierung der Sphäre richtig.\\
Aus diesen Aussagen folgt der Satz.



\pagebreak

\subsection {Galoisgruppen von Überlagerungen, normale Überlagerungen}
In diesem Abschnitt untersuchen wir, welche Überlagerungen vollständig durch Gruppentheorie beschrieben sind. Der gesamte Abschnitt ist vollkommen analog zur Galoistheorie von Körpererweiterungen.

Die Standardbegriffe der Gruppentheorie werden wir auch für Automorphismen von Riemannschen Flächen verwenden. Für eine Untergruppe $G \leq \operatorname{Aut}(X)$ heißt $G \cdot x=\{g(x), g \in G\}$ die Bahn von $x$ unter der Aktion von $G$. Die Fixgruppe eines Punktes $x \in X$ ist definiert als $\operatorname{Fix}(x):=\{g \in G: g(x)=x\}$.

Beispielsweise sei $X=\mathbb{C}$. Für alle $\lambda \in \mathbb{C}^{*}$ ist $m_{\lambda}: z \longmapsto \lambda \cdot z$ ein Automorphismus von $X$. Für alle $\mu \in \mathbb{C}$ ist ebenso $t_{\mu}: z \longmapsto z+\mu$ ein Automorphismus von $X$.

Als weiteres Beispiel erinnern wir daran, dass für alle $M=\left(\begin{array}{ll}a & b \\ c & d\end{array}\right) \in \mathrm{SL}_{2}(\mathbb{R})$ die Abbildung $z \longmapsto M \cdot z:=\frac{a z+b}{c z+d}$ ein Automorphismus der oberen Halbebene $\mathbb{H}$ ist. Ein solcher Automorphismus wird Möbiustransformation genannt.

Definitionen 3.36 (Galoisgruppe, normale Überlagerung) Sei $f: X \longrightarrow Y$ eine (nicht notwendigerweise unverzweigte) Überlagerung zwischen zusammenhängenden Riemannschen Flächen.\\
i) Die Gruppe Gal(f):=Gal(X/Y):=\{g $\in \operatorname{Aut}(X): f \circ g=f\}$ heißt Galoisgruppe von $f$ (oder auch Deckgruppe).\\
ii) Ist $f$ unverzweigt, so heißt $f$ normal (oder Galoissch), wenn Gal( $f$ ) für alle $y \in Y$ transitiv auf der Faser $f^{-1}(y)$ operiert, d.h. wenn für alle $x_{1}, x_{2} \in f^{-1}(y)$ ein $g \in \operatorname{Gal}(f)$ existiert mit $g\left(x_{1}\right)=x_{2}$.\\
iii) Sei $X$ zusammenhängend und $A \subset X$ die Menge der Verzweigungspunkte von $f$. Dann heißt allgemeiner $f$ normal (oder Galoissch), wenn die unverzweigte Überlagerung $f: X \backslash f^{-1}(f(A)) \rightarrow Y \backslash f(A)$ galoissch ist.

Obige Definition von Galois verlangt Transitivität in allen Fasern, aber Transitivität in einer Faser genügt schon, wie das folgende Lemma besagt.

Lemma $3.37 f: X \rightarrow Y$ ist normal, falls es einen Punkt $y \in Y \backslash f(A)$ gibt, so dass es zu jedem Paar $x_{i}, x_{j} \in f^{-1}(y)$ ein Element $g$ von $\operatorname{Gal}(f)$ gibt mit $g\left(x_{i}\right)=x_{j}$.

Beweis: Sei $\widetilde{y} \in Y$ beliebig. Verbinde $y$ und $\widetilde{y}$ mit einem Weg $\gamma$. Seien $\widetilde{x}_{i}, \widetilde{x}_{j} \in f^{-1}(\widetilde{y})$ und $\gamma_{i}, \gamma_{j}$ die Liftungen von $\gamma$ nach $X$ beginnend bei $\widetilde{x}_{i}$ bzw. $\widetilde{x}_{j}$ und endend bei $x_{i}$ bzw. $x_{j}$. Dann leistet $g \in \operatorname{Gal}(f)$ mit $g\left(x_{i}\right)=x_{j}$ das Verlangte. \(\square\)

Beispiele 3.38 Wir geben ein Beispiel für eine normale Überlagerung und eine Überlagerung ohne diese Eigenschaft an. Man beachte, dass wir nicht gefordert haben, dass der Grad der Überlagerung endlich ist.\\
i) Betrachte die Exponentialabbildung $\exp (\cdot): X=\mathbb{C} \longrightarrow Y=\mathbb{C}^{*}$. Dann ist

$$
\operatorname{Gal}(X / Y)=\{z \longmapsto z+2 \pi i k, k \in \mathbb{Z}\}
$$

und deswegen ist $\exp$ eine normale Überlagerung.\\
ii) Sei $X^{*}$ die Riemannsche Fläche, die aus $X_{1}$ durch Verkleben der Paare von Seiten mit gleichem Symbol (analog zu Beispiel 2.8 entsteht. Wir werden diese

\begin{figure}[h]
\begin{center}
  \includegraphics[scale=0.2]{2025_10_06_374d4cfa8b1948eb4c41g-034}
 
\caption{Abbildung 3.12: Beispiel für eine nicht-galoissche Überlagerung}
\end{center}
\end{figure}

Konstruktion nochmal formal in Abschnitt Baufgreifen. Es gibt eine offensichtliche Abbildung $f: X^{*} \longrightarrow T_{\{1, i\}}^{*}$ auf den punktierten Torus, welche das Innere jedes Quadrats bijektiv auf das Innere des Quadrat abbildet, welches zur Konstruktion von $T^{*}$ verwendet wurde. Dies ist eine Überlagerung vom Grad 4. Es ist aber $\operatorname{Gal}\left(X^{*} / T_{\{1, i\}}^{*}\right) \cong \mathbb{Z} /{ }_{2 \mathbb{Z}}$, denn wenn man das Quadrat links oben auf eines der vier Quadrate schickt, lässt sich dies nur in zwei der vier Fällen zu einem Biholomorphismus von X* fortsetzen. Somit ist diese Überlagerung nicht Galoissch.

Sei ab sofort $\zeta_{n}$ eine primitive $n$-te Einheitswurzel, d.h. es gilt $\zeta_{n}^{n}=1$ und $\zeta_{n}^{k} \neq 1$ für $1 \leq k<n$.

Lemma 3.39 Sei $f: \triangle_{\varepsilon} \rightarrow \triangle_{\varepsilon^{n}}$ die Überlagerung $z \mapsto z^{n}$. Dann ist $\operatorname{Gal}(f) \cong \mathbb{Z} / n \mathbb{Z}$.

Beweis: Ist $g \in \operatorname{Gal}(f)$, so ist $g(z)^{n}=f(g(z))=z^{n}$, d.h. $g(z)=\zeta_{n}^{i} \cdot z$, wobei der Exponent $i$ a priori vom Punkt $z$ abhängt. Da aber $g$ stetig ist, hängt $i$ de facto nicht von $z$ ab. Umgekehrt ist jede Abbildung $z \mapsto \zeta_{n}^{i} \cdot z$ für $i \in\{1, \ldots, n\}$ in $\operatorname{Gal}(f)$.

Ist $f: X \rightarrow Y$ eine Überlagerung vom Grad $d$ und $y \in Y$ ein Punkt mit $d$ Urbildern $\left\{x_{1}, \ldots, x_{d}\right\}$. Dann definiert man eine Abbildung

$$
\rho(\cdot):\left\{\begin{array}{ccc}
\operatorname{Gal}(f) & \longrightarrow & \operatorname{Perm}\left(\left\{x_{1}, \ldots, x_{d}\right\}\right) \cong S_{d} \\
g & \longmapsto & \left(x_{i} \mapsto g\left(x_{i}\right), i=1, \ldots, d\right)
\end{array}\right.
$$

die Permutationsdarstellung der Galoisgruppe genannt wird.\\
Beispiel 3.40 Die Abbildung $\exp (\cdot)$ ist eine normale Überlagerung, die Abbildung $f$ aus Lemma 3.39 ebenfalls. Die Überlagerung $X^{*} \rightarrow T_{\{1, i\}}^{*}$ aus Beispiel 3.38 ist nicht normal.

Proposition 3.41 Sei $f$ normal und $f^{-1}(y)=\left\{x_{1}, \ldots, x_{k}\right\}$ die Faser von $f$ über einem beliebigen Punkt $y \in Y$. Dann ist $e_{x_{1}}(f)=\ldots=e_{x_{k}}(f)$.

Beweis: Sei $V$ eine Umgebung von $y$ wie in der Definition einer Überlagerung und $f^{-1}(V)=U_{1} \cup \ldots \cup U_{k}$ mit $x_{i} \in U_{i}$. Aus Normalität folgt die Existenz von $g_{i} \in \operatorname{Gal}(f)$ mit $g_{i}\left(x_{i}\right)=x_{1}$ für $i=2, \ldots, k$. Wir erhalten

$$
e_{x_{1}}(f)=e_{x_{1}}\left(f \circ g_{i}\right)=e_{x_{i}}(f)
$$

da $g_{i}$ ein Biholomorphismus ist.

Satz 3.42 Sei $f: X \rightarrow Y$ eine unverzweigte Überlagerung mit $f(a)=b$. Ist $\widetilde{a} \in f^{-1}(b)$ ein weiterer Punkt im Urbild, so ist $f_{*}\left(\pi_{1}(X, \widetilde{a})\right)$ zu $f_{*}\left(\pi_{1}(X, a)\right)$ in $\pi_{1}(Y, b)$ konjugiert. Mit einem solchem Basispunktwechsel erhält man alle zu $f_{*}\left(\pi_{1}(X, \widetilde{a})\right)$ konjugierten Untergruppen von $\pi_{1}(Y, b)$.

Insbesondere gilt: $f$ ist genau dann normal, wenn $f_{*}\left(\pi_{1}(X, a)\right)$ ein Normalteiler in $\pi_{1}(Y, b)$ ist. In diesem Fall ist

$$
\operatorname{Gal}(f) \cong \pi_{1}(Y, b) / f_{*}\left(\pi_{1}(X, a)\right)
$$

Beweis: Sei $\gamma$ ein Weg von a nach $\widetilde{a}$. Dann ist $\bar{\gamma}=f \circ \gamma$ ein geschlossener Weg in $Y$. Die Abbildung $\eta \mapsto \gamma^{-1} \eta \gamma$ definiert einen Isomorphismus $\pi_{1}(X, \widetilde{a}) \rightarrow \pi_{1}(X, a)$. Folglich ist $\bar{\gamma}^{-1} f_{*}\left(\pi_{1}(X, a)\right) \bar{\gamma}=f_{*}\left(\pi_{1}(X, \widetilde{a})\right)$, wie behauptet. Umgekehrt sei die konjugierte Gruppe $\bar{\gamma}^{-1} f_{*}\left(\pi_{1}, a\right) \bar{\gamma}$ gegeben. Lifte $\bar{\gamma}$\\
zu einem Weg in $X$ beginnend mit $a$. Diese endet in einem Punkt $\widetilde{a}$. Dann ist $\bar{\gamma}^{-1} f_{*}\left(\pi_{1}(X, a)\right) \bar{\gamma}=f_{*} \pi_{1}(X, \widetilde{a})$ und wir haben die konjugierte Gruppe als Fundamentalgruppe basiert bei $\widetilde{a}$ geschrieben.

Ist $f_{*}\left(\pi_{1}(X, a)\right)$ ein Normalteiler, so können wir das Liftungslemma 3.21 auf $f$, auf $g=f$ und zwei beliebige Punkte $a, \widetilde{a} \in f^{-1}(b)$ anwenden. Der so erhaltene Lift ist das gesuchte Element von $\operatorname{Gal}(f)$. Damit erhält man auch die Abbildung von rechts nach links im gesuchten Isomorphismus. Die Umkehrabbildung haben wir bereits oben kennengelernt, denn Nebenklassen $\pi_{1}(Y, b) / f_{*}\left(\pi_{1}(X, a)\right)$ entsprechen den Punkten im Urbild von $b$ und $\operatorname{Gal}(f)$ operiert frei auf diesen.



\pagebreak

\subsection {Unverzweigte Überlagerungen zu gegebenen Untergruppen}
Satz 3.43 Zu jeder Untergruppe $H<\pi_{1}(Y, b)$ gibt es (bis auf Isomorphie) genau eine zusammenhängende unverzweigte Überlagerung $f: X \rightarrow Y$ mit gewähltem Basispunkt $a \in X$, sodass $f(a)=b$ und $f_{*}\left(\pi_{1}(X, a)\right)=H$. Weiterhin ist $\left[\pi_{1}(Y, b):\right. H]=\operatorname{deg}(f)$.

Beweis: Die Eindeutigkeit bis auf Isomorphie haben wir bereits in Satz 3.20 gezeigt. Wir nennen zwei Wege $\gamma_{1}, \gamma_{2}$, die in $b$ beginnen, $H$-äquivalent, falls $\gamma_{1}^{-1} \cdot \gamma_{2} \in H$.

Sei $X$ die Menge der $H$-Äquivalenzklassen von Wegen in $Y$ und $f: X \rightarrow Y$ die Abbildung, die jedem Weg den Endpunkt zuordnet. Sei $x=[\gamma]_{H} \in X$ ein Weg von $b$ nach $y$ und $U_{y}$ eine einfach zusammenhängende Umgebung von $y$. Wir definieren

$$
U_{x}=\left\{\left[\gamma_{1} \cdot \gamma\right]_{H}, \gamma_{1} \text { Weg in } U_{y} \text { beginnend bei } y\right\}
$$

Dann ist $\left.f\right|_{U_{x}}: U_{x} \rightarrow U_{y}$ eine Bijektion, da $U_{y}$ zusammenhängend und einfach zusammenhängend ist. Mit Bijektionen dieser Art definieren wir auf $X$ einen Atlas und damit eine komplexe Struktur.

Es bleibt noch zu zeigen, dass $\left|f^{-1}(b)\right|=\left[\pi_{1}(Y, b): H\right]$ ist. Da nach Konstruktion $f^{-1}(b)$ die Menge der $H$-Äquivalenzklassen von Wegen von $b$ nach $b$ ist, folgt dies sofort.

Angewandt auf die triviale Gruppe zeigt dieser Satz zusammen mit Korollar 3.23 sofort:

Korollar 3.44 Jede Riemannsche Fläche besitzt eine universelle Überlagerung.\\
Korollar 3.45 Zu jedem surjektiven Gruppenhomomorphismus $\pi_{1}(Y, a) \rightarrow G$ gibt es bis auf Isomorphie genau eine normale unverzweigte Überlagerung $f: X \rightarrow Y$ mit Deckgruppe $\operatorname{Gal}(f)=G$.

Der Satz gibt auch ein Verfahren an, um Überlagerungen nach Wunsch herzustellen. Bevor wir dies als Verfahren formulieren, analysieren wir noch einmal eine unverzweigte Überlagerung und stellen fest, dass wir in jedem Fall einen Homomorphismus erhalten.

Definition 3.46 (Monodromie) Die Monodromiedarstellung einer unverzweigten Überlagerung $f: X \rightarrow Y$ ist der Homomorphismus

Mon : $\left\{\begin{array}{cll}\pi_{1}(Y, b) & \longrightarrow & \operatorname{Perm}\left(f^{-1}(b)\right) \cong S_{\operatorname{deg}(f)} \\ \gamma & \longmapsto & (a \mapsto \text { Endpunkt des Lifts von } \gamma \text { nach } X \text { beginnend in a })\end{array}\right.$\\
Das Bild dieses Homomorphismus wird Monodromiegruppe der Überlagerung genannt.

Ist $f$ normal, so faktorisiert Mon über die Untergruppe $\pi_{1}(X, a)$ und der Quotientenhomomorphismus ist gerade die Permutationsdarstellung der Galoisgruppen.



\pagebreak

\subsection {Der Satz von Riemann-Hurwitz}
Die Geschlechter der Riemannschen Fläche in einer Überlagerung erfüllen folgende fundamentale Beziehung.

Satz 3.47 (Riemann-Hurwitz-Formel) Sei $f: X \rightarrow Y$ eine nicht-konstante holomorphe Abbildung zwischen kompakten, zusammenhängenden Riemannschen Flächen. Dann gilt

$$
2 g(X)-2=\operatorname{deg}(f) \cdot(2 g(Y)-2)+\sum_{x \in X}\left(e_{x}(f)-1\right)
$$

Beweis: Wir betrachten nochmal den Beweis der Triangulierbarkeit von $Y$. Wir können annehmen, dass jedes Dreieck in einer Standardkreisscheibe (der Definition von Überlagerung) liegt und - nach Verfeinerung - dass das Bild jedes Verzweigungspunkts eine Ecke der Triangulierung ist. Aufgrund dieser Bedingung definieren die Urbilder der Dreiecksflächen und Kanten eine Triangulierung von $X$.\\
Seien $E_{X}, E_{Y}, K_{X}, K_{Y}, F_{X}, F_{Y}$ die Ecken, Kanten und Flächen auf $X$ und $Y$. Dann gilt $\left|K_{X}\right|=\left|K_{Y}\right| \cdot \operatorname{deg}(f)$ und $\left|F_{X}\right|=\left|F_{Y}\right| \cdot \operatorname{deg}(f)$ und es gilt

$$
\left|E_{X}\right|=\left|E_{Y}\right| \cdot \operatorname{deg}(f)+\sum_{x \in X}\left(1-e_{x}(f)\right)
$$

Zusammen mit $2-2 g=E-K+F$ aus Satz 3.35 folgt die Behauptung.



\pagebreak


\section {Differentialformen}
Ziel der nächten Abschnitte ist es, eine holomorphe Funktion $f: X \longrightarrow \mathbb{P}_{\mathbb{C}}^{1}$ für eine kompakte Riemannsche Fläche $X$ zu konstruieren. Dazu müssen wir in der Lage sein, $f$ abzuleiten und wir untersuchen hier, welches Objekt wir auf einer Riemannschen Fläche erhalten. Ableiten wollen und können wir nur kartenweise. Sei dazu $U \subset X$ genügend klein sodass $f(U) \subset \mathbb{C} \subset \mathbb{P}^{1}$ und sei $g: U \longrightarrow \mathbb{C}$ eine Karte von $X$. Dann kennen wir die Definition von $\frac{\partial\left(f \circ g^{-1}\right)}{\partial x}$ und $\frac{\partial\left(f \circ g^{-1}\right)}{\partial y}$. Wir bezeichnen nun die Kartenabbildung nicht mit $g$, sondern suggestiver mit $z: U \longrightarrow \mathbb{C} \cong \mathbb{R}^{2}$ und nennen die Koordinatenabbildungen $x$ und $y$, also $z=x+i y$. Sei $z_{1}=x_{1}+i y_{1}$ eine weitere Karte von $U$. Dann ist nach der Kettenregel

$$
\begin{aligned}
& \frac{\partial\left(f \circ z_{1}^{-1}\right)}{\partial x_{1}}=\frac{\partial\left(x \circ z_{1}^{-1}\right)}{\partial x_{1}} \cdot \frac{\partial\left(f \circ z^{-1}\right)}{\partial x}+\frac{\partial\left(y \circ z_{1}^{-1}\right)}{\partial x_{1}} \cdot \frac{\partial\left(f \circ z^{-1}\right)}{\partial y} \\
& \frac{\partial\left(f \circ z_{1}^{-1}\right)}{\partial y_{1}}=\frac{\partial\left(x \circ z_{1}^{-1}\right)}{\partial y_{1}} \cdot \frac{\partial\left(f \circ z^{-1}\right)}{\partial x}+\frac{\partial\left(y \circ z_{1}^{-1}\right)}{\partial y_{1}} \cdot \frac{\partial\left(f \circ z^{-1}\right)}{\partial y}
\end{aligned}
$$

Auch ohne die Existenz einer Abbildung $f$ kann man die Objekte mit dieser Kartenwechseleigenschaft definieren. Dies führt uns zur Definition von Differentialformen.

\subsection {Meromorphe Funktionen}
Wir wiederholen hier den Meromorphiebegriff aus der Funktionentheorie und setzen ihn auf Riemannsche Flächen fort. Ziel ist der Vergleich von meromorphen Funktionen mit Abbildungen auf den $\mathbb{P}_{\mathbb{C}}^{1}$.

Sei $U \subset \mathbb{C}$ ein Gebiet. Eine Funktion $f$ heißt meromorphe Funktion auf $U$, falls es eine diskrete Teilmenge $A \subset U$ gibt, sodass $f: U \backslash A \rightarrow \mathbb{C}$ holomorph ist und für jeden Punkt $a \in A$ und jede Folge $a_{n} \rightarrow a$ gilt, dass $\lim _{n \rightarrow \infty}\left|f\left(a_{n}\right)\right|=\infty$. Die Menge $A$ wird Polstellenmenge von $f$ genannt. Man zeigt (siehe z.B. [MF, Abschnitt 3.2]), dass $f$ in der Nähe eines Punktes $a \in A$ eine Entwicklung

$$
f(z)=\sum_{n \geq-k} c_{n}(z-a)^{n}, \quad k \in \mathbb{N}, \quad c_{-k} \neq 0
$$

besitzt, welche Laurent-Reihe genannt wird. Die Zahl $k$ wird in dem Fall Polstellenordnung von $f$ bei $a$ genannt. Genauso definieren wir zu jedem $p \in U$ mit $f(p)=0$ die Nullstellenordnung von $f$ bei $p$, als $e_{p}(f)$. Die Menge der meromorphen Funktionen auf $U$ bildet einen Körper.

Definition 4.1 (Meromorphe Funktion) Eine meromorphe Funktion auf einer Riemannschen Fläche $X$ ist eine Abbildung $f: X \backslash A \rightarrow \mathbb{C}$, wobei $A \subset X$ diskret ist, sodass für jede Karte $g: U \rightarrow \mathbb{C}$ von $X$ die Abbildung $f \circ g^{-1}$ eine meromorphe Funktion auf $g(U)$ ist.

Da Kartenwechselabbildungen biholomorph sind, hängt die Menge der Polstellen und die Polordnung nicht von der gewählten Karte ab.

Definition 4.2 (Null- und Polstellenordnung) Zu einer meromorphen Funktion $f$ auf einer Riemannschen Fläche $X$ sei $A \subset X$ die Menge der Polstellen und $Z= f^{-1}(0) \subset X$ die Menge der Nullstellen von $f$. Dann ist die Ordnung von $f$ bei $p$ definiert als

$$
\operatorname{ord}_{p}(f)= \begin{cases}0, & \text { falls } p \notin Z \cup A \\ k, & \text { falls } p \in Z \text { ist } \\ -k, & \text { falls } p \in A \text { ist }\end{cases}
$$

wobei $k$ die Null- bzw. Polstellenordnung bei p ist.\\
Der Begriff mermorphe Funktion ist nützlich, denn Koeffizientenfunktionen z.B. in der Definition einer Differentialform sind oftmals von dieser Art. Andererseits können wir nun die gewünschte Beziehung zu Abbildungen formulieren.

Proposition 4.3 Ist $X$ eine Riemannsche Fläche und $f$ eine meromorphe Funktion auf $X$, so erhält man mit der Definition $f(a)=\infty$ für alle Polstellen von $f$ eine holomorphe Abbildung $f: X \rightarrow \mathbb{P}_{\mathbb{C}}^{1}$. Umgekehrt ist die Einschränkung einer holomorphen Abbildung $f: X \rightarrow \mathbb{P}_{\mathbb{C}}^{1}$ auf $f^{-1}\left(\mathbb{P}_{\mathbb{C}}^{1} \backslash\{\infty\}\right)$ eine meromorphe Funktion auf $X$.

Insbesondere ist also

$$
\operatorname{ord}_{p}(f)= \begin{cases}e_{p}(f), & \text { falls } p \in f^{-1}(0) \text { ist } \\ -e_{p}(f), & \text { falls } p \in f^{-1}(\infty) \text { ist } \\ 0, & \text { sonst }\end{cases}
$$

für eine meromorphe Funktion $f$ aufgefasst als $f: X \rightarrow \mathbb{P}^{1}$.

\section {$4.2 C^{1}$ - und $C^{\infty}$-Differentialformen}

Sei $X$ eine (nicht notwendigerweise) kompakte Riemannsche Fläche. Eine reellwertige (oder komplexwertige) $C^{1}$-Differentialform (erster Ordnung) $\omega$ auf $X$ sind für jede Karte $z=x+i y: U \rightarrow \mathbb{C}$ die Angabe stetig differenzierbarer reellwertiger (oder komplexwertiger) Funktionen $\alpha, \beta$ auf $U$, sodass für die Funktionen $\alpha_{1}, \beta_{1}$ zu einer anderen Karte $z_{1}=x_{1}+i y_{1}$ auf dem gemeinsamen Definitionsbereich gilt

$$
\binom{\alpha_{1}}{\beta_{1}}=J\binom{\alpha}{\beta} \quad \text { mit } \quad J=\left(\begin{array}{cc}
\frac{\partial x}{\partial x_{1}} & \frac{\partial x}{\partial y_{1}} \\
\frac{\partial y}{\partial x_{1}} & \frac{\partial y}{\partial y_{1}}
\end{array}\right) .
$$

Wir schreiben $\omega=\alpha d x+\beta d y$. Differentialformen erster Ordnung werden auch als 1-Formen bezeichnet. Analog sind reellwertige (oder komplexwertige) $C^{\infty}$ Differentialformen durch Angaben von $C^{\infty}$-Funktionen $\alpha, \beta: U \longmapsto \mathbb{R}$ bzw. $U \longmapsto \mathbb{C}$, die der Transformation (4.2) genügen, definiert.

Wir können nun den Begriff der Ableitung modellieren. Gegeben eine Funktion $f: X \rightarrow \mathbb{R}$ oder $f: X \rightarrow \mathbb{C}$ so definieren wir (lokal, in einer Karte)

$$
d f=\frac{\partial f(z)}{\partial x} d x+\frac{\partial f(z)}{\partial y} d y
$$

wobei die Koeffizientenfunktionen reellwertig bzw. komplexwertig sind. Für die Funktionen $x: U \rightarrow \mathbb{R}$ und $y: U \rightarrow \mathbb{R}$, aus denen die Kartenabbildungen zusammengesetzt ist, ist diese Notation konsistent, d.h. $d x=d x$ und $d y=d y$. Diese Zuordnung einer Differentialform zu einer Funktion wird auch äußere Ableitung genannt.

Auch $z$ und $\bar{z}: U \longrightarrow \mathbb{C}$ sind solche Funktionen und es gilt

$$
d z=d x+i d y \quad \text { und } \quad d \bar{z}=d x-i d y
$$

Also kann man jede 1-Form $\omega$ lokal als $\omega=u d z+v d \bar{z}$ schreiben, wobei

$$
u=\frac{1}{2}(\alpha-i \beta) \quad \text { und } \quad v=\frac{1}{2}(\alpha+i \beta)
$$

Insbesondere können wir Gleichung 4.3 in dieser Basis schreiben und erhalten

$$
d f=\frac{1}{2}\left(\frac{\partial f}{\partial x}-i \frac{\partial f}{\partial y}\right) d z+\frac{1}{2}\left(\frac{\partial f}{\partial x}+i \frac{\partial f}{\partial y}\right) d \bar{z}
$$

Dies motiviert die Definition der Operatoren

$$
\frac{\partial}{\partial z}:=\frac{1}{2}\left(\frac{\partial}{\partial x}-i \frac{\partial}{\partial y}\right) \quad \text { und } \quad \frac{\partial}{\partial \bar{z}}:=\frac{1}{2}\left(\frac{\partial}{\partial x}+i \frac{\partial}{\partial y}\right)
$$

Insbesondere implizieren die Cauchy-Riemannschen Differentialgleichungen sofort, dass $\frac{\partial f}{\partial \bar{z}}=0$ ist, wenn $f$ holomorph ist.

Wir nennen eine ( $C^{\infty}$ )-Differentialform $\omega$ eine ( 1,0 )-Form bzw.eine ( 0,1 )Form, falls sie sich in allen Karten einer Überdeckung von $X$ als

$$
\omega=u d z \quad \text { bzw. als } \quad \omega=v d \bar{z}
$$

schreiben lässt. Man überprüft, dass dieser Begriff nicht von der Wahl der Karte abhängt, da die Kartenwechselabbdilungen holomorph sind, die Jacobimatrix $J$ in der Basis $z, \bar{z}$ von $\mathbb{C}$ also Diagonalgestalt hat.

Jede Differentialform ist also Summe einer ( 1,0 )-Form und einer ( 0,1 )Form. Diese ( $p, q$ )-Formen sind ein wichtiges technisches Hilfsmittel, wir wir z.B. in den Beweisen in Abschnitt 6.2 sehen werden. Unter den ( 1,0 )-Formen wollen wir holomorphe Differentialformen besonders herausstreichen.

Definition 4.4 (holomorphe Differentialformen) Eine komplexwertige Differentialform $\omega$ heißt holomorph, falls in jeder Karte $U$ von $X$ die Kartendarstellung $\omega=\alpha d z+\beta d \bar{z}$ die Eigenschaft $\beta=0$ hat und zudem $\alpha: U \longrightarrow \mathbb{C}$ holomorph ist. Den Vektorraum der holomorphen 1-Formen auf $X$ bezeichnen wir mit $\Omega^{1}(X)$.\\
Ein Punkt $x \in X$ heißt Nullstelle der Ordnung $n$ einer holomorphen Differentialform $\omega$, falls sich $\omega$ in einer Karte um $x$ als $\omega=\alpha d z$ mit $e_{x}(\alpha)=n$ schreiben lässt. Wir setzen $\operatorname{ord}_{x} \omega=n$.

In gleicher Weise definieren wir eine meromorphe 1-Form indem wir in obiger Definition von $\alpha$ lediglich 'meromorph' statt 'holomorph' verlangen.

Im nächsten Abschnitt wollen wir Differentialformen integrieren. In Vorbereitung dazu erinnern wir daran, dass eine 1 -Form $\alpha d x+\beta d y$ höchstens integrierbar ist, wenn $\alpha_{y}=\beta_{x}$ ist. Dieses formalisieren wir in dem folgenden Begriff.

Definition 4.5 (Differential zweiter Ordnung) Eine $C^{1}$-Differentialform zweiter Ordnung (oder auch kurz eine 2-Form) $\Omega$ ist eine Zuordnung einer $\mathbb{C}$-wertigen $C^{1}$-Funktion a zu jeder Karte $z$ mit folgender Eigenschaft. Ist $z_{1}$ eine weitere Karte und $\alpha_{1}$ die zugehörige Funktion der 2-Form, dann gilt $\alpha_{1}=\operatorname{det}\left(\operatorname{Jac}\left(z \circ z_{1}^{-1}\right)\right) \cdot \alpha$.

Dabei ist Jac die Jacobi-Matrix einer Funktion $\mathbb{R}^{2} \rightarrow \mathbb{R}^{2}$. Wir können jede 2-Form in einer Karte als Vielfaches von $d x \wedge d y$ schreiben, d.h.

$$
\Omega=\alpha d x \wedge d y
$$

Wir definieren nun auch die äußere Ableitung einer 1-Form $\omega=\alpha d x+\beta d y$ durch

$$
d \omega=\left(\frac{\partial \beta}{\partial x}-\frac{\partial \alpha}{\partial y}\right) d x \wedge d y
$$

Dabei ist zu prüfen, dass dies überhaupt wohldefiniert ist. Ist $\omega=\alpha_{1} d x_{1}+ \beta_{1} d y_{1}$ die Darstellung der Differentialform in einer anderen Karte. Dann ist

$$
\begin{aligned}
\frac{\partial \beta_{1}}{\partial x_{1}}-\frac{\partial \alpha_{1}}{\partial y_{1}} & =\frac{\partial \beta_{1}}{\partial x} \frac{\partial x}{\partial x_{1}}+\frac{\partial \beta_{1}}{\partial y} \frac{\partial y}{\partial x_{1}}-\frac{\partial \alpha_{1}}{\partial x} \frac{\partial x}{\partial y_{1}}-\frac{\partial \alpha_{1}}{\partial y} \frac{\partial y}{\partial y_{1}} \\
& =\frac{\partial}{\partial x}\left(\alpha \frac{\partial x}{\partial y_{1}}+\beta \frac{\partial y}{\partial y_{1}}\right) \frac{\partial x}{\partial x_{1}}+\frac{\partial}{\partial y}\left(\alpha \frac{\partial x}{\partial y_{1}}+\beta \frac{\partial y}{\partial y_{1}}\right) \frac{\partial y}{\partial x_{1}} \\
& -\frac{\partial}{\partial x}\left(\alpha \frac{\partial x}{\partial x_{1}}+\beta \frac{\partial y}{\partial x_{1}}\right) \frac{\partial x}{\partial y_{1}}-\frac{\partial}{\partial y}\left(\alpha \frac{\partial x}{\partial x_{1}}+\beta \frac{\partial y}{\partial x_{1}}\right) \frac{\partial y}{\partial y_{1}} \\
& =\left(\frac{\partial \beta}{\partial x}-\frac{\partial \alpha}{\partial y}\right) \operatorname{det}\left(\begin{array}{ll}
\frac{\partial x}{\partial x_{1}} & \frac{\partial y}{\partial x_{1}} \\
\frac{\partial x}{\partial y_{1}} & \frac{\partial y}{\partial y_{1}}
\end{array}\right)
\end{aligned}
$$

was zu zeigen war.\\
Auch die äußere Ableitung lässt sich in der Basis $d z, d \bar{z}$ schreiben, was manchmal sehr bequem ist.

Lemma 4.6 Ist $\omega=u d z+v d \bar{z}$ eine 1-Form, so ist

$$
d \omega=\left(\frac{\partial v}{\partial z}-\frac{\partial u}{\partial \bar{z}}\right) d z \wedge d \bar{z}
$$

Beweis: Direktes Nachrechnen unter Verwendung von (4.4).\\
Dieses Lemma macht einen Zweck der äusseren Ableitung offensichtlich, sie charakterisiert die holomorphen Differentialformen unter den (1, 0)-Formen.

Korollar 4.7 Eine $(1,0)$-Form $\omega$ is holomorph genau dann wenn $d(\omega)=0$ ist.\\
Sind $\omega_{i}=\alpha_{i} d x+\beta_{i} d y$ für $i=1,2$ zwei Differentialformen erster Ordnung, so ist $\omega_{1} \wedge \omega_{2}=\left(\alpha_{1} \beta_{2}-\alpha_{2} \beta_{1}\right) d x \wedge d y$ eine 2-Form, das äußere Produkt von $\omega_{1}$ und $\omega_{2}$.

Definition 4.8 (Geschlossene und exakte Formen) Eine 1-Form $\omega$ auf $X$ heißt geschlossen, falls $d \omega=0$. Eine 1-Form auf $X$ wird exakt genannt, falls es ein holomorphe Funktion $f$ auf $X$ gibt mit $\omega=d f$.

Lemma 4.9 Ist die Differentialform $\omega$ exakt, so ist $\omega$ geschlossen. Ist $\omega$ eine holomorphe 1-Form, so ist $\omega$ geschlossen.

Beweis: Die erste Aussage kann man Karte für Karte beweisen. Dort ist es Folge von $\frac{\partial^{2}}{\partial y \partial x} f=\frac{\partial^{2}}{\partial x \partial y} f$. Die zweite Aussage folgt direkt aus der Holomorphie und (4.6).

\subsection {Integration von Differentialform}
In der Funktionentheorie definiert man das Integral eines Ausdrucks $f d z$ längs eines (stückweise stetigen) Weges in $\mathbb{C}$ durch Summation der einzelnen Wegintegrale. Auf einer Riemannschen Fläche können wir dies auch tun, denn stückweise befinden wir uns in einer Karte der Fläche.

Definition 4.10 (Wegintegral) Ist $\omega=\alpha d x+\beta d y$ eine 1-Form auf einer Riemannschen Fläche $X$ und $\gamma:[0,1] \rightarrow U \subseteq X$ ein Weg mit Bild in einer offenen Menge $U$ mit Karte z, so definieren wir

$$
\int_{\gamma} \omega=\int_{0}^{1}\left(\alpha(\gamma(s)) \frac{\partial x(\gamma(s))}{\partial s}+\beta(\gamma(s)) \frac{\partial y(\gamma(s))}{\partial s}\right) d s
$$

Für beliebige Wege $\gamma:[0,1] \rightarrow X$ definieren wir $\int_{\gamma} \omega$ durch Zerlegen in endlich viele Stücke, die jeweils in einer offenen Menge mit Kartenabbildung wie oben liegen.

Um die Wohldefiniertheit dieser Definition zu zeigen, ist noch nachzuprüfen, dass das Integral nicht von der Kartenwahl und Zerlegung des Weges abhängt.

Wir definieren nun das Integral von 2-Formen.\\
Ist $\Omega$ eine 2-Form und $D: S \rightarrow X$ ein 2-Simplex, d.h. eine stetige Einbettung des (gefüllten) Dreiecks

$$
S=\left\{(x, y) \in \mathbb{R}^{2}: 0 \leq x, y \leq 1, x+y \leq 1\right\}
$$

nach $X$, dann definieren wir, falls das Bild in einer offenen Menge $U$ mit Kartenabbildung $z: U \rightarrow \mathbb{C}$ enthalten ist

$$
\iint_{D} \Omega=\iint_{z(D)} f(x, y) d x d y
$$

mit $\Omega=f(x, y) d x \wedge d y$ auf $U$. Im Allgemeinen definieren wir $\iint_{D} \Omega$ durch eine Zerteilung von $D$ in 2-Simplices, die in offenen Mengen mit Kartenabbildungen enthalten sind und die Addition der einzelnen Beiträge.

Damit überträgt sich auch die folgende Formel von der zweidimensionalen Analysis auf Riemannsche Flächen.

Satz 4.11 (Stokes-Formel) Ist $\omega$ eine Differentialform auf der Riemannschen Fläche $X$ und ist das Gebiet $D \subseteq X$ relativ kompakt (d.h. $\bar{D} \subseteq X$ ist kompakt) mit stückweise $C^{\infty}$-Rand, dann gilt

$$
\int_{\partial D} \omega=\iint_{D} d \omega
$$

Beweis: Aufgrund der relativen Kompaktheit kann man $D$ in endlich viele Teilmengen mit offenem Rand und Kreisscheibenkarten zerlegen. Die Formel gilt bekanntlich in $\mathbb{C} \cong \mathbb{R}^{2}$ und beide Seiten waren additiv auf Zerlegungen definiert.

Ein Korollar der Stokes-Formel ist folgende Charakterisierung von geschlossenen Differentialformen.

Korollar 4.12 Eine Differentialform $\omega$ auf $X$ ist geschlossen genau dann, wenn für je zwei homotope Wege $\gamma_{1}$ und $\gamma_{2}$ gilt

$$
\int_{\gamma_{1}} \omega=\int_{\gamma_{2}} \omega
$$

Beweis: Gilt für je zwei homotope Wege die Integralgleichheit, dann betrachten wir für jede Kreisscheibe $D \subset X$ zwei Wege $\gamma_{1}$ und $\gamma_{2}$, sodass der Rand $\partial D=\gamma_{1} \cdot \gamma_{2}^{-1}$ aus den zwei Wegen, in verschiedener Richtung durchlaufen, zusammengesetzt ist. Diese Wege sind offenbar homotop zueinander. Hieraus folgt mit der Stokes-Formel, dass $\int_{D} d \omega=0$ für jede Kreisscheibe und damit, dass $d \omega=0$.

Sei ungekehrt $\omega$ geschlossen. Sind $\gamma_{1}$ und $\gamma_{2}$ zwei $C^{\infty}$-Wege, die sich nicht überkreuzen und sodass das Bild der Homotopie, die $\gamma_{1}$ in $\gamma_{2}$ überführt gerade das Innere des von $\gamma_{1}$ und $\gamma_{2}$ berandeten Gebiets überdeckt, so folgt die Behauptung wieder direkt aus der Stokes-Formel. Durch Zerteilen der Wege und der Homotopie in Teilstücke erhält man die Behauptung zunächst für beliebige $C^{\infty}$-Wege. Schliesslich nähert man beliebige Wege durch $C^{\infty}$-Wege an und erhält die Behauptung in voller Allgemeinheit.

Proposition 4.13 Sei X eine einfach zusammenhängende Riemannsche Fläche und $\omega$ eine geschlossene komplexwertige Differentialform. Dann besitzt $\omega$ eine Stammfunktion, d.h. es gibt eine stetig differenzierbare Funktion $f: X \rightarrow C$ mit $d f=\omega$.

Beweis: Den Fall, dass $X$ ein Gebiet in $\mathbb{C}$ ist, übernehmen wir aus der reellen zweidimensionalen Analysis. Sei $\omega=\alpha d x+\beta d y$. Wir definieren

$$
F(x, y)=\int_{0}^{1}(x \alpha(t x, t y)+y \beta(t x, t y)) d t
$$

Diese Funktion ist offenbar differenzierbar. Wir können unter dem Integral ableiten und erhalten aufgrund von $\frac{\partial \beta}{\partial x}=\frac{\partial \alpha}{\partial y}$, dass

$$
\begin{aligned}
\frac{\partial F}{\partial x} & =\int_{0}^{1}\left(t x \frac{\partial \alpha}{\partial x}(t x, t y)+t y \frac{\partial \beta}{\partial x}(t x, t y)+\alpha(t x, t y)\right) d t \\
& =\int_{0}^{1}\left(t \frac{\partial \alpha}{\partial t}(t x, t y)+\alpha(t x, t y)\right) d t \\
& =\int_{0}^{1}\left(\frac{\partial}{\partial t} t \alpha(t x, t y)\right) d t=\alpha(x, y)
\end{aligned}
$$

Ganz analog rechnet man $\frac{\partial F}{\partial y}=\beta$ nach. Zusammen zeigen diese Aussagen die Behauptung im Spezialfall.

Im Allgemeinen sei $a \in X$ ein fixierter Basispunkt und wir definieren $f(x)=\int_{a}^{x} \omega$. Die Wohldefiniertheit dieser Definition folgt unmittelbar aus Korollar 4.12 und offenbar ist $d f=\omega$.

\section {Garbenkohomologie}
Im vorigen Abschnitt haben wir Differentialformen integriert und Funktionen abgeleitet. Holomorphe Funktionen gibt es lokal, d.h. auf jeder Karte einer Riemannschen Fläche, und damit gibt es lokal auch holomorphe Differentialformen. Aber gibt es auf einer Riemannschen Fläche global überhaupt holomorphe Differentialformen? Solche Fragestellungen, ob und wie viele globale Objekte aus lokalen Daten zusammengesetzt werden können, werden vom Begriff der Garbe und deren Kohomologiegruppen formalisiert. Wir führen in diesem Abschnitt Grundbegriffe der Garbentheorie und deren CzechKohomologie ein.

\subsection {Garben}
In der Definition von Differentialformen haben wir jeder offenen Menge $U$ in einer Riemannschen Fläche $X$ eine abelsche Gruppe (eigentlich sogar einen $\mathbb{R}$ -Vektorraum) zugeordnet, die Gruppe der 1-Formen auf $U$. Eine 1-Form auf $U$ gibt durch Einschränken eine 1-Form auf jeder offenen Teilmenge wieder eine

1-Form. Umgekehrt will man auf einer offenen Menge $U \cup V$ etwas als 1-Form zulassen, wenn es eine 1-Form auf $U$ und eine 1-Form auf $V$ ist und die beiden Daten auf dem Durschnitt übereinstimmen.

Genau diese Eigenschaften formalisiert der Begriff Garbe. Wir schreiben diesen allgemein für topologische Räume auf. Sei in diesem Abschnitt $X$ ein topologischer Raum und $I$ die Menge der offenen Mengen in $X$.

Definition 5.1 (Prägarbe) Eine Prägarbe $\mathcal{F}$ von abelschen Gruppen auf $X$ besteht aus folgenden Objekten.

\begin{itemize}
  \item $\mathcal{F}(\emptyset)$ ist die triviale Gruppe
  \item Für jede offene Menge $U \subset X$ ist $\mathcal{F}(U)$ eine abelsche Gruppe.
  \item Zu jedem Paar offener Mengen $U \subset V$ gibt es Restriktionsmorphismen $\rho_{U}^{V}$ : $\mathcal{F}(V) \rightarrow \mathcal{F}(U)$ mit folgenden Kompatibilitätsbedingungen: Es ist $\rho_{U}^{U}$ die Identität auf $\mathcal{F}(U)$ und für je drei offene Mengen $U \subset V \subset W$ gilt $\rho_{U}^{W}=\rho_{U}^{V} \circ \rho_{V}^{W}$.
\end{itemize}

Ist $f \in \mathcal{F}(V)$, so schreiben wir oft $\left.f\right|_{U}$ für $\rho_{U}^{V}(f)$.\\
Definition 5.2 (Garbe) Eine Prägarbe auf $X$ wird Garbe genannt, falls für jede Überdeckung einer offenen Menge $U$ durch $\cup_{i \in J} U_{i}$ folgende Eigenschaften erfüllt sind.

\begin{itemize}
  \item Sind $f, g \in \mathcal{F}(U)$ und gilt $\left.f\right|_{U_{i}}=\left.g\right|_{U_{i}}$ für alle $i \in J$, so ist $f=g$.
  \item Gibt es Elemente $f_{i} \in \mathcal{F}\left(U_{i}\right)$ mit der Kompatibilitätsbedingung $\left.f_{i}\right|_{U_{i} \cap U_{j}}= \left.f_{j}\right|_{U_{i} \cap U_{j}}$ für alle $i, j \in J$, so gibt es ein Element $f \in \mathcal{F}(U)$ mit $f_{i}=\left.f\right|_{U_{i}}$ für alle $i \in J$.
\end{itemize}

Beispiele für Garben ist die Garbe der holomorphen Funktionen $\mathcal{O}_{X}$ auf einer Riemannschen Fläche, welche jeder offenen Menge die holomorphen Funktionen auf dieser Menge zuordnet. Analog bilden meromorphe Funktionen auf $X$ eine Garbe $\mathcal{M}_{X}$. Diese Garben sind nicht nur Garben abelscher Gruppen, sondern Garben von Ringen, d.h. für jedes $U$ ist $\mathcal{O}_{X}(U)$ bzw. $\mathcal{M}_{X}(U)$ ein Ring und die Einschränkungsabbildungen sind Ringhomomorphismen.

Das motivierende Beispiel der Garbe der holomorphen 1-Formen bezeichnen wir mit $\Omega_{X}^{1}$, d.h. $\Omega_{X}^{1}(U)$ sind die holomorphen 1-Formen auf $U \subset X$. Ebenso haben wir im letzten Abschnitt 2-Formen definiert, welche offensichtlich eine Garbe bilden, die wir mit $\Omega_{X}^{2}$ bezeichnen.

Eine weiteres wichtige Beispiel von Garben sind konstante Garben. D.h. wenn man eine abelsche Gruppe $G$ fixiert und zu jeder nicht-leeren offenen Menge $G(U)=G$ definiert, so sind mit der Identität als Restriktionsabbildung alle Axiome einer Prägarbe erfüllt. Ist aber $U$ nicht zusammenhängend und durch seine Zusammenhangskomponenten $U_{i}$ überdeckt, so erkennt man,\\
dass diese Definition keine Garbe liefert. D.h. man sollte $G(U)$ als direkte Summe von $G$ definieren, für jede Zusammenhangskomponente von $U$ eine Kopie. Das gleiche leistet folgende Definition.

Sei $G$ eine abelsche Gruppe und $G(U)$ die Gruppe der lokal konstanten Abbildungen von $U$ nach $G$. (Eine solche Abbildung ist notwendigerweise auf jeder Zusammenhangskomponente von $U$ konstant.) Zu jeder Inklusion $U \subset V$ von offenen Mengen erhalten wir eine Restriktionsabbildung $\rho_{U}^{V}: G(V) \rightarrow G(U)$. Die so erhaltene Prägarbe ist offenbar eine Garbe, genannt die konstante Garbe mit Werten in $G$.

Neben den oben genannten Garben von holomorphen Funktione und Differentialformen bildet auch die unendlich oft differentierbaren Funktionen eine Garbe, die wir mit $\mathcal{E}$ bezeichen. Die Garbe der unendlich oft differenzierbaren 1 -Formen werden wir mit $\mathcal{E}_{X}^{1}$ bezeichnen. Die $(1,0)$-Formen und $(0,1)$ Formen bilden eine Garbe die wir mit $\mathcal{E}^{(1,0)}$ und $\mathcal{E}^{(0,1)}$ bezeichnen. Die unendlich oft differenzierbaren 2-Formen bilden natürlich auch eine Garbe, die mit $\mathcal{E}_{X}^{2}$ bezeichnet wird. Ein Ziel dieses Abschnitts ist es aufzuzeigen, dass sich $\mathcal{E}_{X}^{1}$ und $\Omega_{X}^{1}$ sich in einem Punkt wesentlich verschieden verhalten: ihre Kohomologiegruppen sind verschieden.

\subsection {Komplexe und exakte Sequenzen}
Dieser Abschnitt vermittelt nur eine Sprechweise. Für $i$ in $\mathbb{Z}$ seien $G_{i}$ abelsche Gruppen und $\varphi_{i}: G_{i} \rightarrow G_{i+1}$ Homomorphismen. Diese Situation deuten wir mit

$$
\cdots G_{-1} \rightarrow G_{0} \rightarrow G_{1} \rightarrow G_{2} \rightarrow \cdots
$$

an. In der Praxis sind i.d.R. alle bis auf endlich viele der $G_{i}$ trivial und wir lassen den Rest der Kette weg.

Bei abelschen Gruppen kennen wir den Begriff von Kern und Bild. Wer nennen (5.1) einen Komplex, falls

$$
\operatorname{Bild}\left(\varphi_{i}\right) \subset \operatorname{Ker}\left(\varphi_{i+1}\right) \quad \text { für alle } \quad i \in \mathbb{Z} .
$$

Äquivalent gesagt, ist (5.1) ein Komplex, falls $\varphi_{i+1} \circ \varphi_{i}=0$ für alle $i$ ist.\\
Ein Komplex heißt exakt, falls $\operatorname{Bild}\left(\varphi_{i}\right)=\operatorname{Ker}\left(\varphi_{i+1}\right)$ für alle $i \in \mathbb{Z}$ ist. In diesem Fall wird (5.1) auch eine exakte Sequenz genannt. Sind alle $G_{i}=0$ mit Ausnahme von $G_{0}$ und $G_{1}$, so bedeutet Exaktheit von (5.1) gerade, dass $\varphi_{0}$ ein Isomorphismus ist. Ein exakte Sequenz, bei denen $G_{i}$ für genau drei auf einanderfolgenden Indices von Null verschieden ist, wird kurze exakte Sequenz genannt.

\subsection {Czech-Kohomologie}
Bei einer Garbe $\mathcal{F}$ auf $X$ ist oft die Menge der globalen Schnitte von primärem Interesse. Aber diese sind oft nicht leicht zu bestimmen. Um zumindest die Di-\\
mension auszurechnen, ordnen wir einer Garbe eine ganze Reihe von Vektorräumen (und damit deren Dimensionen) zu. Diese Hilfsgrößen sind wichtige Invarianten einer Garbe und lassen indirekt Schlüsse auf die Dimensions des Raums der globalen Schnitte zu.

Sei dazu $\mathfrak{U}=\left\{U_{i}, i \in I\right\}$ eine beliebige Überdeckung von $X$. Eine Czech-pKokette (oder kurz: $p$-Kokette) einer Garbe $\mathcal{F}$ ist ein Schnitt der Garbe auf dem Produkt aller ( $p+1$ )-fachen Durchschnitten der $U_{i}$. Die Menge aller $p$-Koketten bilden eine Gruppe, bezeichnet mit

$$
C^{p}(\mathfrak{U}, \mathcal{F})=\prod_{\left(i_{0}, \ldots, i_{p}\right) \in I^{p}} \mathcal{F}\left(U_{i_{0}} \cap U_{i_{1}} \cap \cdots \cap U_{i_{p}}\right)
$$

Man beachte hierbei, dass Elemente von $I^{p}$ geordnete Tupel sind, die Reihenfolge der Indizes eine wesentliche Rolle spielt.

Wir wollen nun aus den Gruppen der $p$-Koketten einen mit $p$ indizierten Komplex machen. Dazu definieren wir einen Korand-Operator $d: C^{p}(\mathfrak{U}, \mathcal{F}) \rightarrow C^{p+1}(\mathfrak{U}, \mathcal{F})$ wie folgt. Ist $s \in C^{p}(\mathfrak{U}, \mathcal{F})$, dann ist

$$
d(s)\left(U_{0} \cap \cdots \cap U_{p+1}\right)=\left.\sum_{j=0}^{p+1}(-1)^{j} s\left(U_{0} \cap \cdots \cap \widehat{U_{j}} \cap \cdots \cap U_{p+1}\right)\right|_{U_{0} \cap \cdots \cap U_{p+1}} .
$$

Dabei bedeute das Symbol $\widehat{U_{j}}$, dass die Menge $U_{j}$ bei der Durchschnittsbildung weggelassen wurde.

Lemma 5.3 Mittels des Korandoperators $d$ wird $\left(C^{p}(\mathfrak{U}, \mathcal{F})\right)_{p \in \mathbb{N}_{0}}$ zu einem Komplex.\\
Beweis: Wir beginnen mit dem Fall $p=1$. Sei $d(s) \in d\left(C^{0}(\mathfrak{U}, \mathcal{F})\right)$ durch $s$ gegeben, also $d(s)\left(U_{i} \cap U_{j}\right)=s\left(U_{j}\right)_{U_{i} \cap U_{j}}-s\left(U_{i}\right)_{U_{i} \cap U_{j}}$. Dann ist

$$
\begin{aligned}
d^{2}(s)\left(U_{0} \cap U_{1} \cap U_{2}\right) & =d(s)\left(U_{1} \cap U_{2}\right)-d(s)\left(U_{0} \cap U_{2}\right)+d(s)\left(U_{0} \cap U_{1}\right) \\
& =\left(s\left(U_{2}\right)-s\left(U_{1}\right)\right)-\left(s\left(U_{2}\right)-s\left(U_{0}\right)\right)+\left(s\left(U_{1}\right)-s\left(U_{0}\right)\right)=0
\end{aligned}
$$

wobei wir auf der rechten Seite die Einschränkungen aller Schnitte auf den Durchschnitt $U_{0} \cap U_{1} \cap U_{2}$ der Übersichtlichkeit halber weggelassen haben. Bei dieser Rechnung wird die Bedeutung des Minuszeichens in der Definition von $d$ deutlich und der allgemeine Fall läßt sich analog nachrechnen.

Wir nennen $s$ einen $p$-Kozykel, falls $d(s)=0$ gilt. Die Kozykelbedingung impliziert z.B. bei einem 1-Kozykel $f=\left\{f_{i_{0} i_{1}}\right\}_{\left(i_{0}, i_{1}\right) \in I^{2}} \in Z^{1}(\mathfrak{U}, \mathcal{F})$, dass

$$
f_{i i}=0 \quad \text { und } \quad f_{i_{0} i_{1}}=-f_{i_{1} i_{0}}
$$

gilt. Die Menge aller Kozykel bildet eine Gruppe, die wir mit $Z^{p}(\mathfrak{U}, \mathcal{F})$ bezeichnen. Die Bilder $d\left(C^{p-1}(\mathfrak{U}, \mathcal{F})\right)$ bilden nach dem vorigen Lemma eine Untergruppe $B^{p}(\mathfrak{U}, \mathcal{F})$ von $Z^{p}(\mathfrak{U}, \mathcal{F})$, deren Elemente wir als $p$-Koränder bezeichnen.

Definition 5.4 (Kohomologiegruppen zu Überdeckungen) Sei $\mathcal{F}$ eine Garbe und $\mathfrak{U}$ eine Überdeckung von $X$. Dann heißt

$$
H^{p}(\mathfrak{U}, \mathcal{F})=Z^{p}(\mathfrak{U}, \mathcal{F}) / B^{p}(\mathfrak{U}, \mathcal{F})
$$

die p-te Kohomologiegruppe von $\mathcal{F}$ bezüglich $\mathfrak{U}$.\\
Aus der Definition folgt direkt, dass

$$
H^{0}(\mathfrak{U}, \mathcal{F})=Z^{0}(\mathfrak{U}, \mathcal{F})=\mathcal{F}(X)
$$

für alle Überdeckungen die globalen Schnitte von $\mathcal{F}$ beschreibt.\\
Bisher haben wir bei Garben nur die Gruppeneigenschaft von $\mathcal{F}(U)$ gefordert. Diese Mengen sind aber oft $\mathbb{C}$ - Vektorräume, z.B. bei $\mathcal{E}, \mathcal{E}^{1}, \mathcal{O}_{X}$ und $\Omega_{X}^{1}$. Die Restriktionsabbildungen der Garben sind $\mathbb{C}$-lineare Abbildungen und so sind auch alle Homomorphismen, die wir bei der Definition der Kohomologiegruppen benötigt haben, lineare Abbildungen. In diesem Fall ist $H^{p}(\mathfrak{U}, \mathcal{F})$ also ein $\mathbb{C}$-Vektorraum.

Diese Definition liefert noch keine Gruppen, die der Garbe $\mathcal{F}$ auf dem Raum intrinsisch zugeordnet sind, denn sie hängen noch von der Wahl der Überdeckung $\mathfrak{U}$ ab. Ein intrinsisches Objekt erhält man, wenn man alle Überdeckungen nimmt.

Eine Überdeckung $\mathfrak{V}$ von $X$ wird feiner als $\mathfrak{U}$ genannt, falls es zu jedem Element $V$ von $\mathfrak{V}$ ein Element von $\mathfrak{U}$ gibt, in dem $V$ enthalten ist.

Ist also $\mathfrak{V}$ feiner als $\mathfrak{U}$ und wählt man zu jedem $V \in \mathfrak{V}$ ein $U=U(V) \in \mathfrak{U}$ mit $V \subset U$ so erhält man offenbar durch Restriktion für alle $p \in \mathbb{N}_{0}$ Vergleichsabbildungen zwischen den Koketten $v_{\mathfrak{V}}^{\mathfrak{U}}: C^{p}(\mathfrak{U}, \mathcal{F}) \rightarrow C^{p}(\mathfrak{V}, \mathcal{F})$.

Lemma 5.5 Ist $\mathfrak{V}$ feiner als $\mathfrak{U}$, so induzieren die Abbildungen $v_{\mathfrak{V}}^{\mathfrak{U}}$ Vergleichsabbildungen

$$
v_{\mathfrak{V}}^{\mathfrak{U}}: H^{p}(\mathfrak{U}, \mathcal{F}) \rightarrow H^{p}(\mathfrak{V}, \mathcal{F})
$$

welche unabhängig von der Wahl von $U \supset V$ und wohldefiniert sind. Für $p=1$ ist diese Abbildung zudem injektiv.

Beweis: Sei $\mathfrak{U}=\left\{U_{i}, i \in I\right\}, \mathfrak{V}=\left\{V_{j}, j \in J\right\}$ und $\tau: J \rightarrow I$ eine der Abbildungen, die nach Definition von 'feiner' existieren, sodass $V_{j} \subset U_{\tau(i)}$ gilt.

Wir zeigen die Injektivität. Sei dazu $f=\left\{f_{i_{0} i_{1}}\right\}_{\left(i_{0}, i_{1}\right) \in I^{2}} \in Z^{1}(\mathfrak{U}, \mathcal{F})$ mit der Eigenschaft, dass $v_{\mathfrak{V}}^{\mathfrak{U}}(f)=d(g)$ für ein $g \in C^{0}(\mathfrak{V}, \mathcal{F})$, d.h. für alle $\left(j_{0}, j_{1}\right)$ ist $f_{\tau\left(j_{0}\right), \tau\left(j_{1}\right)}=g_{j_{1}}-g_{j_{0}}$ auf $V_{j_{0}} \cap V_{j_{1}}$. Nach der Kozykeleigenschaft gilt auf $U_{i} \cap V_{j_{0}} \cap V_{j_{1}}$

$$
g_{j_{1}}-g_{j_{0}}=f_{\tau\left(j_{0}\right), \tau\left(j_{1}\right)}=f_{\tau\left(j_{0}\right), i}+f_{i, \tau\left(j_{1}\right)}=f_{\tau\left(j_{0}\right), i}-f_{\tau\left(j_{1}\right), i}
$$

das heißt $f_{\tau\left(j_{0}\right), i}+g_{j_{0}}=f_{\tau\left(j_{1}\right), i}+g_{j_{1}}$. Diese Summen kann man also nach den Garbenaxiomen auf $U_{i}$ zu einem Schnitt $h_{i} \in \mathcal{F}\left(U_{i}\right)$ verkleben, sodass\\
$h_{i}=f_{\tau\left(j_{0}\right), i}+g_{j_{0}}$ auf $U_{i} \cap V_{j_{0}}$. Wir behaupten, dass diese $h_{i}$ das Gewünschte leisten. Dazu rechnen wir auf $U_{i_{0}} \cap U_{i_{1}} \cap V_{j}$ nach, dass

$$
f_{i_{0} i_{1}}=f_{i_{0} \tau(j)}+f_{\tau(j) i_{1}}=\left(f_{\tau(j) i_{1}}+g_{\tau(j)}\right)-\left(f_{\tau(j) i_{0}}+g_{\tau(j)}\right)=h_{i_{1}}-h_{i_{0}} .
$$

Da die $V_{j}$ für $j \in J$ ganz $X$, insbesondere ganz $U_{i_{0}} \cap U_{i_{1}}$ überdecken, gilt die geforderte Gleichheit dort überall.

Den Beweis der Wohldefiniertheit überlassen wir als Übung, um schnell zur endgültigen Definition von Czech-Kohomologie zu gelangen.

Zu je zwei Überdeckungen erhält man durch Durchschnittsbildung eine Überdeckung, die feiner als beide ursprünglichen Überdeckungen ist. Überdeckungen mit 'feiner' als Relation bilden also ein gerichtetes System. Auf der Summe $\bigoplus_{\mathfrak{U}} H^{p}(\mathfrak{U}, \mathcal{F})$ der Kohomologiegruppen definiert man eine Äquivalenzrelation, welche $s \in H^{p}(\mathfrak{U}, \mathcal{F})$ und $t \in H^{p}(\mathfrak{V}, \mathcal{F})$ identifiziert, falls $\mathfrak{V}$ feiner als $\mathfrak{U}$ ist und $t=v_{\mathfrak{V}}^{\mathfrak{U}}(s)$. Die Faktorgruppe wird als direkter Limes über alle Überdeckungen bezeichnet, und $\lim _{\mathfrak{U}} H^{p}(\mathfrak{U}, \mathcal{F})$ geschrieben.

Definition 5.6 (Czech-Kohomologiegruppe) Der direkte Limes der p-ten Kohomologiegruppen zu einer Überdeckung

$$
H^{p}(X, \mathcal{F})=\underset{\mathfrak{U}}{\lim } H^{p}(\mathfrak{U}, \mathcal{F})
$$

heißt p-te Czech-Kohomologiegruppe der Garbe $\mathcal{F}$.\\
Wie auch bei den Kohomologiegruppen bzgl. einer Überdeckung sind die Czech-Kohomologiegruppen $\mathbb{C}$-Vektorräume, falls die Schnitte bzw. Restriktionsabbildungen der Garbe $\mathbb{C}$-Vektorräume bzw. lineare Abbildungen sind. In diesem Fall schreiben wir

$$
h^{p}(X, \mathcal{F})=\operatorname{dim}_{\mathbb{C}} H(X, \mathcal{F})
$$

Das Rechnen mit dieser Definition ist nicht sehr praktisch, wenn man alle Überdeckungen berücksichtigen muss. Wir berechnen einige wichtige Kohomologiegruppen um dann zu einer Berechnungsmethode für CzechKohomologiegruppen zu gelangen, die auf Überdeckungen mit endlich vielen offenen Mengen durchgeführt werden kann.

Proposition 5.7 Die erste Kohomologiegruppe der Garbe unendlich oft differenzierbarer Funktionen verschwindet auf einer Riemannschen Fläche $X$, d.h. $H^{1}(X, \mathcal{E})=$ 0 .

Beweis: Technisches Hilfsmittel hierzu ist die Existenz einer Partition der Eins, welche für stetig differenzierbare Funktionen existiert. (Für einen Beweis der Existenz einer stetigen Partition der Eins siehe z.B. [Jän87, Kapitel 8]. Die dabei verwendeten Funktionen aus dem Uryson-Lemma kann man durch Faltung\\
mit einer $C^{\infty}$-Funktion mit kompaktem Träger glätten und erhält somit auch eine Partition der Eins aus $C^{\infty}$-Funktionen.) Sei $\mathfrak{U}=\left\{U_{i}, i \in I\right\}$ eine Überdeckung von $X$. Dann gibt es eine Kollektion unendlich oft differenzierbarer Funktionen $\psi_{i}, i \in I$ jeweils mit Träger $\operatorname{supp}\left(\psi_{i}\right) \subset U_{i}$, sodass jedes $x \in X$ eine Umgebung hat, welche nur endlich viele der Trägermengen $\operatorname{supp}\left(\psi_{i}\right)$ trifft und sodass $\sum_{i \in I} \psi_{i}=1$ gilt.

Sei $\left(f_{i j}\right) \in Z^{1}(\mathfrak{U}, \mathcal{E})$ ein 1-Kozykel, gegeben durch $f_{i j} \in \mathcal{E}\left(U_{i} \cap U_{j}\right)$. Die Funktion $\psi_{i} f_{i j}$ hat Träger in $U_{i}$. Die Fortsetzung mit Null ausserhalb ihrers Tägers zu einer Funktion auf ganz $U_{i}$ ist also immer noch unendlich oft differenzierbar, d.h. $\psi_{i} f_{i j} \in \mathcal{E}\left(U_{i}\right)$.

Für jedes $i \in I$ setzen wir nun $g_{i}=\sum_{k \in I} \psi_{k} f_{i k}$. Die Endlichkeitsvoraussetzung für die Träger in der Definition der Zerlegung der Eins gestattet diese Definition. Also ist $g=\left(g_{i}\right)_{i \in I} \in C^{0}(\mathfrak{U}, \mathcal{E})$ und

$$
\begin{aligned}
d(g)\left(U_{i} \cap U_{j}\right) & =\sum_{k \in I} \psi_{k} f_{j k}-\sum_{k \in I} \psi_{k} f_{i k} \\
& =\sum_{k \in I} \psi_{k}\left(f_{j k}+f_{k i}\right)=\sum_{k \in I} \psi_{k} f_{j i}=-f,
\end{aligned}
$$

wobei wir für die beiden mittleren Umformungen die Kozykeleigenschaft von $f$ verwendet haben.

Mit dem gleichen Beweis zeigt man auch den analogen Satz für $C^{\infty}$ Differentialformen.

Proposition 5.8 Auf einer Riemannschen Fläche $X$ ist $H^{1}\left(X, \mathcal{E}_{X}^{1}\right)=0$, $H^{1}\left(X, \mathcal{E}_{X}^{(1,0)}\right)=0$ und $H^{1}\left(X, \mathcal{E}_{X}^{(0,1)}\right)=0$.

Für holomorphe Funktionen existiert keine Partition der Eins, wie der Identitätssatz leicht zeigt. Wir sehen dann später im Allgemeinen, dass $H^{1}\left(X, \mathcal{O}_{X}\right) \neq 0$ ist. Zunächst brauchen wir noch einen Verschwindungssatz.

Satz 5.9 Die ersten Kohomologiegruppen mit konstanten Koeffizienten $\mathbb{Z}$ und $\mathbb{C}$ verschwinden auf einer einfach zusammenhängenden Riemannschen Fläche X , d.h.

$$
H^{1}(X, \mathbb{C})=0 \quad \text { und } \quad H^{1}(X, \mathbb{Z})=0
$$

Beweis: Zum Beweis der ersten Aussage sei $c=\left(c_{i j}\right) \in Z^{1}(\mathfrak{U}, \mathbb{C})$ gegeben. Da konstante Funktionen unendlich oft differenzierbar sind, können wir diesen auch als Element von $Z^{1}(\mathfrak{U}, \mathcal{E})$ auffassen. Nach Proposition 5.7 gibt es also $g \in C^{0}(\mathfrak{U}, \mathcal{E})$ mit $d(g)=c$. Die Ableitungen $d g_{i} \in C^{0}\left(\mathfrak{U}, \mathcal{E}_{X}^{1}\right)$ stimmen auf allen Durchschnitten überein (man verwechsle hier nicht Ableitung und Korandoperator), denn

$$
\left.d g_{i}\left(U_{i}\right)\right|_{U_{i} \cap U_{j}}-\left.d g_{j}\left(U_{j}\right)\right|_{U_{i} \cap U_{j}}=d c_{i j}=0
$$

da $c_{i j}$ konstant ist. Aufgrund der Garbeneigenschaft gibt es also eine global definierte $C^{\infty}$-Differentialform $\omega \in \mathcal{E}_{X}^{1}(X)$, deren Restriktionen die $d g_{i}$ sind. Auf jeder Menge $U_{i}$ ist $\left.\omega\right|_{U_{i}}=d g_{i}$, also exakt und damit nach Lemma 4.9 geschlossen. Da $X$ als einfach zusammenhängend vorausgesetzt ist, gibt es nach Proposition 4.13 eine Stammfunktion von $\omega$, d.h. eine unendlich oft differenzierbare Funktion $f$ mit $d f=\omega$. Die Funktion $\left(g_{i}-f\right) \in \mathcal{E}\left(U_{i}\right)$ hat Ableitung Null, ist also konstant und somit ist $\left(g_{i}-f\right)_{i \in I} \in C^{0}(\mathfrak{U}, \mathbb{C})$. Der Korand hiervon ist $d\left(g_{i}-f\right)=d\left(g_{i}\right)=c$, was gesucht war.

Zum Beweis der zweiten Aussage sei $b=\left(b_{i j}\right) \in Z^{1}(\mathfrak{U}, \mathbb{Z})$ gegeben. Da $\mathbb{Z} \subset \mathbb{C}$ gibt es nach der ersten Aussage eine $\mathbb{C}$-Kokette $\left(g_{i}\right) \in C^{0}(\mathfrak{U}, \mathbb{C})$ mit $d\left(\left(g_{i}\right)\right)=$ b. Wir schreiben $\exp (z)$ für $e^{2 \pi i z}$. Die Konstanten $\exp \left(g_{i}\right) \in \mathcal{E}\left(U_{i}\right)$ stimmen auf den Durchschnitten überein, denn die Differenz auf den Durchschnitten ist nach Konstruktion exp angewandt auf eine ganze Zahl. Also ist $\exp \left(g_{i}\right)$ gleich einer globalen konstanten $E$. Sei $f \in \mathbb{C}$ derart, dass $\exp (f)=E$. Dann ist $g_{i}-f$ ganzzahlig, denn exp hiervon ist Eins. Also ist $\left(g_{i}-f\right)_{i \in I} \in C^{0}(\mathfrak{U}, \mathbb{Z})$ und der Korand hiervon ist $b$, was zu zeigen war.

Wir können nun das Kriterium angeben, welche Überdeckung fein genug ist, um die Czech-Kohomologie korrekt auszurechnen.

Satz 5.10 (Leray) Ist $\mathcal{F}$ eine Garbe und $\mathfrak{U}$ eine Überdeckung von $X$ mit $H^{1}\left(U_{i}, \mathcal{F}\right)=0$ für alle $i$, dann ist

$$
H^{1}(X, \mathcal{F})=H^{1}(\mathfrak{U}, \mathcal{F})
$$

Beweis: Siehe z.B. [FG81, Theorem 12.8].\\
Damit können wir in einem Beispiel eine nichtverschwindende Kohomologiegruppe berechnen. Alle weiteren solchen Ergebnisse basieren auf der Kohomologiesequenz im nächsten Abschnitt.

Beispiel 5.11 Wir wollen zeigen, dass

$$
H^{1}\left(\mathbb{C}^{*}, \mathbb{Z}\right)=\mathbb{Z}
$$

ist. Dazu überdecken wir $\mathbb{C}^{*}$ mit $U_{1}=\mathbb{C}^{+}=\mathbb{C} \backslash \mathbb{R}_{\geq 0}$ und $U_{2}=\mathbb{C}^{-}=\mathbb{C} \backslash \mathbb{R}_{\leq 0}$. Diese Mengen sind sternförmig, also erfüllt $\mathfrak{U}=\left\{U_{i}\right\}_{i=1}^{2}$ nach dem Satz 5.9 die Voraussetzungen des Satzes von Leray.

Ist $\left(a_{i j}\right) \in Z^{1}(\mathfrak{U}, \mathbb{Z})$ ein Kozykel, so ist wegen der Kozykeleigenschaft $a_{i i}=0$ und $a_{12}=-a_{21}$, also der Kozykel durch $a_{12}$ eindeutig festgelegt. Da $U_{1} \cap U_{2}$ zwei Komponenten hat, ist $\mathbb{Z}\left(U_{1} \cap U_{2}\right) \cong \mathbb{Z}^{2}$, gegeben durch den Wert des Schnitts auf der oberen und unteren Halbebene. Es ist $C^{0}(\mathfrak{U}, \mathbb{Z}) \cong \mathbb{Z}^{2}$, gegeben durch den Wert auf $U_{1}$ und $U_{2}$. Der Korand von ( $b_{1}, b_{2}$ ) ist in jeder der Komponenten von $U_{1} \cap U_{2}$ gegeben durch $b_{2}-b_{1}$, also ist

$$
H^{1}\left(\mathbb{C}^{*}, \mathbb{Z}\right) \cong \mathbb{Z}^{2} /\left\{\left(b_{2}-b_{1}, b_{2}-b_{1}\right):\left(b_{1}, b_{2}\right) \in \mathbb{Z}^{2}\right\} \cong \mathbb{Z}
$$

\subsection {Die lange exakte Kohomologiesequenz}
Von einzelnen Garben die Kohomologie auszurechnen ist mühsam. Wir wollen also Garben miteinander Vergleichen, d.h. Homomorphismen und exakte Sequenzen definieren und schließlich prüfen, wie diese sich mit Kohomologie vertragen.

Definition 5.12 (Garbenhomomorphismen) Sind $\mathcal{F}$ und $\mathcal{G}$ Garben auf $X$, so besteht ein Garbenhomomorphismus $\alpha$ aus Gruppenhomomorphismen $\alpha(U)$ : $\mathcal{F}(U) \rightarrow \mathcal{G}(U)$ für jede offene Menge $U \subset X$, welche mit den Restriktionsabbildungen kompatibel sind, d.h. sodass wenn immer $U \subset V$, so ist $\rho_{V}^{U} \circ \alpha(U)=\alpha(V) \circ \rho_{V}^{U}$.

Beispiele für Garbenhomomorphismen sind die Inklusion der konstanten Garbe $\mathbb{C}$ in die Garbe $\mathcal{E}_{X}$, welche wir bereits im Beweis von Satz 5.9 implizit verwendet haben, sowie die äussere Ableitung $d: \mathcal{E}^{1} \rightarrow \mathcal{E}^{2}$ oder $d: \Omega_{X}^{1} \rightarrow \Omega_{X}^{2}$, denn diese Abbildungen hatten wir ja bereits für offene Menge einzeln definiert.

Beispiel 5.13 Auf einer Riemannschen Fläche $X$ ist für jedes $U \subset X$ die Menge der holomorphen Funktionen, welche nirgends eine Nullstelle haben, mit der Multiplikation als Verknüpfung eine abelsche Gruppe. Diese abelschen Gruppen bilden offenbar eine Garbe, die wir mit $\mathcal{O}_{X}^{*}$ bezeichnen.\\
Für jede offene Menge $U$ und $f$ holomorph auf $U$ definieren wir die Exponentialabbildung e durch $\mathbf{e}(f)(z)=\exp (2 \pi i f(z))$. Diese definiert offenbar einen Garbenmorphismus

$$
\mathbf{e}: \mathcal{O}_{X} \rightarrow \mathcal{O}_{X}^{*}
$$

welcher uns noch oft begegnen wird.\\
Definition 5.14 (Kern) Ist $\alpha: \mathcal{F} \rightarrow \mathcal{G}$ ein Garbenmorphismus so ist $\operatorname{Ker}(\alpha)$ definiert durch $\operatorname{Ker}(\alpha)(U)=\operatorname{Ker}(\alpha(U))$ eine Garbe, die wir als Kern von $\alpha$ bezeichnen.

Um die implizite Behauptung dieser Definition zu verifzieren, überzeugt man sich, dass das erste Garbenaxiom für $\operatorname{Ker}(\alpha)$ gilt, da dies in $\mathcal{F}$ gilt und dass man aus dem gleichen Grund kompatible Elemente auf offenen Mengen zusammenkleben kann. Der Begriff des Kerns ist also nichts besonders Neues im Vergleich zur Gruppentheorie, der Begriff des Bildes ist jedoch weit schwieriger korrekt zu definieren.

Beispiel 5.15 Sei e: $\mathcal{O}_{X} \rightarrow \mathcal{O}_{x}^{*}$ wieder die Exponentialabbildung. Dann kann man, wie bei jedem Garbenmorphismus $\operatorname{Bild}(\mathbf{e})(U)=\operatorname{Bild}(\mathbf{e}(U)$ : $\left.\mathcal{O}_{X}(U) \rightarrow \mathcal{O}_{X}^{*}(U)\right)$ definieren. Die Restriktionsmorphismen leisten das Gewünschte, Bild(e) ist eine Prägarbe. Aber Bild(e) ist im Allgemeinen keine Garbe. Um dies zu beweisen sei $X=\mathbb{C}^{*}$, es sei $U_{1}=\mathbb{C}^{+}$und $U_{2}=\mathbb{C}^{-}$. Da die $U_{i}$ einfach zusammenhängend sind, ist $\mathbf{e}\left(U_{i}\right)$ für $i=1,2$ surjektiv, also ist\\
zum Beispiel $f_{i}(z)=z$ auf $U_{i}$ für $i=1,2 \mathrm{im}$ Bild des jeweiliegen $\mathbf{e}\left(U_{i}\right)$. Diese Funktionen verkleben sich zu $f(z)=z$, einem Element von $\mathcal{O}_{X}^{*}(X)$, aber dieses liegt nicht im Bild der Exponentialabbildung.

Die Garbeneigenschaft will man beim Begriff des Bildes keinesfalls aufgeben. Stattdessen nennt man einen Garbenmorphismus surjektiv, wenn er dies auf genügend kleinen offenen Mengen ist. Der folgende Begriff liefert diese Eigenschaft präzise, sodass die Exponentialabbildung am Ende wie erwartet ein surjektiver Garbenhomomorphismus ist.

Definition 5.16 (Halm) Sei $x \in X$ ein Punkt und $\mathcal{F}$ eine Garbe auf $X$. Der Halm von $\mathcal{F}$ bei $x$ ist

$$
\mathcal{F}_{x}=\underset{U \ni x}{\underline{\lim }} \mathcal{F}(U)=\bigoplus_{U \ni x} \mathcal{F}(U) / \sim,
$$

der direkte Limes der abelschen Gruppen über alle offenen Mengen, die $x$ enthalten, d.h. die Summe über die Werte auf allen solchen Menge modulo der Äquivalenzrelation $f \sim g$ für $f \in \mathcal{F}(U)$ und $g \in \mathcal{F}(V)$, falls es eine Umgebung $W \subseteq U \cap V$ von $x$ gibt mit $\left.f\right|_{W}=\left.g\right|_{W}$.

Als direkter Limes ist jeder Halm natürlich wiederum eine Gruppe.\\
Beispiel 5.17 Der Halm der Garbe der holomorphen Funktionen $\mathcal{O}_{X} \mathrm{im}$ Punkt $x$ besteht aus allen Potenzreihen, entwickelt um $x$ mit beliebigem positiven Konvergenzradius. In der Tat definiert jede solche Potenzreihe eine holomorphe Funktion auf einer Umgebung von $x$ und alle diese Umgebungen werden bei der Bildung des direkten Limes berücksichtigt. Stimmen umgekehrt zwei Funktionen $f \in \mathcal{O}_{X}(U)$ und $g \in \mathcal{O}_{X}(V)$ auf dem Durchschnitt ihrer Definitionsbereiche überein, so stimmen nach dem Identitätssatz ihre Potenzreihenentwicklungen im Punkt $x$ überein.

Der Halm von $\mathcal{O}_{X}^{*}$ im Punkt $x$ besteht aus dem gleichen Grund aus allen Potenzreihen mit positivem Konvergenzradius und nicht-verschwindendem Konstantglied.

Ist $\alpha: \mathcal{F} \rightarrow \mathcal{G}$ ein Garbenmorphismus, so induziert dieser einen Morphismus $\alpha_{x}: \mathcal{F}_{x} \rightarrow \mathcal{G}_{x}$ der Halme an jedem Punkt $x \in X$, indem man für alle $U \ni x$ das Element $f \in \mathcal{F}(U)$ auf die Klasse von $\alpha(U)(f) \in \mathcal{G}(U)$ abbildet. Die Wohldefiniertheit dieser Definition ist direkte Konsequenz der Verträglichkeit eines Garbenmorphismus mit den Restriktionsabbildungen von $U$ bzw. von $V$ nach $U \cap V$.

Definition 5.18 (Exakte Sequenz von Garben) Eine Kette von Garbenmorphismen $\varphi_{i}: \mathcal{F}_{i} \rightarrow \mathcal{F}_{i+1}$ für $i \in I \subset \mathbb{Z}$ wird exakt genannt, falls für jedes $x$ die Kette von Homomorphismen abelscher Gruppen $\left(\varphi_{i}\right)_{x}:\left(\mathcal{F}_{i}\right)_{x} \rightarrow\left(\mathcal{F}_{i+1}\right)_{x}$ exakt ist.

Dies beinhaltet die angekündigte Surjektivitätsdefinition eines Garbenmorphismus $\alpha: \mathcal{F} \rightarrow \mathcal{G}$ nach der Begriffsbildung aus Abschnitt ??. Dieser wird surjektiv genannt, wenn $\mathcal{F} \rightarrow \mathcal{G} \rightarrow 0$ eine exakte Sequenz ist. Er wird injektiv genannt, wenn die Sequenz $0 \rightarrow \mathcal{F} \rightarrow \mathcal{G}$ exakt ist. Dabei passiert, im Gegensatz zur Surjektivität nichts ungewöhnliches: $\alpha$ ist injektiv genau dann, wenn $\operatorname{Ker}(\alpha)=0$ die Nullgarbe ist.

Beispiel 5.19 Die natürliche Inklusion und die Exponentialabbildung ergeben eine exakte Sequenz

$$
0 \rightarrow \mathbb{Z} \rightarrow \mathcal{O}_{X} \rightarrow \mathcal{O}_{X}^{*} \rightarrow 0
$$

Die Injektivität ist klar. Für die Surjektivität sei ein Element $f$ im $\operatorname{Halm} \mathcal{O}_{X, x}^{*}$ an einem beliebigen Punkt $x \in X$ gegeben. Dieses Element wird durch $f \in \mathcal{O}_{X}^{*}(U)$ für ein $U \ni x$ repräsentiert. Aufgrund der Äquivalenzrelation in der Definition des Halmes können wir zu einer kleineren Umgebung von $x$ übergehen, z.B. zu einer kleinen Kreisscheibe $\Delta_{\varepsilon}(x)$. Diese ist einfach zusammenhängend und daher hat die Restriktion von $f$ auf $\Delta_{\varepsilon}(x)$ ein Urbild unter der Exponentialabbildung. Dessen Äquivalenzklasse ist das gesuchte Urbild in $\mathcal{O}_{X, x}$.

Für die Exaktheit in der Mitte bemerken wir, dass offenbar das Bild der Inklusion auf jeder offenen Menge im Kern der Exponentialabbildung ist und dass, umgekehrt, auf einer kleinen offenen Menge um einen beliebigen Punkt $x \in X$ ein Funktion nur dann im Kern der Exponentialabbildung liegt, wenn sie konstant eine ganze Zahl ist.

Ein Garbenhomomorphismus $\alpha: \mathcal{F} \rightarrow \mathcal{G}$ ergibt offenbar einen Gruppenhomomorphismus $C^{p}(\mathfrak{U}, \mathcal{F}) \rightarrow C^{p}(\mathfrak{U}, \mathcal{G})$ für jedes $p$ und jede Überdeckung $\mathfrak{U}$, denn die Gruppen der Koketten sind ja nichts anderes als die Werte der Garben auf gewissen offenen Mengen, die als Durchschnitte der $U \in \mathfrak{U}$ auftreten. Wir könnten diesen mit $\alpha_{p}$ bezeichnen, werden aber später einfach auch $\alpha$ schreiben. Der Korandoperator ist, nach Definition, eine Linearkombination von Restriktionsabbildungen. Da $\alpha$ mit diesen verträglich ist, folgt $d \circ \alpha_{p}=\alpha_{p+1} \circ d$ und somit haben wir eine eine induzierte Abbildung $\alpha_{p}: Z^{p}(\mathfrak{U}, \mathcal{F}) \rightarrow Z^{p}(\mathfrak{U}, \mathcal{G})$ zwischen den Kozykeln. Ist ein Kozykel $f=d g$ ein Korand, so ist aus dem gleichen Grund $\alpha(f)=d(\alpha(g))$, d.h. sein Bild ist auch ein Korand.

Als Fazit erhalten wir, dass jeder Garbenmorphimus $\alpha: \mathcal{F} \rightarrow \mathcal{G}$ in natürlicher Weise einen Gruppenhomomorphismus zwischen den $p$-ten Kohomologiegruppen bzgl. jeder Überdeckung und damit einen Gruppenhomomorphismus

$$
\alpha_{p}: H^{p}(X, \mathcal{F}) \rightarrow H^{p}(X, \mathcal{G})
$$

zwischen den Czech-Kohomologiegruppen induziert. Diese wollen wir nun in eine exakte Sequenz packen.

Satz 5.20 Seien $\alpha: \mathcal{F}_{1} \rightarrow \mathcal{F}$ und $\beta: \mathcal{F} \rightarrow \mathcal{F}_{2}$ Homomorphismen von Garben, sodass

$$
0 \rightarrow \mathcal{F}_{1} \rightarrow \mathcal{F} \rightarrow \mathcal{F}_{2} \rightarrow 0
$$

eine exakte Sequenz ist. Dann gibt es Verbildungshomomorphismenn

$$
\delta: H^{0}\left(X, \mathcal{F}_{2}\right) \rightarrow H^{1}\left(X, \mathcal{F}_{1}\right)
$$

die zusammen mit den induzierten Abbildung $\alpha_{p}, \beta_{p}$ eine exakte Sequenz

$$
\begin{aligned}
0 \longrightarrow H^{0}\left(X, \mathcal{F}_{1}\right) \xrightarrow{\alpha_{0}} H^{0}(X, \mathcal{F}) \xrightarrow{\beta_{0}} H^{0}\left(X, \mathcal{F}_{2}\right) \\
\delta^{0} \\
H^{1}\left(X, \mathcal{F}_{1}\right) \xrightarrow{\alpha_{1}} H^{1}(X, \mathcal{F}) \xrightarrow{\beta_{1}} H^{1}\left(X, \mathcal{F}_{2}\right)
\end{aligned}
$$

definieren.\\
Das folgende Lemma beweist unter anderem die Exaktheit der ersten Zeile.\\
Lemma 5.21 Ist $0 \rightarrow \mathcal{F}_{1} \rightarrow \mathcal{F} \rightarrow \mathcal{F}_{2} \rightarrow 0$ eine exakte Sequenz, dann ist für jedes $U \subset X$ die Sequenz

$$
0 \rightarrow \mathcal{F}_{1}(U) \rightarrow \mathcal{F}(U) \rightarrow \mathcal{F}_{2}(U)
$$

exakt.\\
Beweis: Ist $f \in \mathcal{F}_{1}(U)$ und $\alpha(U)(f)=0$, dann folgt aus Exaktheit der Sequenz, dass der Halm von $f$ bei jedem Punkt $x$ Null ist. D.h. jeder Punkt hat eine Umgebung, auf der $f$ Null ist. Nach den Garbenaxiomen folgt also, dass $f=0$ ist.

Sei $\mathcal{F}(U) \ni g=\alpha(U)(f)$. Aufgrund der Exaktheit von $\left(\mathcal{F}_{1}\right)_{x} \rightarrow \mathcal{F}_{x} \rightarrow\left(\mathcal{F}_{2}\right)_{x}$ hat jeder Punkt $x \in X$ eine Umgebung $V=V(x)$, sodas $\left.\beta(g)\right|_{V}=0$ ist. Nach den Garbenaxiomen folgt also $\beta(g)=0$, d.h. $\operatorname{Bild}(\alpha(U)) \subset \operatorname{Ker}(\beta(U))$.

Sei nun $g \in \mathcal{F}(U)$ mit $\beta(U)(g)=0$. Wir müssen noch zeigen, dass dieses Element im Bild von $\alpha(U)$ liegt. Wiederum benutzen wir die Exaktheit der Sequenz auf allen Halmen, um $U$ mit Mengen $V_{i}, i \in I$ zu überdecken, auf denen es $f_{i} \in \mathcal{F}_{1}\left(V_{i}\right)$ gibt mit $\alpha\left(V_{i}\right)\left(f_{i}\right)=\left.g\right|_{V_{i}}$. Auf den Schnittmengen $V_{i} \cap V_{j}$ ist $\alpha\left(V_{i} \cap V_{j}\right)\left(f_{i}-f_{j}\right)=0$. Aufgrund der bereits gezeigten Injektivität von $\alpha(U)$ für jedes $U$ ist also $f_{i}=f_{j}$ auf $V_{i} \cap V_{j}$ und die $f_{i}$ setzen sich nach den Garbenaxiomen also zu dem gewünschten Element $f \in \mathcal{F}_{1}(U)$ zusammen.

Beweis von Satz 5.20: Wir beginnen mit der Konstruktion des Verbindungshomomorphismus. Sei $h \in Z^{0}\left(\mathfrak{U}, \mathcal{F}_{2}\right)$ ein Repräsentant von $h \in H^{0}\left(X, \mathcal{F}_{2}\right)$. Wir können die Überdeckung $\mathfrak{U}$ so fein wählen, dass die Surjektivität von allen Halmabbildungen $\beta_{x}$ die Existenz von $g=\left(g_{i}\right) \in C^{0}(\mathfrak{U}, \mathcal{F})$ mit $\beta\left(g_{i}\right)=h_{i}$ garantiert. Um daraus ein Element in $H^{1}\left(X, \mathcal{F}_{1}\right)$ zu basteln, wenden wir den Korandoperator an und hoffen, dass $d(g)$ das Gewünschte leistet. Wegen $d^{2}=0$\\
ist $d(g) \in Z^{1}(\mathfrak{U}, \mathcal{F})$. Wegen $\beta_{1} \circ d=d \circ \beta_{0}$ ist $\beta_{1}(d(g))=0$, da $h$ ein Kozykel war. Nach dem Lemma 5.21 angewandt auf alle 2-fachen Durchschnitte der Mengen $U_{i}$ gibt es also ein Element $f \in C^{1}\left(\mathfrak{U}, \mathcal{F}_{1}\right)$ mit $\alpha_{1}(f)=d(g)$. Aus $0=d^{2}(g)=d\left(\alpha_{1}(f)\right)=\alpha_{2}(d(f))$ und der Injektivitätsaussage in Lemma 5.21 angewandt auf alle 3 -fachen Durchschnitte der Mengen $U_{i}$ folgt $d(f)=0$. Also ist $f \in Z^{1}\left(\mathfrak{U}, \mathcal{F}_{1}\right)$ und dessen Klasse in $f \in H^{1}\left(X, \mathcal{F}_{1}\right)$ definieren wir als $\delta(g)$.

Es ist noch zu prüfen, dass $\delta$ nicht von der getroffenen Wahl der Urbilder $g=\left(g_{i}\right)$ abhängt. Sei $\widetilde{g}=\left(\widetilde{g}_{i}\right)$ auf einer Überdeckung $\widetilde{\mathfrak{U}}$, bestehend aus offenen Mengen $\widetilde{U}_{i}$ ein weiteres solches Urbild. Dann gehen wir zu einer Überdeckung $\mathfrak{V}$ über, die feiner als $\mathfrak{U}$ und $\tilde{\mathfrak{U}}$ ist. Wir können mittels der Vergleichsabbildungen $g$ und $\widetilde{g}$ als Elemente von $C^{0}(\mathfrak{V}, \mathcal{F})$ auffassen. Nach der Definition von Czech-Kohomologie als Limes über alle Überdeckungen genügt es zu zeigen dass das Bild $\delta(h) \in H^{1}\left(\mathfrak{V}, \mathcal{F}_{1}\right)$ unabhängig davon ist, ob man mittels $g$ ein $f$ oder mittels $\widetilde{g}$ ein $\widetilde{f}$ konstruiert hat. Nach Konstruktion ist $\beta_{0}(g-\widetilde{g})=0$. Aufgrund der Exaktheit der Garbensequenzen hat jeder Punkt $x \in X$ eine Umgebung $W(x)$, auf der ein Element $u$ mit $\alpha_{0}(u)=\left.(g-\widetilde{g})\right|_{U}$ existiert. Durch nochmaliges Verfeinern der Überdeckung $\mathfrak{V}$ - wir ändern die Notation dafür allerdings nicht noch - können wir annehmen, dass $\mathfrak{V}$ diese $W(x)$ schon enthält, d.h. dass es ein Element $u \in C^{0}\left(\mathfrak{V}, \mathcal{F}_{1}\right)$ mit $\alpha_{0}(u)=g-\widetilde{g}$ gibt. Dann aber ist

$$
\alpha_{0}(h-\widetilde{h})=d(g-\widetilde{g})=d\left(\alpha_{0}(u)\right)=\alpha_{1}(d(u))
$$

und wegen der Injektivität von $\alpha_{1}$ ist diese Differenz ein Korand, die Klasse von $h$ in $H^{1}\left(X, \mathcal{F}_{1}\right)$ also wohldefiniert. Nach Konstruktion ist $\delta$ offenbar ein Gruppenhomomorphismus.

Wir kommen jetzt zum Beweis der Exaktheit der Sequenz. Die Exaktheit der ersten Zeile haben im bereits im vorangehenden Lemma gezeigt. Dieses zeigt allgemeiner $\operatorname{Bild}\left(\alpha_{p}\right) \subset \operatorname{Ker}\left(\beta_{p}\right)$, indem wir es auf alle $p+1$-fachen Durchschnitte von Mengen in $\mathfrak{U}$ anwenden. Für die umgekehrte Inklusion nehmen wir $g \in \operatorname{Ker}\left(\beta_{p}\right)$ her, repränsentiert durch ein Element in $Z^{p}(\mathfrak{U}, \mathcal{F})$. Dies bedeutet, dass es $h \in C^{p-1}\left(\mathfrak{U}, \mathcal{F}_{2}\right)$ gibt mit $d h=\beta_{p}(g)$. Wieder aufgrund der Surjektivität von $\mathcal{F} \rightarrow \mathcal{F}_{2}$ hat jeder Punkt $x \in X$ eine Umgebung $W(x)$ sodass $\left.h\right|_{W}=\beta\left(u_{W}\right)$ ist. Wir gehen wieder zu einer Verfeinerung $\mathfrak{V}$ von $\mathfrak{U}$ über sodass die entprechenden Umgebungen $W(x)$ in $\mathfrak{V}$ enhalten sind. Dann also gibt es $u \in C^{p-1}(\mathfrak{V}, \mathcal{F})$ mit $\beta(u)=h$. Als Elemente von $H^{p}(X, \mathcal{F})$ sind die Kozykel $g$ und $\widetilde{g}=g-d(u)$ gleich. Nach Konstruktion ist $\beta(\widetilde{g})=d(h)-d(h)=0$. Jetzt könen wir das Lemma 5.21 angewandt auf $p+1$-fache Durchschnitte verwenden, um ein Element $f \in C^{p}\left(\mathfrak{V}, \mathcal{F}_{1}\right)$ zu basteln, dessen Bild $\widetilde{g}$ ist. Wegen $d(\widetilde{g})=0$ ist $d(f)=0$, also $f$ der gewünschte Kozykel.

Für die Inklusion $\operatorname{Bild}\left(\beta_{0}\right) \subset \operatorname{Ker}(\delta)$ genügt es noch einmal den Beweis der Wohldefiniertheit von $\delta$ zu betrachten. Ist umgekehrt $h \in \operatorname{Ker}(\delta)$, so verwenden wir auch die Bezeichnungen in der Konstruktion von $\delta$ wieder. Es gibt also ein $u \in C^{0}\left(\mathfrak{U}, \mathcal{F}_{1}\right)$ mit $d(u)=f$. Wir ändern das dort gewählte

Urbild $g \in C^{0}(\mathfrak{U}, \mathcal{F})$ um $-\alpha(u)$ ab, d.h. $g^{\prime}=g-\alpha_{0}(u)$. Dann ist $d\left(g^{\prime}\right)= d(g)-d\left(\alpha_{0}(u)\right)=d(g)-\alpha_{1}(f)=0$, also $g$ wie gewünscht ein Kozykel und $\beta_{0}\left(g^{\prime}\right)=\beta_{0}(g)-\beta_{0}\left(\alpha_{0}(u)\right)=\beta_{0}(g)=h$.

Die Inklusion $\operatorname{Bild}(\delta) \subset \operatorname{Ker}\left(\alpha_{1}\right)$ ist wiederum direkt aus der Konstruktion von $\delta$ ersichtlich. Für die umgekehrte Inklusion schließlich sei $f \in \operatorname{Ker}\left(\alpha_{1}\right)$ gegeben, d.h. $\alpha_{1}(f)=d(g) \in C^{1}(\mathfrak{U}, \mathcal{F})$ für eine Kokette $g=\in C^{0}(\mathfrak{U}, \mathcal{F})$. Setzt $\operatorname{man} h=\beta_{0}(g)$, so ist $d(h)=\beta_{0}(d(g))=\beta_{0}\left(\alpha_{0}(f)\right)=0$, also ist $h \in Z^{0}\left(\mathfrak{U}, \mathcal{F}_{2}\right)$ und nach Konstruktion des Verbindungshomomorphismus ist $\delta(h)=f$.

\subsection {Das Dolbeault-Lemma}
In diesem Abschnitt beginnen wir die Kohomologie der Strukturgarbe $\mathcal{O}_{X}$ zu untersuchen. Diese wird eine der wichtigsten Invarianten für kompakte Riemannsche Flächen sein. Für Kreisscheiben hingegen trägt sie keine Information, wie der folgende Satz zeigt.

Satz 5.22 Ist $\Delta$ eine Kreisscheibe vom Radius $R$ mit $R>0$ (d.h. $R=\infty$, zugelassen), so ist $H^{1}\left(\Delta, \mathcal{O}_{\Delta}\right)=0$.

Anders gesagt, Kreisscheiben bilden eine genügend feine Überdeckung für den Satz von Leray, wenn wir die Kohomologie der Strukturgarbe bestimmen wollen.

Korollar 5.23 Für die projektive Gerade gilt $H^{1}\left(\mathbb{P}_{\mathbb{C}}^{1}, \mathcal{O}_{\mathbb{P}_{\mathbb{C}}^{1}}\right)=0$.\\
Beweis: Wir nehmen die Standardüberdeckung von $\mathbb{P}_{\mathbb{C}}^{1}$ mit zwei Kreisscheiben von Radius $\infty$ her, d.h. $U_{1}=\mathbb{P}_{\mathbb{C}}^{1} \backslash\{\infty\}$ und $U_{2}=\mathbb{P}_{\mathbb{C}}^{1} \backslash\{0\}$ und $\mathfrak{U}=\left\{U_{1}, U_{2}\right\}$. Nach dem Satz von Leray und Satz 5.22 ist $H^{1}\left(\mathbb{P}_{\mathbb{C}}^{1}, \mathcal{O}_{\mathbb{P}_{\mathbb{C}}^{1}}\right)=H^{1}\left(\mathfrak{U}, \mathcal{O}_{\mathbb{P}_{\mathbb{C}}^{1}}\right)$. Sei $f_{i j} \in Z^{1}\left(\mathfrak{U}, \mathcal{O}_{\mathbb{P}_{\mathbb{C}}^{1}}\right)$ gegeben. Dieser Kozykel ist wiederum durch $f_{12} \in \mathcal{O}_{X}\left(\mathbb{C}^{*}\right)$ vollständig bestimmt. Wir müssen zeigen, dass dies ein Korand ist. Dies ist Konsequenz der Laurententwicklung einer holomorphen Funktion auf einem Kreisring, siehe z.B. [MF, Satz 3.4]. Eine solche Funktion $f$ lässt sich als $f^{+}-f^{-}$ schreiben, wobei $f^{+}$auf $\mathbb{C}$ konvergiert und $f^{-}$auf $\mathbb{C} \backslash\{0\}$ konvergiert und $\lim _{z \rightarrow \infty} f^{-}(z)=0$ ist. Also ist $f^{-}$auf ganz $U_{2}$ holomorph und $g=\left(g_{1}=\right. \left.f^{+}, g_{2}=f^{-}\right) \in C^{0}\left(\mathfrak{U}, \mathcal{O}_{\mathbb{P}_{\mathbb{C}}^{1}}\right)$ ist der gewünschte Kozykel mit $d g=f$.

Den Beweis von Satz 5.22 bereiten wir mit einigen Hilfsaussagen vor. Wir erinnern daran, dass der Differentialoperator

$$
\frac{\partial}{\partial \bar{z}}=\frac{1}{2}\left(\frac{\partial}{\partial x}+i \frac{\partial}{\partial y}\right)
$$

eingeführt wurde, um holomorphe Funktionen zu detektieren, d.h. $f \in \mathcal{E}(\mathbb{C})$ ist holomorph, genau dann wenn $\frac{\partial f}{\partial \bar{z}}=0$.

Lemma 5.24 Sei $g \in \mathcal{E}(\mathbb{C})$ eine unendlich oft differenzierbare Funktion mit kompaktem Träger. Dann gibt es eine unendlich oft differenzierbare Funktion $f \in \mathcal{E}(\mathbb{C})$ mit

$$
\frac{\partial f}{\partial \bar{z}}=g
$$

Beweis: Wir behaupten, dass wir das gewünschte $f$ durch Integration gegen den Cauchy-Kern erhalten, d.h.

$$
f(\zeta)=\frac{1}{2 \pi i} \iint_{\mathbb{C}} \frac{g(z)}{z-\zeta} d z \wedge d \bar{z}
$$

Wir zeigen gleichzeitig, dass das Integral wohldefiniert ist und die gewünschte Eigenschaft hat und führen dazu Polarkoordinaten zentriert bei $\zeta$ ein, d.h. $z= \zeta+r e^{i \theta}$. Dann ist

$$
d z \wedge d \bar{z}=-2 i d x \wedge d y=-2 i r d r \wedge d \theta
$$

also

$$
f(\zeta)=\frac{-1}{\pi} \iint_{\mathbb{C}} g\left(\zeta+r e^{i \theta}\right) e^{-i \theta} d r \wedge d \theta
$$

An diesem Ausdruck ist die Endlichkeit des Integrals offensichtlich. Wegen des kompakten Trägers und majorisierter Konvergenz genügt es über ein genügend großes Quadrat $B$ (zentiert bei Null mit Kantenlänge $2 R$ ) zu integrieren und wir können unter dem Integral differenzieren. Sei $B_{\varepsilon}$ dieses Quadrat ohne den Innenkreis vom Radius $\varepsilon$ zentriert um Null. Dann ist

$$
\begin{aligned}
\frac{\partial f}{\partial \bar{\zeta}}(\zeta) & =\frac{-1}{\pi} \iint_{B} e^{-i \theta} \frac{\partial}{\partial \bar{\zeta}} g\left(\zeta+r e^{i \theta}\right) d r \wedge d \theta \\
& =\lim _{\varepsilon \rightarrow 0} \frac{-1}{\pi} \iint_{B_{\varepsilon}} e^{-i \theta} \frac{\partial}{\partial \bar{\zeta}} g\left(\zeta+r e^{i \theta}\right) d r \wedge d \theta \\
& =\lim _{\varepsilon \rightarrow 0} \frac{1}{2 \pi i} \iint_{B_{\varepsilon}} \frac{1}{z} \frac{\partial}{\partial \bar{\zeta}} g(\zeta+z) d z \wedge d \bar{z} \\
& =\lim _{\varepsilon \rightarrow 0} \frac{1}{2 \pi i} \iint_{B_{\varepsilon}} \frac{\partial}{\partial \bar{z}}\left(\frac{g(\zeta+z)}{z}\right) d z \wedge d \bar{z} \quad=-\lim _{\varepsilon \rightarrow 0} \iint_{B_{\varepsilon}} d \omega
\end{aligned}
$$

wobei wir beim Übergang zur dritten Zeile $z=r e^{i \theta}$ rücksubstituiert haben, wobei

$$
\omega=\frac{1}{2 \pi i} \frac{g(\zeta+z)}{z} d z
$$

ist und von $d \omega$ nur die $\frac{\partial}{\partial \bar{z}}$-Ableitung einen Beitrag liefert. Nach dem Satz von Stokes (Satz 4.11) ist also

$$
\frac{\partial f}{\partial \bar{\zeta}}(\zeta)=-\lim _{\varepsilon \rightarrow 0} \iint_{B_{\varepsilon}} d \omega=-\lim _{\varepsilon \rightarrow 0} \int_{\partial B_{\varepsilon}} \omega=\lim _{\varepsilon \rightarrow 0} \int_{|z|=\varepsilon} \omega
$$

da wir annehmen können, dass die Aussenkanten von $B_{\varepsilon}$ ausserhalb des Trägers von $g$ liegen und wobei $\{|z|=\varepsilon\}$ im letzten Ausdruck mathematisch positiv durchlaufen wird, sich also die Umlaufrichtung und damit das Vorzeichen\\
geändert hat. Also ist schliesslich

$$
\frac{\partial f}{\partial \bar{\zeta}}(\zeta)=\lim _{\varepsilon \rightarrow 0} \frac{1}{2 \pi} \int g\left(\zeta+\varepsilon e^{i \theta}\right) d \theta
$$

Dies ist der Mittelwert von $g$ auf einem Kreisrand von Radius $\varepsilon$ um $\zeta$. Da $g$ stetig ist, ist der Limes hiervon gerade $g(\zeta)$, was zu zeigen war.

Wir arbeiten jetzt daran, die Voraussetzung der Kompaktheit des Trägers fallenzulassen.

Proposition 5.25 (Dolbeault-Lemma) Sei $g \in \mathcal{E}(\mathbb{C})$ eine unendlich oft differenzierbare Funktion. Dann gibt es eine unendlicht oft differenzierbare Funktion $f \in \mathcal{E}(\mathbb{C})$ mit

$$
\frac{\partial f}{\partial \bar{z}}=g
$$

Beweis: Wir führen die Aussage durch einen Limesprozess auf das vorige Lemma zurück. Sei $r_{n}$ eine streng monoton wachsende unbeschränkte Folge von Radien und $\Delta_{n}=\Delta\left(0, r_{n}\right)$. Sei $\psi_{n}$ eine unendlich oft differenzierbare Funktion mit $\psi_{n}=1$ auf $\Delta_{n}$ und $\psi_{n}=0$ außerhalb von $\Delta_{n+1}$. Dann können wir also das Lemma 5.24 auf $g_{n}=\psi_{n} g$ anwenden. Seien $f_{n}$ die somit erhaltenen Funktionen. Wir möchten erreichen, dass sich die Funktionen $f_{n}$ und $f_{n+1}$ auf $\Delta_{n+1}$ nur um $2^{-n}$ in der Supremumsnorm unterscheiden und modifizieren dazu induktiv $f_{n}$ durch Addition eines Polynoms in $z$. Dies ändert die $\bar{z}$-Ableitung nicht.

Am Induktionsanfang setzen wir $\widetilde{f}_{1}=f_{1}$. Seien $\widetilde{f}_{1}, \ldots, \widetilde{f_{n}}$ bereits konstruiert. Nach Konstruktion der $f_{n}$ ist $f_{n+1}-\widetilde{f_{n}}$ auf $\Delta_{n}$ holomorph, besitzt also eine Potenzreihenentwicklung. Davon nehmen wir eine polynomiale Approximation $P_{n}$, welche genügend viele Terme hat, sodass

$$
\left\|f_{n+1}-\widetilde{f_{n}}-P_{n}\right\|_{\infty}<2^{-n} \quad \text { auf } \quad \Delta_{n}
$$

Wir setzen also $\widetilde{f_{n+1}}=f_{n+1}-P_{n}$. Dann gilt immer noch $\frac{\partial}{\partial \bar{z}} \widetilde{f_{n+1}}=g_{n+1}$ auf $\Delta_{n+1}$ und nach Konstruktion

$$
\left\|\widetilde{f_{n+1}}-\widetilde{f_{n}}\right\|_{\infty}<2^{-n} \quad \text { auf } \quad \Delta_{n}
$$

Wir setzen für $z \in \Delta_{n}$

$$
f(z)=\widetilde{f_{n}}(z)+\sum_{k=n}^{\infty}\left(\widetilde{f_{k+1}}(z)-\widetilde{f}_{k}(z)\right)
$$

Dieser Limes existiert, da die Reihe gleichmäßig konvergiert. Da jede Differenz in der Reihe auf $\Delta_{n}$ eine holomorphe Funktion ist, muss die Reihe auch wieder eine auf $\Delta_{n}$ holomorphe Funktion, also insbesondere unendlich oft differenzierbar sein. Deswegen ist $f$ unendlich oft differenzierbar und

$$
\frac{\partial}{\partial \bar{z}} f=\frac{\partial}{\partial \bar{z}} \widetilde{f_{n}}=g_{n}=g
$$

auf $\Delta_{n}$. Schliesslich ist noch festzuhalten, dass die verschiedenen Definitionen von $f$ für $z \in \Delta_{n} \subset \Delta_{n+1}$ übereinstimmen.

Beweis von Satz 5.22 Sei $\mathfrak{U}$ eine beliebige Überdeckung von $X$ und $f= \left(f_{i j}\right) \in Z^{1}\left(\mathfrak{U}, \mathcal{O}_{X}\right)$. Wir folgen die Beweisprinzip von Satz 5.9 und fassen $Z^{1}\left(\mathfrak{U}, \mathcal{O}_{X}\right) \subset Z^{1}(\mathfrak{U}, \mathcal{E})$ auf. Da $H^{1}(\mathfrak{U}, \mathcal{E})=0$ gibt es $g=\left(g_{i}\right) \in C^{0}(\mathfrak{U}, \mathcal{E})$ mit $d g=f$. Da die $f_{i j}$ holomorph sind, ist $\frac{\partial}{\partial \bar{z}} f_{i j}=0$, d.h. die $\frac{\partial}{\partial \bar{z}}$-Ableitungen der $g_{i}$ stimmen auf zweifachen Durchschnitten überein. Also gibt es ein globales $h \in \mathcal{E}(X)$ mit $\left.h\right|_{U_{i}}=\frac{\partial}{\partial \bar{z}} g_{i}$. Auf dieses $h$ wenden wir das Dolbeault-Lemma Proposition 5.25 an und erhalten $g_{X} \in \mathcal{E}(X)$ mit $h=\frac{\partial}{\partial \bar{z}} g_{X}$.

Sei nun $g_{i}^{\prime}=g_{i}-\left.g_{X}\right|_{U_{i}}$. Dann ist $\frac{\partial}{\partial \bar{z}} g_{i}^{\prime}=\frac{\partial}{\partial \bar{z}} g_{i}-\left.h\right|_{U_{i}}=0$, also $g^{\prime}=\left(g_{i}^{\prime}\right) \in C^{0}\left(\mathfrak{U}, \mathcal{O}_{X}\right)$. Ausserdem ist $d\left(g^{\prime}\right)=d(g)-d\left(g_{X}\right)=d(g)=f$, was zu zeigen war.

\subsection {Garben, holomorphe Abbildungen und Kohomologie}
Bisher haben wir die Riemannsche Fläche $X$ festgehalten und Garben auf $X$ sowie deren Garbenmorphismen betrachtet. Wir betrachten nun den einfachsten Fall einer holomorphen Abbildung, nämlich die Inklusionsabbildung $i$ : $V \rightarrow X$ von einer offenen Menge $V$.

Ist $\mathcal{F}$ eine Garbe auf $X$, so definiert $\left.\mathcal{F}\right|_{V}(U)=\mathcal{F}(U \cap V)$ eine Garbe auf $V$, die wir mit $\left.\mathcal{F}\right|_{V}$ oder auch $i^{*} \mathcal{F}$ bezeichnen.

Ist $\mathfrak{U}=\left\{U_{i}, i \in I\right\}$ eine offene Überdeckung von $X$, so ist $\mathfrak{U}_{V}=\left\{U_{i} \cap V, i \in I\right\}$ eine offene Überdeckung von $V$. Die Restriktionsabbildungen ergeben also Abbildungen $i^{*}: C^{p}(\mathfrak{U}, \mathcal{F}) \rightarrow C^{p}\left(\mathfrak{U}_{V}, i^{*} \mathcal{F}\right)$, die Kozykel auf Kozykel und Koränder auf Koränder abbilden. Da dies für alle offenen Überdeckungen funktioniert und mit Einschränkungen verträglich ist, erhält man Abbildungen auf dem induktiven Limes, d.h.

$$
i^{*}: H^{p}(X, \mathcal{F}) \rightarrow H^{p}\left(V, i^{*} \mathcal{F}\right)
$$

für alle $p \in \mathbb{N}$.

\subsection {Der Endlichkeitssatz für $H^{1}\left(X, \mathcal{O}_{X}\right)$}
Wir haben im vorigen Abschnitt bereits begonnen, die erste Kohomologie der Garbe der holomorphen Funktionen auszurechnen. Sie verschwindet für Kreisscheiben, in der Tat für alle nicht-kompakten Riemannschen Flächen ([IFG81, Theorem 26.1]), und für die projektive Gerade. Für kompakte Riemannsche Flächen ist es aufgrund des folgenden Satzes eine wichtige Invariante.

Satz 5.26 Sei $X$ eine kompakte Riemannsche Fläche. Dann ist $\operatorname{dim} H^{1}\left(X, \mathcal{O}_{X}\right)< \infty$.

Definition 5.27 (Geschlecht) Ist $X$ eine kompakte Riemannsche Fläche. so wird die Dimension $g=\operatorname{dim} H^{1}\left(X, \mathcal{O}_{X}\right)$ das Geschlecht von $X$ genannt.

Im Moment ist dieser Begriff doppelt belegt und wir bezeichnen vorübergehend das so definierte Geschlecht auch als kohomologisches Geschlecht $g_{\mathrm{koh}}= \operatorname{dim} H^{1}\left(X, \mathcal{O}_{X}\right)$, um es vom topologisch definierten Geschlecht zu unterscheiden. Wir werden in Abschnit 6.2 sehen, dass beide Begriffe in der Tat übereinstimmen.

Bis wir ein Beispiel ausrechnen, bei dem diese Invariante auch tatsächlich nicht verschwindet, werden wir noch einige theoretische Resultate abhandeln. Selbst für den Torus ist das direkte Ausrechnen von der Definition weg mit der Überdeckung aus Beispiel 2.8 aufgrund der Vielzahl von Schnittmengen ziemlich lästig.

\subsection {Der Existenzsatz für Abbildungen}
Als erste Anwendung des Endlichkeitssatzes zeigen wir die Existenz von meromorphen Funktionen.

Satz 5.28 Sei $X$ eine kompakte Riemannsche Fläche und $x \in X$. Dann gibt es eine meromorphe Funktion auf $X$, die genau im Punkt $x$ einen Pol hat.

Beweis: Sei $g=\operatorname{dim} H^{1}\left(X, \mathcal{O}_{X}\right)$. Sei $\left(U_{1}, z\right)$ eine Karte einer Umgebung $U_{1}$ von $x$ mit $z(x)=0$. Sei $U_{2}=X \backslash\{x\}$. Dann ist $\mathfrak{U}=\left\{U_{1}, U_{2}\right\}$ eine offene Überdeckung von $X$. Die Funktionen $z^{-j}$ für $j \in \mathbb{N}$ sind auf $U_{1} \cap U_{2}$ holomorph, repräsentieren also ein Element in $Z^{1}\left(\mathfrak{U}, \mathcal{O}_{X}\right)$, das wir mit $z_{j}$ bezeichnen. Da $H^{1}\left(\mathfrak{U}, \mathcal{O}_{X}\right) \rightarrow H^{1}\left(X, \mathcal{O}_{X}\right)$ nach Lemma 5.5 injektiv ist, sind die Elemente $z_{1}, \ldots, z_{g+1}$ in $H^{1}\left(\mathfrak{U}, \mathcal{O}_{X}\right)$ linear abhängig. Das heisst, dass es $f=\left(f_{1}, f_{2}\right) \in C^{0}\left(\mathfrak{U}, \mathcal{O}_{X}\right)$ und komplexe Zahlen $c_{1}, \ldots, c_{g+1}$ gibt, sodass

$$
f_{2}-f_{1}=\sum_{j=1}^{g+1} c_{j} z^{-j} \quad \text { auf } \quad U_{1} \cap U_{2}
$$

gilt. Umgeschrieben bedeutet diese Gleichung, dass die Funktion definiert duch $f_{2}$ auf $U_{2}$ und $f_{1}+\sum_{j=1}^{g+1} c_{j} z^{-j}$ auf $U_{1} \backslash\{x\}$ auf $X \backslash\{x\}$ holomorph ist und bei $x$ einen Pol (der Ordnung höchstens $g+1$ ) hat.

Diese Satz besagt, dass jede Riemannsche Fläche Überlagerung der projektiven Geraden $\mathbb{P}^{1}$ ist (Nichtkonstante meromorphe Funktionen sind gerade die Überlagerungen von $\mathbb{P}^{1}$ ). Er wird uns gestatten zusammen mit dem Satz von Riemann-Hurwitz und Riemann-Roch im nächsten Abschnitt $g=g_{\text {koh }}$ zu zeigen. Er gestattet auch zu zeigen, dass kompakte Riemannsche Flächen stets algebraische Varietäten der Dimension Eins sind. Der Satz wird erst dann richtig erstaunlich, wenn man im Vergleich dazu gesehen hat, dass komplexe Mannigfaltigkeiten der Dimension zwei nicht immer eine nicht-konstante meromorphe Funktionen gestatten. In Dimension zwei bricht auch der Vergleich zwischen komplexen Mannifaltigkeiten und algebraischen Varietäten im Allgemeinen zusammen.

Statt einer festen Polstelle kann man sich auch Werte einer meromorphen Funktionen an beliebig (aber endlich) vielen Stellen wünschen.

Korollar 5.29 Sei $X$ eine kompakte Riemannsche Fläche, $x_{1}, \ldots, x_{n}$ paarweise verschiedene Punkte in $X$ und $c_{1}, \ldots, c_{n} \in \mathbb{C}$. Dann gibt es eine meromorphe Funktion $f$ auf $X$ mit $f\left(x_{i}\right)=c_{i}$ für $i=1, \ldots, n$.

Beweis: Übung.\\
Aus dem Beweis von Satz 5.28 ist offensichtlich, dass die meromorphe Funktion, die dort konstruiert wurde, eine Überlagerung vom Grad höchstens $g_{\text {koh }}+1$ ist. Daraus erhalten wir unmittelbar folgendes Korollar.

Korollar 5.30 Ist $X$ eine kompakte Riemannsche Fläche mit $g_{\mathrm{koh}}(X)=0$, so gibt es einen Isomorphismus $X \rightarrow \mathbb{P}_{\mathbb{C}}^{1}$.

\section {Divisoren und der Satz von Riemann-Roch}
Eine fundamentale Frage lautet, wie viele holomorphe Differentialformen es auf einer gegebenen Riemannschen Fläche $X$ gibt, insbesondere im Fall, dass $X$ kompakt ist. Diese Frage beantwortet man am besten, indem man gleich die allgemeinere Frage nach Differentialformen mit vorgegebenen Nullstellen an vorgegebenen Punkten beantwortet. Dies ist mit Nullstellenvielfachheit zu interpretieren und daher kommt die folgende Begriffsbildung.

Definition 6.1 (Divisor) Sei $X$ kompakte eine Riemannsche Fläche. Ein Divisor $D=\sum_{x \in X} a_{x}[x]$ ist eine formale Linearkombination von Punkten $x \in X$ mit Koeffizienten $a_{x} \in \mathbb{Z}$, welche für alle bis auf endlich viele $x \in X$ Null sind. Die Menge der Divisoren bildet eine additive Gruppe Div ( $X$ ), die freie abelsche Gruppe erzeugt von den Punkten $x \in X$.

Ein Divisor wird effektiv genannt, falls $a_{x} \geq 0$ für alle $x \in X$ ist. Die ganze Zahl

$$
\operatorname{deg}(D)=\sum_{x \in X} a_{x}
$$

heißt der Grad des Divisors.\\
Ist $f: X \longrightarrow \mathbb{P}_{\mathbb{C}}^{1}$ eine nicht-konstante holomorphe Abbildung (also eine meromorphe Funktion auf $X$ ), so sei

$$
\operatorname{div}(f)=\sum_{x \in X} \operatorname{ord}_{x}(f) \cdot[x]
$$

In diesem Fall wird $\operatorname{div}(f)$ der zu $f$ gehörige Hauptdivisor genannt.\\
Da die Gradabbildung als Funktion auf $\mathbb{P}_{\mathbb{C}}^{1}$ konstant ist, folgt, dass $\operatorname{deg}(\operatorname{div}(f))=0$ ist für alle holomorphen Abbildungen $f: X \longrightarrow \mathbb{P}_{\mathbb{C}}^{1}$. Hauptdivisoren haben also Grad Null. Für konstante holomorphe Abbildungen $f$ setzt $m a n \operatorname{div}(f)=0$.

Ist $U \subset X$ offen und $D$ ein Divisor auf $X$, so definieren wir $\left.D\right|_{U}= \sum_{x \in U} a_{x}[x]$.

Neben Hauptdivisoren sind die wichtigste Quelle von Divisoren diejenigen, die von Differentialformen kommen. Der Leser prüfe die Kartenunabhängigkeit der folgenden Definition.

\section {Definition 6.2 (Divisor einer Differentialform und kanonischer Divisor)}
Sei $\omega$ eine meromorphe Differentialform auf einer Riemannschen Fläche $X$. Ist $x \in X$ und $\omega=f d z$ in einer Karte ( $U, z$ ) um $x$, so setzen wir $\operatorname{ord}_{x}(\omega)=\operatorname{ord}_{x}(f)$. Schließlich definieren wir

$$
\operatorname{div}(\omega)=\sum_{x \in X} \operatorname{ord}_{x}(\omega)[x]
$$

Ist $\omega$ holomorph, oder äquivalent $\operatorname{div}(\omega)$ effektiv, so wird $\operatorname{div}(\omega)$ ein kanonischer Divisor genannt.

Proposition 6.3 Sei $\omega \not \equiv 0$ eine meromorphe Differentialform auf einer kompakten Riemannschen Fläche $X$ vom Geschlecht $g$. Dann ist $\operatorname{deg}(\operatorname{div}(\omega))=2 g-2$.

Beweis: Sind $\omega$ und $\omega_{2}$ zwei solche Differentialformen auf $X$. Dann ist

$$
\omega / \omega_{2}: X \longrightarrow \mathbb{P}_{\mathbb{C}}^{1}
$$

eine holomorphe Funktion, also $\operatorname{div}(\omega)=\operatorname{div}\left(\omega_{2}\right)+\operatorname{div}\left(\omega / \omega_{2}\right)$, insbesondere

$$
\operatorname{deg} \operatorname{div}(\omega)=\operatorname{deg} \operatorname{div}\left(\omega_{2}\right)
$$

Also genügt es, die Aussage für eine meromorphe Differentialform auf $X$ zu zeigen. Für $X=\mathbb{P}_{\mathbb{C}}^{1}$ nehmen wir $\omega=d z$ in einer der Karten. In der Karte $\widetilde{z}$ mit $\widetilde{z}=\frac{1}{z}$ ist $\omega=d\left(\frac{1}{z}\right)=-\frac{1}{z^{2}} d \widetilde{z}$, also $\operatorname{deg}(\operatorname{div}(\omega))=-2$, was zu zeigen war.

Für $X$ beliebig sei $f: X \longrightarrow \mathbb{P}_{\mathbb{C}}^{1}$ eine holomorphe und nicht-konstante Abbildung, welche nach Satz 5.28 existiert. Wir können durch Möbiustransformation annehmen, dass $f$ über $\infty$ unverzweigt ist und setzen $\omega=d(f)$. Dann ist lokal $\omega=f^{\prime} d z$, also falls $f(x) \neq \infty$, so ist $\operatorname{ord}_{x}(\omega)=\operatorname{ord}_{x}(f)-1$.\\
Ist $f(x)=\infty$, so folgt aus der Unverzweigtheit, dass $f$ in einer Kreisscheibenkarte ( $U, t$ ) um $x$ mit $t(x)=0$ die Form $f(t)=c_{1} t^{-1}+c_{0}+c_{1} t+\ldots$ mit $c_{1} \neq 0$ hat. Also ist die Nullstellenordnung der Ableitung bei $t=0$ gleich -2 und somit $\operatorname{ord}_{x}(\omega)=2$. Da es $\operatorname{deg}(f)$ solche Punkte gibt, folgt insgesamt

$$
\operatorname{deg} \omega=\sum_{x \in X}\left(\operatorname{ord}_{x}(f)-1\right)-2 \operatorname{deg}(f)=2 g(X)-2
$$

nach der Riemann-Hurwitz-Formel Satz 3.47

Der Satz von Riemann-Roch kann auch motiviert werden durch die Frage nach Funktionen mit gegebenem Polstellenverhalten. In erster Näherung haben wir diese Frage mit Satz 5.28 beantwortet. Für eine genauere Antwort machen wir folgende Definitionen.

Sind $D_{1}$ und $D_{2}$ zwei Divisoren, so schreiben wir $D_{1} \geqslant D_{2}$, falls $D_{1}-D_{2}$ effektiv ist.

Definitionen 6.4 Zu einem Divisor $D$ auf $X$ sei

$$
\mathcal{L}(D):=\{0\} \cup\left\{f: X \longrightarrow \mathbb{P}_{\mathbb{C}}^{1} \text { holomorph und } \operatorname{div}(f)+D \geqslant 0\right\}
$$

und

$$
\ell(D)=\operatorname{dim}_{\mathbb{C}} \mathcal{L}(D)
$$

Beispiel 6.5 Ist $X$ kompakt und $D=0$, so besteht $\mathcal{L}(D)$ nur aus den konstanten Funktionen nach dem Maximumsprinzip. In diesem Fall ist also $\ell(D)=1$. Ist $D \geqslant D^{\prime}$, so ist $\mathcal{L}\left(D^{\prime}\right)$ in $\mathcal{L}(D)$ enthalten. Ist $\operatorname{deg}(D)<0$, so folglich $\ell(D)=0$, $\operatorname{denn} \operatorname{deg}(\operatorname{div}(f))=0$ für jede holomorphe Funktion $f$.

Das Ziel dieses Abschnitts ist der folgende Satz.\\
Satz 6.6 (Riemann-Roch) Sei $\omega \not \equiv 0$ eine meromorphe Differentialform und $K= \operatorname{div}(\omega)$. Dann gilt für jeden Divisor $D$ auf $X$ :

$$
\ell(D)<+\infty \quad \text { und } \quad \ell(D)-\ell(K-D)=\operatorname{deg}(D)+1-g_{\mathrm{koh}}
$$

Ist also $\operatorname{deg}(D)>2 g_{\text {koh }}-2$ so erhält man mit $\ell(D)=\operatorname{deg}(D)+1-g_{\text {koh }}$ direkt die Antwort auf die motivierende Frage nach Funktionen mit vorgegebenem Verschwindungsverhalten. Im Allgemeinen liefert Riemann-Roch direkt eine Abschätzung $\ell(D) \geq \operatorname{deg}(D)+1-g_{\text {koh }}$, aber der Beitrag $\ell(K-D)$ ist nicht einfach zu bestimmen.

\subsection {Kohomologische Version von Riemann-Roch}
Wir wollen die Vektorräume $\mathcal{L}(D)$ nun kohomologisch interpretieren. Dazu definieren wir zwei wichtige Garben. Für einen Divisor $D$ sei die Garbe $\mathcal{O}_{X}(D)$ definiert durch

$$
\mathcal{O}_{X}(D)(U)=\left\{f: U \rightarrow \mathbb{P}_{\mathbb{C}}^{1} \text { holomorph und } \operatorname{div}(f) \geq-\left.D\right|_{U}\right\} \cup\{0\}
$$

Die Garbenaxiome prüft man wie bei $\mathcal{O}_{X}$ nach. Nach Definition ist

$$
\mathcal{L}(D)=\mathcal{O}_{X}(D)(X)=H^{0}\left(X, \mathcal{O}_{X}(D)\right)
$$

Ist $\omega$ eine meromorphe Differentialform mit $K=\operatorname{div}(\omega)$ und ist $f \in \mathcal{L}(K)$, so ist $f \omega$ eine meromorphe Differentialform ohne Pole, also eine holomorphe

Differentialform. Ist umgekehrt $\omega_{1} \in H^{0}\left(X, \Omega_{X}^{1}\right)$, also eine globale holomorphe Differentialform, so ist $f=\omega_{1} / \omega \in \mathcal{L}(K)$. Auf diese Weise haben wir einen Isomorphismus

$$
\mathcal{L}(K) \cong H^{0}\left(X, \Omega_{X}^{1}\right)
$$

von $\mathbb{C}$-Vektorräumen definiert. Wir erweitern diese Korrespondenz auf Verschwindungsbedingungen wie bei einem Divisor. Dazu definieren wir die Garbe

$$
\Omega_{X}^{1}(-D)(U)=\left\{\omega \in \Omega_{X}^{1}(U): \operatorname{div}(\omega)-D \geq 0\right\}
$$

und erhalten mit dem gleichen Argument einen Isomorphismus

$$
\mathcal{L}(K-D) \cong H^{0}\left(X, \Omega_{X}^{1}(-D)\right)
$$

(Bei dieser Definition ist nicht $D$ effektiv vorausgesetzt. Das Minuszeichen sorgt lediglich dafür, dass der Term direkt wie im Satz von Riemann-Roch aussieht.)

Der Satz von Riemann-Roch beinhaltet im Spezialfall $D=0$ auch die Antwort auf die zweite motivierende Fragestellung nach globalen Differentialformen.

Korollar 6.7 Sei $X$ eine kompakte Riemannsche Fläche. Dann hat der Vektorraum $H^{0}\left(X, \Omega_{X}^{1}\right)$ der holomorphen Differentialformen auf $X$ die Dimension $g$.

Wir bereiten nun den Beweis des Satzes von Riemann-Roch vor. Offenbar gibt es für einen effektiven Divisor $D$ einen injektiven Garbenmorphismus $\mathcal{O}_{X} \rightarrow \mathcal{O}_{X}(D)$. Wir wollen diesen zu einer exakten Sequenz ergänzen und die zugehörige Kohomologiesequenz betrachten.

Definition 6.8 (Wolkenkratzergarbe) Sei $D=\sum_{x \in X} a_{x}[x]$ ein effektiver Divisor auf $X$. Durch

$$
\mathbb{C}_{D}(U)=\oplus_{x \in U} \mathbb{C}^{a_{x}}
$$

und Restriktionsabbildungen gegeben durch die natürlichen Vergissabbildungen $\oplus_{x \in U} \mathbb{C}^{a_{x}} \rightarrow \oplus_{x \in V} \mathbb{C}^{a_{x}}$ für $V \subset U$ wird eine Garbe $\mathbb{C}_{D}$ auf $X$ beschrieben, die Wolkenkratzergarbe zu $D$ genannt wird.

Ist $f \in \mathcal{O}_{X}(D)(U)$ und $x \in U$ beliebig, so hat $f$ eine Potenzreihenentwicklung um $x$ der Form $f=\sum_{i \geq-a_{x}} c_{i}(z-x)^{i}$. Indem man $f$ auf seinen Polarteil $\sum_{i=-a_{x}}^{-1} c_{i}(z-x)^{i}$ abbildet, erhält man einen Garbenmorphismus $\mathcal{O}_{X}(D) \rightarrow \mathbb{C}_{D}$. Man rechnet nach, dass damit $0 \rightarrow \mathcal{O}_{X} \rightarrow \mathcal{O}_{X}(D) \rightarrow \mathbb{C}_{D} \rightarrow 0$ und allgemeiner für alle Divisoren $D_{1}$

$$
0 \rightarrow \mathcal{O}_{X}\left(D_{1}\right) \rightarrow \mathcal{O}_{X}\left(D+D_{1}\right) \rightarrow \mathbb{C}_{D} \rightarrow 0
$$

wobei der Pfeil nach $\mathbb{C}_{D}$ auf den um die Koeffizienten von $D_{1}$ verschobenen Polarteil abbildet, zu einer exakten Sequenz von Garben wird.

Lemma 6.9 Ist $\mathbb{C}_{D}$ eine Wolkenkratzergarbe zum effektiven Divisor $D$, so ist

$$
h^{0}\left(X, \mathbb{C}_{D}\right)=\operatorname{deg}(D) \quad \text { und } \quad h^{1}\left(X, \mathbb{C}_{D}\right)=0
$$

Beweis: Die erste Aussage ist klar und zum Beweis der zweiten Aussage starten wir mit einer Überdeckung $\mathfrak{V}$ und verfeinern sie derart zu einer Überdeckung $\mathfrak{U}$ von $X$, sodass jeder Punkt $x \in X$ mit $a_{x} \neq 0$ in nur einer Menge $U_{i}$ enthalten ist. Dann enthält also ein Durchschnitt von Mengen in $\mathfrak{U}$ keinen solchen Punkt. Ist also $f \in Z^{1}\left(\mathfrak{U}, \mathbb{C}_{D}\right)$ in Kozykel, so ist $f=0$ nach Definition.

Satz 6.10 (Riemann-Roch, kohomologische Version) Sei $D$ ein Divisor auf $X$. Dann sind $H^{0}\left(X, \mathcal{O}_{X}(D)\right)$ und $H^{1}\left(X, \mathcal{O}_{X}(D)\right)$ endichdimensionale $\mathbb{C}$ Vektorräume und es gilt

$$
h^{0}\left(X, \mathcal{O}_{X}(D)\right)-h^{1}\left(X, \mathcal{O}_{X}(D)\right)=\operatorname{deg}(D)+1-g_{\mathrm{koh}}
$$

Beweis: Für den Divisor $D=0$ folgt dies nach der Bemerkung in Beispiel 6.5 und der Definition von $g_{\text {koh }}$. Den allgemeinen Fall zeigen wir, indem wir die Situation $D^{\prime}=D+[x]$ betrachten und zeigen, dass die Behauptung für $D$ die Behauptung für $D^{\prime}$ und umgekehrt impliziert. Dazu betrachten wir die lange exakte Sequenz\\
\includegraphics[scale=0.2, center]{2025_10_06_374d4cfa8b1948eb4c41g-066}\\
zu (6.1) unter Verwendung von Lemma 6.9. Sei $V=\operatorname{Bild}\left(\beta_{0}\right)$. Dann ist die Exaktheit der langen Sequenz äquivalent zur Exaktheit der beiden folgenden kurzen exakten Sequenzen

$$
0 \rightarrow H^{0}\left(X, \mathcal{O}_{X}(D)\right) \rightarrow H^{0}\left(X, \mathcal{O}_{X}\left(D^{\prime}\right)\right) \rightarrow V \rightarrow 0
$$

und

$$
0 \rightarrow \mathbb{C} / V \rightarrow H^{1}\left(X, \mathcal{O}_{X}(D)\right) \rightarrow H^{1}\left(X, \mathcal{O}_{X}\left(D^{\prime}\right)\right) \rightarrow 0
$$

Dies bedeutet

$$
h^{0}\left(X, \mathcal{O}_{X}\left(D^{\prime}\right)\right)=\operatorname{dim}_{\mathbb{C}} V+h^{0}\left(X, \mathcal{O}_{X}(D)\right)
$$

und

$$
-h^{1}\left(X, \mathcal{O}_{X}\left(D^{\prime}\right)\right)=-h^{1}\left(X, \mathcal{O}_{X}(D)\right)+\left(1-\operatorname{dim}_{\mathbb{C}} V\right) .
$$

Durch Addition der beiden Zeilen folgt die Behauptung in beiden Fällen.\\
Daraus kann man direkt eine grobe, aber oft nützliche Abschätzung ableiten.

Korollar 6.11 Für einen beliebigen Divisor $D$ gilt $h^{0}\left(X, \mathcal{O}_{X}(D)\right) \geq \operatorname{deg}(D)+1- g_{\text {koh }}$.

\subsection {Serre-Dualität}
Für den Beweis des Riemann-Roch-Satzes 6.6 fehlt, in Anbetracht des Satzes 6.10 und der kohomologischen Interpretation von $\mathcal{L}(D)$ und $\mathcal{L}(K)$ genau noch folgende Aussage.

Satz 6.12 (Serre-Dualität) Es gibt einen Isomorphismus

$$
i_{D}: H^{0}\left(X, \Omega_{X}^{1}(-D)\right) \rightarrow \operatorname{Hom}_{\mathbb{C}}\left(H^{1}\left(X, \mathcal{O}_{X}(D), \mathbb{C}\right)\right.
$$

Insbesondere ist $h^{0}\left(X, \Omega_{X}^{1}(-D)\right)=h^{1}\left(X, \mathcal{O}_{X}(D)\right)$.\\
Korollar 6.13 Die Definitionen des Geschlechts für eine kompakte Riemannsche Fläche aus Abschnitt 3.6 und aus Abschnitt 5.7stimmen überein, d.h. $g=g_{\mathrm{koh}}$.

Beweis: Wir wenden Riemann-Roch auf $D=K$ sowie die implizit bereits eingeführten Garbenisomorphismen $\mathcal{O}_{X}(K) \cong \Omega_{X}^{1}$, allgemeiner $\mathcal{O}_{X}(K-D) \cong \Omega_{X}^{1}(-D)$ und somit $\mathcal{O}_{X} \cong \Omega_{X}^{1}(-K)$ sowie Proposition 6.3 an und erhalten

$$
h^{0}\left(X, \Omega_{X}^{1}\right)-h^{1}\left(X, \Omega_{X}^{1}\right)=(2 g-2)+1-g_{\mathrm{koh}} .
$$

Zusammen mit den Konsequenzen der Serre-Dualität $h^{0}\left(X, \Omega_{X}^{1}\right)= h^{1}\left(X, \mathcal{O}_{X}\right)=g_{\text {koh }}$ sowie $h^{1}\left(X, \Omega_{X}^{1}\right)=h^{0}\left(X, \mathcal{O}_{X}\right)=1$ folgt die Behauptung. \(\square\)

Jede Bilinearform $\langle\cdot, \cdot\rangle: V_{1} \times V_{2} \rightarrow \mathbb{C}$ liefert eine lineare Abbildung $V_{1} \rightarrow \operatorname{Hom}_{\mathbb{C}}\left(V_{2}, \mathbb{C}\right)$ ) (und umgekehrt liefert solch ein Homomorphismus eine Bilinearform). Angewandt auf die in Satz 6.12 vorkommenden Räume verwenden wir also zunächst die Bilinearform

$$
\Omega_{X}^{1}(-D)(U) \times \mathcal{O}_{X}(D)(U) \rightarrow \Omega_{X}^{1}(U), \quad(\omega, f) \mapsto f \omega
$$

für jede offene Menge $U \subset X$. Diese induziert für jede offene Überdeckung $\mathfrak{U}$ eine Bilinearform $C^{0}\left(\mathfrak{U}, \Omega_{X}^{1}(-D)\right) \times C^{1}\left(\mathfrak{U}, \mathcal{O}_{X}(D)\right) \rightarrow C^{1}\left(\mathfrak{U}, \Omega_{X}^{1}\right)$, welche $Z^{0} \times Z^{1}$ nach $Z^{1}$ abbildet. Man prüft nach, dass das Bild von $d\left(C^{0}\left(\mathfrak{U}, \mathcal{O}_{X}(D)\right)\right)$ dabei nach $d\left(C^{0}\left(\mathfrak{U}, \Omega_{X}^{1}\right)\right)$ abgebildet wird und man somit eine bilineare Abbildung

$$
H^{0}\left(X, \Omega_{X}^{1}(-D)\right) \times H^{1}\left(X, \mathcal{O}_{X}(D)\right) \rightarrow H^{1}\left(X, \Omega_{X}^{1}\right)
$$

erhält. Der erste Schritt ist also, eine natürliche Abbildung des Bildes dieser bilinearen Abbildung nach $\mathbb{C}$ zu konstruieren.

In der Garbensprache interpretiert besagt Korollar 4.7, dass die Sequenz

$$
0 \rightarrow \Omega_{X}^{1} \rightarrow \mathcal{E}_{X}^{(1,0)} \rightarrow \mathcal{E}_{X}^{2} \rightarrow 0
$$

exakt ist. Aufgrund von Proposition 5.8 wird die lange exakte Kohomologiesequenz zu einer kurzen exakten Sequenz

$$
\cdots \rightarrow H^{0}\left(X, \mathcal{E}_{X}^{(1,0)}\right) \rightarrow H^{0}\left(X, \mathcal{E}_{X}^{2}\right) \rightarrow H^{1}\left(X, \Omega_{X}^{1}\right) \rightarrow 0 .
$$

Wir erhalten also die gewünschte Abbildung durch die Definition

$$
\text { Res : } H^{1}\left(X, \Omega_{X}^{1}\right) \rightarrow \mathbb{C}, \quad \operatorname{Res}(\Omega)=\frac{1}{2 \pi i} \iint_{X} \Omega
$$

auf einem Repräsentanten $\Omega \in \mathcal{E}_{X}^{2}(X)$ von $H^{1}\left(X, \Omega_{X}^{1}\right)$. Nach dem Satz von Stokes (Satz 4.11) ist $\operatorname{Res}(d \omega)=0$, die Abbildung also in der Tat auf $H^{1}\left(X, \Omega_{X}^{1}\right)$ wohldefiniert. Der Ursprung des Begriffs Residuum wird sogleich klar.

Mit dem Ziel die Injektivität der Serre-Dualität zu zeigen, repräsentieren wir Elemente in $H^{1}\left(X, \Omega_{X}^{1}\right)$ durch sogenannte Mittag-Leffler-Koketten $\mu=\left(\mu_{i}\right)$ bezüglich einer Überdeckung $\mathfrak{U}$ von $X$, d.h. $\mu_{i}$ ist eine meromorphe Differentialform of $U_{i}$ und die Differenzen $\mu_{i}-\mu_{j}$ sind auf den Überlappungsmengen holomorph. Mit der üblichen Definition des Korandoperators ist also $d\left(\left(\mu_{i}\right)_{i \in I}\right) \in H^{1}\left(X, \Omega_{X}^{1}\right)$. Der Vorteil hiervon ist, dass man diese leicht angeben kann (siehe Lemma 6.15)). Auch kann man einer Mittag-Leffler-Kokette leicht ein Residuum zuordnen. Ist $x \in U_{i}$ so sei $\operatorname{Res}_{x} \mu_{i}$ das übliche Residuum. Ist $x \in U_{i} \cap U_{j}$, so ist $\operatorname{Res}_{x} \mu_{i}=\operatorname{Res}_{x} \mu_{j}$ nach Definition einer Mittag-LefflerKokette. Da diese Residuen nur in endlich vielen Punkten von Null verschieden sind, definieren wir mit

$$
\operatorname{Res}\left(\left(\mu_{i}\right)_{i \in I}\right)=\sum_{x \in X} \operatorname{Res}_{x} \mu_{i(x)}
$$

wobei $i(x)$ ein Index mit $x \in U_{i}$ ist. Die Arbeit besteht nun darin zu zeigen, dass die zwei Residuendefinitionen übereinstimmen.

Lemma 6.14 Ist $\left(\mu_{i}\right)_{i \in I}$ eine Mittag-Leffler-Kokette, so ist

$$
\operatorname{Res}\left(\left(\mu_{i}\right)_{i \in I}\right)=\operatorname{Res}\left(d\left(\left(\mu_{i}\right)_{i \in I}\right)\right)
$$

Beweis: Wir verfolgen die Definition des Verbindungshomomorphismus $\delta^{0}$ um eine 2 -Form $\Omega \in \mathcal{E}^{(2)}(X)$ mit $\delta^{0}(\Omega)=d(\mu)$ zu konstruieren. Da $d(\mu)$ ein 1-Kozykel von $\Omega_{X}^{1}$ und somit auch von $\mathcal{E}^{(1,0)}$ ist und $H^{1}\left(X, \mathcal{E}^{(1,0)}\right)=0$ ist, gibt es $g=\left(g_{i}\right) \in C^{0}\left(\mathfrak{U}, \mathcal{E}^{(1,0)}\right)$ mit $d g=d \mu$. Die äussere Ableitung auf $\mu_{i}-\mu_{j}$ angewandt ist Null, da diese Differenz holomorph ist. Diese stimmt wegen der vorigen Bemerkung mit der äusseren Ableitung von $g_{i}-g_{j}$ überein. Also verkleben sich diese äusseren Ableitungen zu einer globlen 2-Form $\Omega \in \mathcal{E}^{(2)}(X)$. Aus der Konstruktion des Verbindungshomomorphismus folgt nun unmittelbar $\delta^{0}(\Omega)=d(\mu)$. Es ist nun also

$$
\operatorname{Res}\left(\left(\mu_{i}\right)_{i \in I}\right)=\frac{1}{2 \pi i} \iint_{X} \Omega
$$

zu zeigen.\\
Wir zeigen zunächst, dass nur die Umgebungen der Pole von $\mu$ zum Integral beitragen. Seien $\left\{x_{1}, \ldots, x_{n}\right\}$ diese endlich vielen Pole, sei $Y=X \backslash$\\
$\left\{x_{1}, \ldots, x_{n}\right\}$. Auf $Y \cap U_{i} \cap U_{j}$ ist $g_{i}-\mu_{i}=g_{j}-\mu_{j}$. Also verkleben sich die $h_{i}=g_{i}-\mu_{i}$ zu einem globalen $h \in \mathcal{E}_{X}^{(1,0)}(Y)$ mit $\left.h\right|_{U_{i}}=h_{i}$.

Für jeden der Punkte $x_{j}$ wählen wir einen Index $i(j)$, sodass $x_{j} \in U_{i(j)}$. Um diese Punkte wählen wir Karten $\left(V_{j}, z_{j}\right)$ mit $V_{j} \subset U_{i(j)}$ und $z_{j}\left(x_{j}\right)=0$ und wir verkleinern diese Umgebungen gegebenenfalls, sodass sie paarweise disjunkt sind. Wir basteln außerdem $f_{j} \in \mathcal{E}(X)$, welche auf einer verkleinerten Umgebung $V_{j}^{\prime} \subset V_{j}$ identisch Eins sind und welche Träger in $V_{j}$ haben. Sei $f^{c}=1-\sum_{j=1}^{n} f_{j}$. Da der Träger von $f^{c}$ disjunkt zu den Polen ist, können wir den Satz von Stokes anwenden und erhalten $\iint_{X} d\left(f^{c} h\right)=0$.

Wir betrachten nun die Umgebungen $V_{j}^{\prime}$. Außerhalb von $x_{j}$ ist dort $d\left(f_{j} h\right)= d(h)=d g_{i(j)}$. Da die $g_{i}$ auf ganz $V_{j}^{\prime}$ definiert sind, gilt das auch für die Bilder unter dem Operator $d$. Wir können demnach wahlweise auf ganz $X$ über $\Omega$ oder $d\left(f_{j} h\right)$ integrieren. Also ist\\
$\iint_{X} \Omega=\sum_{j=1}^{n} \iint_{V_{j}} d\left(f_{j} h\right)=\sum_{j=1}^{n} \iint_{V_{j}} d\left(f_{j} g_{i(j)}\right)-d\left(f_{j} \mu_{i(j)}\right)=\sum_{j=1}^{n} \iint_{V_{j}}-d\left(f_{j} \mu_{i(j)}\right)$,\\
wobei wir nochmal den Satz von Stokes verwendet haben. Dieser Satz zusammen mit dem Residuensatz zeigt\\
$\iint_{V_{j}}-d\left(f_{j} \mu_{i(j)}\right)=\iint_{V_{j}^{\prime}}-d\left(f_{j} \mu_{i(j)}\right)=\iint_{V_{j}^{\prime}}-d\left(\mu_{i(j)}\right)=\int_{\partial V_{j}^{\prime}}-\mu_{i(j)}=2 \pi i \operatorname{Res}_{x_{j}} \mu_{i(j)}$.\\
Durch Addition dieser Beiträge folgt nun die Behauptung.

Lemma 6.15 Die Abbildung $i_{D}$ aus Satz 6.12 ist für alle $D$ injektiv.\\
Beweis: Wir müssen zu jedem $\omega \in H^{0}\left(X, \Omega_{X}^{1}(-D)\right)$ eine geeignete Überdeckung $\mathfrak{U}$ und einen Kozykel $f=\left(f_{i j}\right) \in C^{1}\left(\mathfrak{U}, \mathcal{O}_{X}(D)\right)$ finden, sodass $\operatorname{Res}(f \omega) \neq 0$ ist. Sei $x \in X$ ein Punkt, der nicht im Träger des Divisors $D$ liegt und $(U, z)$ eine Karte einer Umgebung von $x$ mit $z(x)=0$. In dieser Umgebung schreiben wir $\omega=f(z) d z$ mit einer holomorphen Funktion $f$. Wir können $U$ nun so verkleinern, dass $f$ außer möglicherweise bei $z=0$ keine Nullstellen hat. Sei $f_{0}=(z f(z))^{-1}$, also $f_{o}(z) \omega=d z / z$ auf $U_{0}$.

Sei $U_{1}=X \backslash\{x\}$. Dann ist $\mu=(d z / z, 0)$ eine Mittag-Leffler-Kokette bezüglich $\mathfrak{U}=\left\{U_{0}, U_{1}\right\}$ und es ist $\operatorname{Res}(\mu)=1$. Wir suchen also $\left(f_{i j}\right) \in C^{1}\left(\mathfrak{U}, \mathcal{O}_{X}(D)\right)$, sodass $f \omega=d \mu$ ist. Dies leistet offenbar der Kozykel gegeben durch $f_{01}= \left.f_{0}\right|_{U_{0} \cap U_{1}}$.

Wir bewegen uns nun in Richtung Surjektivität der Abbildungen $i_{D}$. Der Beweis läuft über den Vergleich von Divisoren $D$ und $D^{\prime}$ mit $D \geq D^{\prime}$ und wir treffen zunächst einige Vorbereitungen.

Sei $B$ ein beliebiger Divisor und $0 \neq g \in H^{0}\left(X, \mathcal{O}_{X}(B)\right)$. Dann liefert die Multiplikation $f \mapsto f g$ einen Garbenmorphismus $\mathcal{O}_{X}(D-B) \rightarrow \mathcal{O}_{X}(D)$ und somit einen Homomorphismus $\psi: H^{1}\left(X, \mathcal{O}_{X}(D-B)\right) \rightarrow H^{1}\left(X, \mathcal{O}_{X}(D)\right)$ der zugehörigen Kohomologiegruppen

\section {Lemma 6.16 Die induzierte Abbildung auf den Dualräumen}
$\psi^{\vee}: \operatorname{Hom}_{\mathbb{C}}\left(H^{1}\left(X, \mathcal{O}_{X}(D)\right), \mathbb{C}\right) \rightarrow \operatorname{Hom}_{\mathbb{C}}\left(H^{1}\left(X, \mathcal{O}_{X}(D-B)\right), \mathbb{C}\right), \quad \ell \mapsto \ell \circ \psi$ ist injektiv.

Beweis: Ist $A=\operatorname{div}(g)$, so ist per Definition $A \geq-B$ und die Multiplikation mit $g$ definiert einen Isormorphimus $\mathcal{O}_{X}(D+A) \rightarrow \mathcal{O}_{X}(D)$ und $\psi$ is die Verkettung der Inklusion $\mathcal{O}_{X}(D-B) \rightarrow \mathcal{O}_{X}(D+A)$ mit diesem Isomorphismus. Da Garbenisomorphismen auch Isomorphismen der Kohomologiegruppen induzieren, genügt es also zu zeigen, dass die Abbildung $H^{1}\left(X, \mathcal{O}_{X}(D-B)\right) \rightarrow H^{1}\left(X, \mathcal{O}_{X}(D+A)\right)$ surjektiv ist, denn dann ist die duale Abbildung injektiv. Diese Surjektivitätsaussage folgt aber direkt aus Lemma 6.9. \(\square\)

Für den Speziallfall $B \geq 0$ und $g=$ id d.h. für $D \geq D^{\prime}$ bezeichen wir diese injektive Vergleichsabbildung zwischen der ersten Kohomologiegruppen mit $c_{D^{\prime}}^{D}$. Andererseits haben wir eine Injektion $w_{D^{\prime}, D^{\prime}}: H^{0}\left(X, \Omega_{X}^{1}(-D)\right) \rightarrow H^{0}\left(X, \Omega_{X}^{1}\left(-D^{\prime}\right)\right.$ und die injektiven Abbildungen $i_{D}$ bzw. $i_{D}^{\prime}$ zwischen diesen Räumen, zusammengefasst in dem Diagramm\\
\includegraphics[scale=0.2, center]{2025_10_06_374d4cfa8b1948eb4c41g-070}

Man überprüft mit Hilfe der Definition der Dualitätsabbildungen, dass dieses Diagramm kommutiert.

Lemma 6.17 Ist $\ell \in \operatorname{Hom}_{\mathbb{C}}\left(H^{1}\left(X, \mathcal{O}_{X}(D)\right), \mathbb{C}\right)$ sowohl im Bild von $c_{D^{\prime}}^{D}$ also auch im Bild von $i_{D^{\prime}}$ so liegt es im Bild von $H^{0}\left(X, \Omega_{X}^{1}(-D)\right)$, d.h. ist

$$
\ell=c_{D^{\prime}}^{D}\left(\ell_{D}\right)=i_{D^{\prime}}(\omega),
$$

so liegt $\omega \in H^{0}\left(X, \Omega_{X}^{1}(-D)\right)$.\\
Beweis: Sei $D=\sum a_{x}[x]$ und $D^{\prime}=\sum a_{x}^{\prime}[x]$. Wir erinnern daran, dass per Definition $\ell(\xi)=\langle\omega, \xi\rangle$ für alle $\xi \in H^{1}\left(X, \mathcal{O}_{X}\left(D^{\prime}\right)\right)$ gilt.

Wenn die Behauptung falsch ist, dann gibt es ein $x \in X$ mit $\operatorname{ord}_{x}(\omega)< a_{x}$. Es gilt aber $\operatorname{ord}_{x}(\omega) \geq a_{x}^{\prime}$ nach Definition von $\omega$. Unser Ziel ist es, eine Kohomologieklasse $\xi \in H^{1}\left(X, \mathcal{O}_{X}\left(D^{\prime}\right)\right)$ zu konstruieren, die sich auf Null in $H^{1}\left(X, \mathcal{O}_{X}(D)\right)$ abbildet, sodass also $\ell_{D}(\xi)=0$ ist, sodass aber andererseits $\ell(\xi)=\langle\omega, \xi\rangle \neq 0$ gilt. Dann haben wir einen Widerspruch zu $\ell=i_{D^{\prime}}(\omega)$ und $\ell=c_{D^{\prime}}^{D}\left(\ell_{D}\right)$.

Wir verwenden dazu die gleiche Überdeckungskonstruktion wie beim Beweis der Injektivität. Sei ( $U_{0}, z$ ) eine Umgebung dieses Punktes mit $z(x)=0$ und wir schreiben $\omega=f d z$ mit einer meromorphen Funktion $f$ auf $U_{0}$. Wir\\
können $U_{0}$ so verkleinern, dass der Träger von $D$ und $D^{\prime}$ in $U_{0}$ nur aus dem Punkt $x$ besteht und ausserdem annehmen, dass in $U_{0}$ die Funktion $f$ ausser bei $x$ keine weiter Null- oder Polstelle besitzt.

Wir betrachten die Überdeckung $\mathfrak{U}=\left\{U_{0}, U_{1}=X \backslash\{x\}\right\}$ von $X$ und $f_{0}=(z f)^{-1}$. Dann ist $\operatorname{ord}_{x}\left(f_{0}\right)+a_{x}=-1-\operatorname{ord}_{x}(\omega)+a_{x} \geq 0$, also ist $\left(f_{0}, 0\right) \in C^{0}\left(X, \mathcal{O}_{X}(D)\right)$ (aber es gilt nicht $\left(f_{0}, 0\right) \in C^{0}\left(X, \mathcal{O}_{X}\left(D^{\prime}\right)\right)$ ). Der Korand hiervon liegt sowohl in $Z^{1}\left(X, \mathcal{O}_{X}(D)\right)$ als auch in $Z^{1}\left(X, \mathcal{O}_{X}\left(D^{\prime}\right)\right)$, da $U_{0} \cap U_{1}$ disjunkt zum Träger beider Divisoren ist. Sei also $\xi \in H^{1}\left(X, \mathcal{O}_{X}(D)\right)$ seine Kohomologieklasse. Nach Konstruktion ist also $i_{D^{\prime}}^{D}(\xi)=0$, wie gefordert.

Schließlich ist $\langle\omega, \xi\rangle=\operatorname{Res}(d z / z, 0)=1$, womit wir den gewünschten Widerspruch haben.\\
Beweis der Surjektivität von $i_{D}$ in Satz 6.12: Eine von Null verschiedene Linearform $\ell \in \operatorname{Hom}_{\mathbb{C}}\left(H^{1}\left(X, \mathcal{O}_{X}(D)\right), \mathbb{C}\right)$ sei vorgegeben und wir verwenden die Notationen von obigem Lemma 6.16 mit $B_{n}=n[x]$, wobei wir $n$ später geeignet wählen. Sei außerdem $D_{n}=D-B_{n}$. Im Gegensatz zum Vorgehen in diesem Lemma fixieren wir aber nun das in Element $\ell \in \operatorname{Hom}_{\mathbb{C}}\left(H^{1}\left(X, \mathcal{O}_{X}(D)\right), \mathbb{C}\right)$ und lassen $g$ variieren. Wir betrachten nun die Menge

$$
\Lambda_{n}=\left\{\ell \circ g, g \in H^{0}\left(X, \mathcal{O}_{X}\left(B_{n}\right)\right)\right\} \subset \operatorname{Hom}_{\mathbb{C}}\left(H^{1}\left(X, \mathcal{O}_{X}\left(D-B_{n}\right)\right), \mathbb{C}\right)
$$

Angenommen wir wüssten bereits, dass $\Lambda_{n} \cap \operatorname{Bild}\left(i_{D_{n}}\right) \neq 0$, d.h. es gibt ein $g$ sodass $\ell \circ g=i_{D_{n}}(\omega)$ für ein $\omega \in H^{0}\left(X, \Omega_{X}^{1}\left(-D_{n}\right)\right.$. Sei $A=\operatorname{div}(g)$ der Divisor von $g$, d.h. $1 / g \in H^{0}\left(X, \mathcal{O}_{X}(A)\right)$ und sei $D^{\prime}=D_{n}-A=D-\left(B_{n}+A\right)$. Dann ist

$$
c_{D^{\prime}}^{D}(\ell)=\frac{1}{g}(\ell \circ g)=\frac{1}{g} i_{D_{n}}(\omega)=i_{D^{\prime}}\left(\frac{1}{g} \omega\right)
$$

und wir sind in der Situation von Lemma 6.17. Dieses Lemma besagt, dass $\frac{1}{g} \omega$, welches a priori in $H^{0}\left(X, \Omega_{X}^{1}\left(-D^{\prime}\right)\right)$ liegt, in der Tat in $H^{0}\left(X, \Omega_{X}^{1}(-D)\right)$ liegt und dass $\ell=i_{D}\left(\frac{1}{g} \omega\right)$ ist. Damit ist die Surjektivität gezeigt.

Wir müssen also noch den nichtleeren Durchschnitt erzwingen. Dazu schätzen wir die Dimensionen ab. Nach Korollar 6.11 ist

$$
\operatorname{dim} \Lambda_{n}=h^{0}\left(X, \mathcal{O}_{X}\left(B_{n}\right)\right) \geq n+1-g_{\mathrm{koh}} .
$$

Mithilfe des gleichen Korollars schätzen wir

$$
\operatorname{dim} \operatorname{Bild}\left(i_{D_{n}}\right)=h^{0}\left(X, \Omega_{X}^{1}\left(-D_{n}\right)\right) \geq n-\operatorname{deg}(D)+2 g+2+1-g_{\mathrm{koh}}
$$

ab. Der Gesamtraum, in dem wir diese Untervektorräume schneiden, hat, wenn wir $n$ so groß wählen, dass $\operatorname{deg}\left(D-B_{n}\right)<0$ ist, die Dimension

$$
h^{1}\left(X, \mathcal{O}_{X}\left(D-B_{n}\right)\right)=g-1-\operatorname{deg}(D)+n
$$

Alle diese Dimensionen wachsen also linear in $n$ und der Koeffizient des linearen Terms ist Eins. Für $n$ groß genug müssen sich also $\Lambda_{n}$ und Bild $\left(i_{D_{n}}\right)$ wie gewünscht schneiden.

\section {Der Uniformisierungssatz}
Wir haben in Abschnitt 3.8 festgestellt, dass jede Riemannsche Fläche eine universelle Überlagerung besitzt. Universelle Überlagerungen sind einfach zusammenhängend und somit z.B. bei der Integration von Differentialformen gut handzuhaben.

Ziel dieses Abschnitts ist die Klassifikation aller möglichen universellen Überlagerungen von Riemannschen Flächen. Es wird sich herausstellen, dass es nur drei mögliche Typen von Überlagerungen gibt. Dieser Satz wird auch Uniformisierungssatz genannt.

Wesentliches Hilfmittel hierfür sind harmonische Abbildungen. Mit deren Hilfe erhalten wir auch einen Beweis der Existenz einer holomorphen Abbildung von jeder kompakten Riemannschen Fläche nach $\mathbb{P}_{\mathbb{C}}^{1}$ sowie einen Beweis des Satzes von Riemann-Roch. Diese Sätze haben wir bereits als Satz 5.28 und Satz 6.6 kennengelernt. Der Zugang hier kommt völlig ohne den Garbenbegriff und Kohomologie aus, benötigt jedoch mehr Techniken aus der Analysis.

\subsection {Das Dirichlet-Problem}
Am Ende dieses Abschnitts lösen wir ein Problem der partiellen Differentialgleichungen, das Dirichlet-Problem des Laplace-Operators für Gebiete in $\mathbb{C}$ mit nicht allzu bösartigem Rand. Der Grund für das Auftreten dieses Operators wird sogleich klar: Lösungen des Laplace-Operators sind genau die Realteile holomorpher Funktionen. Und diese Realteile lassen sich, als reellwerte Funktionen, flexibler handhaben als die starren, dem Identitätssatz gehorchenden holomorphen Funktionen. Um genügend Flexibilität zu haben müssen wir den Funktionsbegriff zwischendurch noch etwas aufweichen und von harmonischen zu subharmonischen Funktionen übergehen.

Definition 7.1 (harmonische Funktion) Gegeben sei ein Gebiet $U \subseteq \mathbb{C}$. Eine Funktion $u: U \rightarrow \mathbb{R}$ heißt harmonisch, falls $u \in C^{2}(U)$ und $\Delta u=0$, wobei $\Delta=\frac{\partial^{2}}{\partial x^{2}}+\frac{\partial^{2}}{\partial y^{2}}$ der Laplace-Operator ist.

Wichtigste Beispiele für diesen Begriff erhält man aus holomorphen Funktionen.

Lemma 7.2 Ist $f$ holomorph so sind $\operatorname{Re} f$ und $\operatorname{Im} f$ harmonisch.\\
Ist umgekehrt $u: \Delta \rightarrow \mathbb{R}$ eine harmonische Funktion auf der Einheitskreisscheibe, so ist $f=u+i \cdot \int_{0}^{y} \frac{\partial u}{\partial x}$ holomorph und $\operatorname{Re} f=u$.

Beweis: Aufgrund der Definition von harmonisch und der Cauchy-RiemannDifferentialgleichungen gilt $\frac{\partial \operatorname{Im} f}{\partial x}=\frac{\partial^{2} \int u}{\partial x^{2}}=-\frac{\partial^{2} \int u}{\partial y^{2}}=-\frac{\partial u}{\partial y}=-\frac{\partial \operatorname{Re} f}{\partial y} \quad$ und $\frac{\partial \operatorname{Im} f}{\partial y}=\frac{\partial u}{\partial x}=\frac{\partial \operatorname{Re} f}{\partial x}$.

Aus diesem Lemma folgt, dass eine harmonische Funktion $u$ auf $U$ automatisch unendlich oft differenzierbar ist. In der Tat können wir unendlich oft differenzierbar in einer kleinen Umgebung, also auf einer Kreisscheibe, nachprüfen. Dort ist dann nach dem obigen Lemma $u$ Realteil einer holomorphen, also unendlich oft differenzierbaren Funktion und somit gilt die Behauptung.

Proposition 7.3 Ist die Funktion $u$ harmonisch, dann gilt für jede abgeschlossene Kreisscheibe $\bar{\Delta}_{r}(z) \subseteq U$ die Mittelwertformel

$$
u(z)=\frac{1}{2 \pi} \int_{0}^{2 \pi} u\left(z+r e^{i \theta}\right) d \theta
$$

Beweis: Aus der Cauchy-Integralformel $f(z)=\frac{1}{2 \pi i} \int_{\partial \Delta} \frac{f(\zeta)}{\zeta-z} d \zeta$ für holomorphe Funktionen erhält man durch die Parametrisierung $\zeta=z+r e^{i \theta}$ des Randes der Kreisscheibe die Formel

$$
f(z)=\frac{1}{2 \pi} \int_{0}^{2 \pi} f\left(z+r e^{i \theta}\right) d \theta
$$

Ist $u$ harmonisch, also der Realteil einer solchen Funktion $f$, so folgt die Behauptung durch Nehmen des Realteils auf beiden Seiten.

Proposition 7.4 (Maximumsprinzip) Ist die Funktion $u$ harmonisch, $V \subseteq U$ ein Gebiet und gilt $u\left(x_{0}\right)=\sup _{x \in V} u(x)$, so ist $u \equiv u\left(x_{0}\right)$ auf $V$.

Beweis: Holomorphe Funktionen sind konstant oder offen. Dies gilt auch für deren Realteile, also nach obiger Bemerkung für alle harmonischen Funktionen. Die Offenheit von $u(V)$ für nicht-konstantes $u$ impliziert die Behauptung.

Das Dirichlet-Problem auf einem beschränkten Gebiet $U$ sucht zu gegebener, stetiger Funktion $\alpha: \partial U \rightarrow \mathbb{R}$ eine harmonische Funktion $u: U \rightarrow \mathbb{R}$, die sich stetig auf den Rand fortsetzen lässt und dort mit $\alpha$ übereinstimmt. Für Gebiete $U$ mit glattem Rand ist die Lösung relativ einfach. Schlüssel zur Lösung ist die Integration gegen den Poisson-Kern.

Definition 7.5 (Poisson-Kern) Der Poisson-Kern ist die Funktion

$$
P(z, \zeta)=\frac{|\zeta|^{2}-|z|^{2}}{|z-\zeta|^{2}}=\operatorname{Re} \frac{\zeta+z}{\zeta-z}
$$

Den Poisson-Kern findet man in der Literatur oft in Polarkoordinaten

$$
P\left(r e^{i \varphi}, R e^{i \Phi}\right)=\frac{R^{2}-r^{2}}{R^{2}-2 R r \cos (\Phi-\varphi)+r^{2}}
$$

Die Motivation, diesen Kern einzuführen, ist die folgende Reproduktionsformel für holomorphe Funktionen. Sie ähnelt der Cauchy-Integralformel. Der

Cauchy-Kern, die Funktion $\frac{1}{z-\zeta}$ mit der $f(\zeta)$ im Integrand multipliziert wird, ist jedoch komplexwertig. Der Poisson-Kern hingegen ist reellwertig und besser geeignet für den Übergang zum Realteil, zu einer harmonischen Funktion.

Sei $U=\Delta_{r}(0)$ eine Kreisssscheibe vom Radius $r$.\\
Lemma 7.6 Ist $f$ holomorph auf $\bar{U}$, so gilt für $z \in U$

$$
f(z)=\frac{1}{2 \pi i} \int_{\partial U} P(z, \zeta) f(\zeta) \frac{d \zeta}{\zeta}=\frac{1}{2 \pi} \int_{0}^{2 \pi} P\left(z, r e^{i \varphi}\right) f\left(r e^{i \varphi}\right) d \varphi
$$

Beweis: Wir betrachen $f(z) /\left(r^{2}-\bar{w} z\right)$ als Funktion von $z$. Für $w \in U$ ist diese auf $U$ holomorph. Daher folgt aus der Cauchy-Integralformel

$$
\frac{f(z)}{r^{2}-\bar{w} z}=\frac{1}{2 \pi i} \int_{\partial U} \frac{f(\zeta)}{r^{2}-\bar{w} \zeta} \frac{d \zeta}{\zeta-z} \quad \text { für alle } w, z \in U
$$

Nun spezialisiert man auf $w=z$ und beachtet, dass $|\zeta|^{2}=r^{2}$. Nach Durchmultiplizieren mit dem Nenner der linken Seite folgt die Behauptung.

Der Realteil davon zeigt sofort folgende Verallgemeinerung von Proposition 7.3.

Korollar 7.7 (Poissonsche Integralformel) Ist die Funktion u harmonisch auf der Kreisscheibe $U=\Delta_{r}(0)$, dann gilt

$$
u(z)=\frac{1}{2 \pi} \int_{0}^{2 \pi} P\left(z, r e^{i \varphi}\right) u\left(r e^{i \varphi}\right) d \varphi
$$

Wir behandeln nun einen ersten Fall des Dirichlet-Problems.\\
Satz 7.8 (Dirichlet-Problem für Kreisscheiben) Zur einer gegebenen stetigen Funktion $\alpha: \partial U \rightarrow \mathbb{R}$ ist

$$
u(z)=\frac{1}{2 \pi} \int_{0}^{2 \pi} P\left(z, r e^{i \varphi}\right) \alpha\left(r e^{i \varphi}\right) d \varphi
$$

harmonisch und eine Lösung des Dirichlet-Problems auf der Kreisscheibe $U$.\\
Beweis: Die Definition des Poissonkerns impliziert, dass $u$ Realteil der offensichtlich holomorphen Funktion

$$
f(z)=\frac{1}{2 \pi i} \int_{\partial U} \frac{\zeta+z}{\zeta-z} \alpha(\zeta) \frac{d \zeta}{\zeta}
$$

und somit harmonisch ist.\\
Es bleibt zu zeigen, dass obige Definition für $u$ im Inneren von $U$ und $u(z)= \alpha(z)$ für $z \in \partial U$ stetig ist. Wir zeigen Stetigkeit bei $a \in \partial U$ und durch Drehen\\
aller beteiligten Funktion können wir $a=r \in \mathbb{R}$ annehmen. Dann erhält man durch Anwenden der Poissonschen Integralformel

$$
u(z)-\alpha(a)=\frac{1}{2 \pi} \int_{0}^{2 \pi} P\left(z, a e^{i \theta}\right)\left(\alpha\left(a e^{i \theta}\right)-\alpha(a)\right) d \theta
$$

Wir zerlegen die rechte Seite in $v_{1}(z)$, das Integral von $-\varphi<\theta<\varphi$ und $v_{2}$, das Integral $\varphi<\theta<2 \pi-\varphi$. Wegen Stetigkeit der Integranden ist für $\varphi$ klein genug $\left|v_{1}(z)\right|<\varepsilon / 2$. Wähle $\delta$ so klein, dass die abgeschlossene Kreisscheibe $\bar{\Delta}_{\delta}(a)$ um $a$ und das restliche Kreissegment $K=\left\{z=a e^{i \theta}: \varphi \leq \theta \leq 2 \pi-\varphi\right\}$ disjunkt sind. Sei $d$ der Abstand von $\bar{\Delta}_{\delta}(a)$ und $K$. Dann gilt für $z \in \overline{\Delta_{\delta}(a)} \cap \Delta_{r}(0)$ und $\zeta \in K$ :

$$
P(z, \zeta)=\frac{|a|^{2}-|z|^{2}}{|\zeta-z|^{2}} \leq \frac{2 r}{d^{2}}(|a|-|z|) \leq \frac{2 r}{d^{2}}|z-a|
$$

Damit ist $\left|v_{2}(z)\right| \leq \frac{2 r \cdot M}{d^{2}} \cdot|z-a|$, wobei $M=\max \left\{|\alpha(\zeta)-\alpha(a)|: \zeta \in \partial \Delta_{r}(0)\right\}$, und folglich ist $\left|v_{2}(z)\right|$ kleiner als $\varepsilon / 2$, falls $\delta$ klein genug gewählt ist.

Wir erweitern nun den Begriff harmonisch von Gebieten in $\mathbb{C}$ zu Riemannschen Flächen. Folgende Definition ist gerechtfertigt aufgrund von Lemma 7.2 und da die Übergangsfunktionen auf Riemannschen Flächen holomorph sind.

Definition 7.9 (harmonische Funktionen auf Riemannschen Flächen) Sei $X$ eine Riemannsche Fläche. Eine Funktion $u: X \rightarrow \mathbb{R}$ heißt harmonisch, falls für einen (und damit für jeden) Atlas $\left(U_{i}, g_{i}\right)$, der $X$ überdeckt, jede der Abbildungen $u \circ g_{i}^{-1}$ harmonisch ist.

Das Dirichlet-Problem lässt sich nun für jedes Gebiet in einer Riemannschen Fläche formulieren. Bei der Lösung helfen subharmonische Funktionen, welche sich flexibler verkleben lassen als harmonische.

Definition 7.10 (subharmonische Funktionen) Eine stetige Funktion $u: X \rightarrow \mathbb{R}$ heißt subharmonisch, falls für alle Gebiete $U \subseteq X$ und alle harmonischen Funktionen $v: U \rightarrow \mathbb{R}$ die Funktion $u-v$ konstant oder ohne lokale Maxima ist.

Wir halten fest, dass offenbar harmonische Funktionen nach Proposition 7.4 subhamonisch sind. Ausserdem ist die Summe (und auch das Maximum) zweier subharmonischer Funktionen wieder subharmonisch. Auch subharmonische Funktionen kann man durch ein sub-Maximumsprinzip wie folgt charakterisieren.

Proposition 7.11 Für eine stetige Funktion $u: X \rightarrow \mathbb{R}$ sind äquivalent:\\
i) $u$ ist subharmonisch\\
ii) Für alle Karten $g: U \rightarrow \mathbb{C}$, deren Bild eine Kreisscheibe $\Delta_{r}(0)$ enthält, gilt $u \leq u_{\Delta_{g, r}}$ auf $\overline{\Delta_{r}(0)}$, wobei $u_{\Delta_{g, r}}$ die Lösung des Dirichletproblems auf $\Delta_{r}(0)$ mit Randfunktion $\left.u \circ g^{-1}\right|_{\partial \Delta_{r}(0)}$ ist.\\
iii) Für alle Karten $g: U \rightarrow \mathbb{C}$, deren Bild eine Kreisscheibe $\Delta_{r}(0)$ enthält, gilt $u \circ g^{-1}(0) \leq \int_{\partial \Delta_{r}} u \circ g^{-1}$.

Aus der Eigenschaft iii) folgt, dass auch subharmonische Funktionen dem Maximumsprinzip genügen.\\
Beweis: Wir zeigen der Reihe nach folgende Implikationen\\
i) $\Rightarrow$ ii) Die Funktion $u \circ g^{-1}-u_{\Delta_{g, r}}$ ist Null auf $\partial \Delta_{r}(0)$, subharmonisch auf $U$, also dort konstant oder ohne Maximum. Da sie am Rand von $U$ gleich null ist, muss in jedem der zwei Fälle $u-u_{\Delta_{g, r}} \leq 0$ auf $U$ gelten.\\
ii) $\Rightarrow$ iii) Nach Voraussetzung in ii) und aufgrund von Proposition 7.3 ist

$$
\left(u \circ g^{-1}\right)(0) \leq u_{\Delta_{g, r}}(0)=\int_{\partial \Delta_{r}} u_{\Delta_{g, r}}=\int_{\partial \Delta_{r}} u \circ g^{-1} .
$$

iii) $\Rightarrow$ i) Sei $v$ eine beliebige harmonische Funktion auf $U$, für $u$ gelte iii) und $u-v$ habe ein Maximum $M$ auf einer offenen Menge $U$. Die Menge $U_{M}=\{x \in U: u(x)-v(x)=M\}$ ist nicht leer und abgeschlossen. Es bleibt zu zeigen, dass sie offen ist. Für $x \in U_{M}$ und $\varepsilon$ klein genug, sodass $\Delta_{\varepsilon}(x) \subseteq U$, gilt

$$
M=(u-v)(x) \leq \int_{\partial \Delta_{\varepsilon}(x)}(u-v) \leq M
$$

da $v$ harmonisch ist und aufgrund von iii) angewandt auf die Karte $\left.g\right|_{\Delta_{\varepsilon}(x)}$. Also ist $u-v=M$ auf der gesamten Umgebung $\Delta_{\varepsilon}(x)$ von $x$ und $U_{M}$ damit offen.

Wir nennen die Karten der Riemannschen Fläche, deren Bild wie in der vorangehenden Proposition eine Kreisscheibe ist, auch kurz Kreisscheibenkarten.

Die folgende Aussage gibt eine Technik zur Konstruktion von harmonischen Funktionen aus subharmonischen Funktionen.

Proposition 7.12 (Perron-Argument) Sei $\mathcal{F} \neq \emptyset$ eine Familie von subharmonischen Funktionen auf der Riemannschen Fläche X mit folgenden Eigenschaften:\\
i) Für alle Kreisscheibenkarten $g: U \rightarrow \Delta$ und $u \in \mathcal{F}$ ist auch $u_{\Delta_{g}} \in \mathcal{F}$.\\
ii) Zu je zwei Funktionen $u, v \in \mathcal{F}$ gibt es eine Funktion $w \in \mathcal{F}$ mit $w \geq \max (u, v)$.\\
Ist $u^{*}=\sup _{\mathcal{F}} u$ auf ganz $X$ endlich, so ist $u^{*}$ harmonisch.\\
Beweis: 'Endlich' und 'harmonisch' sind lokale Eigenschaften, also genügt es, eine Karte $g: U \rightarrow \Delta \subset \mathbb{C}$ der Riemannschen Fläche $X$ zu analysieren. Sei also $\mathcal{F}$ so, dass $u^{*}$ endlich ist, sei $x_{0} \in U$. Wir wählen eine Folge $u_{n} \in \mathcal{F}$ mit $u^{*}\left(x_{0}\right)=\lim u_{n}\left(x_{0}\right)$. Wir basteln aus $u_{n}$ zuerst $\widetilde{u}_{n}=\max _{i \leq n} u_{i}$ und dann $\widehat{u_{n}}=\left(\widetilde{u}_{n}\right)_{\Delta_{g}}$ und definieren $\widehat{u}=\lim \widehat{u_{n}}$. Aufgrund von i) und ii) haben wir\\
dabei das Maximum in $x_{0}$ nicht verändert, d.h. es ist $\widehat{u}\left(x_{0}\right)=u^{*}\left(x_{0}\right)$. Der Satz 7.8 zeigt, dass $\widehat{u_{n}}$ eine wachsende Folge harmonischer Funktionen ist, die auf ganz $U$ gegen ein endliche Funktion konvergiert. Aus der Integraldarstellung in Satz 7.8 folgt, dass auch $\widehat{u}=\lim \widehat{u_{n}}$ harmonisch ist. (Dieser Schluss wird auch Harnack-Prinzip genannt).

Es bleibt zu zeigen, dass $\widehat{u}=u^{*}$ ist, was wir im Moment nur in $x_{0}$ wissen. Ist $x \in U$ ein anderer Punkt, so schreiben wir $u^{*}(x)=\lim v_{n}(x)$. Wieder können wir ohne Einschränkung $v_{u} \in \mathcal{F}$ wachsend und harmonisch annehmen. Wir können weiterhin $v_{n}$ durch $\max \left(v_{n}, u_{n}\right)$ ersetzen, anders gesagt ohne Einschränkung $v_{n} \geq u_{n}$. Sei nun $v=\lim v_{n}$. Folglich ist $\widehat{u} \leq v \leq u^{*}$. Also ist $\widehat{u}-v$ harmonisch, nicht-positiv und Null bei $x_{0}$, also identisch Null nach dem Maximumsprinzip. Daraus folgt $\widehat{u}(x)=v(x)=u^{*}(x)$, was zu zeigen war.

Zu $U \subset X$ offen und einer stetigen und beschränkten Funktion

$$
\alpha: \partial U \rightarrow \mathbb{R}
$$

findet man leicht eine Familie, die den Voraussetzungen dieser Proposition genügt, nämlich\\
$\mathcal{F}_{\alpha}=\left\{u: \bar{U} \rightarrow \mathbb{R}\right.$ stetig $, u \leq \sup _{\partial U} \alpha, u$ subharmonisch auf $U$ und $u \leq \alpha$ auf $\left.\partial U\right\}$.

Sei $u_{\alpha}=\sup _{u \in \mathcal{F}_{\alpha}} u$ die harmonische Funktion zu $\mathcal{F}_{\alpha}$ aus der Proposition 7.12, Um diese als Lösung des Dirichlet-Problems allgemein zu verwenden bleibt zu untersuchen, ob sich $u_{\alpha}$ stetig auf den Rand $\partial U$ fortsetzen lässt.

Definition 7.13 Ein Randpunkt $x_{0} \in \partial U$ heißt regulär, falls es eine Barrierenfunktion in $x_{0}$ gibt, d.h. eine Umgebung $V$ von $x_{0}$ und eine Funktion $\beta: V \longrightarrow \mathbb{R}$, die stetig auf $V \cap U$ ist, subharmonisch auf $\overline{V \cap U}$, die Null in $x_{0}$ ist und die negativ auf $\overline{V \cap U} \backslash\left\{x_{0}\right\}$ ist.

Sei beispielsweise $X \subset \mathbb{C}$ und $U \subset X$ eine offene Menge mit Randpunkt $x_{0}=0$. Ist $\mathbb{R}_{<0} \subseteq X \backslash \bar{U}$, so ist $x_{0}$ regulär, denn (in Polarkoordinaten $z=r e^{i \theta}$ )

$$
\beta(z)=-r^{1 / 2} \cos (\theta / 2)=-\operatorname{Re}(\sqrt{z})
$$

hat auf $V=\Delta_{\varepsilon}(0)$ die verlangten Eigenschaften, die die Winkel $\theta$, bei denen $\cos (\theta / 2)$ verschwindet liegen auf der negativen rellen Achse, also ausserhalb von $\bar{U}$.

Allgemeiner kann man zeigen, dass ein Randpunkt regulär ist, wenn es einen unendlich oft differenzierbaren Weg im Komplement von $\bar{U}$ gibt, der an an diesem Randpunkt endet.

Irreguläre Punkte sind also „Spitzen, die nach innen zeigen", wie in nebenstehender Abbildung dargestellt.

Wir passen nun die oberen und unteren Schranken einer Barrierefunktion so an, wie wir sie im Beweis von Satz 7.15 benötigen.

Lemma 7.14 Seien $x_{0}$ ein regulärer Randpunkt von $U$ und $V$ wie in der Definition 7.13. Zu vorgegebenen reellen Konstanten $m$ und $C$ mit $m \leq C$ gibt es eine stetige Funktion $v: \bar{U} \rightarrow \mathbb{R}$, die subharmonisch auf $U$ ist, die beschränkt durch $C$ auf $\bar{U}$ ist, und sodass $v\left(x_{0}\right)=C$ und $v \equiv m$ auf $\bar{U} \backslash V$ gelten.

Beweis: Sei $\beta$ eine Barrierenfunktion, welche negativ auf dem Kompaktum $\partial V \cap \bar{U}$ ist. Durch Reskalieren kann man $\beta<m-c$ auf $\partial V \cap \bar{U}$ erreichen. Dann leistet

$$
v=\left\{\begin{array}{lll}
\max (m, \beta+c) & \text { auf } & \bar{U} \cap V \\
m & \text { auf } & \bar{U} \backslash V
\end{array}\right.
$$

das Verlangte, da am Rand von $V$ das Maximum durch $m$ angenommen wird und $v$ somit stetig ist.

Satz 7.15 (Dirichlet-Problem auf $U$ mit regulärem Rand) Sei $\alpha: \partial U \rightarrow \mathbb{R}$ eine vorgegebene stetiger Randfunktion und $u_{\alpha}$ die harmonische Funktion auf $U$, die mittels $\mathcal{F}_{\alpha}$ aus (7.1) und dem Perron-Argument konstruiert wurde.

Ist der Randpunkt $x_{0}$ regulär, so ist $\lim _{x \rightarrow x_{0}} u_{\alpha}(x)=\alpha\left(x_{0}\right)$. Insbesondere ist $u_{\alpha}$ auf $\bar{U}$ stetig, falls jeder Punkt von $\partial U$ regulär ist.

Beweis: Seien $k \leq K$ Schranken für $\alpha$, d.h. $k \leq \alpha(x) \leq K$ für alle $x \in \partial U$. $\mathrm{Zu} \varepsilon>0$ und einem regulären Randpunkt $x_{0} \in \partial U$ wählen wir eine Umgebung von $V$ von $x_{0}$ so klein, dass es eine Barrierenfunktion auf $V$ gibt und sodass zudem $\left|\alpha\left(x_{0}\right)-\alpha(\alpha)\right|<\varepsilon$ für alle $x \in V \cap \partial U$ ist.

Das Lemma 7.14 angewandt auf $m=k-\varepsilon$ und $C=\alpha\left(x_{0}\right)-\varepsilon$ gibt eine Funktion $v \in \mathcal{F}_{\alpha}$ und folglich ist

$$
\liminf _{x \in \bar{U}, x \rightarrow x_{0}} u_{\alpha}(x) \geq v\left(x_{0}\right)=\alpha\left(x_{0}\right)-\varepsilon
$$

Das Lemma 7.14 angewandt auf $m=-K$ und $C=-\alpha\left(x_{0}\right)$ liefert eine subharmonische Funktion $w$. Für $u \in \mathcal{F}_{\alpha}$ beliebig gilt auf $\partial U \cap V$ nach Definition $u(x) \leq \alpha(x) \leq \alpha\left(x_{0}\right)+\varepsilon$. Also gilt $u_{\alpha}+w<\varepsilon$ am Rand und nach dem Maximumsprinzip für subharmonische Funktionen auf ganz $\bar{U} \cap V$. Wie in Lemma 7.14 bemerkt, gilt $w \leq-K$ auf dem ganzen Rand von $U$ und somit gilt dort die Abschätzung $u+w \leq K-K=0<\varepsilon$. Nach dem Maximumsprinzip gilt dies auch im Inneren. Folglich ist für alle $x \in U$

$$
u_{\alpha}(x) \leq \varepsilon-w(x) \leq \alpha\left(x_{0}\right)+\varepsilon
$$

Beide Abschätzungen zusammen ergeben die Behauptung indem man $\varepsilon$ gegen Null gehen lässt.

\subsection {Greensche Funktionen}
Definition 7.16 Eine Riemannsche Fläche $X$ heißt reich, falls es auf $X$ eine nichtkonstante, beschränkte subharmonische Funktion gibt. Andernfalls heißt $X$ arm.

Reiche Flächen werden manchmal auch hyperbolisch genannt. Wir vermeiden diesen Begriff, da er auch mit anderer Bedeutung verwendet wird.

Lemma 7.17 Sei $X$ reich, $U \subseteq X$ ein Gebiet mit regulärem Rand, dessen Komplement $K=X \backslash U$ nichtleer und kompakt ist. Dann gibt es eine stetige Funktion $w: \bar{U} \rightarrow \mathbb{R}$ mit $w \equiv 1$ auf $\partial U, w$ harmonisch und $0<w<1$ auf $U$.

Beweis: Sei $v=-u$ mit $u$ aus der Definition 7.16. Durch Addition einer Konstanten können wir $\min _{K} v=1$ erreichen. Es gibt einen Punkt $x \in U$ mit $v(x)<1$, da $u$ kein Maximum auf $X$ annimmt und somit $v$ kein Minimum auf $X$ annimmt. Wenn wir $v_{1}:=\min (1, v)$ setzen, dann ist $-v_{1}$ eine subharmonische Funktion und $v_{1}$ ist identisch gleich Eins auf $K$. Wir setzen

$$
\mathcal{F}=\left\{\widetilde{v}: \bar{U} \rightarrow \mathbb{R} \text { stetig, subharmonisch und } \widetilde{v} \leq v_{1} \text { auf } U\right\}
$$

Wir wollen das Perron-Argument auf $\mathcal{F}$ anwenden. Die Eigenschaft ii) ist offensichtlich erfüllt. Zum Nachprüfen von i) sei $g: W \rightarrow \Delta$ eine Kreisscheibenkarte und $\widetilde{v} \in \mathcal{F}$. Wir wollen zeigen, dass auch $\widetilde{v}_{\Delta_{g}} \in \mathcal{F}$ ist. Auf $\partial W$ gilt $\widetilde{v}_{\Delta_{g}}-v_{1}=\widetilde{v}-v_{1} \leq 0$ und da $\widetilde{v}_{\Delta_{g}}-v_{1}$ subharmonisch ist, gilt die Ungleichung auf ganz $U$, was die Behauptung zeigt.

Sei also $w=\sup _{\widetilde{v} \in \mathcal{F}} \widetilde{v}$. Wir wollen zeigen, dass sie die gewünschten Eigenschaften hat. Sei $\widetilde{U}$ eine endliche Vereinigung von Kreisscheibenkarten, die $K$ überdecken. Dann ist $K \subseteq \widetilde{U} \varsubsetneqq X$ und $\widetilde{U}$ und somit aufgrund der Voraussetzung auch $\widetilde{U} \backslash K$ hat regulären Rand. Sei $\widehat{v}$ eine Lösung des Dirichlet-Problems auf $\widetilde{U} \backslash K$ mit den Randwerten 1 auf $\partial K$ und 0 auf $\partial \widetilde{U}$. Diese setzen wir mit Null auf $\bar{U} \backslash \widetilde{U}$ fort und Dann ist $\widehat{v}-v_{1}$ subharmonisch, $\leq 0$ auf $X \backslash \widetilde{U}$, also auch $\leq 0$ auf $\widetilde{U}$ nach dem Maximumsprinzip. Also ist $\widehat{v} \in \mathcal{F}$ und folglich $\widehat{v} \leq \omega \leq v_{1}$.\\
Damit ist $0 \leq w \leq 1$ auf $\bar{U}$ und ist $w \equiv 1$ auf $\partial K$. Da $w(x)<1$ im speziellen Punkt $x$ ist, kann $w$ nicht konstant sein und daher ist $0<w<1$ auf $U$, da $w$ weder Minimum noch Maximum auf $U$ annimmt.

Definition 7.18 (Greensche Funktion) Eine Greensche Funktion auf $X \mathrm{im}$ Punkt $x$ ist eine positive harmonische Funktion $u_{x}: X \backslash\{x\} \rightarrow \mathbb{R}_{>0}$, sodass $u_{x}(z)+\log |z|$ harmonisch in einer Umgebung von $x$ ist und sodass $u_{x} \leq \widetilde{u}$ gilt für jede positive harmonische Funktion $\widetilde{u}: X \backslash\{x\} \rightarrow \mathbb{R}_{>0}$, die auch die Eigenschaft hat, dass $\widetilde{u}(z)+\log |z|$ harmonisch in einer Umgebung von $x$ ist.

Dabei ist $\log |z|$ bezüglich einer Karte $g: U(x) \rightarrow \mathbb{C}$ mit $g(x)=0$ definiert.

Übungen: Die Eigenschaft, dass $g(z)+\log |z|$ harmonisch ist, hängt nicht von der Kartenwahl bei $x$ ab. Die Minimalität macht die Greensche Funktion im Punkt $x$ eindeutig, falls sie existiert.

Satz 7.19 Die folgenden Eigenschaften sind äquivalent:\\
i) Die Riemannsche Fläche $X$ ist reich.\\
ii) Es gibt einen Punkt $x \in X$, sodass es eine Green'sche Funktion $u_{x}$ in $x$ gibt.\\
iii) Zu jedem Punkt $x \in X$ gibt es eine Green'sche Funktion $u_{x}$.

Beweis: Sei $u_{x}$ die Greensche Funktion im Punkt $x$. Sei $m>0$ und $u=-\min \left(m, u_{x}\right)$. Die Funktion $u$ ist subharmonisch auf $X$, denn in einer Umgebung von $x$ ist $u=-m$ und in allen anderen Punkten ist $u$ endlich und harmonisch. Die Minimalität von $u_{x}$ impliziert, dass $u_{x}-\frac{m}{2}$ nicht überall positiv ist. Also ist $u$ mancherorts grrößer als $-m / 2$, auf jeden Fall nicht konstant.\\
Sei umgekehrt $x \in X$ beliebig und $g: U \rightarrow \mathbb{C}$ eine Karte mit $g(x)=0$. Sei

$$
\mathcal{F}=\{v: X \rightarrow \mathbb{R} \cup \infty, \text { subharmonisch und } \geqslant 0 \text { auf } X \backslash\{x\}, \text { mit kompaktem }
$$

$$
\text { Träger, } \left.\quad v \circ g^{-1}(z)+\log |z| \text { ist subharomonisch auf }|z|<1\right\} \text {. }
$$

Wir prüfen die Voraussetzungen für das Perron-Lemma:\\
Die Menge $\mathcal{F}$ ist nicht leer, denn die Funktion $v(z)=-\log |z|$ auf $|z|<1$ und $v=0$ außerhalb von $|z|<1$ liegt in $\mathcal{F}$. Die Maxima zweier Funktionen aus $\mathcal{F}$ und die Funktionen $v_{\Delta_{g}}$ zu $v \in \mathcal{F}$ liegen offenbar auch in $\mathcal{F}$ und damit ist i) und ii) erfüllt.\\
Weiterhin ist zu zeigen, dass $v^{*}:=\sup _{v \in \mathcal{F}} v$ auf $X \backslash\{x\}$ endlich ist. Dazu verwenden wir, dass $X$ reich ist in Form der Funktionen aus Lemma 7.17 angewandt auf $K_{r}=\{|z|<r\}$ mit $0<r<1$ und $U_{r}=X \backslash K_{r}$. Wir bezeichnen die so erhaltene Funktion mit $w_{r}$. Zu $v \in \mathcal{F}$ sei $v_{r}:=\max _{|z|=r} v(z)$. Wir behaupten, dass

$$
v_{r}(z) \cdot w_{r}(z)-v(z) \geqslant 0 \quad \text { außerhalb von } \quad K_{r}
$$

Die Aussage ist offenbar richtig auf $\partial K_{r}$ und am Rand des Trägers von $v$. Da die linke Seite das Negative einer subharmonischen Funktion ist, trifft die Aussage überall außerhalb von $K_{r}$ zu. Aus dieser Gleichung und der Subharmonizität von $v+\log |z|$ folgt

$$
v_{r}+\log r \leq v_{1} \leq v_{r} \cdot \lambda_{r} \quad \text { mit } \quad \lambda_{r}=\max _{|z|=1} \omega_{r}(z) .
$$

Folglich ist

$$
v(z) \leq \frac{\log r}{\lambda_{r}-1} \quad \text { auf } \quad|z|=r
$$

und mit dem gleichen Schluß wie oben gilt dies aufgrund des kompakten Trägers auch für $|z| \geqslant r$. Mit dieser universellen Schranke ist auch das Supremum aller solchen $u$ beschränkt.

Es bleibt zu zeigen, dass die so konstruierte Funktion $v^{*}$ die Greensche Funktion im Punkt $x$ ist. Nach obiger Abschätzung ist

$$
v(z)+\log (z) \leq \frac{\log r}{\lambda_{r}-1}+\log (r) \quad \text { auf } \quad|z|=r
$$

also ist die auf $X \backslash\{x\}$ harmonische Funktion $v^{*}(z)+\log |z|$ in einer Umgebung von $x$ beschränkt, also harmonisch wie man mit Hilfe der Darstellung von harmonischen Funktionen als Realteil holomorpher Funktionen und dem Riemannschen Fortsetzungssatz leicht sieht. Weiter ist $u^{*} \geqslant 0$ nach Definition von $\mathcal{F}$ und $v^{*} \geqslant-\log |z|$, da $-\log |z| \in \mathcal{F}$, also ist $v^{*}>0$.

Bleibt also die Minimalität zu zeigen. Ist $\widetilde{u}$ ein weiterer Kandidat und $v \in \mathcal{F}$, so ist $\widetilde{u}-v \geqslant 0$ außerhalb von supp $v$ und, da $v-\widetilde{u}$ subharmonisch ist, auch in supp $v$. Daraus folgt $\widetilde{u} \geqslant \sup _{v \in \mathcal{F}} v=u^{*}$.

\subsection {Der Existenzsatz für Abbildungen}
Mit diesen Vorbereitungen können wir nun die zentralen Sätze formulieren und beweisen. Der Existenzsatz für Abbildungen ist uns bereits als Satz 5.28 begegnet. Der Beweis hier verwendet keine der Methoden aus Abschnitt 5 . insbesondere weder Garbenkohomologie noch den Endlichkeitssatz 5.26.

Der folgende Satz ist der Schlüssel sowohl zum Existenzsatz, also auch zum Beweis des Satzes von Riemann-Roch weiter unten.

Satz 7.20 Sei $X$ arm, $x \in U \subseteq X$ und $f: U \backslash\{x\} \longrightarrow \mathbb{C}$ holomorph. Dann gibt es genau eine Funktion $u: X \backslash\{x\} \longrightarrow \mathbb{R}$, die harmonisch ist, die im Komplement jeder Umgebung von $x$ beschränkt ist und sodass $u-\operatorname{Re}(f)$ harmonisch auf $U$ und Null in $x$ ist.

Mit Hilfe dieses Satzes beweisen wir nun den Existenzsatz im armen und reichen Fall simultan.

Satz 7.21 Sei $X$ eine Riemannsche Fläche und $x_{1} \neq x_{2}$ Punkte auf $X$. Dann gibt es eine holomorphe Funktion $f: X \longrightarrow \mathbb{P}_{\mathbb{C}}^{1}$ mit $f\left(x_{1}\right)=0$ und $f\left(x_{2}\right)=\infty$.

Beweis von Satz 7.21, Ist $X$ reich, so sei $u_{i}$ die Greensche Funktion im Punkt $x_{i}$. Ist $X$ arm, so sei $g_{i}: U_{i} \rightarrow \Delta \subseteq \mathbb{C}$ eine Kreisscheibenkarte einer Umgebung von $x_{i}$. Wir wenden den Satz 7.20 auf $U_{i}$ und $f_{i}=\frac{1}{z} \circ g_{i}$ an, um eine harmonische Funktion $u_{i}$ zu erhalten.\\
Wir schreiben $u_{x}$ und $u_{y}$ für die partiellen Ableitungen einer Funktion $\mathbb{C} \rightarrow \mathbb{R}$. Sei nun

$$
f=\frac{u_{2, x}-i u_{2, y}}{u_{1, x}-i u_{1, x}}
$$

Außerhalb von $x_{1}$ und $x_{2}$ sind Zähler und Nenner holomorph, da $u_{1}$ und $u_{2}$ harmonisch sind und die Cauchy-Riemann-Differentialgleichungen erfüllt\\
sind. Wir untersuchen eine Umgebung von $x_{1}$. Dort ist der $u_{2, x}-i u_{2, y} \neq 0$ und $u_{1, x}-i u_{1, y}$ verhält sich bis auf eine beschränkte Funktion wie $\left(\frac{1}{z}\right)^{\prime}=-\frac{1}{z^{2}}$ im armen Fall und $\log (z)^{\prime}=\frac{1}{z} \mathrm{im}$ reichen Fall. Der Kehrwert hiervon geht für $x \rightarrow x_{1}$ also gegen Null. Folglich läßt sich $f$ stetig - und nach dem Riemannschen Hebbarkeitssatz holomorph - nach 0 mit $f\left(x_{1}\right)=0$ fortsetzen. Analoges gilt mit $f\left(x_{2}\right)=\infty$.

Für Satz 7.20 führen wir zunächst folgenden Begriff ein.\\
Definition 7.22 (Hodge-Stern) Ist $\omega=\alpha d x+\beta d y$ eine Differentialform, so setzen wir $^{*} \omega=-\beta d x+\alpha d y$.

Dieser Operator hat im Kontext von harmonischen Funktionen seine Berechtigung, da $\Delta u=d^{*} d u$ gilt. Daraus leiten wir einen einfache Integralformel ab.

Lemma 7.23 Sind $u, v: U \longrightarrow \mathbb{R}$ zweimal stetig differenzierbare Funktionen und $U \subset X$ relative kompakt und mit stückweise $C^{\infty}$-Rand, so gilt

$$
\int_{\partial D} u^{*} d v-v^{*} d u=\iint_{U} u \Delta v-v \Delta u
$$

Beweis: Setze $\omega=u^{*} d v-v^{*} d u$. Dann folgt die Behauptung aus dem Satz von Stokes und der oben genannten Faktorisierung des Laplace-Operators.

Lemma 7.24 Sei $X$ eine arme Riemannsche Fläche, sei $U \subset X$ eine offen Menge mit Kreisscheibenkarte $g: U \rightarrow \Delta(1,0) \subseteq \mathbb{C}$ und für $r<1$ sei $u: X \backslash g^{-1}(\{|z|<r\}) \rightarrow \mathbb{R}$ harmonisch und beschränkt. Dann ist $\int_{\partial U}{ }^{*} d u=0$.

Der Satz von Stokes und $\Delta u=d d^{*} u=0$ genügen als Beweis nicht, da $X \backslash U$ nicht notwendig relativ kompakt ist.\\
Beweis: Die Addition einer Konstanten erlaubt $0 \leq u \leq M$ anzunehmen. Sei $\mathcal{D}$ die Menge aller Gebiete $\hat{U}$ mit $U$ relative kompakt in $\hat{U}$ und $\hat{U}$ relative kompakt in $X$ und sodass $\hat{U}$ analytischen Rand hat. $\mathrm{Zu} \hat{U} \in \mathcal{D}$ sei

$$
u_{\hat{U}}=\left\{\begin{array}{rcl}
0 & \text { auf } & \partial \hat{U} \\
u & \text { auf } & |z|=r \\
\text { Lösung des Dirichlet-Problems } & \text { auf } & \hat{U} \cap\{|z| \geqslant r\} \\
0 & \text { auf } & X \backslash \hat{U}
\end{array}\right.
$$

Nach dem Perron-Argument ist $u^{*}=\sup _{\hat{U} \in \mathcal{D}} u_{\hat{U}}$ harmonisch auf $|z| \geqslant r$. Weiterhin ist $0 \leq u-u^{*} \leq M$ auf $\{|z| \geqslant r\}$ und $u-u^{*}=0$ auf $|z|=r$. Wir wollen zeigen, dass $u-u^{*}$ auf $|z| \geqslant r$ sogar identisch verschwindet. Zu gegebenem $0<\varepsilon<M$ sei

$$
u_{\varepsilon}=\left\{\begin{array}{lll}
\max \left(u-u^{*}, \varepsilon\right)-M & \text { auf } & |z|>r \\
\varepsilon-M & \text { auf } & |z| \leqslant r
\end{array}\right.
$$

Diese Funktion ist subharmonisch und $u_{\varepsilon}<0$ auf ganz $X$. Da $X$ arm ist, muss $u_{\varepsilon}$ konstant sein. Also ist $u-u^{*} \leqslant \varepsilon$ auf $|z|>r$. Da $\varepsilon$ beliebig ist, folgt die Zwischenbehauptung. Wir wenden die gleiche Argumentkette auf $\mathcal{D}$ und auf $v_{\hat{U}}$ definiert als die Lösung des Dirichlet-Problems mit $v_{\hat{U}} \equiv 1$ auf $|z|=1$ und sonst wie oben an. Damit erhalten wir, dass $v^{*}=\sup _{\hat{U} \in \infty} v_{\hat{U}}$ harmonisch ist und $v^{*} \equiv 1$ auf $|z| \geqslant r$.\\
Wir schreiben $u=\lim u_{j}$ und $v=\lim v_{j}$ mit $u_{j}=u_{\hat{U}_{j}}$ und $v_{j}=v_{\hat{U}_{j}}$. Dabei können wir in der Tat die gleichen Gebiete $\hat{U}_{j}$ für beide Folgen ohne Einschränkung wählen, da für $\hat{U}_{j} \subseteq \hat{V}_{j} \in \mathcal{D}$ gilt: $v_{\hat{U}_{j}} \leqslant v_{\hat{V}_{j}}$. Dazu erinnern wir uns an die Lösung des Dirichlet-Problems als Supremum über Funktionen in einer Familie und diese Funktionsfamilien sind für $\hat{U}_{j} \subseteq \hat{V}_{j}$ ineinander enthalten. Die Konvergenz der Näherungen für $u$ und $v$ ist gleichmäßig auf jedem Kompaktum aufgrund der Integraldarstellung für harmonische Funktionen.\\
Da $u_{j}=v_{j}=0$ auf dem Rand von $\hat{U}_{j}$, folgt

$$
\int_{\partial U}^{*} d u=\lim _{j \rightarrow \infty} \int_{\partial U} v_{j}^{*} d u_{j}-u_{j}^{*} d v_{j}=-\lim _{j \rightarrow \infty} \int_{\partial\left(\hat{U}_{j} \backslash U\right)} v_{j}^{*} d u_{j}-u_{j}^{*} d v_{j}
$$

Da $u_{j}$ und $v_{j}$ harmonisch sind, verschwindet dieses Integral nach dem Korollar zur Stokes Formel.

Beweis des Satzes 7.20. Zum Beweis der Eindeutigkeit nehmen wir an, dass $u_{1}$ und $u_{2}$ zwei solche Funktionen sind. Dann ist $u_{1}-u_{2}$ beschränkt auf ganz $X$ und Null in $x$. Da $X$ arm ist, muss $u_{1}=u_{2}$ sein.

Zum Beweis der Existenz lösen wir zunächst das Dirichlet-Problem zu $\operatorname{Re}(f)$ auf einer Karte bei $x$. Genauer sei $z$ eine Karte bei $x$ mit $\{|z| \leq 2\} \subseteq U$. Für $0<\rho<1$ sei $u_{\rho}$ die Lösung des Dirichlet-Problems auf $|z| \geqslant \rho$ mit $u_{\rho}=\operatorname{Re}(f) \operatorname{auf}|z|=\rho$.\\
Untenstehendes Lemma besagt, dass für relativ kleine Radien $0<\rho<r<\frac{1}{20}$ die Maxima der $u_{\rho}$ außerhalb von $\Delta_{r}$ unabhängig von $\rho$ beschränkt sind. Wir wählen also $r_{1}<1 / 20$ und unter den $u_{\rho}$ mit $\rho<r_{1}$ eine Teilfolge $u_{j}$ die auf $\left\{r_{1} \leqslant|z| \leqslant 1\right\}$ gleichmäßig konvergiert.\\
Wir definieren $r_{n}=\frac{r_{1}}{n}$ und wählen aus $u_{j}$ sukzessive Teilfolgen $u_{j}^{(n)}$ aus, die auf $\left\{r_{n} \leqslant|z| \leqslant 1\right\}$ gleichmäßig konvergieren. Die Diagonalfolge $u_{j}^{(j)}$ konvergiert gleichmäßig auf jedem Kompaktum, das $x$ nicht enthält, denn nach dem Maximumsprinzip setzt sich diese Eigenschaft auch außerhalb von $\{|z| \leqslant 1\}$ fort. Sei $u=\lim u_{j}^{j}$. Diese Funktion ist harmonisch und auf dem Komplement jeder Umgebung von $x$ beschränkt. Es bleibt noch zu zeigen, dass $u-\operatorname{Re}(f)$ sich durch Null zu einer harmonischen Funktion in $x$ fortsetzen läßt.

Lemma 7.25 Sei $0<r<1 / 20$. Dann gibt es $c(r)$, sodass für alle $\rho<r$ gilt:

$$
\max _{|z| \geqslant r}\left|u_{\rho}(z)\right| \leqslant c(r)
$$

Beweis: Es genügt nach dem Maximumsprinzip das Maximum bei $|z|=r$ abzuschätzen. Für jedes $\left|z_{0}\right| \in(\rho, 2)$ ist nach dem Vorbereitungslemma

$$
\int_{|z|=z_{0}} * d u_{\rho}=0
$$

also für $z_{0}$ mit $\rho<\left|z_{0}\right|<2$ ist

$$
F_{\rho}(z)=u_{\rho}(z)+i \cdot \int_{z_{0}}^{z} d u_{\rho}
$$

wohldefiniert, da wegunabhängig, analytisch, wie man mit den Cauchy-Riemann-Differentialgleichungen leicht nachprüft. Eine analytische Funktion auf einem Kreisring hat eine Laurent-Entwicklung im Nullpunkt, also $F_{\rho}-f=\sum_{n \in \mathbb{Z}} c_{n} \cdot z^{n}$, d.h.

$$
\left(u_{\rho}-\operatorname{Re} f\right)\left(t e^{i \theta}\right)=\sum_{n \in \mathbb{Z}}\left(\alpha_{n} \cos n \theta+\beta_{n} \sin n \theta\right) t^{n}
$$

wobei die $\alpha_{n}=\operatorname{Re}\left(c_{n}\right)$ und $\beta_{n}=\operatorname{Im}\left(c_{n}\right)$ nicht von $\rho$ abhängen.\\
Durch Integration dieser Fourier-Reihe gegen $\cos k \theta$ bzw. gegen $\sin k \theta$ erhält man die $\alpha_{k}$ bzw. $\beta_{k}$, d.h. für $t \in(\rho, 2)$ gilt,

$$
\begin{aligned}
& \frac{1}{\pi} \int_{0}^{2 \pi}\left(u_{\rho}-\operatorname{Re} f\right)\left(t e^{i \theta}\right) \cos k \theta d \theta=\alpha_{k} t^{k}+\alpha_{-k} t^{-k} \\
& \frac{1}{\pi} \int_{0}^{2 \pi}\left(u_{\rho}-\operatorname{Re} f\right)\left(t e^{i \theta}\right) \sin k \theta d \theta=\beta_{k} t^{k}+\beta_{-} t^{-k}
\end{aligned}
$$

Aus $t=\rho$ folgt $\alpha_{-k}=\rho^{-2 k} \alpha_{k}$ und aus $t=1$ folgt

$$
\left|\alpha_{k}\right|\left(1-\rho^{2 k}\right) \leqslant 2 M_{\rho} \quad \text { und } \quad\left|\beta_{K}\right|\left(1-\rho^{2 k}\right) \leqslant 2 M_{\rho}
$$

wobei

$$
M_{\rho}=\max _{|z|=1}\left|u_{\rho}\right|+\max _{|z|=1}|\operatorname{Re} f| .
$$

Folglich erreichen wir für $\rho<1 / 2$ die Abschätzung

$$
\begin{aligned}
\max _{|z| \geqslant r}\left|u_{\rho}(z)\right| & \leqslant \max _{|z|=r} \operatorname{Re} f+4 M_{\rho} \sum_{n=0}^{\infty}\left(r^{n}+\rho^{2 n} r^{-n}\right) \\
& \leqslant \max _{|z|=1} \operatorname{Re} f+4 M_{\rho} \frac{2 r}{1-r}
\end{aligned}
$$

Wir müssen noch eine Schranke für $M_{\rho}$ unabhängig von $\rho$, nur abhängig von $r$ finden. Da $u_{\rho}$ harmonisch ist, folgt mit $r<1 / 20$

$$
\max _{|z|=1}\left|u_{\rho}\right| \leqslant \max _{|z|=r}|\operatorname{Re} f|+8 M_{\rho} \frac{r}{1-r} \leqslant \max _{|z|=r}|\operatorname{Re} f|+\frac{1}{2} M_{\rho}
$$

also

$$
\frac{1}{2} M_{\rho} \leqslant \max _{|z|=r}|\operatorname{Re} f|+\max _{|z|=1}|\operatorname{Re} f|
$$

Beweis des Satzes 7.20, letzter Teil: Wir betrachten die Fourierentwicklung des Realteils von $F_{\rho}-f$. Sei diese

$$
(u-\operatorname{Re} f)\left(t e^{i \theta}\right)=\sum_{n=1}^{\infty}\left(a_{u} \cos n \theta+b_{n} \sin n \theta\right) t^{n}
$$

Mit $\alpha_{n}(\rho)$ und $\beta_{u}(\rho)$, wie im Lemma 7.25zu $u_{\rho}$ definiert, gilt

$$
a_{n}=\lim _{\rho \rightarrow 0} \alpha_{n}(\rho) \leqslant 4 \cdot\left(\max _{|z|=1}|u|+\max _{|z|=1}|\operatorname{Ref}|\right)=: C
$$

und die gleiche Abschätzung für $b_{n}$. Also folgt

$$
\left|(u-\operatorname{Re} f)\left(t e^{i \theta}\right)\right| \leqslant \sum_{n=1}^{\infty} 2 C|t|^{n}=2 C \frac{t}{1-t} \longrightarrow 0 \text { für } t \rightarrow 0
$$

Die Fortsetzbarkeit der harmonischen Funktion folgt nun aus dem Riemannschen Fortsetzungssatz und Lemma 7.2

\subsection {Der Uniformisierungssatz}
Ziel dieses Abschnittes ist es zu zeigen, dass bis auf Biholomorphie $\Delta, \mathbb{C}$ und $\mathbb{P}_{\mathbb{C}}^{1}$ die einzigen einfach zusammenhängenden Riemannschen Flächen sind.

Satz 7.26 Ist $X$ eine reiche, einfach zusammenhängende Riemannsche Fläche, so ist $X$ isomorph zur Einheitskreisscheibe $\Delta$.

Beweis: Sei $x \in X$ und $g_{x}$ die Greensche Funktion bei $x$. Zunächst zeigen wir, dass $g_{x}=-\log \left|f_{x}\right|$ für eine holomorphe Funktion $f_{x}: X \longrightarrow \mathbb{P}_{\mathbb{C}}^{1}$ ist. Da $X$ einfach zusammenhängend ist, müssen wir dies nur lokal zeigen. Außerhalb von $x$ ist $g_{x}$ harmonisch, also gibt es eine Funktion $\widetilde{f}_{x}$, sodass

$$
\begin{aligned}
g_{x} & =\operatorname{Re}\left(\widetilde{f}_{x}\right)=\operatorname{Re}\left(\log \exp \left(\widetilde{f}_{x}\right)\right)=\frac{1}{2}\left(\log \left(\exp \left(\widetilde{f}_{x}\right)\right)+\log \overline{\left(\exp \left(\widetilde{f}_{x}\right)\right)}\right) \\
& =\frac{1}{2} \log \left|\exp \left(\widetilde{f}_{x}\right)\right|^{2}=-\log \left|\exp \left(-\widetilde{f}_{x}\right)\right|
\end{aligned}
$$

Im Punkt $x$ selbst ist $g_{x}+\log |z|$ harmonisch, also gibt es eine Funktion $\widetilde{f}_{x}$, sodass wie in der vorigen Rechnung

$$
g_{x}(z)+\log |z|=-\log \left|\exp \left(-\widetilde{f}_{x}(z)\right)\right|
$$

gilt. Also leistet $f_{x}=\exp \left(-\widetilde{f}_{x}(z)\right) / z$ das Verlangte. Da $g_{x}>0$ ist $\left|f_{x}\right|<1$, also ist $f_{x}$ eine Abbildung in die Einheitskreisscheibe.

Zur Injektivität von $f_{x}$ : Seien $x$ und $y$ fest und sei

$$
\phi(z)=\frac{f_{x}(y)-f_{x}(z)}{1-f_{x}(y) f_{x}(z)}
$$

Dies ist holomorph in $z, \phi(y)=0$ und $|\phi|<1$, denn

$$
\begin{aligned}
\mid \text { Nenner }\left.\right|^{2}-\mid \text { Zähler }\left.\right|^{2} & =\left|1-f_{x}(y) f_{x}(z)\right|^{2}-\left|f_{x}(y)-f_{x}(z)\right|^{2} \\
& =\left(1-\left|f_{x}(y)\right|^{2}\right)\left(1-\left|f_{x}(z)\right|^{2}\right)>0
\end{aligned}
$$

Sei $n=\operatorname{ord}_{y} \phi$. Dann ist $u=-\frac{1}{n} \log |\phi|$ harmonisch und positiv. Wir vergleichen $u$ mit der Greenschen Funktion $u_{y}$ bei $y$. Sei $\zeta$ Koordinate einer Kreisscheibenkarte um $y$. Da $\phi(\zeta)=\zeta^{n} \cdot P(\zeta)$ mit einer Potenzreihe $P(\zeta)$ mit $P(0) \neq 0$ ist, ist $u+\log |\zeta|$ harmonisch in der Nähe von $y$.\\
Aufgrund der Minimalität der Greenschen Funktion ist $u \geqslant u_{y}$. Sei analog zu oben $u_{y}=-\log \left|f_{y}\right|$. Dann gilt

$$
\left|f_{y}(z)\right| \geqslant|\phi(z)|^{1 / n} \geqslant|\phi(z)|
$$

Für $z=x$ folgt $\left|f_{x}(y)\right| \leqslant\left|f_{y}(x)\right|$ und durch Rollentausch von $y$ und $x$ in obigem Argument folgt Gleichheit. Also ist die holomorphe Funktion $h(z)=\phi(z) / f_{y}(z)$ im Betrag durch 1 beschränkt und $|h(x)|=1$. Folglich ist $h$ konstant. Die Funktion $f_{y}(z)$ ist nur in $z=y$ Null, also auch $\phi=h \cdot f_{y}$.

Zur Surjektivität von $f_{x}$ : Angenommen es gibt $a^{2} \in \Delta \backslash f_{x}(X)$. Notwendigerweise ist $a \neq 0$, da $f_{x}(x)=0$. Also ist $h_{2}(z)=\left(z-a^{2}\right) /\left(1-a^{-2} z\right)$ holomorph und ohne Nullstelle auf $f_{x}(X)$. Nach obiger Injektivität ist $f_{x}(X) \cong X$, also insbesondere einfach zusammenhängend. Also hat $h_{2}(z)$ eine Quadratwurzel $h(z)$ mit $h(0)=i \cdot a$. Sei

$$
F=(h-i a) /(1+i \bar{a} h): f_{x}(X) \longrightarrow \mathbb{C}
$$

Ziel ist es zu zeigen, dass $-\log \left|F \circ f_{x}\right|$ alle Eigenschaften einer Greenschen Funktion in $x$ hat, von der Minimalität abgesehen. Daraus erhalten wir dann

$$
-\log \left|F \circ f_{x}\right| \geqslant g_{x}=-\log \left|f_{x}\right|
$$

Also folgt $|F(z)| \leqslant|z|$ in der Nähe von Null, was $\left|F^{\prime}(0)\right|>1$ widerspricht. Mit Hilfe der Kettenregel und $h^{\prime}(0)=\frac{\left(h^{2}\right)^{\prime}(0)}{2 \cdot}$ berechnet man $F^{\prime}(0)=\frac{1+|a|^{2}}{2 i a}$, insbesondere folgt $\left|F^{\prime}(0)\right|>1$.\\
Die Funktion $h^{2}$ ist von der selben Termkomposition wie $\phi$. Obiger Beweis zeigt also, dass $\left|h^{2}(z)\right|<1$ und folglich $|F(z)|<1$ ist.

Es ist noch zu zeigen, dass $F$ nur bei $z=0$ eine Nullstelle hat. Offenbar ist $F(0)=0$. Umgekehrt, falls $F(z)=0$, d.h. falls $h=i a$, so ist

$$
h(z)^{2}=-a^{2}=\frac{z-a^{2}}{1-a^{-2} z}
$$

woraus unmittelbar $z=0$ folgt. Da $f_{x}$ nach dem ersten Teil injektiv ist, folgt dass $F \circ f_{x}$ nur bei $z=0$ eine Nullstelle hat.

Der gleiche Quadratwurzeltrick zeigt auch:\\
Korollar 7.27 Ein einfach zusammenhängendes Gebiet $X$ in $\mathbb{P}_{\mathbb{C}}^{1}$ ist isomorph zu $\mathbb{C}, z u \mathbb{P}_{\mathbb{C}}^{1}$ oder zu $\Delta$.

Beweis: Falls $X \nsupseteq \mathbb{P}_{\mathbb{C}}^{1}$ und $X \neq \mathbb{C}$, so ist ohne Einschränkung $0, \infty \notin X$. Also gibt es eine holomorphe Funktion $f: X \rightarrow \mathbb{C}$ mit $f^{2}(z)=z$. Die Funktion $f$ nimmt niemals die Werte $y$ und $-y$ beide an. Da Bild $(f)$ offen ist, also eine Kreisscheibe enthält, gibt es auch eine Kreisscheibe $\Delta_{r_{0}}\left(z_{0}\right)$, die zum $\operatorname{Bild}(f)$ disjunkt ist. Also ist

$$
h(z)=\frac{1}{f(z)-z_{0}}
$$

holomorph und beschränkt auf $X$. Folglich ist $X$ reich und der vorige Satz anwendbar.

Satz 7.28 Eine kompakte, einfach zusammenhängende Riemannsche Fläche ist isomorph zu $\mathbb{P}_{\mathbb{C}}^{1}$. Eine arme, nicht-kompakte, einfach zusammenhängende Riemannsche Fläche ist isomorph zu $\mathbb{C}$.

Wir sagen, dass eine holomorphe Funktion $f: X \longrightarrow \mathbb{P}_{\mathbb{C}}^{1}$ zulässig bei $x \in X$ ist, falls sie einen einfachen Pol bei $x$ hat und außerhalb jeder Umgebung von $x$ beschränkt ist. Als Vorbereitung zeigen wir der Reihe nach

Lemma 7.29 1.) $Z u$ jedem Punkt $x \in X$ gibt es zulässige Funktionen. Je zwei zulässige Funktionen bei $x$ gehen durch Postkomposition mit einer Möbiustransformation ineinander über.\\
2.) Eine zulässige Funktion in $x$ und eine zulässige Funktion in $y$ gehen durch Postkomposition mit einer Möbiustransformation ineinander über.\\
3.) Zulässige Funktionen sind injektiv.

Beweis des Satzes 7.28, Sei $f: X \longrightarrow \mathbb{P}_{\mathbb{C}}^{1}$ zulässig. Falls $X$ kompakt ist, so ist das Bild offen und kompakt, also ganz $\mathbb{P}_{\mathbb{C}}^{1}$. Aus Lemma 7.29-3) folgt die Isomorphie.\\
Falls $X$ nicht kompakt ist, so ist $f$ nicht surjektiv. Nach Postkomposition mit einer Möbiustransformation ist also $f(X) \subseteq \mathbb{C}$. Wäre diese Inkusion echt, so wäre $X$ reich nach Korollar 7.27. Mit Lemma 7.29-3) folgt, dass $f$ ein Isomorphismus von $X$ nach $\mathbb{C}$ ist.

Das Argument von Beweis von Lemma 7.29-1) ist analog zu Satz7.20. Wegen der Zusatzaussage über Möbiustransformation führen wir es noch einmal aus.\\
Beweis von Lemma 7.29-1): Wie in Satz7.20sei $U=\{|z| \leq 1\}$ eine Kreisscheibenkarte um $X$, sei $u: X \backslash\{x\} \rightarrow \mathbb{R}$ harmonisch, $u-\operatorname{Re}\left(\frac{1}{z}\right)$ harmonisch auf $U$ und Null bei $x$. Schreibe $u=\operatorname{Re} f$ mit $f$ holomorph. Da $X$ einfach zusammenhängend ist, können wir dieses $f$ global wählen. Es ist zunächst zu zeigen, dass der Imaginärteil von $f$ auch beschränkt außerhalb jeder Umgebung von $x$ ist. Wähle dazu $\widetilde{u}$ wieder wie in Satz 7.20, nun aber sodass $\widetilde{u}-\operatorname{Re}\left(\frac{i}{z}\right)$ harmonisch bei $x$ ist. Ist $\widetilde{f}$ holomorph so gewählt, dass $\widetilde{u}=\operatorname{Re} \widetilde{f}$, dann ist der Realteil von $\widetilde{f}$ offenbar beschränkt außerhalb jeder Umgebung von $U$. Es ist also $\widetilde{f}=i f$ zu zeigen.\\
Eine holomorphe Funktion in einer Umgebung von $x$ mit einfachem Pol bei $x$ ist injektiv in einer Umgebung von $X$, denn falls $f=\frac{1}{z}+a_{0}+a_{2} z+\ldots$, so ist

$$
f\left(z_{1}\right)-f\left(z_{2}\right)=\left(z_{1}-z_{2}\right)\left(\frac{-1}{z_{1} z_{2}}+h\right)
$$

wobei $h$ in einer Umgebung von $x$ beschränkt ist. Also können wir durch Verkleinern der Umgebung annehmen, dass $f$ und $\widetilde{f}$ injektiv sind. Wähle $m \geqslant \max _{|z| \geqslant 1}(|u(z)|+|\widetilde{u}(z)|)$ und $x_{2}$ so nahe bei $x$, dass $\left|u\left(x_{2}\right)\right|>2 m$ und $\left|\widetilde{u}\left(x_{2}\right)\right|>2 m$. Sei

$$
g(z)=\frac{1}{f(z)-f\left(x_{2}\right)} \quad \text { und } \quad \widetilde{g}(z)=\frac{1}{\widetilde{f}(z)-\widetilde{f}\left(x_{2}\right)}
$$

Die Injektivität impliziert, dass $g$ und $\widetilde{g}$ holomorph auf $U \backslash\left\{x_{2}\right\}$ sind, beide bei $x_{2}$ einen einfachen Pol haben, und nach Wahl von $m$ außerhalb von $U$ beschränkt sind, da gilt

$$
\left|f(z)-f\left(x_{2}\right)\right| \geqslant\left|u(z)-u\left(x_{2}\right)\right| \geqslant m
$$

und ebenso für $\widetilde{f}$. Folglich ist eine Linearkombination $a g+\widetilde{a} \widetilde{g}$ mit $a, \widetilde{a} \in \mathbb{C}$ holomorph auf $X$ und konstant, da $X$ arm ist. Löst man $a g+\widetilde{a} \widetilde{g}=$ const auf, so erhält man $\widetilde{f}=\frac{\alpha f+\beta}{\gamma f+\delta}$ mit $\alpha, \beta, \gamma, \delta \in \mathbb{C}$. Da $f=\frac{1}{z}+O(z)$ und $\widetilde{f}=\frac{i}{z}+O(z)$ folgt aus dieser Darstellung $\widetilde{f}=i f$.\\
Außerdem können wir das Argument für eine beliebige andere Funktion anstelle von $\widetilde{f}$ wiederholen und erhalten somit die Aussage über die Existenz einer Möbiustransformation zwischen zwei zulässigen Funktionen bei $x$.

Beweis von Lemma 7.29-2): Sei $\Sigma_{x}$ die Menge bei denen die Aussage wahr ist. $x \in \Sigma_{x}$ nach Lemma 7.29-1. Die Menge $\Sigma_{x}$ ist offen, denn für $x_{2}$ nahe bei $x$ und $f$ zulässig bei $x$ liefert das Argument von Lemma 7.29 -1, dass

$$
g(z)=\frac{1}{f(z)-f\left(x_{2}\right)}
$$

zulässig bei $x_{2}$ ist. Angenommen $\Sigma_{x} \subsetneq X$, d.h. es gibt ein $y \in X \backslash \Sigma_{x}$. Mit dem gleichen Argument ist $\Sigma_{y}$ offen, nicht-leer und disjunkt zu $\Sigma_{x}$, da die Verkettung zweier Möbiusformationen wieder eine Möbiusformation ist. Das widerspricht dem Zusammenhang von $X$.

Beweis von Lemma 7.29-3): Falls $f\left(x_{1}\right)=f\left(x_{2}\right)$, sei $g$ eine bei $x_{1}$ zulässige Funktion. Dann geht $g$ auf $f$ durch eine Möbiustransformation hervor nach Lemma 7.29-2. Insbesondere hat $g$ auch einen Pol bei $x_{2}$. Falls $x_{1} \neq x_{2}$, so widerspricht dies der Definition von Zulässigkeit bei $x_{2}$.

\subsection {Der Satz von Riemann-Roch}
Der Zugang dieses Kapitels liefert uns eine weitere Art, den Satz von RiemannRoch zu beweisen, den wir mittels Garbenkohomologie und Serre-Dualität als Satz 6.6 bewiesen haben. Wir setzen in diesem Abschnitt lediglich die Grundbegriffe über Divisoren aus den einleitenden Paragraphen von Abschnitt 6 voraus. Zur logischen Unabhängigkeit sei noch bemerkt, dass wir im Beweis von Proposition 6.3 statt auf den Existenzsatz 5.28 zu veweisen, auch auch den Existenzsatz 7.21verweisen können.

Satz 7.30 (Riemann-Roch) Sei $X$ eine kompakte Riemannsche Fläche vom Geschlecht $g$. Sei $\omega \not \equiv 0$ eine meromorphe Differentialform und $K=\operatorname{div}(\omega)$. Dann gilt für jeden Divisor $D$ auf $X$ :

$$
\ell(D)<+\infty \quad \text { und } \quad \ell(D)-\ell(K-D)=\operatorname{deg}(D)+1-g
$$

Mit Hilfe dieses Satzes von Riemann-Roch gilt unverändert der Beweis des im folgenden immer wieder verwendeten Korollars 6.7, dass der Vektorraum der holomorphen Differentialformen auf einer kompakten Riemannschen Fläche vom Geschlecht $g$ die Dimension $g$ besitzt.\\
Beweis: Schritt 1a: Für jeden positiven Divisor $D$ gilt $\ell(D) \geqslant \operatorname{deg}(D)+1-g$. Da $X$ kompakt ist, können wir Satz 7.20 anwenden. Dieser impliziert, dass es zu gegebenem $k_{j}$ und $x_{j} \in X$ eine harmonische Funktion gibt, die bei $x_{j}$ wie $\operatorname{Re}\left(1 / z_{j}^{k_{j}}\right)$ wächst, wobei $z_{j}$ eine Koordinate bei $x_{j}$ ist. Gleiches gilt für $\operatorname{Im}\left(1 / z_{j}^{k_{j}}\right)$. Wir schreiben den Divisor als $D=\sum_{j \in J} n_{j}\left[x_{j}\right]$ und es sei $H$ der Vektorraum der harmonischen Funktion auf $X \backslash \bigcup_{j \in J} x_{j}$ mit Wachstum höchstens der Ordnung $n_{i}$ bei $x_{i}$. Dann ist

$$
\operatorname{dim}_{\mathbb{R}} H \geqslant 2\left(\sum_{j \in J} n_{j}+1\right)
$$

Wir untersuchen nun, welche Funktionen in $H$ sich als $\operatorname{Re}(f)$ für $f$ holomorph schreiben lassen. Die Fundamentalgruppe $\pi_{1}\left(X \backslash \bigcup_{j \in J}\left\{x_{j}\right\}\right)$ ist erzeugt von Schleifen $\gamma_{j}$ um $x_{j}$ und den $a_{j}, b_{j}, j=1 \ldots g$ aus der kanonischen Zerschneidung nach Satz 3.30. Wir zeigen, dass $\int_{\gamma_{j}}{ }^{*} d u$ für alle $u \in H$ und alle $j \in J$\\
verschwindet. Für $k_{j}>0$ ist die Funktion $u-\operatorname{Re}\left(z_{i}^{-k_{i}}\right)$ harmonisch. Folglich ist

$$
d u+i^{*} d u-z_{j}^{-k_{j}-1} \cdot\left(-k_{j}\right) d z_{j}
$$

holomorph, also ohne Residuum bei $x_{j}$. Da $-k_{j}-1<-1$ ist, beeinflußt der letzte Term das Residuum nicht.\\
Also stellen wir nur die $a_{i}$ und $b_{i}$ Bedingungen, die $u=\operatorname{Re} f$ verhindern könnten, also höchstens $2 g$ viele. Insgesamt ist

$$
\operatorname{dim}_{\mathbb{R}}(\mathcal{L}(D)) \geq 2 \cdot\left(\sum_{j \in J} n_{j}+1\right)-2 g=2 \cdot(\operatorname{deg} D+1-g)
$$

Schritt 1b: Für jeden effektiven Divisor $D$ gilt $\ell(D) \geqslant \operatorname{deg} D+1-g$.\\
Wir schreiben $D=\sum_{i \in I} u_{i}\left[x_{i}\right]-\sum_{j \in J} m_{j}\left[y_{j}\right]$ und es sei $\widetilde{D}=\sum_{i \in I} u_{i}\left[x_{i}\right]$. Sei weiterhin $\varphi: \mathcal{L}(\widetilde{D}) \longrightarrow \mathbb{C}^{\sum m_{j}}$ die lineare Abbildung, die jeder Funktion die Taylorentwicklungen an den Stellen $y_{j}$ bis zur Ordnung $m_{j}$ zuordnet. Dann ist $\operatorname{Ker}(\varphi)=\mathcal{L}(D)$ und $\operatorname{dim}_{\mathbb{C}} \mathcal{L}(D) \geqslant \operatorname{dim}_{\mathbb{C}} \mathcal{L}(\widetilde{D})-\sum_{j \in J} m_{j} \geqslant \operatorname{deg} D+1-g$ nach Schritt 1a.

Schritt 2: Für alle $x \in X$ gilt $\ell(D) \leq \ell(D+[x]) \leq \ell(D)+1$.\\
Die erste Ungleichung ist offensichtlich. Außerdem hat für $D=\sum_{x \in X} n(x) \cdot[x]$ die Abbildung

$$
\varphi:\left\{\begin{array}{lll}
\mathcal{L}(D+[x]) & \longrightarrow & \mathbb{C} \\
f=\sum_{i \in \mathbb{Z}} a_{i} \cdot z^{i} & \longmapsto & a_{n(x)-1}
\end{array}\right.
$$

den Ker $\mathcal{L}(D)$. Aus Schritt 2 folgt induktiv auch die erste Behauptung des Satzes.

Schritt 3: Die Differenzen $\ell(D+[x])-\ell(D)$ und $\ell(K-D)-\ell(K-D-[x])$ sind nicht beide gleichzeitig eins.\\
Angenommen die Behauptung ist falsch und seien $f_{1}$ bzw. $f_{2}$ Funktionen in den Differenzvektorräumen. Dann ist

$$
\begin{gathered}
\operatorname{div}\left(f_{1}\right)+D=-x+\sum_{y \neq x} n(y) \cdot[y] \\
\operatorname{div}\left(f_{2}\right)+K-D=\sum_{y \neq x} m(y) \cdot[y]
\end{gathered}
$$

mit $m \geqslant 0$ und $n \geqslant 0$.\\
Dann aber ist $f_{1} f_{2} \omega$ eine meromorphe Differentialform mit Divisor

$$
\operatorname{div}\left(f_{1}\right)+\operatorname{div}\left(f_{2}\right)+\operatorname{div}(\omega)=-x+\sum_{y \neq x}(m(y)+n(y)) \cdot[y]
$$

Das aber widerspricht dem Residuensatz, denn die Summe der Residuen ist dieser Differentialform ist -1 .

Schritt 4: Wir betrachten die Differenzfunktion

$$
\psi(D)=\ell(D)-\ell(K-D)-\operatorname{deg}(D)
$$

Um eine untere Schranke für $\psi(D)$ zu finden, können wir $D$ ohne Einschränkung vergrößern, da $\psi(D+[x]) \geqslant \psi(D)$ nach Schritt 3 gilt. Sobald $\operatorname{deg}(D) \geqslant 2 g-2$ ist $\ell(K-D)$ und dann ist nach Schritt 1b

$$
\psi(D) \geqslant 1-g \text { für alle } D \in \operatorname{Div}(X) \text {. }
$$

Mit dem gleichen Argument ist $\psi(K-D) \geqslant 1-g$, also

$$
2-2 g \leqslant \psi(D)+\psi(K-D)=-\operatorname{deg}(K)=2-2 g
$$

weswegen überall Gleichheit gelten muss.

\section {Flache Flächen}
In diesem Abschnitt geben wir die formale Definition der Art von Riemannschen Flächen, wie sie in der Motivation 1.2 der Einleitung und später im Beispiel 2.8 für den Torus nochmal aufgetreten sind. Zunächst betrachten wir ein interessanteres Beispiel und definieren dann formal, was Entfaltung bedeutet.

\subsection {Entfalten von Polygonen}
Beispiel 8.1 Ein immer wieder nützliches Beispiel in diesem Kontext ist die L-Fläche, wir skalieren sie mit den Schenkellängen a und b und Schenkelhöhen eins. Sie besteht im Wesentlichen aus den drei Rechtecken $X_{1}, X_{2}$ und $X_{3}$

$$
\begin{aligned}
& X_{1}=\{z=x+i y \in \mathbb{C}: x \in(0,1), y \in(0, a-1)\} \\
& X_{2}=\{z \in \mathbb{C}: x \in(0,1), y \in(0,1)\} \\
& X_{3}=\{z \in \mathbb{C}: x \in(0, b-1), y \in(0,1)\}
\end{aligned}
$$

wie in Abbildung 8.1.\\
Dann verkleben wir linke und rechte Seiten durch horizontale Translationen und obere und untere Seiten durch vertikale Translationen. Diesen offensichtlichen Vorgang formal korrekt zu beschreiben ist notationell etwas lästig. Also definieren wir

$$
\begin{aligned}
& X_{4}=X_{6}=\{z \in \mathbb{C}: x \in(0,1), y \in(-\varepsilon, \varepsilon)\} \\
& X_{5}=X_{7}=\{z \in \mathbb{C}: x \in(-\varepsilon, \varepsilon), y \in(0,1)\} \\
& X_{8}=\{z \in \mathbb{C}: x \in(0, b-1), y \in(-\varepsilon, \varepsilon)\} \\
& X_{9}=\{z \in \mathbb{C}: x \in(-\varepsilon, \varepsilon), y \in(0, a)\}
\end{aligned}
$$

und die zugehörigen Kartenwechselabbildungen wie in Abbildung 8.1 angedeutet. D.h. $X_{4}$ ist verklebt mit $X_{1}$ und $X_{2}$ vermöge\\
$U_{1}^{+}=\{z \in \mathbb{C}: x \in(0,1), y \in(0, \varepsilon)\}$ und $U_{1}^{-}=\{z \in \mathbb{C}: x \in(0,1), y \in(-\varepsilon, 0)\}$

\begin{figure}[h]
\begin{center}
  \includegraphics[scale=0.2]{2025_10_06_374d4cfa8b1948eb4c41g-092}
 
\caption{Abbildung 8.1: Die $L$-Fläche mit Schenkellängen $a$ und $b$}
\end{center}
\end{figure}

und den Verklebeabbildungen

$$
\begin{aligned}
& \varphi_{1}:\left\{\begin{array}{lll}
U_{1}^{+} & \longrightarrow & X_{1} \\
z & \longmapsto & z
\end{array}\right. \\
& \varphi_{1}:\left\{\begin{array}{l}
U_{1}^{-} \\
z
\end{array} \longrightarrow \begin{array}{l}
\longrightarrow \\
z
\end{array} X_{4}\right. \\
& z
\end{aligned} \quad \varphi_{2}:\left\{\begin{array} { r l l } 
{ U _ { 1 } ^ { + } } & { \longrightarrow } & { X _ { 4 } } \\
{ z } & { \longmapsto } & { z }
\end{array} \left\{\begin{array}{rll}
U_{1}^{-} & \longrightarrow & X_{2} \\
z & \longmapsto & 1+z
\end{array}\right.\right.
$$

Dabei bleibt wieder „nur die Ecke frei", d.h. wir nennen $X^{*}$ die Verklebung von $X_{1}, \ldots, X_{9}$ wie oben.

Wir wollen nun die "Löcher" der Riemannschen Fläche schließen, d.h. wir wollen $X^{*}$ mit Kreisscheiben vom Radius $\varepsilon$ (mit $\varepsilon<\frac{\min \{1, a, b\}}{2}$ ) verkleben, sodass die entstehende Fläche kompakt ist. Dazu gibt es ein naheliegendes Prinzip: Wir stellen uns gedanklich in den Ursprung von $X_{2}$, schauen zunächst in Richtung der reellen Achse und drehen uns nach links. Sobald wir in Richtung der imaginären Achse blicken, können wir unseren Standpunkt auch nach ( $b-1,0$ ) in $X_{3}$ verlagern und sehen, aufgrund der Verklebung mit $X_{7}$, denselben Ausschnitt der Fläche wie vorher. Nun drehen wir uns weiter nach links. Diesem Prinzip folgend verkleben wir $X_{\varepsilon}=\triangle_{\varepsilon}$ und $X^{*}$ wie folgt.

Seien $V_{k}=\left\{z \in \triangle_{\varepsilon}: \arg (z) \in\left[(k-1) \frac{\pi}{6}, k \frac{\pi}{6}\right]\right\}$ für $k \in\{1, \ldots, 12\}$,

\begin{figure}[h]
\begin{center}
  \includegraphics[scale=0.2]{2025_10_06_374d4cfa8b1948eb4c41g-093}
 
\caption{Abbildung 8.2: Verkleben des Loches der $L$-Fläche}
\end{center}
\end{figure}

$$
\begin{aligned}
& \varphi_{2}: V_{k} \quad \longrightarrow \triangle_{\varepsilon} \quad \text { die Inklusion und } \\
& \varphi_{1}: V_{1} \quad \longrightarrow X_{2}, \quad z \longmapsto z^{3}, \\
& \varphi_{1}: V_{2} \quad \longrightarrow X_{3}, \quad z \longmapsto z^{3}+(b-1) \\
& \varphi_{1}: V_{3} \quad \longrightarrow X_{3}, \quad z \longmapsto z^{3}+(b-1)+i \\
& \varphi_{1}: V_{4} \quad \longrightarrow X_{2}, \quad z \longmapsto z^{3}+i \\
& \varphi_{1}: V_{5} \quad \longrightarrow X_{1}, \quad z \longmapsto z^{3} \\
& \varphi_{1}: V_{6} \quad \longrightarrow X_{1}, \quad z \longmapsto z^{3}+1 \\
& \varphi_{1}: V_{7} \quad \longrightarrow X_{2}, \quad z \longmapsto z^{3}+1+i \\
& \varphi_{1}: V_{8} \quad \longrightarrow X_{3}, \quad z \longmapsto z^{3}+i \\
& \varphi_{1}: V_{9} \quad \longrightarrow X_{3}, \quad z \longmapsto z^{3} \\
& \varphi_{1}: V_{10} \quad \longrightarrow X_{2}, \quad z \longmapsto z^{3}+1 \\
& \varphi_{1}: V_{11} \quad \longrightarrow X_{1}, \quad z \longmapsto z^{3}+1+(a-1) \cdot i \\
& \varphi_{1}: V_{12} \quad \longrightarrow X_{1}, \quad z \longmapsto z^{3}+(a-1) \cdot i
\end{aligned}
$$

Die obige Verklebungsbeschreibung ergibt mengentheoretisch das richtige Objekt, aber sie ist nicht formal korrekt im Sinne der Definition in Abschnitt 2.1. da dort Verkleben stets entlang offener Mengen definiert war, hier aber die $V_{k}$ abgeschlossen sind. Würde man die $V_{k}$ offen definieren, blieben die horizontalen und vertikalen Kantensegmente in $V_{k}$ verdoppelt übrig. Eine vollständig korrekte Definition erhält man, indem man die offenen $V_{k}$ verklebt und zudem noch Umgebungen der horizontalen und vertikalen Kantensegmente mit geeigneten offenen Mengen in den $X_{4}, \ldots, X_{9}$. Wir führen dies nicht im Detail aus.

Dieses Beispiel verallgemeinern wir. Gegeben seien Polygone $P_{i}, i \in I$, in der komplexen Ebene $\mathbb{C}$, nicht notwendigerweise endlich viele. (Polygone haben hier per Definition nur endlich viele Kanten.) Wir nehmen zusätzlich an,

\begin{figure}[h]
\begin{center}
  \includegraphics[scale=0.2]{2025_10_06_374d4cfa8b1948eb4c41g-094}
 
\caption{Abbildung 8.3: Verkleben mit Hilfe der \$U\_\{K}\$\}\end{center}
\end{figure}

dass es eine Paarung $\sigma$ (d.h. eine fixpunktfreie Involution) auf der Menge der Kanten $\mathcal{K}$ aller dieser Polygone gibt und dass zudem $K$ und $\sigma(K)$ gleich lang und parallel sind für alle $K \in \mathcal{K}$. Zu jedem solchen Kantenpaar $K \in P_{i}$ und $\sigma(K) \in P_{j}$ wählen wir eine Umgebungen $U_{K, i} \in P_{i}$ und $U_{\sigma(K), j} \in P_{j}$. Die Vereinigung $U_{K}$ von $U_{K, i}$, einer geeigneten Translation von $U_{\sigma(K), j} \in P_{j}$ and die Kante $K$ ist eine Umgebung von $K \subset \mathbb{C}$ nach der Voraussetzung an die gepaarten Kanten.

Es bezeichne $P_{i}^{\circ}=P_{i} \backslash \partial P_{i}$ das offene Polygon. Sei also

$$
X^{*}=\left(\coprod_{i \in I} P_{i}^{\circ} \sqcup \coprod_{K \in \mathcal{K} / \sigma} U_{K}\right) / \sim
$$

wobei wir für jedes Kantenpaar $(K, \sigma(K))$ die Inklusionsabbildungen $\varphi_{1}$ : $U_{K, i} \rightarrow P_{i}$ und $\varphi_{1}: U_{\sigma(K), j} \rightarrow P_{j}$ verwenden und $\sim$ die Äquivalenzrelation ist, die durch Verkleben der offenen Polygone und Übergangsflächen $U_{K}$ entlang all dieser offenen Mengen $U_{K, i}$, wie in Abschnitt 2.1 definiert, entsteht.

Wir wollen nun die Löcher von $X^{*}$ wie im vorigen Beispiel stopfen. Dazu bestimmen wir zu jeder Ecke $E_{i j}$ jedes der $P_{i}$ das Kreissegment $S_{i j}$ von maximalem Radius $r_{i j}$, das in $P_{i}$ enthalten ist. In einem festen Drehsinn gesehen (sagen wir gegen den Uhrzeigersinn) definert die Kantenidentifikation $\sigma$ zu jedem $S_{i j}$ einen Nachfolger.

Wir betrachten zunächst den Fall, dass die Bahn unter dieser Nachfolgerabbildung endlich ist. Dann sind die Radien $r_{i j}$ entlang dieser Bahn nach unten durch eine positive Zahl beschränkt und wir ersetzen die $S_{i j}$ durch Kreisbogensemgente des minimalen Radius $r_{\text {min }}$. Diese Summe der Winkel der Kreisbogensegmente in solch einer Bahn ist ein ganzzahliges Vielfaches $2 k \pi$ von $2 \pi$,\\
da $\sigma$ eine Bijektion paralleler Kanten ist. In diesem Fall fügen wir zu $X^{*}$ einen Punkt $x_{0}$ hinzu, indem wir $X^{*}$ mit einer Kreisscheibe $U=\Delta_{r_{\text {min }}}$ längs der Vereinigung der Kreisbogensegmente $S_{i j}$ via geeigneter Translationen der Abbildung $z \mapsto z^{k}$ verkleben. (Genauer gesagt, müssen wir auch noch auf den Einschränkungen der Überlappungsflächen verkleben, wie in der vorangegangenen Bemerkung zu Beispiel 8.1. Wir nennen $x_{0}$ einen Punkt mit konischem Winkel $2 k \pi$.

Sind alle Bahnen unter der Nachfolgerabbildung endlich, so bezeichne $X$ die Riemannsche Fläche, die man aus $X^{*}$ erhält, indem man für jede Bahn einen Punkt wie oben beschrieben hinzufügt. Wir nennen $X^{*}$ die offene Verklebungsfläche der $P_{i}$ und $X$ die (abgeschlossene) Verklebungsfläche der $P_{i}$. Startet man mit endlich vielen Polygonen $P_{i}$, so ist die Verklebungsfläche eine kompakte Riemannsche Fläche.

Ist eine Bahn unter der Nachfolgerabbildung nicht endlich, so sind wir in der Situation der Riemannschen Fläche des Logarithmus, wie im ersten motivierenden Beispiel besprochen. In diesem Fall ist es nicht möglich, $X^{*}$ zu kompaktifizieren, sodass das Resultat eine Riemannsche Fläche ist.

Sei $P \subseteq \mathbb{C}$ ein Polygon. Eine Reflexionskopie von $P$ ist ein Polygon, das aus $P$ durch Spiegelung an einer Kante von $P$ entsteht. Die Menge $R(P)$ aller Reflexionskopien von $P$ sei die kleinste Menge von Polygonen, die $P$ enthält, und die zu jedem Element $Q \in R(P)$ auch ein Translat jeder Reflexionskopie von $Q$ enthält.

Die Menge der Kanten aller Polygone in $R(P)$ hat eine natürliche Involution $\sigma$ mit den Eigenschaften wie in der Definition von Verklebungsflächen gefordert, indem man die Fixgeraden jeder Spiegelung miteinander identifiziert.

Ist $P$ ein Polygon, dessen Winkel allesamt rationale Vielfache von $\pi$ sind, so ist $R(P)$ endlich. Um dies einzusehen genügt es zu beobachten, der Winkel einer festen Kante von $P$ (sagen wir gegen die Horizontale) bei einer Reflektion nur um ein Vielfaches von $\frac{1}{\operatorname{kgV}\left(q_{i}\right)} \pi$ abgeändert wird, falls die Eckenwinkel von $P$ gleich $\frac{p_{i}}{q_{i}}$ sind.

Definition 8.2 (Entfaltungsfläche) Die Entfaltungsfläche $X(P)$ eines Polygons, dessen Winkel rationale Vielfache von $\pi$ sind, ist die Verklebungsfläche der Reflexionskopien von $P$.

Beispielsweise ist die Menge der Reflexionskopien der L-Fläche aus Beispiel 8.1 (d.h. genauer $L$ ist der Abschluss von $X_{1} \cup X_{2} \cup X_{3}$ ) vierelementig. Im Bild 8.4 sind die Entfaltungsflächen der rechtwinkligen Dreiecke mit einem Winkel $\pi / 5$ bzw. $\pi / 10$ dargestellt. Jeweils parallel Aussenkanten sind verklebt.

\begin{figure}[h]
\begin{center}
  \includegraphics[scale=0.2]{2025_10_06_374d4cfa8b1948eb4c41g-096}
 
\caption{Abbildung 8.4: Zwei Entfaltungsflächen von Dreiecken}
\end{center}
\end{figure}

\subsection {Translationsstrukturen}
Die Entfaltungsflächen aus dem vorigen Abschnitt sind Spezialfälle der folgenden zwei Definitionen, die sich als im Wesentlichen äquivalent herausstellen werden.

Definition 8.3 (flache Fläche) Eine flache Fläche ( $X, \omega$ ) ist eine kompakte Riemannsche Fläche $X$ zusammen mit einer von Null verschiedenen holomorphen 1Form $0 \neq \omega \in \Omega_{X}^{1}(X)$.

Zwei flache Flächen $\left(X_{i}, \omega_{i}\right), i=1,2$ werden als isomorph bezeichnet, wenn es einen Isomorphismus $\varphi: X_{1} \rightarrow X_{2}$ gibt, der die Differentialformen ineinander überführt, d.h. sodass $\varphi^{*}\left(\omega_{2}\right)=\omega_{1}$ ist. Wir erinnern an die Definition von $\varphi^{*}$. Ist $z$ eine lokale Karte einer Umgebung eines Punktes $P_{2} \in X_{2}$, so ist $t=z \circ \varphi$ eine lokale Karte einer Umgebung des Punktes $P_{1}=\varphi^{-1}\left(P_{2}\right) \in X_{1}$. Ist auf diesen Karten $\omega_{2}=\alpha_{2}(z) d z$ und $\omega_{1}=\alpha_{1}(t) d t$, so ist $\varphi^{*}\left(\omega_{2}\right)= \alpha_{2}(\varphi(t)) \varphi^{\prime}(t) d t$ und die Bedingung an $\varphi$ ist gerade $\alpha_{1}(t)=\alpha_{2}(\varphi(t)) \varphi^{\prime}(t)$.

Definition 8.4 (Translationsstruktur) Ein Translationsatlas auf einer Riemannschen Fläche $X$ ist ein komplexer Atlas ( $U_{i}, g_{i}: U_{i} \longrightarrow \mathbb{C}$ ), dessen Übergangsfunktionen Translationenen sind. Zwei Translationsatlanten auf $X$ werden äquivalent genannt, wenn die Übergangsfunktionen auf ihrem Definitionsbereich Translationen sind. Eine Riemannsche Fläche X mit einer Äquivalenzklasse von Translationsatlanten wird Fläche mit Translationsstruktur oder Translationsfläche genannt.

Zwei Flächen mit Translationsstrukturen heißen isomorph, wenn es einen Isomorphismus der zugrundeliegenden Riemannschen Flächen gibt, sodass die induzierten Abbildungen auf den Karten der beiden Translationsatlanten Translationen sind.

Vor dem Beweis der Äquivalenz der Definitionen führen wir einige Begriffe ein, die auf flachen Flächen und Flächen mit Translationsstrukturen definiert werden können, aber nicht auf Riemannschen Flächen im Allgemeinen.

Definition 8.5 (Weglänge, Abstand) Ein Weg $\gamma$ hat eine Länge $\ell(\gamma)$ : Ist $\gamma$ in einer Karte des Translationsatlas enthalten, so ist dies die übliche Bogenlänge. Im All-\\
gemeinen definieren wir die Länge von $\gamma$ durch Zerschneiden in Stücke, die in einer Karte des Translationsatlas liegen.

Zwei Punkte $x_{1}, x_{2} \in X$ haben einen Abstand

$$
d\left(x_{1}, x_{2}\right):=\inf _{\substack{\gamma \text { Weg von } \\ x_{1} \text { nach } x_{2}}} \ell(\gamma)
$$

Wir bemerken, dass wir auf einer Fläche mit Translationstruktur einen wohldefinierten Winkelbegriff zwischen zwei sich schneidenden Geraden haben.

Eine flache Fläche hat ein Volumen

$$
\operatorname{vol}(X, \omega)=\frac{i}{2} \iint_{X} \omega \wedge \bar{\omega}
$$

Ist $X$ kompakt, so ist $\operatorname{vol}(X, \omega)$ offenbar endlich.\\
Proposition 8.6 i) Ist ( $X, \omega$ ) eine flache Fläche und $Z(\omega)$ die Menge der Nullstellen von $\omega$, so ist $X \backslash Z(\omega)$ eine Fläche mit Translationsstruktur, indem man $X \backslash Z(\omega)$ mit einfach zusammenhängenden Mengen $U_{i}$ überdeckt, so dass die Vereinigungen $U_{i} \cup U_{j}$ ebenfalls einfach zusammenhängend sind, $x_{i} \in U_{i}$ wählt und Kartenabbildungen wie folgt definiert:

$$
g_{i}:\left\{\begin{array}{rlc}
U_{i} & \longrightarrow & \mathbb{C} \\
z & \longmapsto & \int_{x_{i}}^{z} \omega .
\end{array}\right.
$$

ii) Umgekehrt, sei $X$ eine kompakte Fläche, $Z \subseteq X$ endlich und $X \backslash Z$ habe eine Translationsstruktur. (Wir setzen nicht voraus, dass $X$ eine Riemannsche Fläche ist.) Wir nehmen zudem an, dass der Abschluß von $X \backslash Z$ bezüglich dem Abstandsbegriff der Translationsstruktur homöomorph zu $X$ ist. Dann gibt es eine eindeutige komplexe Struktur auf $X$ und bis auf skalare Vielfache genau eine holomorphe 1-Form $\omega$ auf $X$ mit $Z(\omega) \subseteq Z$, so dass die gegebene Translationsstruktur aus $\omega$ wie im ersten Teil wiedergewonnen werden kann.

Beweis: Zum Beweis von i) seien $U_{1}$ und $U_{2}$ mit Basispunkten $x_{1}$ und $x_{2}$ gegeben.\\
\includegraphics[scale=0.2, center]{2025_10_06_374d4cfa8b1948eb4c41g-097}

Wir wählen einen Weg wie im nebenstehenden Bild. Dann ist

$$
g_{1}(z)=\int_{x_{1}}^{z} \omega=\int_{\gamma} \omega+\int_{x_{2}}^{z} \omega=\int_{\gamma} \omega+g_{2}(z)
$$

d.h. die Karten unterscheiden sich um Translation durch $\int_{\gamma} \omega$.

Wir beweisen nun die Aussage ii). Auf $X \backslash Z$ definieren wir lokal auf $U_{i}$ das Differential $\omega=d z_{i}$, wobei $z_{i}$ die Koordinate auf $g_{i}\left(U_{i}\right) \subset \mathbb{C}$ ist. Dies ist wohldefiniert, denn falls sich $U_{1}$ und $U_{2}$ überlappen, so gilt nach Definition einer Translationsstruktur

$$
z_{2}=z_{1}+c, \quad \text { also } \quad d z_{2}=d\left(z_{1}+c\right)=d z_{1}
$$

Zu einem Punkt $x \in Z$ sei $V=V(x) \subset X$ eine Umgebung von $x$. Zu jedem Punkt im Rand dieser Umgebung ist der Abstand zu $x$ (im Sinne von Translationsstrukturen) wohldefiniert und endlich aufgrund der Voraussetzung über den Abschluß von $X \backslash Z$. Sei

$$
m:=\min _{p \in \partial(V)} d(p, x)
$$

Dieses Minimum ist positiv und wir betrachten den geschlossenen Weg $\gamma$, der $x$ einmal umläuft und stets den Abstand $m / 2$ zu $x$ hat. Sei $U$ das Innere des Weges $\gamma$, d.h. die Komponente von $X \backslash \gamma([0,1])$, die $x$ enthält.

Unser Ziel ist es $U$ mit Kreisscheibensegmenten analog zu den $V_{k}$ in Beispiel 8.1 zu überdecken. Dazu müssen wir zunächst den Mittelpunkt finden. Da dieser nicht in $X \backslash Z$ liegt und somit nicht von einer Karte des Translationsatlas erfasst wird, wenden wir folgende Konstruktion an. Sei $y \in \partial U$ ein beliebiger Punkt und $\ell$ eine Gerade in $U$, die senkrecht auf dem Rand von $U$ steht. Wir können alle Translationskarten um einen festen Winkel rotieren und so ohne Einschränkung annehmen, dass $\ell$ horizontal ist. In einem maximalen Translationsatlas von $X \backslash Z$ ist auch eine Karte ( $V, g$ ) enthalten, bei dem der Punkt auf $\ell$ im Abstand von $m / 2$ von $y$ (in Richtung des Innern von $U$ abgetragen) der Nullpunkt ist. Nach Rotation können wir zudem annehmen, dass $y$ auf der positiven reellen Achse liegt.

Sei $\left(V_{0}, g_{0}\right)$ eine Karte im Translationsatlas $Y o n U \mathbb{K}_{2} X$, sopass $V_{0} \supset(V \backslash \left.g^{-1}(\ell)\right) \sqrt{\text { sodass }} g_{0}\left(V_{Q}\right) \cap \mathbb{R}_{\leq 0}=\emptyset$ und sodass der De Veich Von $V_{0}$ maximal is. Aure requéfinition von $U$ ist $g_{g}\left(V_{0}\right)-$ I $\leftarrow d$ ejne geschlitzteK eisscheibe Ausgenerget von ( $V_{0}, g_{0}$ ) definieren wir nun weitere Karten, welche my $g_{0}$ auf der oberen (bzw. unteren) Halbebene übereinstimmen und maximal unter allen Karten im Translationsatlas von $U$ sind, sodass der Bildbereich nicht die positive reelle Achse trifft. Diese Karten sind wieder geschlitzte Kreisscheiben, ohne die positive reelle Achse.

Iteriert man diesen Vorgang, erhält man nach geeigneter Numerierung Karten $\left(V_{j}, g_{j}\right)$, sodass $g_{j}\left(V_{j}\right)$ geschlitzte Kreisscheiben sind, wobei der Schlitz für gerade $j$ die negative und für ungerade $i$ die positive reelle Achse ist und sodass $g_{j}^{-1}\left(\mathbb{H}^{+}\right)=g_{j+1}^{-1}\left(\mathbb{H}^{+}\right)$für gerade $j$ sowie $g_{j}^{-1}\left(\mathbb{H}^{-}\right)=g_{j+1}^{-1}\left(\mathbb{H}^{-}\right)$für ungerade $j$ ist. Die Menge dieser Karten überdeckt $U$, da nach der Homöomorphievoraussetzung $U \backslash\{x\}$ zusammenhängend ist. Diese Voraussetzung zusammen mit der Kompaktheit von $X$ besagt auch, dass $V_{0}=V_{j_{0}}$ für ein gewisses $j_{0} \in \mathbb{N}$ ist, denn sonst wären die Punkte $g_{2 k}^{-1}\left(\frac{m}{4} i\right)$ eine Folge ohne Häufungs-\\
punkt in $X$. Notwendigerweise ist $j_{0}=2 k$ gerade. (Der Punkt $x$ hat also einen Winkel $2 k \pi$ in der Sprechweise aus der Konstruktion von Verklebungsflächen.)

Damit können wir um $x$ eine Kreisscheibenkarte vom Radius $m / 2$ mit den $V_{i}$ vermöge der Einschränkung von $z \mapsto z^{k}$ auf Kreisscheibensegmente verkleben.

Wenn wir dies für alle Punkte in $Z$ durchgeführt haben, erhalten wir auf $X$ einen komplexen Atlas, sodass die Differentialform $\omega$ definiert durch $d z$ auf den Translationskarten von $X \backslash Z$ und durch $d\left(z^{k}\right)=k z^{k-1} d z$ auf der Kreisschreibenkarte um einen Punkt $x$, falls dieser den Winkel $2 k \pi$ hat, global zu einer holomorphen Differentialform zusammenpasst. Da wir auf den Translationskarten von $X \backslash Z$ die Differentialform $d z$ verwendet haben, ist dies in der Tat die Umkehrung der Konstruktion aus i).

Die Voraussetzung über den Abschluß oder ein Ersatz hierfür ist in der Tat notwendig. Um dies zu zeigen, konzentrieren wir uns auf die Umgebung (homöomorph zu $\Delta$ ) eines Punktes $x$ in $Z \subseteq X$. Für alle reellen Zahlen $r<R$ ist $\Delta \backslash\{x\}$ homöomorph zu dem Kreisring $\Delta_{r, R}=\{z: r<|z|<R\}$. Die Einbettung des Kreisrings in die komplexe Ebene definiert also einen Translationsatlas bestehend aus einer Karte von $\Delta \backslash\{x\}$. Wir werden später sehen, dass $\Delta_{r, R}$ für $0<r<R<\infty$ und die punktierte Kreisscheibe nicht isomorph als Riemannsche Flächen sind. Wenn wir dies verwenden, wird klar, dass es keine komplexe Struktur auf $\Delta$ gibt, sodass $\Delta \backslash \Delta_{r, R}$ aus nur einem Punkt besteht. Natürlich ist in diesem Kreisringbeispiel die Voraussetzung über den metrischen Abschluß nicht erfüllt.

Wir enden diesen Abschnitt mit drei wichtigen Beispielen von flachen Flächen. Zwei davon erhalten wir aus der vorangehenden Proposition.

Korollar 8.7 Sind alle Winkel des Polygons rationale Vielfache von $\pi$, so ist die Entfaltungsfläche eine flache Fläche.

Beweis: Aus der Endlichkeit der Zahl der Reflexionskopien folgt die Kompaktheit der Entfaltungsfläche. Die Existenz einer Translationsstruktur auf $X^{*}$ folgt direkt aus der Definition der Verklebungsabbildungen und die metrische Voraussetzung ist offenbar erfüllt.

Beispiel 8.8 Die Riemannschen Flächen aus Beispiel 8.1 sind ebenfalls flache Flächen, da $X \backslash X^{*}$ einpunktig ist und die metrische Voraussetzung von Proposition 8.6 erfüllt ist.

Das letzte Beispiel betrifft die Konstruktion von flachen Flächen durch Überlagerungen. Ist $f: X \longrightarrow Y$ eine surjektive holomorphe Abbildung und $\omega \in \Omega^{1}(Y)$ eine holomorphe 1-Form, so definiert man die holomorphe 1-Form $f^{*} \omega$ wie folgt. Ist $\omega=P(z) d z$ lokal in einer Karte auf $Y$ und $z_{x}$ eine Koordinate auf $X$, so ist $f^{*} \omega=P\left(f\left(z_{x}\right)\right) d f\left(z_{x}\right)$.

Da eine Riemannsche Fläche $E$ mit $g(E)=1$ bis auf skalare Vielfache genau eine holomorphe 1-Form $\omega_{E}$ besitzt, wird eine Überlagerungsfläche $f: X \longrightarrow E$ vermöge $\omega=f^{*} \omega_{E}$ in natürlicher Weise zu einer flachen Fläche.

\subsection {Die affine Gruppe}
Wir definieren nun eine Gruppe, die in natürlicher Weise einer Fläche mit Translationsstruktur zugeordnet ist, die aber im Allgemeinen sehr schwer zu berechnen sein wird. Wir beschränken uns auf den Fall kompakter Riemannscher Flächen.

Definition 8.9 (affiner Homöomorphismus) Sei ( $X, \omega$ ) eine flache Fläche und $Z=Z(\omega)$ die Nullstellenmenge von $\omega$. Ein Homöomorphismus $\varphi: X \longrightarrow X$ heißt affin bzgl. $\omega$, falls für jeden Punkt $x \in X \backslash Z$ und für Karten $(U, g)$ um $x$ und $(V, h)$ um $\varphi(x)$ im Translationsatlas die Verkettung

$$
h \circ \varphi \circ g^{-1}: \quad \mathbb{R}^{2} \cong \mathbb{C} \longrightarrow \mathbb{C} \cong \mathbb{R}^{2}
$$

auf ihrem Definitionsbereich eine affine Abbildung $\binom{x}{y} \longmapsto A\binom{x}{y}+b, A \in G L_{2}(\mathbb{R}), b \in \mathbb{R}^{2}$ von $\mathbb{R}^{2}$ ist. Die Gruppe

$$
\operatorname{Aff}(X, \omega)=\{\varphi: X \longrightarrow X \text { orientierungserhaltend und affin bzgl. } \omega\}
$$

heißt Gruppe affiner Homöomorphismen von $(X, \omega)$.\\
Aus der Definition eines Translationsatlasses folgt leicht folgende Beobachtung:

Lemma 8.10 Seien $\varphi$ affin und $g$ und $h$ Kartenabbildungen bei $x$ und $\varphi(x)$. Dann hängt der Matrixanteil

$$
D(\varphi):=A
$$

(mit der Notation aus obiger Definition) nicht von $x$ und nicht von der Kartenwahl ab. Man erhält also einen Homomorphismus

$$
D: \operatorname{Aff}(X, \omega) \longrightarrow G L_{2}(\mathbb{R})
$$

Definition 8.11 (Veechgruppe) Das Bild $D(\operatorname{Aff}(X, \omega))$ des obigen Homomorphismus wird affine Gruppe oder Veechgruppe genannt und mit SL( $X, \omega$ ) bezeichnet.

Diese Bezeichnung ist durch folgendes Lemma gerechtfertigt.\\
Lemma 8.12 Ist $X$ kompakt, so ist die affine Gruppe in $\mathrm{SL}_{2}(\mathbb{R})$ enthalten.\\
Beweis: Seien $(U, g)$ und $(V, h)$ wie in der Definition 8.9. Ist $h \circ \varphi \circ g^{-1}$ auf der Menge $W \subseteq g(U)$ definiert, so ist $\operatorname{vol}(\varphi(W))=\operatorname{det}(D(\varphi)) \operatorname{vol}(W)$. Da $D(\varphi)$ unabhängig vom Punkt $x \in X \backslash Z$ ist, folgt aus der Endlichkeit des Volumens und der Bijektivität von $\varphi$ die Behauptung.

Beispiel 8.13 Wir bestimmen nun diese Gruppen im Fall eines Torus T. Es ist

$$
T=\mathbb{C} / \Lambda, \quad \Lambda \subseteq \mathbb{C} \quad \text { Gitter }
$$

und $d z$ auf $\mathbb{C}$ definiert eine holomorphe 1-Form auf $\omega$ auf $T$. Nach Korollar 6.7 ist dies bis auf skalare Vielfache die einzige holomorphe 1-Form. Der Torus ist der Ausnahmefall mit $\# Z(\omega)=2 g-2=0$. Folglich sind alle Translationen in $\operatorname{Aff}(T, \omega)$ und haben $D=1$. Jede affine Abbildung läßt sich eindeutig als Komposition einer Abbildung, die einen gewählten Punkt $O \in T$ festhält, und einer Translation schreiben. Eine affine Abbildung, die $O \in T$ festhält, lässt sich zu einer affinen Abbildung von $\mathbb{C}$ hochheben, die $O$ als Punkt und $\Lambda$ mengenweise festhält. Diese Abbildung liegt bei geeigneter Basiswahl in $\mathrm{SL}_{2}(\mathbb{Z})$ und umgekehrt leistet jedes Element von $\mathrm{SL}_{2}(\mathbb{Z})$ das Verlangte. Wir haben also gezeigt, dass

$$
\begin{aligned}
\operatorname{Aff}(T, \omega) & =\left\langle\text { Translationen, } \mathrm{SL}_{2}(\mathbb{Z})\right\rangle \\
\operatorname{SL}(T, \omega) & =\mathrm{SL}_{2}(\mathbb{Z})
\end{aligned}
$$

gilt.\\
Es gibt ein allgemeines Kriterium zum Auffinden von Elementen in der Veechgruppe. Um dieses anzugeben führen wir einige wichtige topologische Begriffe auf flachen Flächen ein und erinnern an grundlegende Begriffe über Fuchssche Gruppen.

Ein Weg auf einer Translationsfläche, der in jeder Karte eine Gerade ist, wird als Gerade oder Geodäte bezeichnet. Ist der Weg geschlossen, so sprechen wir von geschlossenen Geodäten. Im Beispiel der $L$-Fläche können alle horizontalen Geraden mit $y$-Koordinate ungleich $0,1, b$ zu geschlossenen Geodäten verlängert werden.

Definition 8.14 (Sattelverbindung) Das Bild eines Weges $\gamma$ in einer flachen Fläche $X$ wird eine Sattelverbindung genannt, falls der Anfangs- und der Endpunkt von $\gamma$ in $Z(\omega)$ liegen und falls $\gamma$ ausserhalb diese beiden Punkte eine Gerade ist.

Beispielsweise sind in der $L$-Fläche alle Ränder von $X_{1}, X_{2}$ und $X_{3}$ Sattelverbindungen und auch die Diagonalen dieser Teilflächen sind Sattelverbindungen. Betrachtet man Geraden der Steigung $1 / k a$ für $k \in \mathbb{N}$ ausgehend von der linken unteren Ecke dieser $L$-Fläche, so ist offensichtlich, dass diese unendlich viele Sattelverbindungen in unendlich vielen Richtungen hat.

Ist $\gamma$ eine Gerade auf einer Translationsfläche, so wird $\operatorname{hol}(\gamma)=\int_{\gamma} \omega \in \mathbb{C}$ der Holonomievektor von $\gamma$ genannt.

Zwei geschlossene Geodäten auf einer Translationsfläche werden homotop (manchmal auch frei homotop genannt, falls es eine Homotopie im Sinne von Abschnitt 3.2 gibt, die die Geodäten ineinander überführt. Der Unterschied zum dortigen Begriff ist, dass die Geodäten nicht an einem festen Basispunkt beginnen müssen. Aus der Homotopieinvarianz des Wegintegrals folgt sofort die folgende Beobachtung.

Lemma 8.15 Homotope geschlossene Geodäten haben den gleichen Holonomievektor, sind also insbesondere parallel.

Ist $R_{a, b}$ das Rechteck $[0, a] \times[0, b]$, so wird die Verklebungsfläche von $R_{a, b}$ längs der vertikalen Seiten $0 \times[0, b]$ und $a \times[0, b]$ ein (horizontaler) Zylinder der Breite $a$ und Höhe $b$ genannt. Das Verhältnis $m=b / a$ wird Modulus des Zylinders genannt. Die Geraden $[0, a] \times y$ sind geschlossene Geodäten auf dem Zylinder, die auch Kernkurven genannt werden. Man beachte, dass wir unter einem Zylinder eine (nicht-kompakte) Riemannsche Fläche (und gleichzeitig eine Translationsfläche) verstehen. Daher sind die Ober- und Unterkanten von $R_{a, b}$ nicht Teil der Verklebungsfläche. Allgemein nennen wir eine Teilmenge einer flachen Fläche einen Zylinder, falls sie als Translationsfläche zu solch einem Zylinder isomorph ist.

Die Bedeutung von Zylindern ergibt sich aus der folgenden Bemerkung.\\
Lemma 8.16 Ist $X$ eine flache Fläche und $g(X) \geq 2$, so ist eine maximale Vereinigung homotoper geschlossener Geodäten ein Zylinder. Der Rand dieses Zylinders ist eine Vereinigung von Sattelverbindungen, die parallel zu den geschlossenen Geodäten sind.

Die affine Gruppe eines horizontalen Zylinders enthält stets die Translationen $z \mapsto z+c, c \in \mathbb{R}$. Wichtiger ist, dass sie auch nichttriviale Elemente enthält, die zudem den oberen und unteren Rand des Zylinders fixieren.

Lemma 8.17 Ist $C$ ein horizontaler Zylinder mit Modulus $m$, so besteht die Untergruppe der affinen Homöomorphismen, die den Rand punktweise fixieren, genau aus Elementen $\varphi$ mit

$$
D(\varphi)=\left(\begin{array}{cc}
1 & k / m \\
0 & 1
\end{array}\right)
$$

Beweis: Der affine Homöomorphismus

$$
\varphi(x+i y)= \begin{cases}\left(x+m^{-1} y\right)+i y & \text { falls } x<a-y m^{-1} \\ \left(x+m^{-1} y\right)+i y-a & \text { falls } x \geq a-y m^{-1}\end{cases}
$$

und seine Potenzen geben die geforderten Elemente. Für die Umkehrung sei $\varphi \in \operatorname{Aff}(C)$ gegeben. In der Nähe der reellen Achse ist $\varphi$ linear. Setzt man $\varphi$ nach oben fort, so folgt, dass $D(\varphi)$ den Punkt $(0, b)$ auf einen Punkt $(k a, b)$ für $k \in \mathbb{Z}$ schicken muss. Daraus folgt die Behauptung.

Man nennt von Null verschiedene reelle Zahlen $x_{1}, \ldots, x_{n}$ kommensurabel, falls es $q_{j} \in \mathbb{Q}$ gibt, sodass $x_{j}=q_{j} x_{1}$ für $j=2, \ldots, n$. In diesem Fall heißt $x$ ein gemeinsames Vielfaches, falls $x / x_{j} \in \mathbb{Z}$ ist für alle $j$.

Proposition 8.18 Sei ( $X, \omega$ ) eine flache Fläche, sodass alle horizontalen Geodäten geschlossen oder Sattelverbindungen sind, d.h. $X$ ist die Vereinigung von Abschlüssen\\
von disjunkten Zylindern $C_{1}, \ldots, C_{n}$. Sind die Moduli $m_{i}=m\left(C_{i}\right)$ dieser Zylinder kommensurabel, und ist $\ell$ ein gemeinsames Vielfaches der inversen Moduli $m_{i}^{-1}$, so enthält $\operatorname{Aff}(X, \omega)$ ein Element $\varphi$ mit

$$
D(\varphi)=\left(\begin{array}{ll}
1 & \ell \\
0 & 1
\end{array}\right)
$$

Man nennt eine Richtung wie die horizonale Richtung in dieser Proposition auch vollständig periodisch.\\
Beweis: Nach der Definition von $\ell$ und Lemma 8.17 gibt es auf jedem der Zy linder $C_{i}$ einen affinen Homöomorphismus $\varphi_{i}$ mit $D(\varphi)=\left(\begin{array}{ll}1 & \ell \\ 0 & 1\end{array}\right)$, welcher sich auf dem Rand des Zylinders stetig mit Identität fortsetzt. Die so zusammengesetzte Abbildung $\varphi: X \rightarrow X$ fixiert also alle horizontale Sattelverbindungen (und damit alle Nullstellen von $\omega$ ) Damit ist offensichtlich, dass $\varphi$ auch in den Translationskarten, welche die Sattelverbindungen am Rand der Zylinder überdecken, affin ist.

Wir behaupten nicht, dass jeder affine Homöomorphismus $\varphi$, sodass $D(\varphi)$ eine obere Dreiecksmatrix ist, notwendigerweise alle horizontale Sattelverbindungen und damit alle Nullstellen von $\omega$ fixiert.

Wir erinnern an die Klassifikation von Elementen in $\mathrm{SL}_{2}(\mathbb{R})$ nach ihrer Jordan-Normalform.\\
Hat $\gamma$ einen Eigenwert vom Betrag größer 1, so hat der andere Eigenwert einen Betrag kleiner 1 (und umgekehrt!). Der Betrag der Spur ist größer als 2 und $\gamma$ über $\mathbb{R}$ diagonalisierbar. In diesem Fall wird $\gamma$ hyperbolisch genannt.\\
Hat $\gamma$ Eigenwert vom Betrag gleich 1, so gibt hat der andere Eigenwert auch Betrag 1, die Jordan-Normalform ist eine obere Dreiecksmatrix und $\gamma$ wird parabolisch genannt. In diesem Fall hat die Spur von $\gamma$ den Betrag zwei.\\
Hat $\gamma$ schließlich keine reellen Eigenwerte, so ist der Betrag der Spur von $\gamma$ echt kleiner als zwei und $\gamma$ zu einer Drehmatrix konjugiert. In diesem Falls wird $\gamma$ elliptisch genannt.

Definition 8.19 (Fuchssche Gruppe) Eine Untergruppe von $\mathrm{SL}_{2}(\mathbb{R})$, die diskret ist, wird Fuchssche Gruppe genannt.

Beispiele hierfür sind 'grosse Gruppen' wie $\mathrm{SL}_{2}(\mathbb{Z})$ oder auch 'kleine' wie z.B. zyklische Untergruppen erzeugt von einem parabolischen oder hyperbolischen Element.

Satz 8.20 Ist $(X, \omega)$ eine flache Fläche, so ist die Veechgruppe eine Fuchssche Gruppe.\\
Dazu definieren wir für eine flache Fläche ( $X, \omega$ ) mit $g(X) \geq 2$, für welche also $Z(\omega)$ nicht-leer ist,

$$
\operatorname{hol}(X, \omega)=\{\operatorname{hol}(\gamma), \gamma \subset X \text { Sattelverbindung }\}
$$

die Menge aller Holonomievektoren von Sattelverbindungen.

Proposition 8.21 Sei $(X, \omega)$ eine flache Fläche mit $g(X) \geq 2$. Dann ist die Menge der Richtungen von Sattelverbindungen dicht im Einheitskreis. Weiterhin ist die Menge $\operatorname{hol}(X, \omega)$ diskret in $\mathbb{C}$.

Da jeder Holonomievektor endlich nur oft auftritt, höchstens $|Z(\omega)|$-mal, ist die letzte Aussage unabhängig davon, ob $\operatorname{hol}(X, \omega)$ mit Vielfachheit definiert ist oder nicht.

Wir zeigen zunächst, wofür wir diese Aussagen benötigen.\\
Beweis von Satz 8.20. Für den Torus haben wir die Aussage bereits in Beispiel 8.13 behandelt. Wir betrachen ab sofort also den Fall $g(X) \geq 2$. Angenommen $\operatorname{SL}(X, \omega)$ enthält eine Folge $\gamma_{n}$, die gegen ein Element $\gamma \in S L_{2}(\mathbb{R})$ konvergiert. Durch Postkomposition mit $\gamma^{-1}$ können wir annehmen, dass $\gamma$ die Identität ist. Für jedes $v \in \operatorname{hol}(X, \omega)$ ist also $\gamma_{n}(v) \in \operatorname{hol}(X, \omega)$ und $\gamma_{n}(v)$ konvergiert gegen $\gamma(v)=v$. Nach Proposition 8.21 muss also die Folge $\gamma_{n}(v)$ irgendwann stationär werden, d.h. $v$ ist ein Eigenvektor von $\gamma_{n}(v)$ für alle $n \geq N_{0}$.

Ebenfalls aus Proposition 8.21 folgt die Existenz von drei Holonomievektoren, von denen keine zwei proportional sind. Das orbige Argument auf diese drei Holonomievektoren angewandt zeigt, dass es irgendwann ein $N \in \mathbb{N}$ gibt, sodass $\gamma_{n}$ die Identität für $n \geq \mathbb{N}$ ist.

\subsection {Intervallaustauschtransformationen}
Zunächst Mit dem Ziel Proposition 8.21 zu beweisen, führen wir einen Begriff ein, der bei Fragen des dynamischen Systems 'flache Fläche' dutzendfach Anwendung findet.

Definition 8.22 (Intervallaustauschtransformation) Sei $I=(a, b]$ ein Intervall. Eine bijektive Abbildung $T: I \rightarrow I$ heißt Intervallaustauschtransformation, falls es eine Zerlegung von $I$ in Teilintervalle $I=\cup_{i=1}^{n}\left(x_{i}, x_{i+1}\right]$ mit $x_{i}<x_{i+1}$ gibt, sodass $\left.T\right|_{\left(x_{i}, x_{i+1}\right]}$ eine Translation ist.

Die Anzahl der Zerlegungspunkte ist endlich, aber nicht universell festgelegt. Diese enthalten die Unstetigkeitsstellen von $T$, möglicherweise strikt. Intervallaustauschtransformationen bilden mit der Hintereinanderausführung eine Gruppe, offenbar nicht abelsch und 'kompliziert'. Man kann z.B. zeigen, dass sie perfekt ist, d.h. gleich ihrer Kommutatorgruppe.

Sei $v \in \mathbb{R}^{2}$ eine vorgegebene Richtung und $I \subset X$ ein Geradenstück in einer flachen Fläche ( $X, \omega$ ), welches senkrecht auf $v$ steht. Wir werden oft annehmen, dass ein Ende von $I \in Z(\omega)$ ist.

Wir betrachten Geraden die senkrecht auf $I$ stehen, d.h. Geraden in Richtung $v$. Für alle Punkte in $X$ ausserhalb einer Nullmenge (nämlich eine endliche Vereinigung von Halbgeraden, welche in $Z(\omega)$ enden) ist die Verschiebung um $t$ in Richtung $v$ für alle $t \in \mathbb{R}$ definiert. Diese Verschiebung um $t$ definiert\\
einen Fluß auf $X$, eine Abbildung, welchem dem Poincaré-Rekurrenzsatz genügt. Dieser besagt, dass unter dem Fluss jeder Punkt mit Ausnahme einer Nullmenge unendlich oft zu einer vorgegebenen Menge positiven LebesqueMaßes zurückkehrt, sogar in einer unbeschränkten Folge von Zeiten. Angewandt auf eine rechteckige Umgebung von $I$ besagt dies, dass abgesehen von abzählbar vielen Punkten alle Punkte von $I$ unter der Verschiebung in Richtung $v$ in einer unbeschränkten Folge von Zeiten wieder zu $I$ zurückkehren.

Die Menge der Punkte, sodass man unter Verschiebung in Richtung $v$ wieder zu $I$ zurückkehrt, ist anderseits offen: Man kann das Geradenstück von $x$ in Richtung $v$ bis zur Rückkehr nach links und rechts verschieben, bis man ein Punkt in $Z(\omega)$ berührt. Also ist die Menge der Punkte, bei denen diese Abbildung nicht definiert ist, endlich. Zusammengefasst haben wir folgendes gezeigt.

Lemma 8.23 Ist $I \subset X$ ein Intervall in einer flachen Fläche, senkrecht zur Richtung v, so definiert die Erstrückkehrabbildung nach Verschiebung in Richtung v nach linksseitig stetiger Fortsetzung eine Intervallaustauschtransformation auf I.

Beweis von Proposition 8.21, Wir starten mit der Dichtheit der Richtungen von Sattelverbindungen. Sei die Richtung $v \in \mathbb{R}^{2}$ vorgegeben und $x \in Z(w)$. Wir wollen zeigen, dass zu vorgegebenem $\varepsilon>0$ es beginnend bei $x$ eine Sattelverbindung gibt, die mit $v$ einen Winkel kleiner als $\varepsilon$ hat. Sei $\ell$ eine Gerade beginnend bei $x$ in Richtung $v$. Endet $\ell$ in einem Punkt von $Z(\omega)$ sind wir fertig. Andernfalls kommt $\ell$ für jedes vorgegebene $s$ in eine $s$-Umgebung von $x$ zurück, trifft also eines der $2\left(\operatorname{ord}_{x}(\omega)+1\right)$ Intervalle der Länge $s$ beginnend bei $x$, welche senkrecht auf $v$ stehen. Wir nennen dieses Intervall $I$. Nach den obigen Überlegungen trifft $\ell$ also $I$ sogar in einer unendlichen Folge von Punkten mit Abstand $s_{n}$ zu $x$. Diese Punkte seien $\ell_{n}$ von $x$ weg, gemessen entlang $\ell$. Wir betrachten nun Geraden $g_{\delta}$, welche bei $x$ starten, mit $v$ einen kleinen Winkel $0 \leq \delta$ bilden, sodass der Winkel zu $I$ stumpf ist. Für $\delta$ klein genug ist der $n$-te Schnittpunkt mit $I$ bei $\ell_{n} / \cos (\delta)<\ell_{n}+\varepsilon$ (gemessen entlang $g_{\delta}$ ) und im Abstand $\ell_{n} \tan (\delta)$ von $s_{n}$ gemessen entlang $I$. Man vergrößert nun $\delta$, bis $g_{\delta}$ eine Singularität trifft. Dies geschieht spätestens bei $\delta=\arctan \left(s_{n} / \ell_{n}\right)$, da dann der Schnittpunkt auf $I$ gegen $x$ gelaufen ist. Da $s_{n}<s$ universell beschränkt ist und die $\ell_{n}$ eine unbeschränkte Folge bilden, kann man durch genügend große Wahl von $n$ erreichen, dass $\delta<\varepsilon$ ist.

Jeder Punkt in $Z(\omega)$ hat eine Umgebung, die $Z(\omega)$ in nur einem Punkt trifft. Also gibt es ein $\varepsilon>0$, sodass jedes Element in $\operatorname{hol}(X, \omega)$ die Länge mindestens $\varepsilon$ hat.

Das Argument lässt sich nicht nur auf den Nullpunkt anwenden. Sei $v \in \operatorname{hol}(X, \omega)$. Wir tragen von jedem Punkt $x \in Z(\omega)$ den Vektor $v$ in alle $\operatorname{ord}_{x}(\omega)+1$ möglichen Richtungen ab. Die Endpunkte hiervon sind eine endlich Menge von Punkten, höchstens $4 g(X)-4$. Es gibt wiederum ein $\varepsilon_{v}>0$, sodass die\\
$\varepsilon_{v}$-Umgebungen um diese Punkte disjunkt zu $Z(\omega)$ sind. Also enthält die $\varepsilon_{v^{-}}$ Umgebung von $v$ kein Element von $\operatorname{hol}(X, \omega)$ außer $v$.

\subsection {Translationsüberlagerungen}
Als nächstes untersuchen wir das Verhalten der affinen Gruppe beim Übergang zu Überlagerungen. Eine Überlagerung

$$
f:(X, \omega) \longrightarrow(Y, \eta)
$$

von flachen Flächen heißt Translationsüberlagerung, falls $f^{*} \eta=\omega$ gilt. Wir nennen $f$ eine strikte Translationsüberlagerung, falls die Verzweigungspunkte von $f$ nach $Z(\eta)$ abgebildet werden oder, im Fall $g(Y)=1$, falls alle Verzweigungspunkte auf einen Punkt abgebildet werden. Ohne diesen Zusatz müssten alle strikten Translationsüberlagerungen des Torus nach Riemann-Hurwitz wieder Tori sein und die folgenden Überlegungen wären ziemlich gegenstandslos.

Beispiel 8.24 Die Verklebungsflächen aus Beispiel 8.1 mit $a, b \in \mathbb{R}$ sind strikte Translationsüberlagerungen des Torus mit Perioden 1 und i, falls $a, b \in \mathbb{Q}$. Die Überlagerung ist offensichtlich, wenn man die L-Fläche in Quadrate unterteilt, deren Seitenlänge das kleinste gemeinsame Vielfache der Nenner von a und b ist.

Ziel dieses Abschnitts ist folgender Satz, dessen Beweis den Rest des Abschnitts füllt.

Satz 8.25 Sind $X$ und $Y$ kompakt und ist $f:(X, \omega) \longrightarrow(Y, \eta)$ eine strikte Translationsüberlagerung, so sind $\mathrm{SL}(X, \omega)$ und $\mathrm{SL}(Y, \eta)$ kommensurabel, d.h. es gibt eine Untergruppe $\Gamma \subset \mathrm{SL}_{2}(\mathbb{R})$, die Untergruppe von endlichem Index sowohl in $\operatorname{SL}(X, \omega)$ als auch in $\operatorname{SL}(Y, \eta)$ ist.

Aus technischen Gründen führen wir die Gruppen $\operatorname{Aff}^{f}(Y, \eta)$ und $\operatorname{Aff}_{f}(X, \omega)$ ein, die aus den affinen Homöomorphismen $\psi$ (bzw. $\varphi$ ) bestehen, die sich via $f$ liften (bzw. via $f$ herunterdrücken) lassen, d.h. genauer

$$
\operatorname{Aff}^{f}(Y, \eta)=\{\psi \in \operatorname{Aff}(Y, \eta): \text { es existiert } \widetilde{\psi} \in \operatorname{Aff}(X, \omega): f \circ \widetilde{\psi}=\psi \circ f\}
$$

und

$$
\operatorname{Aff}_{f}(X, \omega)=\{\varphi \in \operatorname{Aff}(X, \omega): \text { es existiert } \bar{\varphi} \in \operatorname{Aff}(Y, \eta): f \circ \varphi=\bar{\varphi} \circ f\}
$$

Die jeweiligen Bilder unter $D$ bezeichnen wir mit $\mathrm{SL}^{f}(Y, \eta)$ und $\mathrm{SL}_{f}(X, \omega)$. Wir beginnen den Beweis mit Abschätzungen über die Liftbarkeit von Homöomorphismen.

Proposition 8.26 Der Index $\left[\operatorname{Aff}(Y, \eta): \operatorname{Aff}^{f}(Y, \eta)\right]$ ist endlich.

Beweis: Die Überlagerung $f$ schränkt sich zu einer Überlagerung

$$
f_{0}: X \backslash f^{-1}(Z(\eta)) \longrightarrow Y \backslash Z(\eta)
$$

ein, welche aufgrund der Voraussetzung strikte Translationsüberlagerung unverzweigt ist.\\
Nach Satz 3.43 ist $f_{0}$ eindeutig durch eine Untergruppe von $\pi_{1}(Y \backslash Z(\eta))$ vom Index $d=\operatorname{deg}\left(f_{0}\right)$ gegeben, nämlich durch $f_{*}\left(\pi_{1}\left(X \backslash f^{-1}(Z(\eta))\right)\right)$. Da $\pi_{1}(Y \backslash Z(\eta))$ endlich erzeugt ist, gibt es endlich viele, sagen wir $k$, Untergruppen von Index $d$. (Man betrachte dazu die Untergruppen als Stabilisator der Eins unter einem Homomorphismus $\pi_{1}(Y \backslash Z(\eta)) \rightarrow S_{d}$.)

Ist $\psi$ in $\operatorname{Aff}(Y, \eta)$, so bewahrt $\psi$ die Menge $Z(\eta)$ - wenn auch nicht notwendigerweise punktweise. Also ist

$$
(\psi \circ f)_{0}=\left.(\psi \circ f)\right|_{(\psi \circ f)^{-1}(Y \backslash Z(\eta))}
$$

wieder eine unverzweigte Überlagerung von $Y \backslash Z(\eta)$ von Grad $d$. Bestimmt diese eine zu $f_{*}\left(\pi_{1}\left(X \backslash f^{-1}(Z(\eta))\right)\right)$ konjugierte Untergruppe (also dieselbe Untergruppe nach Wahl eines geeigneten Basispunkts, welche wir hier in der Notation unterdrücken), so gibt es einen Homöomorphismus $\widetilde{\psi}$ nach Lemma 3.21, sodass $f \circ \widetilde{\psi}=\psi \circ f$. Wegen der Voraussetzung Translationsüberlagerung ist $\widetilde{\psi}$ offenbar affin, denn dies ist eine lokal verifizierbare Eigenschaft, also $\psi \in \operatorname{Aff}^{f}(Y, \eta)$. Wir betrachten die Operation von $\operatorname{Aff}(Y, \eta)$ auf den Untergruppen von Index $d$ von $\pi_{1}(Y \backslash Z(\eta))$ via $\psi \mapsto\left(U \mapsto \psi_{*} U\right)$ und stellen fest, dass eine Untergruppe vom Index höchstens $k$ die ursprüngliche Untergruppe $\left(f_{*}\left(\pi_{1}\left(X \backslash f^{-1}(Z(\eta))\right)\right)\right)$ sogar fixiert, also insbesondere konjugiert ist. Daraus folgt die Behauptung über den endlichen Index.

Um die Frage zu klären, welche affinen Homömorphismen auf $X$ via einer Überlagerung zu einem affinen Homöomorphismus von $Y$ absteigen, werden wir aus beweistechnischen Gründen die flache Struktur der Riemannschen Flächen „unten" modifizieren müssen. Dabei tritt zum ersten Mal eine $\mathrm{GL}_{2}(\mathbb{R})$-Aktion auf flachen Flächen auf, die uns später noch oft begegnet.

Satz 8.27 Ist $(X, \omega)$ eine flache Fläche und $A \in G L_{2}(\mathbb{R})$, so gibt es auf dem $X$ zugrundeliegenden topologischen Raum genau eine komplexe Struktur $Y$ und eine holomorphe 1-Form $\eta$, sodass $Z(\eta)=Z(\omega)$ und sodass für einen Translationsatlas $\left(U_{i}, G_{i}\right)$ von $X \backslash Z(\omega)$ die Abbildungen

$$
A \circ g_{i}: U_{i} \xrightarrow{g_{i}} \mathbb{C} \cong \mathbb{R}^{2} \xrightarrow{A \cdot} \mathbb{R}^{2} \cong \mathbb{C}
$$

einen Translationsatlas auf $Y \backslash Z(\eta)$ zu $\eta$ definieren. Diese Zuordnung $A \cdot(X, \omega):= (Y, \eta)$ definiert eine Aktion von $G L_{2}(\mathbb{R})$ auf der Menge der flachen Flächen. Sie läßt das Geschlecht der flachen Fläche, die Zahl der Nullstellen von $\omega$, sowie deren Nullstellenordnung invariant.

Beweis: Wir definieren auf $X \backslash Z(\omega)$ einen komplexen Atlas durch die Karten (8.1). Dies ist ein Translationsatlas, da ( $U_{i}, g_{i}$ ) diese Eigenschaft hat und die Verkettung mit $A=\left(\begin{array}{ll}a & b \\ c & d\end{array}\right) \in \mathrm{GL}_{2}(\mathbb{R})$ eine Translation um $x+i y$ in eine Translation um $(a x+b y)+i(c x+d y)$ überführt. Wir definieren $\eta=d z$ bezüglich der neuen Karten (8.1) und erhalten so eine holomorphe 1-Form auf $Y \backslash Z(\omega)$. Die komplexe Struktur in einer Umgebung jedes Punktes von $Z(\omega)$ definieren wir wie im Beweis von Proposition 8.6.

Die zweite Etappe des Beweises von Satz 8.25 klärt, wie viele Homömorphismen via $f$ absteigen.

Proposition 8.28 Sei $f:(X, \omega) \longrightarrow(Y, \eta)$ eine strikte Translationsüberlagerung. Dann ist $\left[\operatorname{Aff}(X, \omega): \operatorname{Aff}_{f}(X, \omega)\right]$ endlich.

Folgende Beweisidee liegt nahe: Sei $\varphi \in \operatorname{Aff}(X, \omega)$ vorgegeben. Dann ist ein Kandidat für $\bar{\varphi} \in \operatorname{Aff}(Y, \eta)$ gegeben durch die Aktion von $D(\varphi) \in \mathrm{SL}_{2}(\mathbb{R})$. Dies setzt natürlich voraus, dass $D(\varphi) \cdot(Y, \eta)$ zu $(Y, \eta)$ isomorph ist. Jedenfalls kommutiert folgendes Diagramm\\
\includegraphics[scale=0.2, center]{2025_10_06_374d4cfa8b1948eb4c41g-108}\\
wobei die vertikalen Pfeile strikte Translationsüberlagerungen sind. Der Index $\left[\operatorname{Aff}(X, \omega): \operatorname{Aff}_{f}(X, \omega)\right]$ ist also beschränkt durch die Anzahl von strikten Translationsüberlagerungen vom Grad $\operatorname{deg}(f)$ ausgehend von ( $X, \omega$ ) und mit Ziel eine Fläche vom Geschlecht $g(Y)$. Diese Schranke wird nicht unmittelbar durch die Untergruppen von vorgegebenem Index wie im vorigen Lemma geliefert.

Wir gehen also anders vor und definieren zunächst eine Zellzerlegung von flachen Flächen $(X, \omega)$. Es sei $\Sigma \supseteq Z(\omega)$ endlich, falls $g(x) \geqslant 2$ und falls $g(x)=$ 1 , so enthalte $\Sigma$ einen beliebigen fest gewählten Punkt $P$. Wir setzen für $x \in X$

$$
d(x, \Sigma):=\min _{y \in \Sigma} d(x, y)
$$

Das Gerippe von ( $X, \Sigma$ ) ist die Menge der Punkte $x \in X$, so dass es mindestens zwei Geradensegmente gibt, die den Abstand $d(x, \Sigma)$ realisieren, d.h. einen Endpunkt in $x$ und den anderen in einem Punkt von $\Sigma$ haben. Wir bezeichnen es mit $\operatorname{Ger}(X, \Sigma)$. Sei $K^{\prime} \subset \operatorname{Ger}(X, \Sigma)$ die Menge von Punkten, sodass genau 2 Geradensegmente den Abstand realisieren. $K^{\prime}$ zerfällt in die disjunkte Vereinigung von offenen Geradensegmenten mit Randpunkten in $\operatorname{Ger}(X, \Sigma) \backslash K^{\prime}$. Wir nennen diese Geradensegmente Knochen und bezeichnen die Menge der Knochen mit $K$. Mit den Knochen als Kanten wird das Gerippe ein Graph, der in $X$ eingebettet ist. Die Zusammenhangskomponenten von $X \backslash \operatorname{Ger}(X, \Sigma)$ heißen Schirme. Die Voronoi-Zerlegung von $X$ ist die Zellzerlegung, deren 2-Zellen

\begin{figure}[h]
\begin{center}
  \includegraphics[scale=0.2]{2025_10_06_374d4cfa8b1948eb4c41g-109}
 
\caption{Abbildung 8.5: Voronoi-Zerlegung (rot) und Gerippe-Triangulierung (blau)}
\end{center}
\end{figure}

die Schirme sind. Jeder Punkt von $\Sigma$ liegt in genau einem Schirm. Die Knochen zusammen mit den Segmenten von $x \in \Sigma$ zu den Eckpunkten des Gerippes am Rand des Schirms, der $x$ enthält, bilden eine Triangulierung von $X$. Wir nennen sie die Gerippe-Triangulierung $T_{X, \Sigma}$. (Sie ist eine Verfeinerung der Voronoi-Triangulierung - für Leser, die mit diesem Begriff vertraut sind.)

Wir zeigen der Reihe nach:\\
Lemma 8.29 Der Index von $\operatorname{Aff}_{f}(X, \omega)$ in $\operatorname{Aff}(X, \omega)$ ist beschränkt durch die Anzahl der Isomorphieklassen von strikten Translationsüberlagerungen $\hat{f}:(X, \omega) \longrightarrow (\widehat{Y}, \hat{\eta})$ mit $g(\widehat{Y})=g(Y)$, wobei zwei solche Überlagerungen $\hat{f}_{i}:(X, \omega) \longrightarrow\left(\widehat{Y}_{i}, \widehat{\omega}_{i}\right)$ isomorph genannt werden, falls es eine holomorphe Abbildung $g: \widehat{Y_{1}} \longrightarrow \widehat{Y}_{2}$ gibt mit $g^{*} \widehat{\omega}_{2}=\widehat{\omega}_{1}$ und $g \circ \hat{f}_{1}=\hat{f}_{2}$.

Beweis: Ist $\varphi \in \operatorname{Aff}(X, \omega)$, so ist $D(\varphi)^{-1} \circ f \circ \varphi$ eine strikte Translationsüberlagerung. Falls diese zu $f$ in obigem Sinne isomorph ist, so liegt $\varphi$ in $\operatorname{Aff}_{f}(X, \omega)$. Da die Zuordnung $\varphi \longmapsto D(\varphi)^{-1} \circ f \circ \varphi$ eine Gruppenoperation ist, folgt die Behauptung. \(\square\)

Lemma 8.30 Sei $f$ eine strikte Translationsüberlagerung, $\Sigma=Z(\eta), \Sigma_{X}= f^{-1}(Z(\eta))$. Die Abbildung, die jedem solchen $f$ die polygonale Abbildung $\mathcal{T}_{f}$ : $\mathcal{T}_{X, \Sigma_{X}} \longmapsto \mathcal{T}_{Y, \Sigma}$ zwischen den Gerippetriangulierungen zuordnet, ist injektiv.

Beweis: Man überzeugt sich leicht, dass $y$ den Abstand $\left(x, \Sigma_{X}\right)$ genau dann realisiert, wenn $f(y)$ den Abstand $(f(y), \Sigma)$ realisiert. Dabei verwenden wir, dass $\Sigma_{X}=f^{-1}(\Sigma)$ ist. Folglich ist $f^{-1}(\operatorname{Ger}(Y, \Sigma))=\operatorname{Ger}\left(X, \Sigma_{X}\right)$, die Urbilder von Knochen bzw. Knoten sind genau die Knochen bzw. Knoten in $\operatorname{Ger}\left(X, \Sigma_{X}\right)$.\\
Folglich werden die 1-Zellen von $X$ isometrisch auf die 1-Zellen von $Y$ abgebildet. Die Längen eines Dreiecks bestimmen eindeutig eine flache Metrik in\\
seinem Inneren. Also werden auch die 2-Zellen von $T_{X, \Sigma_{X}}$ isometrisch auf die 2-Zellen von $T_{Y, \Sigma}$ abgebildet. Insgesamt ist die Abbildung $f: X \longrightarrow Y$ als Abbildung zwischen metrischen Räumen durch die Abbildung $\mathcal{T}_{f}$ bestimmt und damit auch als Translationsüberlagerung.

Lemma 8.31 Die Anzahl der Isomorphieklassen (im Sinne von Lemma 8.29) von strikten Translationsüberlagerungen $f_{i}:(X, \omega) \longrightarrow\left(Y_{i}, \eta_{i}\right)$ mit fester Ausgangsfläche $(X, \omega)$ und variabler Zielfläche $Y_{i}$ mit $g\left(Y_{i}\right)=g(Y)$ ist endlich.

Beweis: Sei $\Sigma=Z\left(\eta_{i}\right)$ und $\Sigma_{X}=f_{i}^{-1}(\Sigma)$ wie in Lemma 8.30. Es ist $|\Sigma| \leqslant 2 g(Y)-2$ und folglich $\left|\Sigma_{X}\right| \leqslant \operatorname{deg}(f) \cdot(2 g(Y)-2)$. Wir müssen die Anzahl der Abbildungen von Zellkomplexen $\mathcal{T}_{X, \Sigma_{X}} \longrightarrow \mathcal{T}_{Y_{i}, \Sigma}$ beschränken. Wenn wir eine Schranke $B$ für die Anzahl $T(X)$ von Dreiecken auf $X$ kennen, so ist die Anzahl von Dreiecken in $\mathcal{T}_{Y, \Sigma}$ ebenfalls durch $B$ beschränkt. Da es maximal drei orientierungserhaltende isometrische Abbildungen zwischen zwei Dreiecken gibt, ist eine große Abschätzung für die Gesamtzahl der Abbildungen $(3 B)^{B}$.

Für jede Zellenzerlegung von $X$ mit $E$ Ecken, $K$ Kanten und $Z$ Zellen folgt aus der Eulerformel Satz 3.35 und dem Verhältnis von Zellen und Kanten

$$
Z \leqslant 2 \cdot(E+2 g(X)-2)
$$

wobei Gleichheit erreicht wird, falls alle Zellen Dreiecke sind. Geht man zur dualen Zellenzerlegung über (d.h. Zellen werden nun zu Ecken und jede Ecke wird zu einer Zelle) so folgt auch die Ungleichung

$$
E \leqslant 2 \cdot(Z+2 g(X)-2)
$$

für jede Zellenzerlegung.\\
Die Ecken $E_{\tau}$ der Gerippetriangulierung sind die Ecken $E_{V}$ der VoronoiZellzerlegung und $\Sigma_{X}$, was in Bijektion zu den Zellen $Z_{V}$ der VoronoiZellzerlegung steht. Also ist

$$
\begin{aligned}
T(X) & =2 \cdot\left(E_{\tau}+2 g(X)-2\right)=2 \cdot\left(E_{V}+Z_{V}+2 g(X)-2\right) \\
& \leqslant 6 \cdot Z_{V}+12 g(X)-12=6 \cdot\left|\sum_{X}\right|+12 g(X)-12
\end{aligned}
$$

beschränkt, was zu zeigen war.\\
Beweis von Proposition 8.28. Mit der Vorbemerkung unmittelbar nach der Proposition ist dies eine direkte Konsequenz von Lemma 8.29, Lemma 8.30 und Lemma 8.31.\\
Beweis von Satz 8.25. Die Zuordnung $\varphi \longmapsto \bar{\varphi}$ aus Proposition 8.28 ist eine surjektive Abbildung $f_{*}: \operatorname{Aff}_{f}(X, \omega) \longrightarrow \operatorname{Aff}^{f}(Y, \eta)$, denn Proposition 8.26 definiert eine Rechtsinverse. Der Kern von $f_{*}$ sind Homöomorphismen $\varphi$ mit $D(\varphi)=I \in \mathrm{SL}_{2}(\mathbb{R})$. Also induziert $f_{*}$ einen Isomorphismus $f_{*}: \mathrm{SL}_{f}(X, \omega) \longmapsto \mathrm{SL}^{f}(X, \omega)$. Aufgefasst als Untergruppe von $\mathrm{SL}_{2}(\mathbb{R})$ ist dies die gesuchte Gruppe $\Gamma$.\\
\includegraphics[scale=0.2, center]{2025_10_06_374d4cfa8b1948eb4c41g-111}\\
\includegraphics[scale=0.2, center]{2025_10_06_374d4cfa8b1948eb4c41g-111(1)}

\section {Weierstraßpunkte und Automorphismen}
Auf flachen Flächen sind Singularitäten die offensichtlichen speziellen Punkte. Bei Riemannschen Flächen ohne das Zusatzdatum der holomorphen 1-Form gibt es auch besondere Punkte, die mit Hilfe aller 1-Formen und Überlegungen wie im Satz von Riemann-Roch definiert sind. Mit Hilfe dieser Punkte zeigen wir, dass die Automorphismengruppen von kompakten Riemannschen Flächen $X$ mit $g(X) \geqslant 2$ stets endlich ist.

\subsection {Lückenreihe und Wronski-Determinante}
In diesem Abschnitt sind $X$ und $Y$ stets kompakte Riemannsche Flächen vom Geschlecht größer als Null.

Definition 9.1 (Lückenreihe) $Z u x \in X$ sei

$$
\Gamma_{x}=\left\{n \in \mathbb{N}_{>0}: \exists f \in \mathcal{L}(n[x]) \backslash \mathcal{L}((n-1)[x])\right\}
$$

die Halbgruppe (von Polordnungen holomorpher Funktionen) zu x. Das Komplement\\
$\mathbb{N} \backslash \Gamma_{x}$ heißt Lückenreihe von $x$.\\
Dies ist wirklich eine Halbgruppe, wie aus $\operatorname{ord}_{x}(f g)=\operatorname{ord}_{x}(f)+\operatorname{ord}_{x}(g)$ folgt.

Lemma 9.2 Für einen Punkt $x \in X$ sind äquivalent:\\
i) $n \in \Gamma_{x}$.\\
ii) Es gibt keine holomorphe Differentialform auf $X$, die bei $x$ eine Nullstelle der Ordnung genau $n-1$ hat.\\
iii) $\ell(n[x])=\ell((n-1)[x])+1$.

Beweis: Die Äquivalenz von i) und iii) folgt indem man die Abbildung betrachtet, die einer Funktion in $\mathcal{L}((n+1)[x])$ ihren $n$-ten Taylorkoeffizienten bei $x$ zuordnet. Die Äquivalenz von ii) und iii) folgt unmittelbar aus dem Satz von Riemann-Roch. \(\square\)

Die Halbgruppe (oder die Lückenreihe) wird dann zu einer nützlichen Invariante, wenn wir wissen, was bei fast allen Punkten von $X$ einerseits und an speziellen Punkten andererseits passieren kann.

Proposition 9.3 Es gibt für jeden Punkt $x \in X$ genau $g=g(X)$ Lücken. Listen wir sie aufsteigend als $n_{1}, \ldots, n_{g}$, so gilt $n_{1}=1$ und $n_{g} \leqslant 2 g-1$.

Beweis: Da $X$ nicht $\mathbb{P}^{1}$ ist, folgt $\ell(0)=\ell([x])=1$, also $n_{1}=1$. Nach RiemannRoch ist $\ell(n[x])=n+1-g$ für $n \geqslant 2 g-1$. Da die Folge $\ell(n[x])$ um Null oder Eins wächst, muss sie zwischen $n=0$ und $n=2 g-1$ genau $g$-mal um Null wachsen.

Definition 9.4 (Weierstraßpunkt) Ein Punkt $x \in X$ heißt Weierstraßpunkt, falls die Menge der Lücken nicht mit $\{1, \ldots, g\}$ übereinstimmt.

Dies ist äquivalent dazu, dass ein $i \in\{1, \ldots, g\}$ in der Halbgruppe von $x$ liegt, d.h. dass es eine holomorphe 1-Form mit Nullstellenordnung $\geq g$ bei $x$ gibt (denn $\ell(K-g[x])=\ell(g[x])-1$ ). Zu $x \in X$ definieren wir das Gewicht

$$
w(x)=\sum_{i=1}^{g}\left(n_{i}-i\right)
$$

Offenbar ist $w(x) \geqslant 0$ und $w(x)>0$ genau dann, wenn $x$ ein Weierstraßpunkt ist. Hauptziel dieses Abschnitts ist der folgende Satz.

Satz 9.5 Die Summe der Gewichte der Weierstraßpunkte erfüllt die Bedingung

$$
\sum_{x \in X} w(x)=g\left(g^{2}-1\right) .
$$

Insbesondere gibt es höchstens $g\left(g^{2}-1\right)$ Weierstraßpunkte auf $X$.\\
Dazu treffen wir zunächst einige Vorbereitungen. In Abschnitt 4 haben wir Differentiale erster Ordnung oder 1-Formen kartenweise als $\omega=\alpha d x+\beta d y$ eingeführt und gefordert, dass bei Kartenwechsel die naheliegende zweidimensionale Kettenregel gilt. Als wir in Definition 4.4 holomorphe Differentiale erster Ordnung eingeführt haben, haben wir festgestellt, dass für holomorphe Kartenwechsel die eindimensionale Kettenregel (mit einer komplexen Variablen) gilt, d.h. ist $t=\varphi(z)$ die Kartenwechselabbildung und $\omega=\alpha_{1}(t) d t= \alpha_{2}(z) d z$ auf dem gemeinsamen Definitionsbereich, so ist

$$
\begin{gathered}
\alpha_{1}(t) d t=\alpha_{1}(\varphi(z)) d \varphi(z)=\alpha_{1}(\varphi(z)) \frac{d \varphi}{d z}(z) d z=\alpha_{2}(z) d z \\
\text { also ist } \quad \alpha_{1} \cdot \frac{d \varphi}{d z}=\alpha_{2}
\end{gathered}
$$

Dieses Konzept von Differentialformen kann man auf zwei Arten multilinear erweitern. Man kann äußere Potenzen nehmen, aus Dimensionsgründen bei Riemannschen Flächen höchstens zweite, danach landet man beim Nullvektorraum. Damit erhält man das Konzept von Differentialen zweiter Ordnung oder 2-Formen, das bereits mehrfach verwendet wurde. Oder man $\operatorname{kann} q$ fache Tensorpotenzen nehmen für beliebige $q \geqslant 1$ und erhält folgendes Konzept.

Definition 9.6 ( $q$-Differentiale) Ein $q$-Differential $\omega$ auf $X$ ist die Zuordnung einer Funktion $\alpha(z)$ zu jeder Karte mit Koordinate z, geschrieben $\omega=\alpha(z)(d z)^{q}$, sodass sich die $\alpha(z)$ unter Kartenwechsel $t=\varphi(z)$ mit dem Faktor $\left(\frac{d \varphi}{d z}\right)^{q}$ transformieren.

In unserer Notation gibt die Zahl vor „Differential" also eine Tensorpotenz, die Zahl vor „Form" also eine äußere Potenz an. 2-Differentiale heißen auch quadratische Differentiale. Sie spielen später noch eine wichtige Rolle, wenn wir alle Riemannschen Flächen in einer Mannigfaltigkeit parametrisieren wollen. Man beachte, dass die Nullstellenordnung eines $q$-Differentials definiert als Nullstellenordnung von $\alpha$ wohldefiniert ist.\\
Ist $\omega$ ein $q$-Differential, so ist $\operatorname{deg}(\operatorname{div}(\omega))=q \cdot(2 g-2)$, denn für $\omega=\eta^{q}$ mit einer 1-Form $\eta$ ist dies offensichtlich und im allgemeinen Fall ist $\omega / \eta^{q}$ eine holomorphe Funktion $X \longrightarrow \mathbb{P}^{1}$, also

$$
\operatorname{deg}(\omega)=\operatorname{deg}\left(\omega / \eta^{q}\right)+\operatorname{deg}\left(\eta^{q}\right)=0+q \cdot(2 g-2)
$$

Beweis des Satzes 9.5. Sei $q=g(g+1) / 2$. Ziel ist es, ein $q$-Differential zu konstruieren, dessen Divisor gerade $\sum_{x \in X} w(x)[x]$ ist. $\mathrm{Zu} x \in X$ seien $n_{i}(x)$ die Lücken bei $x$ und $\omega_{i, x} \in \Omega_{X}^{1}(X)$ seien 1-Formen der exakten Nullstellenordnung $n_{i}(x)-1$ in $x$. Deren Existenz ist nach Lemma 9.2 gesichert. In einer Karte $(U, z)$ um $x$ schreiben wir $\omega_{i, x}=f_{i}(z) d z$. Wir behaupten, dass die Wronski-Determinante

$$
W_{x}\left(f_{1}, \ldots, f_{g}\right)(z)=\operatorname{det}\left(f_{i}^{(j-1)}(z)\right)_{i, j=1}^{g}
$$

ein $q$-Differential ist.\\
Dazu sei $(V, t)$ eine weitere Karte und wir untersuchen den Kartenwechsel $z=\psi(t)$. Falls $\omega_{i, x}=h_{i}(t) d t$, so ist $h_{i}(t)=f_{i}(\psi(t)) \frac{d \psi}{d t}$. Für die höheren Ableitungen folgt also

$$
h_{i}^{(j)}(t)=\left(f_{i} \circ \psi\right)^{(j)}(t) \cdot \frac{d \psi}{d t}+\left\langle\text { Linearkombination von }\left(f_{i} \cdot \psi\right)^{(k)}, k<j,\right\rangle
$$

wobei die Koeffizienten der Linearkombination Funktionen in $t$, aber unabhängig von $i$ sind. In der Wronski-Determinante sind die Terme niedrigerer Ordnung also Linearkombination anderer Zeilen und es gilt

$$
W_{x}\left(h_{1}, \ldots, h_{g}\right)=\operatorname{det}\left(\left(f_{i} \circ \psi\right)^{(j-1)} \cdot \frac{d \psi}{d t}\right)_{i, j=1}^{g}
$$

Mit dem selben Gedanken können wir die höheren Ableitungen von $f_{i} \circ \psi$ berechnen. Es gilt\\
$\left(f_{i} \circ \psi\right)^{(j)}=\left(f_{i}^{(j)} \circ \psi\right) \cdot\left(\frac{d \psi}{d t}\right)^{j}+\left\langle\right.$ Linearkombination von $\left.\left(f_{i}^{(k)} \circ \psi\right) \cdot\left(\frac{d \psi}{d t}\right)^{k}, k<j\right\rangle$,\\
und daher insgesamt\\
$W_{x}\left(h_{1}, \ldots, h_{g}\right)(t)=\operatorname{det}\left(\left(f_{i}^{(j-1)} \circ \psi\right) \cdot\left(\frac{d \psi}{d t}\right)^{j}\right)_{i, j=1}^{g}=\left(\frac{d \psi}{d t}\right)^{q} W_{x}\left(f_{1}, \ldots, f_{g}\right)(\psi(t))$.\\
Also ist die Wronski-Determinante ein $q$-Differential. Es ist nicht Null, denn die $f_{i}$ haben verschiedene Nullstellenordnungen. Es bleibt der Divisor von $\omega_{x}$ zu bestimmen.

Wir zeigen zuerst, dass $\operatorname{ord}_{x} W_{x}=w(x)$ für jedes $x \in X$ gilt. Durch Multiplizieren der $\omega_{i, x}$ mit einer von Null verschiedenen Konstante erlaubt uns $f_{i}(z)=z^{n_{i}-1}+$ Terme höherer Ordnung zu schreiben. Also ist

$$
\begin{aligned}
W_{x}\left(f_{1}, \ldots, f_{g}\right)(z) & =\operatorname{det}\left(\left(\left(z^{n_{i}-1}\right)^{(j-1)}\right)_{i, j=1}^{g}\right)+\ldots \\
& =\sum_{\tau \in S_{g}} \operatorname{sign}(\tau) \prod_{i=1}^{g} \frac{\left(n_{i}-1\right)!}{\left(n_{i}-\tau(i)\right)!} z^{n_{i}-\tau(i)}+\ldots \\
& =c \cdot z^{\sum_{i=1}^{g}\left(n_{i}-i\right)}+\ldots
\end{aligned}
$$

Wir behaupten, dass die Konstante $c$ nicht Null ist, denn die führenden Terme $z^{n_{i}-1}$ der $f_{i}$ sind verschieden. In der Tat ist

$$
c=\operatorname{det}\left(\frac{\left(n_{i}-1\right)!}{\left(n_{i}-j\right)!}\right)_{i, j=1, \ldots, g}
$$

und durch Zeilenoperationen formt man diese Determinante zu einer vom VandermondŽschen Typ um.

Die Abbildung

$$
\left\{\begin{aligned}
\Omega_{X}^{1}(X)^{g} & \longrightarrow q \text {-Differentiale } \\
\left(\omega_{1}, \ldots, \omega_{g}\right) & \longmapsto W\left(\omega_{1}, \ldots, \omega_{g}\right)
\end{aligned}\right.
$$

ist alternierend, faktorisiert also über $\wedge^{g} \Omega_{X}^{1}(X)$ und nach Korollar 6.7 ist dieser Vektorraum isomorph zu $\mathbb{C}$. Deswegen sind für zwei beliebige Punkte $x, y \in X$ die Wronski-Determinanten $\mathbb{C}$-linear abhängig und es gilt ord ${ }_{y} W_{x}= \operatorname{ord}_{y} W_{y}=w(y)$, was noch zu zeigen war.

Anstatt über die Lücken kann man das Gewicht eines Punktes auch über die Dimensionen $\ell(n[x])$ berechnen. Aus Lemma 9.2 folgert man leicht:

Lemma 9.7 Ist $x \in X$ ein beliebiger Punkt, so gilt $w(x)=1-\frac{g(g+1)}{2}+ \sum_{n=1}^{2 g-2} \ell(n[x])$.

\subsection {Hyperelliptische Kurven}
Wir studieren hier eine wichtige Klasse von Beispielen für Riemannsche Flächen. Bisher hatten wir Riemannsche Flächen intrinsisch, d.h. als Verklebung\\
von Karten, oder als Überlagerung erhalten. Dieser Abschnitt gibt Gelegenheit, die extrinsische Sichtweise, d.h. Untermannigfaltigkeiten eines umgebenden Raumes, kurz einzuführen. Sie wird üblicherweise in der algebraischen Geometrie, insbesondere bei algebraischen Kurven verwendet. Daher auch der Begriff hyperelliptische Kurven, obwohl hyperelliptische Riemannsche Flächen auch üblich ist und vielleicht konsequenter wäre.

Definition 9.8 (hyperelliptisch) Eine kompakte Riemannsche Fläche $X$ heißt hyperelliptisch, falls es eine holomorphe Abbildung $f: X \longrightarrow \mathbb{P}^{1}$ vom Grad 2 gibt.

Ist $X$ hyperelliptisch, so sagt der Satz von Riemann-Hurwitz, dass die Anzahl der Verzweigungspunkte gleich $2 g+2$ ist.

Proposition 9.9 Jede Riemannsche Fläche $X$ vom Geschlecht $g(X) \leqslant 2$ ist hyperelliptisch.

Beweis: Ist $g(X)=0$, so ist $f(z)=z^{2}$ eine Überlagerung vom Grad 2. Ist $g(X)=1$, und $D$ ein beliebiger effektiver Divisor vom Grad 2, so ist nach Riemann-Roch $\ell(D)=2$. Also gibt es eine nicht-konstante Funktion in $\mathcal{L}(D)$. Diese hat Grad 2, definiert also die gewünschte Überlagerung.

Ist schließlich $g(X)=2$, so ist $\ell(K)=2$ für jeden kanonischen Divisor $K$. Wiederum gibt es eine nicht-konstante Funktion in $\mathcal{L}(K)$ und diese hat den gewünschten den Grad 2, da $K$ effektiv ist.

Wir wollen nun alle hyperelliptischen Kurven konstruieren. Seien dazu nun $z_{1}, \ldots, z_{2 g+2}$ paarweise verschiedene Punkte in $\mathbb{C}$. Wir betrachen

$$
X^{*}=\left\{(x, y) \in \mathbb{C}^{2}: y^{2}=\prod_{i=1}^{2 g+2}\left(x-z_{i}\right)\right\}
$$

und die Projektionsabbildungen

$$
p_{x}^{*}: X^{*} \rightarrow \mathbb{C},(x, y) \mapsto x \quad \text { sowie } \quad p_{y}^{*}: X^{*} \rightarrow \mathbb{C},(x, y) \mapsto y
$$

Satz 9.10 Die Menge $X^{*}$ kann mit der Struktur einer Riemannschen Fläche $X^{*} \subset \mathbb{C}^{2}$ versehen werden, sodass die Projektionsabbildungen holomorph sind.

Es gibt genau eine Kompaktifizierung von $X^{*}$ zu einer kompakten Riemannschen Fläche $X$, sodass sich $p_{x}^{*}$ zu einer Abbildung $p_{x}: X \longrightarrow \mathbb{P}_{\mathbb{C}}^{1}$ vom Grad zwei fortsetzt. Insbesondere ist $X$ also hyperelliptisch.\\
Umgekehrt kann man zu jeder hyperelliptischen Kurve $p^{\prime}: Y^{\prime} \longrightarrow \mathbb{P}^{1}$ den offenen Teil $p^{\prime-1}(\mathbb{C})$ durch eine Gleichung wie in (9.1) beschreiben, falls $p^{\prime}$ über $\infty$ unverzweigt ist.

Beweis: Wir betrachten $X^{*} \subset \mathbb{C}^{2}$ mit der induzierten Topologie. Ist $x_{0} \neq z_{i}$, so ist $p_{x}$ eine Kartenabbildung am Punkt $\left(x_{0}, y_{0}\right)$, denn

$$
x \longmapsto\left(x, \pm \sqrt{\prod_{i=1}^{2 g+2}\left(x-z_{i}\right)}\right)
$$

ist, für die richtige Wahl des Vorzeichens, eine stetige Umkehrabbildung. Insbesondere ist $\left.p_{x}^{*}\right|_{\left(p_{x}^{*}\right)^{-1}\left(\mathbb{C} \backslash\left\{z_{1}, \ldots, z_{2 g+2}\right\}\right)}$ eine unverzweigte Überlagerung. In den übrigen Punkten ( $z_{i}, 0$ ) behaupten wir, dass $p_{y}$ eine Karte ist. Wir bezeichnen die rechte Seite von 9.1 mit $g(x)$. Dann ist $\frac{d g}{d x}\left(z_{i}\right) \neq 0$, da die Punkte $z_{i}$ paarweise verschieden sind. Der Satz über implizite Funktionen liefert also die gesuchte lokale Umkehrabbildung. Verkettungen dieser Kartenabbildungen und der lokalen Umkehrabbildungen sind offenbar holomorph.\\
Zur Kompaktifizierung von $X^{*}$ schreiben wir die Gleichung (9.1) in Termen eines Parameters $\widetilde{x}=\frac{1}{x}$ um $\infty \in \mathbb{P}^{1}$ um. Wir nehmen dazu an, dass alle $z_{i} \neq 0$ sind, sonst führen wir noch vorher eine Koordinatenverschiebung $z \longmapsto z+c$ auf $\mathbb{C}$ durch. Dann ist $y^{2}=\prod_{i=1}^{2 g+2}\left(\frac{1}{\widetilde{x}}-z_{i}\right)$ äquivalent zu

$$
\left(y x^{-(g+1)}\left(\prod_{i=1}^{2 g+2} \frac{1}{z_{i}}\right)^{1 / 2}\right)^{2}=\prod_{i=1}^{2 g+2}\left(\widetilde{x}-\frac{1}{z_{i}}\right)
$$

Setzen wir also $\widetilde{y}=y x^{-(g+1)}\left(\prod_{i=1}^{2 g+2} \frac{1}{z_{i}}\right)^{1 / 2}$ für eine beliebige Wahl der Quadratwurzel, so definiert (9.2) eine Riemannsche Fläche $X_{2}^{*}$ und die Abbildung $(\widetilde{x}, \widetilde{y}) \longmapsto \widetilde{x}$ eine zweiblättrige Überlagerung $p_{x, 2}^{*}: X_{2}^{*} \longrightarrow \mathbb{C}$. Die Verklebung von $X^{*}$ und $X_{2}^{*}$ auf der offenen Menge $\left\{(x, y) \in X^{*}: x \neq 0\right\}$ via $(x, y) \longmapsto(\widetilde{x}, \widetilde{y})$ definiert die gewünschte Kompaktifizierung $X$.\\
Umgekehrt, gegeben $p^{\prime}: Y^{\prime} \longrightarrow \mathbb{P}^{1}$ vom Grad $s$ mit $z_{1}, \ldots, z_{2 g+2}$ die Bilder der Verzweigungspunkte, so definiert obige Kompaktifizierung von $y^{2}= \prod_{i=1}^{2 g+2}\left(x-z_{i}\right)$ eine zweiblättrige Überlagerung $p_{x}: X \longrightarrow \mathbb{P}^{1}$. Um zu zeigen, dass $X \cong Y^{\prime}$ behaupten wir, dass es genügt einen Homöomorphismus $\psi: X \longrightarrow Y^{\prime}$ anzugeben, sodass $p^{\prime} \circ \psi=p_{x}$ ist. In der Tat ist so ein $\psi$ auf $\left(p_{x}\right)^{-1}\left(\mathbb{C} \backslash\left\{z_{1}, \ldots, z_{2 g+2}\right\}\right)$ holomorph, denn auf dieser Menge ist $p_{x}$ bzw. $p^{\prime}$ eine Kartenabbildung. Nimmt man nun einen der endlich vielen verbleibenden Punkte $P=p_{x}^{-1}\left(z_{i}\right)$ bzw. $P^{\prime}=\left(p^{\prime}\right)^{-1}\left(z_{i}\right)$ her, so ist bzgl. Karten ( $U, g$ ) um $P$ und $(V, h)$ um $P^{\prime}$ die Abbildung $h \circ \psi \circ g^{-1}$ stetig und holomorph ausserhalb von $g(P)$, also überall holomorph.

Es genügt sogar $\left.\psi\right|_{p_{x}^{-1}\left(\mathbb{C} \backslash\left\{z_{1}, \ldots, z_{2 g+2}\right\}\right)}$ anzugeben, da die Fortsetzung auf die diskrete Menge von Verzweigungspunkten von $p_{x}$ automatisch (eindeutig) möglich ist.\\
Verzweigte Überlagerungen vom Grad 2 von $\mathbb{P}^{1} \backslash\left\{z_{1}, \ldots, z_{2 g+2}\right\}$ bis auf Homöomorphie entsprechen bijektiv den Untergruppen vom Index 2 von

$$
\pi_{1}\left(\mathbb{P}^{1} \backslash\left\{z_{1}, \ldots, z_{2 g+2}\right\}\right)=\left\langle\gamma_{1}, \ldots, \gamma_{2 g+2}\right\rangle=: F_{2 g+1}
$$

Da Untergruppen von Index 2 stets normal sind, kann man diese eindeutig Homomorphismen

$$
F_{2 g+1} \longrightarrow \mathbb{Z} / 2 \mathbb{Z}
$$

zuordnen. Der Homomorphismus zu $p_{x}$ und der zu $p^{\prime}$ muss jede Schleife $\gamma_{i}$ auf das nichttriviale Element in $\mathbb{Z} / 2 \mathbb{Z}$ schicken, da $p_{x}$ (und $p^{\prime}$ ) über den Punkten $z_{i}$\\
wirklich verzweigt ist. Das legt aber den Homomorphismus und folglich die Untergruppe von $F_{2 g+1}$ eindeutig fest.

Das Argument im obigen Beweis besagt auch, dass hyperelliptische Kurven stets einen Automorphismus besitzen, der die Deckgruppe der Grad-2Überlagerung auf $\mathbb{P}_{\mathbb{C}}^{1}$ erzeugt. Dieser wird hyperelliptische Involution genannt, sie Abbildung 9.1

\begin{figure}[h]
\begin{center}
  \includegraphics[scale=0.2]{2025_10_06_374d4cfa8b1948eb4c41g-117}
 
\caption{Abbildung 9.1: Hyperelliptische Involution}
\end{center}
\end{figure}

Mit der konkreten Gleichung können wir auch holomorphe 1-Formen auf hyperelliptischen Kurven konkret angeben.

Proposition 9.11 Auf der hyperelliptischen Kurve $X$ mit Gleichung (9.1) ist eine Basis von holomorphen Differentialen gegeben durch

$$
\frac{d x}{y}, x \frac{d x}{y}, \ldots, x^{g-1} \frac{d x}{y}
$$

Beweis: Seien $\infty^{+}$und $\infty^{-}$die zwei Urbilder von $\infty \in \mathbb{P}^{1}$ und $f(x)$ die rechte Seite der Gleichung in (9.1). Es ist

$$
\operatorname{div}(y)=\left[\left(z_{1}, 0\right)\right]+\ldots+\left[\left(z_{2 g+2}, 0\right)\right]-(g+1)\left(\left[\infty^{+}\right]+\left[\infty^{-}\right]\right)
$$

wobei die Beschreibung bei $\infty$ ohne Kartenwechselrechnungen aus Grad- und Symmetriegründen (Invarianz unter $y \longmapsto-y$ ) folgt. Ebenso ist

$$
\begin{gathered}
\operatorname{div}(x)=[(0, \sqrt{f(0)})]+[(0,-\sqrt{f(0)})]-\left[\infty^{+}\right]-\left[\infty^{-}\right] \\
\text {bzw. } \operatorname{div}(x)=2 \cdot([0,0])-\left[\infty^{+}\right]-\left[\infty^{-}\right]
\end{gathered}
$$

je nachdem ob Null unverzweigt oder verzweigt ist.\\
Außerhalb der Verzweigungspunkte ist $x$ eine Karte, also $\operatorname{ord}(d x)=0$. Bei ( $z_{i}, 0$ ) ist $y$ eine Karte und $d x=d\left(y^{2}\right) h(x)$, wobei $h$ in der Nähe von ( $0, z_{i}$ ) holomorph ist. Aus Grad- und Symmetriegründen ist also

$$
\operatorname{div}(d x)=\left[\left(z_{1}, 0\right)\right]+\ldots\left[\left(z_{2 g+2}, 0\right)\right]-2 \cdot\left(\left[\infty^{+}\right]+\left[\infty^{-}\right]\right)
$$

Daraus folgt, dass die angegebenen 1-Formen holomorph und offensichtlich linear unabhängig sind.

Korollar 9.12 Ist $X$ hyperelliptisch, so sind die Weierstraßpunkte genau die Verzweigungspunkte. Die Lückenreihe jedes solchen Punktes ist $\{1,3,5, \ldots, 2 g-1\}$ und folglich ist das Gewicht jedes solchen Punktes gleich $g(g-1) / 2$.

Beweis: Wir nehmen an, dass $X$ durch die Gleichung (9.1) gegeben ist. Die Differentiale $\left(x-z_{i}\right)^{j} \frac{d x}{y}$ haben für $j=0,1, \ldots, g-1$ bei ( $z_{i}, 0$ ) Nullstellen der Ordnung $0,2,4, \ldots, 2 g-2$. Daraus folgt die Behauptung über Lückenreihen. Das Gesamtgewicht dieser ( $2 g+2$ ) Weierstraßpunkte ist $(2 g+2) \frac{g(g-1)}{2}=g\left(g^{2}-\right.$ 1). Also haben wir alle Weierstraßpunkte gefunden.

Daraus folgt auch, dass hyperelliptische Kurven die Dimensionen $\ell(n[x])$ maximieren.

Korollar 9.13 Für einen beliebigen Punkt $x$ auf $X$ und $n \leq 2 g$ gilt $\ell(n[x]) \leq 1+ \left\lfloor\frac{n}{2}\right\rfloor$ und die Gleichheit gilt genau für Weierstraßpunkte auf hyperelliptischen Kurven.

Beweis: Sei $n_{0}$ eine Zahl für die die Aussage falsch ist. Es genügt zu zeigen, dass alle Lücken kleiner gleich $n_{0}$ sind, denn dann gilt die für große $n$ wahre Konsequenz von Riemann-Roch $\ell(n[x])=n+1-g$ auch für $n_{0}$ und sie impliziert $n_{0}+1=g+\ell\left(n_{0}[x]\right)>g+1+\left\lfloor\frac{n_{0}}{2}\right\rfloor$ und damit $n_{0}>2 g$.

Um dies zu zeigen, nehmen wir wiederum im Widerspruchsbeweis an, dass $m=n_{0}+1$ eine Lücke sei. Dann ist für jedes $i \leq\left\lfloor\frac{m}{2}\right\rfloor$ eine der beiden Zahlen $i$ oder $m-i$ auch eine Lücke, denn sonst wäre $m=i+(m-i) \in \Gamma_{x}$ nach der Halbgruppeneigenschaft. Also gibt es mindestens $\left\lfloor\frac{m}{2}\right\rfloor$ Lücken $\leq m-1$. Daraus folgt, dass $\ell\left(n_{0}[x]\right) \leq 1+\left\lfloor\frac{n_{0}}{2}\right\rfloor$, im Widerspruch zur ersten Annahme. Also ist $m$ keine Lücke. Es folgt, dass auch $\ell(m[x])>1+\left\lfloor\frac{m}{2}\right\rfloor$ gilt. Wir können also denselben Widerspruchsbeweis für die größere Zahl $m$ durchführen, um so induktiv die Behauptung zu zeigen, dass alle Lücken kleiner gleich $n_{0}$ sind.

Ist $X$ hyperelliptisch und $x$ ein Weierstraßpunkt, so impliziert die Lückenreihe in Korollar 9.12, dass $\ell(n[x])=1+\left\lfloor\frac{n}{2}\right\rfloor$ für $n \leq 2 g$ ist. Ist andererseits bereits $\ell(2[x])=2$, so ist $X$ hyperelliptisch und $x$ ein Weierstraßpunkt, da es somit auf $X$ eine nichtkonstante holomorphe Abbildung nach $\mathbb{P}^{1}$ gibt, die genau bei $x$ eine doppelte Polstelle hat, also von Grad 2 mit $x$ als Verzweigungspunkt ist.

\subsection {Die Endlichkeit der Automorphismengruppe}
Wir haben alle Begriffe zusammen, die wir im Beweis des folgenden Satzes benötigen.

Satz 9.14 Ist $X$ eine kompakte Riemannsche Fläche mit $g(X) \geq 2$, so ist die Automorphismengruppe endlich.

Dazu einige Vorüberlegungen, die von unabhängigem Nutzen sind.

Proposition 9.15 Die hyperelliptische Involution $(y, z) \longmapsto(-y, z)$ ist der einzige Automorphismus $\varphi$ (abgesehen von der Identität) einer hyperelliptischen Kurve $X$ vom Geschlecht $g(X) \geq 2$, der mindestens 5 Punkte fixiert.

Beweis: Sei $f: X \longrightarrow \mathbb{P}^{1}$ eine holomorphe Abbildung vom Grad 2. Jede weitere Abbildung $\widetilde{f}: X \longrightarrow \mathbb{P}^{1}$ vom Grad 2 ist die Verkettung von $f$ mit einer Möbiustransformation $A: \mathbb{P}^{1} \longrightarrow \mathbb{P}^{1}$, denn ist $x \in X$ ein Weierstraßpunkt, so sind die Funktionen $1 /(f-f(x)), 1 /(\widetilde{f}-\widetilde{f}(x))$ und die Konstanten allesamt in $\mathcal{L}(2[x])$ und $\ell(2[x])=2$. Aus der Existenz einer nichttrivialen Linearkombination folgt die Behauptung.

Sei also $A$ diese Möbiustransformation mit der Eigenschaft $A \circ f=f \circ \varphi$. Notwendigerweise fixiert $A$ die Bildpunkte der Fixpunkte von $\varphi$. Davon gibt es mindestens $5 / \operatorname{deg}(f)>2$ und folglich ist $A=i d$. Daraus folgt, dass $\varphi$ die Identität ist oder die Blätter von $f$ vertauscht. In diesem Fall ist $\varphi$ die hyperelleptische Involution.

Proposition 9.16 Sei $X$ eine Riemannsche Fläche mit $g(X) \geqslant 2$. Dann hat ein nichttrivialer Automorphismus von $X$ höchstens $2 g+2$ Fixpunkte.

Beweis: Sei $x \in X$ ein Punkt, der von dem Automorphismus $\varphi$ nicht festgehalten wird. Riemann-Roch impliziert, dass

$$
\ell((g+1)[x]) \geqslant 2 .
$$

Wir nehmen also $f \in \mathcal{L}((g+1) \cdot[x])$ eine nicht-konstante Funktion und setzen $h=f-f \circ \varphi$. Da $h^{-1}(\infty)=\left\{x, \varphi^{-1}(x)\right\}$ und bei jedem dieser Punkte die Abbildung höchstens die Ordnung $g+1$ hat, ist $\operatorname{deg}(h) \leq 2 g+2$. Also hat Null höchstens $2 g+2$ Urbilder unter $h$ und jeder Fixpunkt von $\varphi$ ist eine Nullstelle von $h$.

Hieraus können wir die angekündigte Endlichkeit der Automorphismengruppe folgern.\\
Beweis von Satz 9.14. Ein Automorphismus bewahrt Polordnungen von Funktionen (und auch Nullstellenordnungen von 1-Formen). Also permutiert er die Menge $W$ der Weierstraßpunkte. Da $W$ endlich ist, genügt es zu zeigen, dass die Untergruppe $\operatorname{Aut}(X)_{W} \leq \operatorname{Aut}(X)$, die trivial auf den Weierstraßpunkten operiert, endlich ist.\\
Ist $X$ hyperelliptisch, so fixiert so ein Automorphismus mindestens $2 g+2>5$ Punkte, ist also die Identität oder die hyperelliptische Involution.\\
Ist $X$ nicht hyperelliptisch, so ist $w(x)<\frac{g(g-1)}{2}$ für jeden Weierstraßpunkt $x$ nach untenstehendem Lemma 9.17. Also gibt es mindestens $2 g+3$ Weierstraßpunkte und $\operatorname{Aut}(X)_{W}=\{\mathrm{id}\}$ in diesem Fall.

Lemma 9.17 Ist $X$ nicht hyperelliptisch, so ist $w(x)<\frac{g(g-1)}{2}$ für jeden Punkt $x \in X$.

Beweis: Aus Lemma 9.7 und Korollar 9.13 folgt, dass $w(x) \leq \frac{g(g-1)}{2}$ und dass diese Ungleichung strikt ist, sobald eine der Ungleichungen in Korollar 9.13 strikt ist. Dies ist auf einer nicht-hyperelliptischen Kurve wegen $\ell(2[x])=1$ für alle $x$ schon für $n=2$ der Fall.

\section {Perioden und die Jakobische}
Sei $X$ eine kompakte Riemannsche Fläche vom Geschlecht $g \geqslant 1$. Im Abschnitt über flache Flächen haben wir bereits mehrfach mit den Abständen $d(x, y)$ argumentiert, insbesondere falls $x$ und $y$ Nullstellen von $\omega$ sind. Nach Definition ist

$$
d(x, y)=\inf \left|\int_{x}^{y} \omega\right|
$$

wobei das Infimum über alle Wege von $x$ nach $y$ genommen wird. Allgemein spielen die Integrale von $\omega$ über geschlossene Wege oder Wege zwischen Nullstellen eine wichtige Rolle bei der Parametrisierung des Raums aller flachen Flächen.

In diesem Abschnitt wollen wir zunächst keine Differentialform auf $X$ auszeichnen und allgemein die Rolle von Perioden, d.h. von Integralen holomorpher 1-Formen entlang geschlossener Wege, untersuchen.\\
Sei $x_{0} \in X$ ein Referenzpunkt und $\omega_{1}, \ldots, \omega_{g}$ eine Basis von $\Omega^{1}(X)$. Die Abbildung

$$
u:\left\{\begin{aligned}
X & \longrightarrow \mathbb{C}^{g} \\
x & \longrightarrow\left(\int_{x_{0}}^{x} \omega_{1}, \ldots, \int_{x_{0}}^{x} \omega_{g}\right)^{T}
\end{aligned}\right.
$$

ist nicht wohldefiniert, da die Integrale von der Wahl des Weges zwischen $x_{0}$ und $x$ abhängen und sich damit um Integrale entlang geschlossener Kurven bei $x_{0}$ unterscheiden. Wir fixieren eine kanonische Zerschneidung von $X$ in ein Polygon $P$ mit Seiten $a_{i}, b_{i}, a_{i}^{-1}, b_{i}^{-1}$, jeweils $i=1 \ldots g$. Wir definieren die Gruppe der Perioden von $X$ als

$$
\Lambda=\left\{\left(\int_{\gamma} \omega_{1}, \ldots, \int_{\gamma} \omega_{g}\right)^{T}, \gamma \in \pi_{1}\left(X, x_{0}\right)\right\} \subset \mathbb{C}^{g}
$$

Diese hängt offenbar nicht vom Basispunkt ab und $\Lambda$ ist nach Satz 3.32erzeugt von

$$
\left(\int_{a_{i}} \omega_{1}, \ldots, \int_{a_{i}} \omega_{g}\right)^{T} \text { und }\left(\int_{b_{i}} \omega_{1}, \ldots, \int_{b_{i}} \omega_{g}\right)^{T}, \quad i=1 \ldots g
$$

Ist $\widehat{\omega}_{1}, \ldots, \widehat{\omega}_{g}$ eine weitere Basis von $\Omega^{1}(X)$ und $M \in \mathrm{GL}_{g}(\mathbb{C})$ die Matrix des Basiswechsels von $\omega_{1}, \ldots, \omega_{g}$ zu $\widehat{\omega}_{1}, \ldots, \widehat{\omega}_{g}$, so gilt

$$
\widehat{\Lambda}:=\left\{\left(\int_{\gamma} \widehat{\omega}_{1}, \ldots, \int_{\gamma} \widehat{\omega}_{g}\right)^{T}, \gamma \in \pi_{1}\left(X, x_{0}\right)\right\}=M \cdot \Lambda
$$

Die Quotientengruppen $\mathbb{C}^{g} / \Lambda$ und $\mathbb{C}^{g} / \widehat{\Lambda}$ sind also isomorph, der Isomorphismus wird durch Linksmultiplikation mit $M$ gegeben. Die Isomorphieklasse des Quotienten $\mathbb{C}^{g} / \Lambda$, versehen mit einer Zusatzstruktur "komplexer Torus", werden wir später als die Jacobische $\operatorname{Jac}(X)$, von $X$ bezeichnen. Zunächst vergleichen wir $\operatorname{Jac}(X)$ mit einer weiteren Gruppe, die in natürlicher Weise $X$ zugeordnet ist. Sei $\operatorname{Div}_{0}(X)$ die Untergruppe von $\operatorname{Div}(X)$ von Divisoren vom Grad null und $\operatorname{HD}(X) \subset \operatorname{Div}_{0}(X)$ die Gruppe der Hauptdivisoren, d.h. der Divisoren von meromorphen Funktionen auf $X$. Die Quotientengruppe

$$
\operatorname{Pic}_{0}(X):=\operatorname{Div}_{0}(X) / \operatorname{HD}(X)
$$

heißt Picardgruppe von $X$. Die oben definierte Abbildung $u$ setzt sich zu einer Abbildung

$$
u:\left\{\begin{aligned}
\operatorname{Div}_{0}(X) & \longrightarrow \operatorname{Jac}(X) \\
\sum_{x \in X} n_{x} \cdot[x] & \longmapsto \sum_{x \in X} n_{x} \cdot u(x)
\end{aligned}\right.
$$

fort. Ziel dieses Abschnitts ist der Beweis des folgenden Satzes.\\
Satz 10.1 (Abel-Jacobi) Die Abbildung u ist unabhängig von der Wahl eines Basispunktes und definiert einen Isomorphismus

$$
u: \operatorname{Pic}_{0}(X) \Longrightarrow \operatorname{Jac}(X)
$$

Dem Beweis geht folgender Hilfssatz voraus, der Residuen von meromorphen Differentialen und Perioden in Verbindung bringt.

Lemma 10.2 Sei $\psi$ eine meromorphe Differentialform, $P$ ein kanonisches Polygon, $x_{0} \in P$ ein Basispunkt, $\omega \in \Omega^{1}(X)$ und $u_{\omega}(x)=\int_{x_{0}}^{x} \omega$, wobei der Weg in $P$ gewählt ist. Dann gilt

$$
\begin{gathered}
2 \pi i \cdot \sum_{x \in X} \operatorname{res}_{x}\left(u_{\omega} \cdot \psi\right)=-\sum_{i=1}^{g} \tau_{i} \beta_{i}+\sum_{i=1}^{g} \sigma_{i} \alpha_{i} \\
\text { wobei } \quad \alpha_{i}=\int_{a_{i}} \omega, \quad \beta_{i}=\int_{b_{i}} \omega, \quad \tau_{i}=\int_{a_{i}} \psi \quad \text { und } \quad \sigma_{i}=\int_{b_{i}} \psi
\end{gathered}
$$

Beweis: Sei $P^{\prime}$ ein Polygon, das in $P$ einbeschrieben ist und alle Pole von $\psi$ enthält. Die Seiten von $P^{\prime}$ seien mit $a_{i}^{\prime}, b_{i}^{\prime}, c_{i}^{\prime}$ und $d_{i}^{\prime}$ wie in Abbildung 10.1 bezeichnet.

Sei $U_{i}$ eine einfach zusammenhängende Umgebung von $a_{i} \backslash\left\{x_{\infty}\right\}$, wobei $x_{\infty} \in X$ der Eckpunkt des Polygons ist. Wir fixieren zudem Punkte $x_{1}, x_{2} \in U_{i}$ auf gegenüberliegenden Seiten von $a_{i}$. Für einen Punkt $x \in U_{i}$ sei $u_{k, i}(x)= \int_{x_{0}}^{x_{k}} \omega+\int_{x_{k}}^{x} \omega$ für $k=1,2$, wobei die Integrale über $\omega$ in $P$ bzw. in $U_{i}$ definiert und daher wohldefiniert sind. Dann gilt

$$
u_{1, i}-u_{2, i}=\int_{b_{i}} \omega=\beta_{i}
$$

\begin{figure}[h]
\begin{center}
  \includegraphics[scale=0.2]{2025_10_06_374d4cfa8b1948eb4c41g-122}
 
\caption{Abbildung 10.1: Das Polygon \$\textbackslash mathrm\{P}\^{}\{\textbackslash prime\}\$ in der kanonischen Zerschneidung P\}\end{center}
\end{figure}

Definiert man analog Umgebungen $V_{i}$ um $b_{i} \backslash\left\{x_{\infty}\right\}$ und fixiert $y_{1}, y_{2} \in V_{i}$ zu beiden Seiten von $b_{i}$ und setzt $v_{k, i}=\int_{x_{0}}^{y_{k}} \omega+\int_{y_{k}}^{x} \omega$ für $k=1,2$ so ist

$$
v_{1, i}-v_{2, i}=\int_{a_{i}^{-1}} \omega=-\alpha_{i}
$$

Folglich ist im Limes, wenn $P^{\prime}$ sich an $P$ annähert

$$
\begin{aligned}
2 \pi i \sum \operatorname{res}\left(u_{\omega} \psi\right) & =\lim \sum_{i=1}^{g}\left(\int_{a_{i}^{\prime}} u_{\omega} \psi+\int_{b_{i}^{\prime}} u_{\omega} \psi+\int_{c_{i}^{\prime}} u_{\omega} \psi+\int_{d_{i}^{\prime}} u_{\omega} \psi\right) \\
& =\sum_{i=1}^{g}\left(\int_{a_{i}}\left(u_{1, i}-u_{2, i}\right) \psi+\int_{b_{i}}\left(v_{1, i}-v_{2, i}\right) \psi\right) \\
& =\sum_{i=1}^{g}\left(\beta_{i} \cdot \tau_{i}-\sigma_{i} \cdot \alpha_{i}\right)
\end{aligned}
$$

Der Satz von Riemann-Roch 7.30 kann zum Existenzbeweis von Funktionen oder Differentialformen mit speziellen Null- oder Polstellenvorgaben verwendet werden.

Proposition 10.3 Ist $\psi$ eine meromorphe 1-Form auf $X$, so ist $\sum_{x \in X} \operatorname{res}_{x}(\psi)=0$. Umgekehrt, seien $x_{1}, \ldots, x_{m} \in X$ Punkte und $\Gamma_{1}, \ldots, \Gamma_{m} \in \mathbb{C}$ mit $\sum_{i=1}^{m} \Gamma_{i}=0$ vorgegeben. Dann gibt es eine meromorphe 1-Form $\psi$ auf $X$, holomorph außerhalb der $x_{i}$, höchstens einfache Pole bei den $x_{i}$ und mit $\operatorname{res}_{x_{i}}(\psi)=\Gamma_{i}$.

Beweis: Für die erste Aussage nehmen wir wieder die Polygone $P^{\prime}$ und $P$ wie im vorigen Beweis. Nach dem Residuensatz für Gebiete in $\mathbb{C}$ ist $\sum_{x \in X} \operatorname{res}_{x}(\psi)=\frac{1}{2 \pi i} \int_{\partial P^{\prime}} \psi$, wenn wir $P^{\prime}$ dicht genug bei $P$ wählen. Zum anderen ist $\partial P^{\prime}$ aber eine Schleife um den Eckpunkt $x_{\infty}$, in deren Inneren keine Pole liegen. Also ist $\int_{\partial P^{\prime}} \psi=0$.

Die Umkehrung ist eine Anwendung des Satzes von Riemann-Roch. Sei $D_{0}= \sum\left[x_{i}\right]$ und $\Omega_{D_{0}}:=H^{0}\left(X, \Omega_{X}^{1}\left(D_{0}\right)\right.$ der Raum der 1-Formen mit höchstens einfachen Polen bei den $x_{i}$ und holomorph außerhalb. Ist $\omega \in \Omega^{1}(X)$ und $K=\operatorname{div}(\omega)$, so liefert $f \longmapsto f \cdot \omega$ ein Isomorphismus auf $\mathcal{L}\left(K+D_{0}\right) \Longrightarrow \Omega_{D_{0}}$. Ist $m=\operatorname{deg}\left(D_{0}\right)$, so ist

$$
\operatorname{Res}:\left\{\begin{aligned}
\Omega_{D_{0}} & \longrightarrow \mathbb{C}^{m} \\
\psi & \longmapsto\left(\operatorname{res}_{x_{1}}(\psi), \ldots, \operatorname{res}_{x_{m}}(\psi)\right)
\end{aligned}\right.
$$

eine lineare Abbildung, deren Kern genau $\Omega^{1}(X)$ ist, und deren Bild in der Hyperebene $H$ mit Koordinatensumme Null liegt. Es ist

$$
\operatorname{dim} \Omega_{D_{0}}=\ell\left(K-\left(-D_{0}\right)\right)=m-1+g
$$

Da der Kern von Res $g$-dimensional ist, muss Res surjektiv auf $H$ abbilden.

\section {Beweis des Satzes von Abel-Jacobi:}
Schritt 1: Hauptdivisoren liegen im Kern von $u$. Zum Beweis dieser Behauptung sei $D=\operatorname{div}(f)=\sum n_{x} \cdot[x]$ ein Hauptdivisor und $\psi=\frac{d f}{f}$. Nach dem Lemma 10.2 ist für $i \in\{1 \ldots g\}$ die $i$-te Komponente

$$
\begin{aligned}
u(D)_{i} & =\sum_{x} n_{x} \cdot u_{\omega_{i}}(x)=\sum_{x} \cdot \operatorname{Res}_{x}\left(u_{\omega_{i}} \cdot \psi\right) \\
& =\frac{1}{2 \pi i} \cdot \sum_{j}\left(-\left(\int_{b_{j}} \omega_{i}\right) \cdot\left(\int_{a_{j}} \frac{d f}{f}\right)+\left(\int_{a_{j}} \omega_{i}\right) \cdot\left(\int_{b_{j}} \frac{d f}{f}\right)\right)
\end{aligned}
$$

Die Integrale $\int_{a_{i}} \frac{d f}{f}, \int_{b_{i}} \frac{d f}{f}$ liegen in $2 \pi i \cdot \mathbb{Z}$, denn der Logarithmus ist genau $2 \pi i \cdot \mathbb{Z}$-vieldeutig. Daraus folgt die Behauptung.\\
Schritt 2: Die Spalten von $\Lambda$ in obiger Basis $\left\{a_{1}, \ldots, a_{g}, b_{1}, \ldots, b_{g}\right\}$ von $\pi_{1}\left(X, x_{0}\right)^{a b}$ erzeugen $\mathbb{C}^{g}$ als $\mathbb{C}$-Vektorraum. Wir bezeichnen zum Beweis dieser Behauptung diese Spalten mit $S_{i}$ für $i=1, \ldots, 2 g$. Falls die Behauptung falsch ist, gibt es eine nicht-triviale Linearform $\mu=\left(\mu_{1}, \ldots, \mu_{g}\right) \in \mathbb{C}^{g}$, sodass das Erzeugnis der $S_{i}$ im Kern von $\mu$ liegt. In diesem Fall ist $\omega=\sum_{i=1}^{g} \mu_{i} \omega_{i}$ eine holomorphe 1-Form, deren Perioden alle verschwinden. Also ist $\int_{x_{0}}^{x} \omega$ eine wohldefinierte holomorphe Funktion auf $X$, notwendigerweise konstant, da $X$ kompakt ist. Dann aber ist $\omega=0$, was der Definition von $\omega_{1}, \ldots, \omega_{g}$ als Basis von $\Omega^{1}(X)$ widerspricht. In Vorbereitung auf den nächsten Schritt nummerieren wir die Spalten um und setzen $T_{i}=-S_{i+g}$ für $i \in\{1, \ldots, g\}$ und $T_{i}=S_{i-g}$ für $i \in\{g+1, \ldots, 2 g\}$.\\
Schritt 3: Der Vektor $\lambda=\left(\lambda_{1}, \ldots, \lambda_{2 g}\right) \in \mathbb{C}^{2 g}$ erfüllt $\sum_{i=1}^{2 g} \lambda_{i} S_{i}=0$ genau dann, wenn es ein $\mu=\left(\mu_{1}, \ldots, \mu_{g}\right) \in \mathbb{C}^{g}$ gibt mit der Eigenschaft $\lambda_{i}=\sum_{j=1}^{g} \mu_{j} \cdot\left(T_{i}\right)_{j}$ für alle $i=1, \ldots, 2 g$. Zum Beweis hiervon nehmen wir an, dass es so ein $\mu$ gibt. Dann ist $\psi=\sum_{j=1}^{g} \mu_{j} \omega_{j}$ holomorph. Nach dem Lemma 10.2 leistet also $\lambda_{i}=-\int_{b_{i}} \psi$ für $i=1, \ldots, g$ und $\lambda_{i}=\int_{a_{i}} \psi$ für $i=g+1, \ldots, 2 g$ das Verlangte. Umgekehrt ist nach Schritt 2 die Abbildung

$$
\begin{aligned}
& \mathbb{C}^{2 g} \longrightarrow \mathbb{C}^{g} \\
& \lambda \\
& \lambda \sum_{i=1}^{2 g} \lambda_{i} S_{i}
\end{aligned}
$$

surjektiv. Nach der ersten Implikation enthält der Kern das Bild von

$$
K: \begin{aligned}
& \mathbb{C}^{g} \quad \longrightarrow \mathbb{C}^{2 g} \\
& \mu \quad \longmapsto\left(\sum_{j=1}^{g} \mu_{j} \cdot\left(T_{i}\right)_{j}\right)_{i} .
\end{aligned}
$$

Genau wie in Schritt 2 zeigt man, das $K$ injektiv ist.\\
Schritt 4: Auf Pic ${ }_{0}$ ist $u$ injektiv. Sei zum Beweis hiervon also $D=\sum n_{x}[x]$ ein Divisor vom Grad Null mit $u(D)=0$ in $\operatorname{Jac}(X)$. Wir wollen zeigen, dass es $f: X \longrightarrow \mathbb{P}^{1}$ gibt mit $D=\operatorname{div}(f)$. Sei $\psi$ ein meromorphes Differential mit einfachen Polen und Residuen $n_{x}$ bei $x$ und holomorph sonst. Sobald wir gezeigt haben, dass wir $\int_{a_{i}} \psi \in 2 \pi i \mathbb{Z}$ und $\int_{b_{i}} \psi \in 2 \pi i \mathbb{Z}$ erreichen können, leistet $f(z)=\exp \left(\int_{x_{0}}^{z} \psi\right)$ das Verlangte: Die Wohldefiniertheit von $f$ ist klar und in der Nähe eines Pols $x$ von $\psi$ ist $\psi(z)=n_{x} \cdot \frac{d z}{z}$, also

$$
\begin{aligned}
f(z) & =\exp \left(n_{x} \cdot\left(\log (z)-\log \left(x_{0}\right)\right)\right) \\
& =z^{n_{x}} \cdot \exp \left(-n_{x} \log \left(x_{0}\right)\right)
\end{aligned}
$$

Damit läßt sich $f$ durch $f(x)=0$ falls $n_{x}>0$ und $f(x)=\infty$ falls $n_{x}<0$ stetig und damit holomorph fortsetzen und wir erhalten $\operatorname{div}(f)=D$. Zum Nachweis der Integralbedingungen bemerken wir, dass nach Voraussetzung und Lemma 10.2

$$
\begin{aligned}
\Lambda \ni u(D) & =\sum_{x} n_{x} \cdot u(x)=\sum_{x}\left(\operatorname{res}_{x}\left(u_{\omega_{1}} \psi\right), \ldots, \operatorname{res}_{x}\left(u_{\omega_{g}} \psi\right)\right) \\
& =\frac{1}{2 \pi i} \sum_{k=1}^{g}\left(\sigma_{k} S_{k}-\tau_{k} S_{k+g}\right)
\end{aligned}
$$

gilt, wobei wie oben $\tau_{k}=\int_{a_{k}} \psi$ und $\sigma_{k}=\int_{b_{k}} \psi$. Also ist

$$
\frac{1}{2 \pi i} \sum_{k=1}^{g}\left(\sigma_{k} S_{k}-\tau_{k} S_{k+g}\right)=\sum_{k=1}^{2 g} m_{k} S_{k}
$$

mit $m_{k} \in \mathbb{Z}$. Nach Schritt 3 gibt es also $\mu=\left(\mu_{1}, \ldots, \mu_{g}\right)$ mit $\widetilde{\tau}_{k}-2 \pi i m_{k}= \sum_{j=1}^{g} \mu_{j} \cdot\left(T_{k}\right)_{j}$ für alle $k \in\{1, \ldots, 2 g\}$, wobei wir jetzt $\widetilde{\tau}_{k}=\sigma_{k}$ und $\widetilde{\tau}_{k+g}=-\tau_{k}$ für $k \in\{1, \ldots, g\}$ gesetzt haben. Also hat $\psi_{2}=\psi+\sum_{j} \mu_{j} \omega_{j}$ die gleichen Pole und Residuen wie $\psi$ und nun wirklich Perioden in $2 \pi i \cdot \mathbb{Z}$.\\
Schritt 5: Wir zeigen die Surjektivität von $u$. Wir nehmen dazu ein Tupel $x_{1}, \ldots, x_{g}$ von Basispunkten in $X$ her, das wir später noch genauer festlegen. Da der Definitionsbereich von $u$ Divisoren vom Grad Null sind, stimmt für $y_{i}$ in einer Umgebung $V_{i}$ der $x_{i}$ die oben definierte Abbildung $u$ mit

$$
\left.u\right|_{\left(V_{i}\right)_{i}}:\left\{\begin{array}{lll}
\prod_{i=1}^{g} V_{i} & \longrightarrow & \mathbb{C}^{g} \\
\left(y_{1}, \ldots, y_{g}\right) & \longmapsto & \left(\sum_{i=1}^{g} \int_{x_{i}}^{y_{i}} \omega_{k}\right)_{k=1, \ldots g}
\end{array}\right.
$$

überein. Da das Bild von $u$ eine Gruppe ist, genügt es zu zeigen, dass $\left.u\right|_{\left(V_{i}\right)_{i}}$ surjektiv auf eine Umgebung der Null ist. Ist $z_{i}$ eine Koordinate bei $x_{i}$ und\\
$\omega_{k}=f_{i k}\left(z_{i}\right) d z_{i}$ in dieser Koordinate, so besagt der Satz über implizite Funktionen, das Surjektivität auf eine Umgebung aus

$$
\operatorname{det}\left(f_{i k}(0)\right)_{i, k=1 \ldots g} \neq 0
$$

folgt. Also ist zu zeigen, dass die Auswertungsabbildung

$$
\mathrm{ev}:\left\{\begin{array}{lll}
\Omega^{1}(X) & \longrightarrow & \mathbb{C}^{g} \\
\omega=f\left(z_{i}\right) d z_{i} & \longmapsto & \left(f\left(x_{1}\right), \ldots, f\left(x_{g}\right)\right)
\end{array}\right.
$$

ein Isomorphismus ist. Offenbar genügt es dafür, $\operatorname{Ker}(\mathrm{ev})=0$ zu zeigen. Um das sicher zu stellen, wählen wir nun die Punkte $x_{i}$ geeignet. Sei $\omega_{1} \neq 0$ beliebig und $x_{1}$ mit $\omega_{1}\left(x_{1}\right) \neq 0$. Sei $\omega_{2} \neq 0$ mit $\omega_{2}\left(x_{1}\right)=0$ gewählt. (Falls es so ein $\omega_{2}$ nicht gibt, sind wir sowieso fertig). Wir nehmen $x_{2} \neq x_{1}$ mit $\omega_{2}\left(x_{2}\right) \neq 0$. Schließlich nehmen wir $\omega_{3} \neq 0$ mit $\omega_{3}\left(x_{1}\right)=\omega_{3}\left(x_{2}\right)=0$ (falls es das noch gibt), wählen $x_{3} \notin\left\{x_{1}, x_{2}\right\}$ mit $\omega_{3}\left(x_{3}\right) \neq 0$ usw. Falls wir den Prozess nicht bis zum $g$-ten Schritt durchführen können, sind wir (wie oben) sowieso fertig. Anderfalls ist Matrix $\left(\mathrm{ev}\left(\omega_{1}\right), \ldots, \mathrm{ev}\left(\omega_{g}\right)\right.$ nach Konstruktion eine Dreiecksmatrix und wir haben durch diese Konstruktion eine Basis $\omega_{1}, \ldots, \omega_{g}$ von $\Omega^{1}(X)$ erhalten. Verschwindet nun die Evaulation einer Linearkombination $\sum_{i=1}^{g} \mu_{i} \omega_{i}$ bei allen Punkten $x_{j}, j=1, \ldots, g$, so argumentiert man sukzessive mit der Dreiecksgestalt, dass $\mu_{i}=0$ für alle $i$.

Lemma 10.4 Die Spaltenvektoren $S_{1}, \ldots, S_{2 g}$ sind $\mathbb{R}$-linear unabhängig.\\
Beweis: Wir nehmen an, dass dies nicht so ist. Dann ist $\Lambda$ in einer reellen Hyperebene von $\mathbb{R}^{2 g}$ enthalten, also $\mathbb{C}^{g} / \Lambda \cong \mathbb{R} \times\left(\mathbb{R}^{2 g-1} / \Lambda\right)$ ist nicht kompakt. Wir betrachten die stetige Abbildung

$$
v:\left\{\begin{array}{lll}
X^{g} & \longrightarrow \operatorname{Pic}_{0}(X) & \longleftrightarrow \mathbb{C}^{g} / \Lambda \\
\left(x_{i}\right) & \longmapsto \sum_{i=1}^{g}\left(x_{i}\right)-\left(x_{0}\right) & \longmapsto\left(\sum_{i=1}^{g} \int_{x_{0}}^{x_{i}} \omega_{1}, \ldots, \sum_{i=1}^{g} \int_{x_{0}}^{x_{i}} \omega_{g}\right)
\end{array}\right.
$$

Da $X^{g}$ kompakt ist, ist auch das Bild kompakt. Andererseits behaupten wir, dass $v$ surjektiv ist. Von $u$ haben wir die Surjektivität bereits gezeigt. Sei $D \in \operatorname{Div}_{0}(X)$ beliebig. Nach Riemann-Roch ist

$$
\ell\left(g \cdot\left[x_{0}\right]+D\right) \geqslant g+1-g=1
$$

also gibt es eine Funktion $f$ mit $\operatorname{div}(f)=\sum_{i=1}^{g}\left[x_{i}\right]-g \cdot\left[x_{0}\right]-D$. Das bedeutet, dass modulo Hauptdivisoren $D=\sum_{i=1}^{g}\left(\left[x_{i}\right]-\left[x_{0}\right]\right)$ ist, und damit im Bild von $X^{g}$ liegt.

Definition 10.5 Eine Untergruppe $L \subseteq \mathbb{R}^{d}$ heißt Gitter, falls $L \cong \mathbb{Z}^{d}$ ist und $L$ diskret in $\mathbb{R}^{d}$ liegt.

Satz 10.6 $L \subseteq \mathbb{R}^{d}$ ist Gitter genau dann, wenn $L \cong \mathbb{Z}^{d}$ und $L$ nicht in einer Hyperebene liegt.

Korollar 10.7 Das Periodengitter $\Lambda$ ist ein Gitter in $\mathbb{R}^{2 g}$. Wir definieren $\{x\}=x- [x] \in[0,1)$ für ein $x \in \mathbb{R}$.

Beweis des Satzes: Seien $b_{1}, \ldots, b_{d}$ Erzeugende von $L$. Wenn $L$ nicht in einer Hyperebene liegt, so ist $\left\{b_{1}, \ldots, b_{d}\right\}$ eine $\mathbb{R}$-Basis von $\mathbb{R}^{d}$. Daher ist

$$
\left\{\alpha_{1} b_{1}+\ldots+\alpha_{d} b_{d}:-1<\alpha_{i}<1\right\} \cap L=\{0\}
$$

Die Menge $\left\{\alpha_{1} b_{1}+\ldots \cdot \alpha_{d} b_{d}:-1<\alpha_{i}<1\right\}$ ist ein Parallelotop in $\mathbb{R}^{d}$, sie enthält also eine Kugel um Null von einem Radius $\rho$. Wäre $L$ nicht diskret, so gäbe es eine Folge $x \neq \ell_{n} \in L$, die gegen ein $x \in \mathbb{R}^{d}$ konvergiert. Dann ist aber $\lim _{n \rightarrow \infty}\left\|\ell_{n}-\ell_{n+1}\right\|=0$. Insbesondere ist für ein $n$ der Abstand $\left\|\ell_{n}-\ell_{n+1}\right\|<\rho$ und $\ell_{n} \neq \ell_{n+1}$. Folglich liegt $\ell_{n+1}-\ell_{n}$ in der $\rho$-Kugel um Null. Umgekehrt, sei $d^{\prime}=\operatorname{dim}\left\langle b_{1}, \ldots, b_{d}\right\rangle_{\mathbb{R}}<d$. Ohne Einschränkung können wir annehmen, dass $b_{1}, \ldots, b_{d^{\prime}}$ diesen Unterraum $L_{\mathbb{R}}$ als $\mathbb{R}$-Vektorraum aufspannen. Es gilt $b_{d}=\sum_{i=1}^{d^{\prime}} \mu_{i} b_{i}$ und die $\mu_{i}$ liegen nicht alle in $\mathbb{Q}$, sonst wäre $L$ nicht isomorph zu $\mathbb{Z}^{d}$, Also liegen für alle $n \in \mathbb{N}$ die unendlich vielen Punkte $\sum_{i=1}^{d^{\prime}}\left\{n \mu_{i}\right\} b_{i}$ einerseits im Gitter, andererseits in der beschränkten Menge

$$
\left\{\alpha_{1} b_{1}+\ldots+\alpha_{d^{\prime}} b_{d^{\prime}}: 0 \leqslant \alpha_{i}<1 \quad \text { für } \quad i \in\left\{1, \ldots, d^{\prime}\right\} .\right.
$$

Also ist $L$ nicht diskret. \(\square\)

\subsection {Ein Beispiel}
Gegeben sei die Überlagerung $f: X \longrightarrow E$ mit der Monodromiedarstellung $\varphi(a)=\left(\begin{array}{lll}1 & 2 & 3\end{array}\right)\left(\begin{array}{lll}4 & 5 & 6\end{array}\right) \in S_{6}$ und $\varphi(b)=\left(\begin{array}{lll}1 & 4 & 2\end{array}\right)\left(\begin{array}{lll}3 & 6 & 5\end{array}\right) \in S_{6}$. Als flache Fläche mit $\omega=f^{*} \omega_{E}$ ist diese untenstehend abgebildet. Wir nehmen nun den Fall, dass $X$ den Torus mit Perioden 1 und $\zeta=e^{2 \pi i / 6}$ überlagert.

\begin{figure}[h]
\begin{center}
  \includegraphics[scale=0.2]{2025_10_06_374d4cfa8b1948eb4c41g-126}
 
\caption{Abbildung 10.2: Die flache Fläche $X$ zur Monodromiedarstellung $\varphi$}
\end{center}
\end{figure}

Was ist die Periodenmatrix von $X$ ? Die geschlossenen Wege $a_{1}, a_{2}, b_{1}, b_{2}$ bilden eine symplektische Basis von $H_{1}(X, \mathbb{Z})$, d.h. die vorzeichenbehaftete Schnittpaarung (siehe einen folgenden Abschnitt) erfüllt $\left\langle a_{i}, b_{j}\right\rangle= \delta_{i j},\left\langle a_{i}, a_{j}\right\rangle=\left\langle b_{i}, b_{j}\right\rangle=0$ für alle ( $i, j$ ). Wir nehmen für den Moment ohne Beweis an, dass wir die Einträge der Periodenmatrix statt mittels der Wege einer kanonischen Zerschneidung auch mit Hilfe einer symplektischen Basis berechnen können. Ist $\left\{\omega, \omega_{2}\right\}$ also eine Basis von $\Omega^{1}(X)$, so sehen wir leicht, dass die erste Zeile der Periodenmatrix ( $3,3, \zeta+1, \zeta+1$ ) ist. Aber wie kommt man an die zweite Zeile heran und für welches $\omega_{2}$ ist diese günstig zu berechnen? Im Allgemeinen ist diese Aufgabe (fast) hoffnungslos und auch hier gelingt sie nur aufgrund vieler Zufälle. Zunächst stellen wir fest, dass $X$ als Translationsfläche gleich den folgenden Flächen ist.

\begin{figure}[h]
\begin{center}
  \includegraphics[scale=0.2]{2025_10_06_374d4cfa8b1948eb4c41g-127}
 
\caption{Abbildung 10.3: Weitere Zerschneidungen von $X$}
\end{center}
\end{figure}

Aus der Eulercharakteristik sehen wir, dass $g(X)=2$ ist, also hyperelliptisch sein muss. Die Rotation der Sechsecke um $180^{\circ}$ gefolgt von Vertauschen der zwei Sechsecke ist kompatibel mit den Verklebungen der Kanten und hat genau 6 Fixpunkte, die Kantenmitten. Es ist also die hyperelliptische Involution $h: X \rightarrow X$. Wir bezeichnen die zugehörige zweiblättrige Überlagerung mit $\pi: X \longrightarrow \mathbb{P}^{1}$.\\
Die Rotation jedes Sechsecks um $60^{\circ}$ ist ein weiterer Automorphismus $\varphi$ von $X$. Man beachte, dass $\varphi^{3}$ die Blätter nicht vertauscht, also nicht gleich $h$ ist. Offenbar gilt aber $\varphi \circ h=h \circ \varphi$. Daher gibt es einen Automorphismus $\bar{\varphi}$ von $\mathbb{P}^{1}$, sodass $\pi \circ \varphi=\bar{\varphi} \circ \pi=: \bar{\pi}$ ist. Betrachtet man die Verzweigungspunkte, so wird klar, dass $\bar{\pi}$ wieder Grad 6 hat. Legt man die Fixpunkte von $\bar{\varphi}$ nach 0 und $\infty$ und normiert mit Hilfe einer geeigneten Möbiustransformation einen Verzweigungspunkt von $\pi$ nach 1, so müssen die anderen Verzweigungspunk-\\
te $\zeta_{6}^{i}, i=1, \ldots, 5$ sein. Also hat $X$ die Gleichung

$$
y^{2}=z^{6}-1
$$

und es gilt

$$
h(y, z)=(-y, z), \quad \varphi(y, z)=\left(y, \zeta_{6} \cdot z\right) .
$$

Da $\varphi$ die flache Struktur respektiert, gilt $\varphi^{*} z=\lambda \cdot z$ für ein $\lambda \in \mathbb{C}^{*}$. Anhand der Perioden sieht man sofort, dass $\lambda=\zeta_{6}$ sein muss. Nun wird auch die natürliche Wahl von $\omega_{2}$ klar. Da $\Omega^{1}(X)=\left\langle\frac{d z}{y}, z \frac{d z}{y}\right\rangle$ ist, sind die Eigendifferentiale von $\varphi$ gerade $\frac{d z}{y}$ und $z \frac{d z}{y}$. Das erste hat Eigenwert $\zeta_{6}$, das zweite Eigenwert $\zeta_{6}^{2}$. Das erste muss also $\omega$ sein, das zweite nehmen wir als $\omega_{2}$. Als nächstes wollen wir die Aktion von $\varphi$ auf $H_{1}(X, Z)=\pi_{1}(X)^{a b}$ bestimmen.\\
Im Bild 10.5 ist die sympektische Basis $\left\{a_{1}, a_{2}, b_{1}, b_{2}\right\}$ mit Perioden $(3,3, \zeta+ 1, \zeta+1)$ nochmal in der Sechseckpräsentation der Fläche gezeichnet. Im Bild 10.5 ist\\
ein geschlossener Weg eingezeichnet, der offenbar homotop zu $\varphi\left(a_{1}\right)$ ist. Er setzt sich aus $b_{2}, b_{1},-a_{2}, b_{2}$ und der Schleife um den schwarzen Punkt, welche natürlich nullhomotop ist, zusammen. Also ist $\varphi\left(a_{2}\right)=b_{1}+2 b_{2}-a_{2} \in H_{1}(X, \mathbb{Z})$.\\
Ganz analog ist $\varphi\left(a_{1}\right)=2 b_{1}+b_{2}-a_{1}$ und $\varphi\left(b_{1}\right)=\varphi\left(a_{1}\right)-b_{1}=b_{1}+b_{2}-a_{1}$ sowie analog $\varphi\left(b_{2}\right)=b_{1}+b_{2}-a_{2}$ läßt sich mit gleicher Methode herleiten. Also ist die Darstellungsmatrix von $\varphi$ auf $H_{1}(X, \mathbb{Z})$ in der Basis $\left\{a_{1}, a_{2}, b_{1}, b_{2}\right\}$ gleich

$$
M_{\varphi}=\left(\begin{array}{rrrr}
-1 & 0 & 2 & 1 \\
0 & -1 & 1 & 2 \\
-1 & 0 & 1 & 1 \\
0 & -1 & 1 & 1
\end{array}\right)
$$

Für die Perioden gilt offenbar

$$
\int_{a_{i}} \omega_{j}=\int_{\varphi\left(\omega_{j}\right)}\left(\varphi^{*}\right)^{-1} \omega_{j}
$$

Ist also

$$
\Pi=\left(\int_{a_{i}} \omega_{j} \int_{b_{i}} \omega_{j}\right)_{i, j=1, \ldots g} \in \mathbb{C}^{g \times 2 g}
$$

die Periodenmatrix und $M_{\varphi} \in \mathrm{GL}_{2 g}$ die Darstellung von $\varphi$ auf $H_{1}(X, \mathbb{Z})$ sowie $A_{\varphi} \in \mathrm{GL}_{g}(\mathbb{C})$ die Darstellung von $\varphi$ auf $\Omega^{1}(X)$, so besagt (10.1), dass

$$
A_{\varphi}^{-1} \cdot \Pi \cdot M_{\varphi}^{T}=\Pi
$$

Im konkreten Fall muss also

$$
\left(\begin{array}{cc}
\zeta^{-1} & 0 \\
0 & \zeta^{-2}
\end{array}\right)\left(\begin{array}{cccc}
3 & 3 & 1+\zeta & 1+\zeta \\
\alpha_{1} & \alpha_{2} & \beta_{1} & \beta_{2}
\end{array}\right) M_{\varphi}^{T}=\left(\begin{array}{cccc}
3 & 3 & 1+\zeta & 1+\zeta \\
\alpha_{1} & \alpha_{2} & \beta_{1} & \beta_{2}
\end{array}\right)
$$

\begin{figure}[h]
\begin{center}
  \includegraphics[scale=0.2]{2025_10_06_374d4cfa8b1948eb4c41g-129}
 
\caption{Abbildung 10.4: Symplektische Basis von \$H\_\{1}(X, \textbackslash mathbb\{Z\})\$ und Bestimmung von $\varphi\left(a_{2}\right) \in H_{1}(X, \mathbb{Z})$\}\end{center}
\end{figure}

gelten. Die obere Zeile prüft man leicht nach, aus der unteren Zeile erhält man ein lineares Gleichungssystem für $\alpha_{1}, \alpha_{2}, \beta_{1}, \beta_{2}$, dessen Lösungsraum eindimensional ist. Mehr kann man nicht erwarten, da wir $\omega_{2}$ als zweites Eigendifferential nur bis auf skalare Vielfache spezifiziert haben. Die Periodenmatrix in diesem Beispiel ist also

$$
\Pi=\left(\begin{array}{cccc}
3 & 3 & 1+\zeta & 1+\zeta \\
\zeta^{2} & \zeta^{5} & -1 & 1
\end{array}\right) .
$$

\begin{figure}[h]
\begin{center}
  \includegraphics[scale=0.2]{2025_10_06_374d4cfa8b1948eb4c41g-130}
 
\caption{Abbildung 10.5: Symplektische Basis von \$H\_\{1}(X, \textbackslash mathbb\{Z\})\$ und Bestimmung von $\varphi\left(a_{2}\right) \in H_{1}(X, \mathbb{Z})$\}\end{center}
\end{figure}

\section {Grundbegriffe aus der Topologie}
In diesem Abschnitt bezeichnen wir für eine Menge $X$ mit $\mathcal{P}(X)$ immer die Potenzmenge von $X$.

Definition 1.1 (Topologischer Raum) Ein topologischer Raum ist ein Paar ( $X, \mathcal{O}$ ) bestehend aus einer Menge $X$ und einer Menge $\mathcal{O} \subset \mathcal{P}(X)$, so dass gilt:\\
i) Es sind $\emptyset \in \mathcal{O}$ und $X \in \mathcal{O}$.\\
ii) Für alle $U_{1}, U_{2} \in \mathcal{O}$ ist auch $U_{1} \cap U_{2} \in \mathcal{O}$.\\
iii) Für jede Familie $\left\{U_{i}\right\}_{i \in I} \subset \mathcal{O}$ ist auch $\bigcup_{i \in I} U_{i} \in \mathcal{O}$.

Die Elemente von $\mathcal{O}$ heißen offene Mengen, Komplemente $U^{c}$ von offenen Mengen $U$ heißen abgeschlossene Mengen. $\mathcal{O}$ selbst heißt Topologie auf $X$.\\
i) Ist $U \subset X$ offen und $x \in U$, dann heißt $U$ offene Umgebung von $x$. Ist $M \subset X$ eine beliebige Menge, die eine offene Umgebung von $x$ enthält, dann heißt $M$ Umgebung von $x$.\\
ii) Für $Y \subset X$ ist $\stackrel{\circ}{Y}:=\bigcup_{U \in \mathcal{O}, U \subset Y} U$ die größte offene Menge, die in $Y$ enthalten ist und heißt Inneres von $Y$.\\
iii) Für $Y \subset X$ ist $\bar{Y}:=\bigcap_{A^{c} \in \mathcal{O}, Y \subset A} A$ die kleinste abgeschlossene Menge, die $Y$ enthält und heißt Abschluss von $Y$.

Sind zwei Topologien $\mathcal{O}_{1}, \mathcal{O}_{2}$ auf der gleichen Menge $X$ gegeben mit $\mathcal{O}_{1} \subset \mathcal{O}_{2}$, dann heißt $\mathcal{O}_{1}$ gröber als $\mathcal{O}_{2}$ und $\mathcal{O}_{2}$ feiner als $\mathcal{O}_{1}$.

Beispiel 1.2 Sei $X$ eine beliebige Menge. Dann sind sowohl ( $X,\{\emptyset, X\}$ ) als auch ( $X, \mathcal{P}(X)$ ) topologische Räume. Wir nennen $\{\emptyset, X\}$ die triviale und $\mathcal{P}(X)$ die diskrete Topologie auf $X$.

Beispiel 1.3 Wie aus der Analysis bekannt definieren wir auf dem Raum $\mathbb{R}^{n}$ eine Topologie $\mathcal{O}$ durch folgende Vorschrift: Eine Menge $U \subset \mathbb{R}^{n}$ liegt genau dann in $\mathcal{O}$, wenn für jeden Punkt $x \in U$ ein $\varepsilon>0$ existiert mit

$$
B_{\varepsilon}(x):=\left\{y \in \mathbb{R}^{n}:|x-y|<\varepsilon\right\} \subset U .
$$

Die Topologie $\mathcal{O}$ werden wir in der Notation oft unterdrücken, also nur von einem topologischen Raum $X$ sprechen.

Die Morphismen in der Kategorie der topologischen Räume sind die stetigen Abbildungen:

Definition 1.4 (Stetigkeit) Seien $X, Y$ topologische Räume. Eine Abbildung $f$ : $X \rightarrow Y$ heißt stetig, wenn gilt: Für jede offene Menge $U \subset Y$ ist auch das Urbild $f^{-1}(U) \subset X$ offen.\\
Eine stetige Abbildung $f: X \rightarrow Y$ heißt Homöomorphismus, wenn es eine stetige Abbildung $f^{-1}: Y \rightarrow X$ mit $f^{-1} \circ f=i d_{X}$ und $f \circ f^{-1}=i d_{Y}$ gibt.

Die Komposition von zwei stetigen Abbildungen ist wieder stetig und die Identität $\mathrm{id}_{X}: X \rightarrow X$ ist ebenfalls stetig.

Bei der letzten Definition ist etwas Vorsicht geboten: Anders, als man es aus der linearen Algebra für lineare Abbildungen gewohnt ist, ist eine stetige Bijektion nicht automatisch ein Homöomorphismus, die Umkehrabbildung also nicht automatisch ebenfalls stetig.

Die grundlegendsten Methoden, aus gegebenen topologischen Räumen neue Räume zu basteln, sind die der Teilraum-, Summen-, Produkt- und Quotientenbildung:

Definition und Proposition 1.5 (Teilraum) Sei ( $X, \mathcal{O}$ ) ein topologischer Raum, $T \subset X$ eine Teilmenge. Wir definieren

$$
\left.\mathcal{O}\right|_{T}:=\{U \cap T: U \in \mathcal{O}\}
$$

Dann ist $\left(T, \mathcal{O}_{T}\right)$ ein topologischer Raum und $\mathcal{O}_{T}$ heißt die von $\mathcal{O}$ induzierte Teilraumtopologie auf $T$.

Insbesondere kann ein Teilraum $T \subset X$ offene Mengen haben, die in $X$ nicht offen sind, so ist z.B. für $X=\mathbb{R}$ und $T=[0,1]$ die Menge $\left[0, \frac{1}{2}\right)=\left(-348, \frac{1}{2}\right) \cap T$ offen in $T$.

Definition und Proposition 1.6 (Summe) Seien ( $X_{1}, \mathcal{O}_{1}$ ), ( $X_{2}, \mathcal{O}_{2}$ ) topologische Räume, und $X_{1}+X_{2}:=X_{1} \coprod X_{2}$ die disjunkte Vereinigung von $X_{1}$ und $X_{2}$. Wir definieren

$$
\mathcal{O}:=\left\{U_{1} \cup U_{2}: U_{i} \in \mathcal{O}_{i}\right\}
$$

Dann ist ( $X_{1}+X_{2}, \mathcal{O}$ ) ein topologischer Raum, die topologische Summe von $X_{1}$ und $X_{2}$.

Definition und Proposition 1.7 (Produkt) Seien $\left(X_{1}, \mathcal{O}_{1}\right),\left(X_{2}, \mathcal{O}_{2}\right)$ topologische Räume. Wir definieren auf dem kartesischen Produkt eine Topologie $\mathcal{O}$ wie folgt: Es sei $U \subset \mathcal{O}$ genau dann, wenn es zu jedem Punkt $x \in U$ offene Mengen $U_{1} \subset \mathcal{O}_{1}$ und $U_{2} \subset \mathcal{O}_{2}$ gibt mit $x \in U_{1} \times U_{2} \subset U$. Dann ist ( $X_{1} \times X_{2}, \mathcal{O}$ ) ein topologischer Raum, das topologische Produkt von $X_{1}$ und $X_{2}$.

Das Produkt trägt damit die gröbste Topologie, bezüglich der die beiden Projektionsabbildungen

$$
\pi_{i}: X_{1} \times X_{2} \rightarrow X_{i}
$$

stetig sind.

Definition und Proposition 1.8 (Quotient) Sei ( $X, \mathcal{O}$ ) ein topologischer Raum, \~{} eine Äquivalenzrelation auf $X$ und

$$
\pi: X \rightarrow X / \sim, x \mapsto[x]
$$

die kanonische Projektion auf die Menge der Äquivalenzklassen. Mit

$$
\mathcal{O}_{\sim}:=\left\{U \in \mathcal{P}(X / \sim): \pi^{-1}(U) \in \mathcal{O}\right\}
$$

ist dann ( $X / \sim, \mathcal{O}_{\sim}$ ) ein topologischer Raum, der Quotientenraum von $X$ modulo $\sim$.

Der Quotient trägt damit die feinste Topologie, bezüglich der die kanonische Projektion stetig ist.

Im Folgenden einige wichtige Eigenschaften, die ein topologischer Raum haben kann, sowie Aussagen über die Verträglichkeit mit den bisherigen Konstruktionen:

Definition 1.9 (Zusammenhang) Sei $X$ ein topologischer Raum.\\
i) $X$ heißt zusammenhängend, wenn gilt:

Für alle $U_{1}, U_{2} \subset X$ offen mit $U_{1} \cap U_{2}=\emptyset$ und $U_{1} \cup U_{2}=X$ gilt $U_{1}=\emptyset$ oder $U_{2}=\emptyset$.\\
ii) $X$ heißt wegzusammenhängend, wenn gilt:

Für alle $x, y \in X$ existiert ein Weg von $x$ nach $y$, d.h. eine stetige Abbildung $\gamma:(0,1) \rightarrow X$ mit $\gamma(0)=x$ und $\gamma(1)=y$.

Proposition 1.10 Sei $f: X \rightarrow Y$ stetig und $f(X)$ mit der Teilraumtopologie von $Y$ ausgestattet. Dann gilt:\\
i) Ist $X$ zusammenhängend, dann ist auch $f(X)$ zusammenhängend.\\
ii) Ist $X$ wegzusammenhängend, dann ist auch $f(X)$ wegzusammenhängend.\\
iii) Jeder wegzusammenhängende Raum ist zusammenhängend.\\
iv) Für Mannigfaltigkeiten (wie z.B. Riemannsche Flächen) gilt sogar: Jede zusammenhängende Mannigfaltigkeit ist wegzusammenhängend.

Definition 1.11 (Hausdorffraum) $X$ heißt Hausdorffsch, wenn gilt:\\
Für alle $x, y \in X$ existiert eine offene Umgebung $U_{x}$ von $x$ und eine offene Umgebung $U_{y}$ von $y$ mit $U_{x} \cap U_{y}=\emptyset$.

Teilräume von Hausdorffräumen sind Hausdorffsch. Zwei nicht-leere topologische Räume sind genau dann Hausdorffsch, wenn ihre Summe Hausdorffsch ist und genau dann, wenn ihr Produkt Hausdorffsch ist.

Definition 1.12 (Kompaktheit) $X$ heißt kompakt, wenn gilt:\\
Jede offene Überdeckung hat eine endliche Teilüberdeckung, soll heißen: Ist $\left\{U_{i}\right\}_{i \in I}$ eine Familie von offenen Mengen in $X$ mit $\bigcup_{i \in I} U_{i}=X$, dann existieren endlich viele $i_{1}, \ldots, i_{n} \in I$ mit $\bigcup_{k=1}^{n} U_{i_{k}}=X$.

Proposition 1.13 i) Ist $X$ kompakt und $f: X \rightarrow Y$ stetig, dann ist $f(X)$ ebenfalls kompakt.\\
ii) Ist $X$ kompakt und $A \subset X$ abgeschlossen, dann ist $A$ ebenfalls kompakt.\\
iii) Ist $X$ ein Hausdorffraum und $K \subset X$ kompakt, dann ist $K$ in $X$ abgeschlossen.\\
iv) Zwei nicht-leere topologische Räume sind genau dann kompakt, wenn ihre Summe kompakt ist und genau dann, wenn ihr Produkt kompakt ist.

Zum Schluss seien hier noch die Begriffe Grenzwerte und Häufungspunkte erklärt:

Definition 1.14 (Limes) Sei $\left(x_{n}\right)_{n \in \mathbb{N}}$ eine Folge in einem topologischen Raum $X$. Ein Punkt $x \in X$ heißt ein Grenzwert oder Limes dieser Folge, wenn es zu jeder offenen Umgebung $U$ von $x$ ein $n_{0} \in \mathbb{N}$ gibt, so dass $x_{n} \in U$ für alle $n \geq n_{0}$ gilt.

Folglich hat in einem Hausdorffraum jede Folge höchstens einen Grenzwert.\\
Definition 1.15 (Häufungspunkt) Sei $Y \subset X$. Ein Punkt $x \in X$ heißt Häufungspunkt von $Y$, wenn für jede offene Umgebung $U \subset X$ von $x$ die Menge $U \cap Y$ nicht endlich ist.

Proposition 1.16 Insbesondere ist jeder Grenzwert einer Folge aus $Y$ ein Häufungspunkt von $Y$. Umgekehrt ist bei Mannigfaltigkeiten (wie z.B Riemannschen Flächen) auch jeder Häufungspunkt von $Y$ Grenzwert einer Folge aus $Y$.

\section {Grundbegriffe aus der Gruppentheorie}
\subsection {Erzeuger und Relationen}
Elemente von Gruppen werden durch Symbole dargestellt, z.B. in der multiplikativen Gruppe $(\mathbb{Z} / 7 \mathbb{Z} \backslash\{0\}, \cdot)$ durch $1,2, \ldots, 6$, manchmal mit einem Querstrich um Restklassen anzudeuten. Aus diesen Gruppenelementen können wir duch Aufmultiplizieren neue bilden, z.B. das Gruppenelement $5 \cdot 3 \cdot 2$. Jedes Element hat ein inverses Element, das Inverse zu 5 bezeichnen wir mit $5^{-1}$. In dieser Gruppe gelten Relationen zwischen den Symbolen, z.B. ist $2 \cdot 3 \cdot 6^{-1}$ das neutrale Element und auch $2 \cdot 4$ ist das neutrale Element. Diese Idee, Gruppen mit Erzeugern und Relationen zu beschreiben werden wir hier knapp formalisieren. Eine ausführliche Quelle hierzu ist [MKS76].

Sei $A$ eine Menge von Symbolen, in diesem Kontext auch Alphabet genannt. Ein Wort in $A$ ist eine endliche Folge von Symbolen

$$
w=\left(w_{1} w_{2} w_{3} \ldots w_{n}\right)
$$

wobei jedes $w_{i}$ entweder ein Element von $A$ oder von der Form $a^{-1}$ mit $a \in A$ ist. Die Zahl $n$ wird auch Länge des Worts genannt. Wir definieren das inverse Wort von $w$ als $w^{-1}=w_{n}^{-1} \ldots w_{2}^{-1} w_{1}^{-1}$. Sind $w=\left(w_{1} w_{2} w_{3} \ldots w_{n}\right)$ und $x=\left(x_{1} x_{2} x_{3} \ldots x_{s}\right)$ Worte, so definiert man die Verkettung als $w x= w_{1} \ldots w_{n} x_{1} \ldots x_{s}$.

Ist $A$ eine Teilmenge einer Gruppe, so gibt es eine offensichtliche Auswertungsabbildung ev, die jedem Wort ein Gruppenelement zuordnet, sodass $\mathrm{ev}\left(w_{1} w_{2}\right)=\mathrm{ev}\left(w_{1}\right) \mathrm{ev}\left(w_{2}\right)$ ist. Jedes Wort, dass unter ev auf das neutrale Element abgebildet wird, nennt man Relation in der Gruppe. (Der Begriff Relator anstelle von Relation ist auch gebräuchlich.) Die Relationen $a a^{-1}$ und $a^{-1} a$ für $a \in A$ gibt es in jeder Gruppe und sie werden triviale Relationen genannt. Ist die Auswertungsabbildung surjektiv, so wird $A$ erzeugend genannt.

Wir starten zunächst von einer gegebenen Gruppe $G$. Es gibt offenbar verschiede Arten ein Gruppenelement als Wort zu schreiben. Dies formalisiert der folgende Begriff.

Sei $w$ ein Wort im Alphabet $A$. Dann heißt das Wort $w^{\prime}$ aus $w$ mit Hilfe der Relationenmenge $R$ abgeleitet, wenn es eine endliche Folge der folgenden zwei Operationen gibt, die $w$ in $w^{\prime}$ überführt. Die erste Ableitungsoperationen ist das Einfügen eines des Worte $r$ oder $r^{-1}$ für $r \in R$ oder $r$ einen der trivialen Relationen. Die zweite Ableitungsoperation ist das Herausstreichen von $r$ oder $r^{-1}$ für $r \in R$ oder für $r$ eine der trivialen Relationen. Kann man jede Relation in $G$ auf diese Art aus einer Relation in $R$ ableiten, so wird die Relationenmenge $R$ definierend genannt. Ein $\operatorname{Paar}(A, R)$ bestehend aus einem Erzeugendensystem und einer definierenden Relationenmenge wird Präsentation der Gruppe $G$ genannt.

Offenbar besitzt jede Gruppe eine Präsentation, indem man $A=G$ nimmt und für $R$ alle Relationen, die in $G$ gelten.

Definition 2.1 (Endlich erzeugt und präsentiert) Eine Gruppe G ist endlich erzeugt, falls es ein endliches Erzeugendensystem gibt. Sie ist endlich präsentierbar, falls sie endliche erzeugt ist und es eine endliche definierende Relationenmenge gibt.

Wir wollen nun umgekehrt Gruppen aus einem gegebenen Alphabet $A$ und einer Relationenmenge $R$ konstruieren. Dazu betrachten wir die Menge der Worte in $A$ und führen darauf die Relation $\sim_{R}$ ein, indem wir sagen, dass $w \sim_{R} w^{\prime}$, falls $w^{\prime}$ in obigem Sinne mit Hilfe von $R$ aus $w$ abgeleitet ist. offenbar ist $\sim_{R}$ eine Äquivalenzrelation.

Definition und Proposition 2.2 (Freie Gruppen) Die Menge der Äquivalenzklassen von Worten bezüglich der Relation $\sim_{R}$ ist bzgl. Verkettung von Worten eine

Gruppe mit dem leeren Wort als neutralem Element. Wir bezeichnen sie als die Gruppe mit der Präsentation ( $A, R$ ). Die Gruppe ( $A, \emptyset$ ) heißt die freie Gruppe $F_{A}$ über dem Alphabet $A$.

Beweis: Offenbar ist das leere Wort neutral bzgl. der Verkettung und das inverse Wort ist ein inverses Element. Die Verkettung von Worten ist assoziativ und damit ist auch die Verknüpfung auf dem Niveau von Äquivalenzklassen assoziativ.

Freie Gruppen sind in gewissem Sinne die einfachsten Gruppen überhaupt. Sie sind definiert durch Worte in Symbolen und ausser der Existenz eines Inversen gibt es keine Relationen. Dadurch sind sie andererseits 'alles andere' als abelsch (zumindest wenn es mehr als ein Symbol gibt) und somit 'kompliziert'.

Der Rang einer freien Gruppe ist die Mächtigkeit von $A$. Offenbar sind zwei freie Gruppen von gleichem Rang isomorph. Die freie Gruppe vom Rang 1 ist isomorph $\mathrm{zu} \mathbb{Z}$.

Die Untergruppe von $\mathrm{SL}_{2}(\mathbb{Z})$ erzeugt von $\left(\begin{array}{ll}1 & 2 \\ 0 & 1\end{array}\right)$ und $\left(\begin{array}{ll}1 & 0 \\ 2 & 1\end{array}\right)$, ist eine freie Gruppe vom Rang zwei. (Die natürliche Abbildung nach SL ${ }_{2}(\mathbb{Z})$ bildet sie surjektiv auf die Gruppe aller Matrizen $\Gamma$ (2) kongruent zu 1 modulo zwei auf den Diagonalelementen und zu 0 modulo zwei auf den Nebendiagonalelementen ab.) Der Beweis hiervon ist eine Anwendung des 'Ping-Pong-Lemmas'.

In einer Gruppe, die gegeben ist durch Erzeuger und Relationen ist es nicht einfach zu entscheiden, ob zwei Elemente gleich sind, da es nicht klar ist, in welcher Reihenfolge die Ableitungsschritte anzuwenden sind. Gleiches gilt für die Frage, ob zwei Elemente konjugiert zueinander sind. Diese Fragen sind als Wortproblem und Konjugationsproblem bekannt, für manche Gruppen bekannt und für viele Gruppen offen. In der freien Gruppe hat das Konjugationsproblem eine einfache Lösung. Dazu nennen wir ein Wort $w$ reduziert wenn $w$ keine Symbolpaar der Form $a a^{-1}$ oder $a^{-1} a$ für $a \in A$ enthält.

Proposition 2.3 In einer freien Gruppe ist jedes Wort von genau einem reduzierten Wort abgeleitet.

Beweis: Ist $w$ nicht reduziert, so kann man an der ersten sich anbieteten Stelle kürzen, d.h. eine trivialen Relation $a a^{-1}$ oder $a^{-1} a$ entfernen. Auf diese Weise verkürzt sich die Länge des Wortes und man endet nach endlich vielen Schritten bei einem reduzierten Wort $\rho(w)$.

Es ist nur noch zu zeigen, das dieses reduzierte Wort das einzige reduzierte Wort ist, dass aus $w$ abgeleitet werden kann. Dazu stellt man fest, dass die Kürzungsprozedur offenbar die Eigenschaften $\rho\left(w_{1} w_{2}\right)=\rho\left(\rho\left(w_{1}\right) w\right)$, und daher $\rho\left(w_{1} a^{\varepsilon} a^{-\varepsilon} w_{2}\right)=\rho\left(w_{1} w_{2}\right)$ für alle $a \in A, \varepsilon \in \pm 1, w_{1}, w_{2} \in F_{A}$ hat. Angenommen zwei reduzierte Worte $w$ und $v$ sind voneinander abgeleitet. Dann gibt es eine Folge von Worten $w=w_{1}, w_{2}, \ldots, w_{k}=v$, sodass $w_{i+1}$\\
aus $w_{i}$ oder $w_{i}$ aus $w_{i+1}$ durch Kürzen von $a^{\varepsilon} a^{-\varepsilon}$ hervorgeht. Dann aber ist $w=\rho(w)=\rho\left(w_{1}\right)=\cdots \rho\left(w_{n}\right)=\rho(v)=v$, was zu zeigen war.

\subsection {Freie und amalgamierte Produkte}
Freie Gruppen werden hier als Lösung eines universellen Abbildungproblems noch einmal auftreten. Als solches sind sie Spezialfall eines freien Produkts und das wieder um ist der Spezialfall eines amalgamierten Produkts. Da amalgamierte Produkte bei Riemannschen Flächen zum Beispiel im Zusammenhang mit dem Satz von Seifert-van-Kampen auftreten, fangen wir mit diesem allgemeinen Fall an.

Seien $B_{0}, B_{1}$ und $B_{2}$ Gruppen und $s_{i}: B_{0} \rightarrow B_{i}$ Gruppenhomomorphismen.\\
Definition 2.4 (Freies und amalgamiertes Produkt) Eine Gruppe $G$ mit Homomorphismen $\alpha_{i}: B_{i} \rightarrow G, i=1,2$, heißt amalgamiertes Produkt $G=B_{1} * B_{0} B_{2}$ von $B_{1}$ und $B_{2}$ entlang $s_{i}: B_{0} \rightarrow B_{i}$, falls $\alpha_{1} \circ s_{1}=\alpha_{2} \circ s_{2}$ gilt und folgende universelle Eigenschaft erfüllt ist: Zu jeder Gruppe $G^{\prime}$ mit Homomorphismen $\alpha_{i}^{\prime}: B_{i} \rightarrow G^{\prime}$, sodass $\alpha_{i}^{\prime} \circ s_{1}=\alpha_{2} \circ s_{2}$ gilt, gibt es genau einen Homomorphismus $\varphi: G \rightarrow G^{\prime}$ mit $\varphi \circ \alpha_{i}=\alpha_{i}^{\prime}$ für $i=1,2$.\\
Ist $B_{0}=\{1\}$, so heißt $G=B_{1} * B_{2}$ freies Produkt.\\
Man kann die Definition von Amalgamen auch auch $B_{i}, i \in I$ eine beliebigen Indexmenge, und $\varphi: B_{0} \rightarrow B_{i}$ ausdehnen, siehe dazu und für weitere Anwendungen von Amalgamen [MKS76] oder [Ser03, Kapitel 1].

Die Eindeutigkeit bis auf Isomorphie eines amalgamierten Produkts ist eine Standard-Schlussweise für universelle Abbdildungseigenschaften. Seien $G_{1}$ und $G_{2}$ zwei amalgamierte Produkte. Dann liefert die Definition Homomorphismen $\varphi: G_{1} \rightarrow G_{2}$ und $\psi: G_{1} \rightarrow G_{2}$. Wenn wir die Eindeutigkeitseigenschaft in der Definition auf $G=G_{1}, G^{!}=G_{1}$ und die Abbildungen $\psi \circ \varphi$ sowie $\operatorname{id}_{G_{1}}$ anwenden, erhalten wir $\psi \circ \varphi=\operatorname{id}_{G_{1}}$. Das gleiche Prinzip auf $G_{2}$ angewandt liefert $\varphi \circ \psi=\operatorname{id}_{G_{2}}$ und somit sind $G_{1}$ und $G_{2}$ isomorph.

Proposition 2.5 Amalgamierte Produkte existieren.\\
Beweis: Wir beschreiben $B_{1} *_{B_{0}} B_{2}$ mit Erzeugern und Relationen. Als Erzeuger nehmen wir alle Elemente von $B_{1}$, alle von $B_{2}$ und alle Elemente von $B_{0}$. Als Relation nehmen wir zunächst alle Element der Form $x y z^{-1}$ falls alle drei Symbole in der selben Gruppe $B_{i}$ liegen und falls in dieser Gruppe die Relation $x y=z$ gilt. Hierzu nehmen wir noch die Relationen $x y^{-1}$ hinzu, falls es ein $x \in B_{0}$ gibt, sodass $y=s_{1}(x) \in B_{1}$ oder sodass $y=s_{2}(x) \in B_{2}$ gilt. Die so definierte Gruppe $G$ kommt mit naürlichen Abbildungen $\alpha_{i}: B_{i} \rightarrow G$ daher, die $x \in B_{i}$ auf die Äquivalenzklassen von $x$ modulo der Relationen abbildet. Die Relationen stellen auch sicher, dass für diese Abbildungen $\alpha_{1} \circ s_{1}=\alpha_{2} \circ s_{2}$ gilt.

Die derart mit Erzeugern und Relationen definierte Gruppe $G$ erfüllt offenbar die universelle Eigenschaft: Sei $G^{\prime}$ mit den Homomorphismen $\alpha_{i}^{\prime}: B_{i} \rightarrow G^{\prime}$ wie in der Definition des amalgamierten Produkts, d.h. mit $\alpha_{1}^{\prime} \circ s_{1}=\alpha_{2}^{\prime} \circ s_{2}$. Falls ein Homomorphismus $\varphi: G \rightarrow G^{\prime}$ mit $\varphi \circ \alpha_{i}=\alpha_{i}^{\prime}$ existiert, so muss $\varphi\left(y_{i}\right)=\alpha_{i}^{\prime}\left(y_{i}\right)$ für alle $y_{i} \in B_{i}$ und $i=1,2$ sein. Da solche $y_{i}$ die Gruppe $G$ per Definition erzeugen, ist $\varphi$ eindeutig, falls es existiert. Zur Existenz von $\varphi$ setzen wir $\varphi\left(y_{i}\right)=\alpha_{i}^{\prime}\left(y_{i}\right)$ für $y_{i} \in B_{i}$ und $i=1,2$ und $\varphi(x)=\alpha_{1}^{\prime} \circ s_{1}(x)$ für $x \in B_{0}$ auf Worte in diesen Erzeugern fort. Damit ist $\varphi$ auf ganz $G$ erklärt und wir müssen nur prüfen, dass $\varphi$ alle Relationen auf das neutrale Element in $G^{\prime}$ schickt um die Wohldefiniertheit sicherzustellen. Für Relationen zwischen drei Elementen in $B_{i}$ ist dies durch die Homomorphie von $\alpha_{i}^{\prime}$ gewährleistet. Für Relationen der Form $y x^{-1}$ mit $y=s_{i}(x)$ ist dies aufgrund von $\alpha_{i}^{\prime}(y)=\alpha_{1}^{\prime}\left(s_{1}(x)\right)=\alpha_{2}^{\prime}\left(s_{2}(x)\right)$ sichergestellt.

Diese Konstruktion von amalgamierten Produkten ist zu konkreten Rechnungen offenbar nicht besonders nützlich, da das Erzeugendensystem unhandlich groß ist. Konkrete Beispiele und praktisches Rechnen in amalgamierten Produkten wird in [MKS76] eingehend behandelt.

Abschließend finden wir so die freien Gruppen wieder:\\
Proposition 2.6 Ist $F_{A_{i}}$ die freie Gruppe über dem Alphabet $A_{i}$ für $i=1,2$, so ist das freie Produkt $F_{A_{1}} * F_{A_{2}} \cong F_{\left\{A_{1} \amalg A_{2}\right\}}$.

Beweis: Die Inklusion von Alphabeten definiert natürliche Abbildungen $\alpha_{i}$ : $F_{A_{i}} \rightarrow F_{\left\{A_{1} \amalg A_{2}\right\}}$. Damit liefert die universelle Abbildungseigenschaft des freien Produkts einen Homomorphismus $\varphi: F_{A_{1}} * F_{A_{2}} \rightarrow F_{\left\{A_{1} \amalg A_{2}\right\}}$. Dieser trifft jeden Erzeuger von $F_{\left\{A_{1} \amalg A_{2}\right\}}$ und ist somit surjektiv. Sei $w$ ein Element im Kern von $\varphi$. Dann ist $w$ nach Konstruktion des amalgamierten Produkts eine Kette von Worten $w=w_{1} w_{2} \ldots, w_{n}$ mit $w_{i} \in F_{A_{1}}$ oder $w_{i} \in F_{A_{2}}$ für alle $i=1, \ldots, n$. Durch Zusammenfassen von Buchstaben in den jeweiligen $F_{A_{i}}$ können wir annehmen, dass zwei aufeinanderfolgende Teilworte $w_{i}$ in verschiedenen $F_{A_{i}}$ liegen. Die Voraussetzung $w \in \operatorname{Ker}(\varphi)$ besagt, dass sich $w$ durch den obigen Kürzungsalgorithmus $\rho$ zum trivialen Wort ableiten lässt. Nach der Voraussetzung über aufeinanderfolgende Teilworte $w_{i}$ und $w_{i+1}$ findet jeder Kürzungsschritt innerhalb einer der Gruppen $F_{A_{i}}$ statt. Also ist $\rho\left(w_{i}\right)$ das triviale Wort für alle $i$ und somit $w$ das neutrale Element in $F_{A_{1}} * F_{A_{2}}$.

\section {Literatur}
[FB06] E. Freitag und R. Busam. Funktionentheorie 1. Berlin: Springer DE, 2006.\\[0pt]
[FG81] O. Forster und B. Gilligan. Lectures on Riemann Surfaces -. 1st ed. 1981. Corr. 4th printing 1999. Berlin, Heidelberg: Springer, 1981.\\[0pt]
[FTDGL] Martin Möller. Funktionentheorie und Differentialgleichungen. 2013. eprint:Vorlesungsskript,Frankfurt/Main.\\[0pt]
[GJ00] E. Gutkin und C. Judge. "Affine mappings of translation surfaces: geometry and arithmetic". In: Duke Mathematical Journal 103.2 (2000), S. 191-214.\\[0pt]
[Hat02] A. Hatcher. Algebraic topology. Cambridge University Press, Cambridge, 2002, S. xii+544. ISBN: 0-521-79160-X; 0-521-79540-0.\\[0pt]
[Jän87] K. Jänich. Topologie. Second. Hochschultext. [University Textbooks]. Berlin: Springer-Verlag, 1987, S. x+215. ISBN: 3-540-17835-X.\\[0pt]
[Lam09] K. Lamotke. Riemannsche Flächen -. Berlin, Heidelberg: Springer, 2009. ISBN: 978-3-642-01711-7.\\[0pt]
[MF] M. Möller. Modulformen. 2013. eprint: Vorlesungsskript, Frankfurt/Main\\[0pt]
[MKS76] W. Magnus, A. Karrass und D. Solitar. Combinatorial group theory. revised. Presentations of groups in terms of generators and relations. New York: Dover Publications Inc., 1976, S. xii+444.\\[0pt]
[MT02] Howard Masur und Serge Tabachnikov. "Rational billiards and flat structures". In: Hrsg. von B. Hasselblatt und A. Katok. Bd. 1, Part A. Handbook of Dynamical Systems. Elsevier Science, 2002, S. 1015 -1089.\\[0pt]
[Rey89] É. Reyssat. Quelques Aspects Des Surfaces de Riemann. Progress in Mathematics - Birkhäuser. Birkhauser Verlag AG, 1989.\\[0pt]
[Ser03] J.P. Serre. Trees. Springer Monographs in Mathematics. Translated from the French original by John Stillwell, Corrected 2nd printing of the 1980 English translation. Berlin: Springer-Verlag, 2003, S. $x+142$.\\[0pt]
[Vor96] Ya. B. Vorobets. "Plane structures and billiards in rational polygons: the Veech alternative". In: Uspekhi Mat. Nauk 51.5(311) (1996), S. 342.\\[0pt]
[ZK75] A. N. Zemljakov und A. B. Katok. "Topological transitivity of billiards in polygons". In: Mat. Zametki 18.2 (1975), S. 291-300.



\pagebreak
\printbibliography
\end{document}